
\section[{Elements Available in All TEI Documents}]{Elements Available in All TEI Documents}\label{CO}\par
This chapter describes elements which may appear in any kind of text and the tags used to mark them in all TEI documents. Most of these elements are freely floating phrases, which can appear at any point within the textual structure, although they should generally be contained by a higher-level element of some kind (such as a paragraph). A few of the elements described in this chapter (for example, bibliographic citations and lists) have a comparatively well-defined internal structure, but most of them have no consistent inner structure of their own. In the general case, they contain only a few words, and are often identifiable in a conventionally printed text by the use of typographic conventions such as shifts of font, use of quotation or other punctuation marks, or other changes in layout.\par
This chapter begins by describing the \hyperref[TEI.p]{<p>} tag used to mark paragraphs, the prototypical formal unit for running text in many TEI modules. This is followed, in section \textit{\hyperref[COPU]{3.2.\ Treatment of Punctuation}}, by a discussion of some specific problems associated with the interpretation of conventional punctuation, and the methods proposed by these Guidelines for resolving ambiguities therein.\par
The next section (section \textit{\hyperref[COHQ]{3.3.\ Highlighting and Quotation}}) describes a number of phrase-level elements commonly marked by typographic features (and thus well-represented in conventional markup languages). These include features commonly marked by font shifts (section \textit{\hyperref[COHQH]{3.3.2.\ Emphasis, Foreign Words, and Unusual Language}}) and features commonly marked by quotation marks (section \textit{\hyperref[COHQQ]{3.3.3.\ Quotation}}) as well as such features as terms, cited words, and glosses (section \textit{\hyperref[COHTGEDR]{3.4.\ Terms and Glosses, Ruby Annotations, and Equivalents and Descriptions}}).\par
Section \textit{\hyperref[COED]{3.5.\ Simple Editorial Changes}} introduces some phrase-level elements which may be used to record simple editorial interventions, such as emendation or correction of the encoded text. The elements described here constitute a simple subset of the full mechanisms for encoding such information (described in full in chapter \textit{\hyperref[PH]{11.\ Representation of Primary Sources}}), which should be adequate to most commonly encountered situations.\par
The next section (section \textit{\hyperref[CONA]{3.6.\ Names, Numbers, Dates, Abbreviations, and Addresses}}) describes several phrase-level and inter-level elements which, although often of interest for analysis or processing, are rarely explicitly identified in conventional printing. These include names (section \textit{\hyperref[CONARS]{3.6.1.\ Referring Strings}}), numbers and measures (section \textit{\hyperref[CONANU]{3.6.3.\ Numbers and Measures}}), dates and times (section \textit{\hyperref[CONADA]{3.6.4.\ Dates and Times}}), abbreviations (section \textit{\hyperref[CONAAB]{3.6.5.\ Abbreviations and Their Expansions}}), and addresses (section \textit{\hyperref[CONAAD]{3.6.2.\ Addresses}}).\par
In the same way, the following section (section \textit{\hyperref[COXR]{3.7.\ Simple Links and Cross-References}}) presents only a subset of the facilities available for the encoding of cross-references or text-linkage. The full story may be found in chapter \textit{\hyperref[SA]{16.\ Linking, Segmentation, and Alignment}}; the tags presented here are intended to be usable for a wide variety of simple applications.\par
Sections \textit{\hyperref[COLI]{3.8.\ Lists}}, and \textit{\hyperref[CONO]{3.9.\ Notes, Annotation, and Indexing}}, describe two kinds of quasi-structural elements: lists and notes. These may appear either within chunk-level elements such as paragraphs, or between them. Several kinds of lists are catered for, of an arbitrary complexity. The section on notes discusses both notes found in the source and simple mechanisms for adding annotations of an interpretive nature during the encoding; again, only a subset of the facilities described in full elsewhere (specifically, in chapter \textit{\hyperref[AI]{17.\ Simple Analytic Mechanisms}}) is discussed.\par
Section \textit{\hyperref[COGR]{3.10.\ Graphics and Other Non-textual Components}} introduces some simple ways of representing graphic or other non-textual content found in a text. A fuller discussion of the multimedia facilities supported by these Guidelines may be found in chapters \textit{\hyperref[FT]{14.\ Tables, Formulæ, Graphics and Notated Music}} and \textit{\hyperref[SA]{16.\ Linking, Segmentation, and Alignment}}.\par
Next, section \textit{\hyperref[CORS]{3.11.\ Reference Systems}}, describes methods of encoding within a text the conventional system or systems used when making references to the text. Some reference systems have attained canonical authority and should be recorded to make the text useable in normal work; in other cases, a convenient reference system should be created by the creator or analyst of an electronic text.\par
Like lists and notes, the bibliographic citations discussed in section \textit{\hyperref[COBI]{3.12.\ Bibliographic Citations and References}}, may be regarded as structural elements in their own right. A range of possibilities is presented for the encoding of bibliographic citations or references, which may be treated as simple phrases within a running text, or as highly-structured components suitable for inclusion in a bibliographic database.\par
Additional elements for the encoding of passages of verse or drama (whether prose or verse) are discussed in section \textit{\hyperref[CODV]{3.13.\ Passages of Verse or Drama}}.\par
The chapter concludes with a technical overview of the structure and organization of the module described here. This should be read in conjunction with chapter \textit{\hyperref[ST]{1.\ The TEI Infrastructure}}, describing the structure of the TEI document type definition.
\subsection[{Paragraphs}]{Paragraphs}\label{COPA}\par
The paragraph is the fundamental organizational unit for all prose texts, being the smallest regular unit into which prose can be divided. Prose can appear in all TEI texts, even those that are primarily of another genre (e.g., verse); thus the paragraph is described here, as an element which can appear in any kind of text.\par
Paragraphs can contain any of the other elements described within this chapter, as well as some other elements which are specific to individual text types. We distinguish \textit{phrase-level} elements, which must be entirely contained within a paragraph or similar structure and cannot appear except within one, from \textit{chunks}, which can appear between, but not within, paragraphs, and from \textit{inter-level} elements, which can appear either within a single paragraph or between paragraphs. The class of phrases includes emphasized or quoted phrases, names, dates, etc. The class of inter-level elements includes bibliographic citations, notes, lists, etc. The class of chunks includes the paragraph itself, and other elements which have similar structural properties, notably the \hyperref[TEI.ab]{<ab>} (anonymous block) element described in \textit{\hyperref[SASE]{16.3.\ Blocks, Segments, and Anchors}}) which may be used as an alternative to the paragraph in some kinds of texts.\par
Because paragraphs may appear in different base or additional tag sets, their possible contents may differ in different kinds of documents. In particular, additional elements not listed in this chapter may appear in paragraphs in certain kinds of text. However, the elements described in this chapter are always by default available in all kinds of text.\par
The paragraph is marked using the \hyperref[TEI.p]{<p>} element: 
\begin{sansreflist}
  
\item [\textbf{<p>}] (paragraph) marks paragraphs in prose.
\end{sansreflist}
\par
If a consistent internal subdivision of paragraphs is desired, the \hyperref[TEI.s]{<s>} or \hyperref[TEI.seg]{<seg>} (‘segment’) elements may be used, as discussed in chapters \textit{\hyperref[SA]{16.\ Linking, Segmentation, and Alignment}} and \textit{\hyperref[AI]{17.\ Simple Analytic Mechanisms}} respectively. More usually, however, paragraphs have no firm internal structure, but contain prose encoded as a mix of characters, entity references, phrases marked as described in the rest of this chapter, and embedded elements like lists, figures, or tables.\par
Since paragraphs are usually explicitly marked in Western texts, typically by indentation, the application of the \hyperref[TEI.p]{<p>} tag usually presents few problems.\par
In some cases, the body of a text may comprise but a single paragraph: \par\bgroup\index{body=<body>|exampleindex}\index{p=<p>|exampleindex}\exampleFont \begin{shaded}\noindent\mbox{}{<\textbf{body}>}\mbox{}\newline 
\hspace*{1em}{<\textbf{p}>}I fully appreciate Gen. Pope's splendid achievements with their\mbox{}\newline 
\hspace*{1em}\hspace*{1em} invaluable results; but you must know that Major Generalships in the\mbox{}\newline 
\hspace*{1em}\hspace*{1em} Regular Army, are not as plenty as blackberries.{</\textbf{p}>}\mbox{}\newline 
{</\textbf{body}>}\end{shaded}\egroup\par \par
This news story shows typically short journalistic paragraphs: \par\bgroup\index{head=<head>|exampleindex}\index{p=<p>|exampleindex}\index{p=<p>|exampleindex}\index{p=<p>|exampleindex}\index{p=<p>|exampleindex}\exampleFont \begin{shaded}\noindent\mbox{}{<\textbf{head}>}SARAJEVO, Bosnia and Herzegovina, April 19{</\textbf{head}>}\mbox{}\newline 
{<\textbf{p}>}Serbs seized more territory in this struggling new country today as\mbox{}\newline 
 the United States Air Force ended a two-day airlift of humanitarian\mbox{}\newline 
 aid into the capital, Sarajevo.{</\textbf{p}>}\mbox{}\newline 
{<\textbf{p}>}International relief workers called on European Community nations\mbox{}\newline 
 to step up their humanitarian aid to the former Yugoslav republic,\mbox{}\newline 
 in conjunction with new American aid flights if necessary.{</\textbf{p}>}\mbox{}\newline 
{<\textbf{p}>}A special envoy from the European Community, Colin Doyle, harshly\mbox{}\newline 
 condemned the decision by Serbs to shell Sarajevo on Saturday night\mbox{}\newline 
 during a visit to the Bosnian capital by a senior American official,\mbox{}\newline 
 Deputy Assistant Secretary of State Ralph R. Johnson.{</\textbf{p}>}\mbox{}\newline 
{<\textbf{p}>}...{</\textbf{p}>}\end{shaded}\egroup\par \par
The following extract from a Russian fairy tale demonstrates how other phrase level elements (in this case \hyperref[TEI.q]{<q>} elements representing direct speech; see section \textit{\hyperref[COHQQ]{3.3.3.\ Quotation}}) may be nested within, but not across, paragraphs: \par\bgroup\index{p=<p>|exampleindex}\index{q=<q>|exampleindex}\index{q=<q>|exampleindex}\index{q=<q>|exampleindex}\index{p=<p>|exampleindex}\index{q=<q>|exampleindex}\index{q=<q>|exampleindex}\index{q=<q>|exampleindex}\index{p=<p>|exampleindex}\index{q=<q>|exampleindex}\index{q=<q>|exampleindex}\index{q=<q>|exampleindex}\exampleFont \begin{shaded}\noindent\mbox{}{<\textbf{p}>}A fly built a castle, a tall and mighty castle.\mbox{}\newline 
 There came to the castle the Crawling Louse. {<\textbf{q}>}Who,\mbox{}\newline 
\hspace*{1em}\hspace*{1em} who's in the castle? Who, who's in your house?{</\textbf{q}>}\mbox{}\newline 
 said the Crawling Louse. {<\textbf{q}>}I, I, the Languishing Fly.\mbox{}\newline 
\hspace*{1em}\hspace*{1em} And who art thou?{</\textbf{q}>}\mbox{}\newline 
\hspace*{1em}{<\textbf{q}>}I'm the Crawling Louse.{</\textbf{q}>}\mbox{}\newline 
{</\textbf{p}>}\mbox{}\newline 
{<\textbf{p}>}Then came to the castle the Leaping Flea. {<\textbf{q}>}Who,\mbox{}\newline 
\hspace*{1em}\hspace*{1em} who's in the castle?{</\textbf{q}>} said the Leaping Flea. {<\textbf{q}>}I,\mbox{}\newline 
\hspace*{1em}\hspace*{1em} I, the Languishing Fly, and I, the Crawling Louse. And\mbox{}\newline 
\hspace*{1em}\hspace*{1em} who art thou?{</\textbf{q}>}\mbox{}\newline 
\hspace*{1em}{<\textbf{q}>}I'm the Leaping Flea.{</\textbf{q}>}\mbox{}\newline 
{</\textbf{p}>}\mbox{}\newline 
{<\textbf{p}>}Then came to the castle the Mischievous Mosquito.\mbox{}\newline 
{<\textbf{q}>}Who, who's in the castle?{</\textbf{q}>} said the Mischievous\mbox{}\newline 
 Mosquito. {<\textbf{q}>}I, I, the Languishing Fly, and I, the\mbox{}\newline 
\hspace*{1em}\hspace*{1em} Crawling Louse, and I, the Leaping Flea. And who art\mbox{}\newline 
\hspace*{1em}\hspace*{1em} thou?{</\textbf{q}>}\mbox{}\newline 
\hspace*{1em}{<\textbf{q}>}I'm the Mischievous Mosquito.{</\textbf{q}>}\mbox{}\newline 
{</\textbf{p}>}\end{shaded}\egroup\par 
\subsection[{Treatment of Punctuation}]{Treatment of Punctuation}\label{COPU}\par
Punctuation marks cause two distinct classes of problem for text markup: the marks may not be available in the character set used, and they may be significantly ambiguous. To some extent, the availability of the Unicode character set addresses the first of these problems, since it provides specific code points for most punctuation marks, and also the second to the extent that it distinguishes glyphs (such as stop, comma, and hyphen) which are used with different functions.  Where punctuation itself is the subject of study, the element \hyperref[TEI.pc]{<pc>} (punctuation character) may be used to mark it explicitly, as further discussed in \textit{\hyperref[AIPC]{17.1.2.\ Below the Word Level}}. Where the character used for a punctuation mark is not available in Unicode, the \hyperref[TEI.g]{<g>} element and other facilities described in chapter \textit{\hyperref[WD]{5.\ Characters, Glyphs, and Writing Modes}} may also be used to mark its presence.
\subsubsection[{Functions of Punctuation}]{Functions of Punctuation}\label{COPU-1}\par
Punctuation is itself a form of markup, historically introduced to provide the reader with an indication about how the text should be read. As such, it is unsurprising that encoders will often wish to encode directly the purpose for which punctuation was provided, as well as, or even instead of, the punctuation itself. We discuss some typical cases below.\par
The \textit{Full stop (period)} may mark (orthographic) sentence boundaries, abbreviations, decimal points, or serve as a visual aid in printing numbers. These usages can be distinguished by tagging S-units, abbreviations, and numbers, as described in sections \textit{\hyperref[SASE]{16.3.\ Blocks, Segments, and Anchors}}, \textit{\hyperref[CONAAB]{3.6.5.\ Abbreviations and Their Expansions}}, and \textit{\hyperref[CONANU]{3.6.3.\ Numbers and Measures}} respectively. However, there are independent reasons for tagging these, whether or not they are marked by full stops, and the polysemy of the full stop itself is perhaps no different from that of any other character in the writing system.\par
The \textit{Question mark} and \textit{exclamation mark} usually mark the end of orthographic sentences, but may also be used as a mid-sentence comment by the author (\textit{!} to express surprise or some other strong feeling, \textit{?} to query a word or expression or mark a sentence as dubious in linguistic discussion). Such usages may be distinguished by marking S-units, in which case the mid-sentence uses of these punctuation marks may be left unmarked, or tagged using the \hyperref[TEI.pc]{<pc>} element discussed in \textit{\hyperref[AILC]{17.1.\ Linguistic Segment Categories}}.\par
\textit{Dashes} are used for a variety of purposes: as a mark of omission, insertion, or interruption; to show where a new speaker takes over (in dialogue); or to introduce a list item. In the latter two cases particularly, it is clearly desirable to mark the function as well as its rendition using the elements \hyperref[TEI.q]{<q>} or \hyperref[TEI.item]{<item>}, on which see section \textit{\hyperref[COHQQ]{3.3.3.\ Quotation}}, and section \textit{\hyperref[COLI]{3.8.\ Lists}}, respectively.\par
\textit{Quotation marks} may be removed from text contained by \hyperref[TEI.q]{<q>} or \hyperref[TEI.quote]{<quote>} elements on editorial grounds, or they may be marked in a variety of ways; see the discussion of quotation and related features in section \textit{\hyperref[COHQQ]{3.3.3.\ Quotation}}.\par
\textit{Apostrophes} should be distinguished from single quote marks. As with hyphens, this disambiguation is best performed by selecting the appropriate Unicode character, though it may also be represented by using appropriate XML markup for quotations as suggested above. However, apostrophes have a variety of uses. In English they mark contractions, genitive forms, and (occasionally) plural forms. Full disambiguation of these uses belongs to the level of linguistic analysis and interpretation.\par
\textit{Parentheses} and other marks of suspension such as dashes or ellipses are often used to signal information about the syntactic structure of a text fragment. Full disambiguation of their uses also belongs to the level of linguistic analysis and interpretation, and will therefore need to use the mechanisms discussed in chapter \textit{\hyperref[AI]{17.\ Simple Analytic Mechanisms}}.\par
Where punctuation marks are disambiguated by tagging their assumed function in the text (for example, quotation), it may be debated whether they should be excluded or left as part of the text. In the case of quotation marks, it may be more convenient to distinguish opening from closing marks simply by using the appropriate Unicode character than to use the \hyperref[TEI.q]{<q>} element, with or without an indication of rendition.\par
Where segmentation of a text is performed automatically, the accuracy of the result may be considerably enhanced by a first pass in which the function of different punctuation characters is explicitly marked. This need not be done for all cases, but only where the structural function of the punctuation markup (for example as a word or phrase delimiter) is ambiguous. Thus, dots indicating abbreviation might be distinguished from dots indicating sentence end, and exclamation or question marks internal to a sentence distinguished from those which terminate one. Furthermore, when encoding historical materials, it may be considered essential to retain the original punctuation, whether by using an appropriate character code, if this is available (or using the \hyperref[TEI.g]{<g>} element where it is not) or by an explicit encoding using \hyperref[TEI.pc]{<pc>}. The particular method adopted will vary depending upon the feature concerned and upon the purpose of the project.
\subsubsection[{Hyphenation}]{Hyphenation}\label{COPU-2}\par
Hyphenation as a phenomenon is generally of most concern when producing formatted text for display in print or on screen: different languages and systems have developed quite sophisticated sets of rules about where hyphens may be introduced and for what reason. These generally do not concern the text encoder, since they belong to the domain of formatting and will generally be handled by the rendition software in use. In this section, we discuss issues arising from the appearance of hyphens in pre-existing formatted texts which are being re-encoded for analysis or other processing. Unicode distinguishes four characters visually similar to the hyphen, including the undifferentiated hyphen-minus (U+002D) which is retained for compatibility reasons. The hard hyphen (U+2010) is distinguished from the minus sign (U+2212) which is for use in mathematical expressions, and also from the soft hyphen (U+00AD) which may appear in ‘born digital’ documents to indicate places where it is acceptable to insert a hyphen when the document is formatted.\par
Historically, the hard hyphen has been used in printed or manuscript documents for two distinct purposes. In many languages, it is used between words to show that they function as a single syntactic or lexical unit. For example, in French, \textit{est-ce que}; in English \textit{body-snatcher}, \textit{tea-party} etc. It may also have an important role in disambiguation (for example, by distinguishing say a \textit{man-eating fish} from a \textit{man eating fish}). Such usages, although possibly problematic when a linguistic analysis is undertaken, are not generally of concern to text encoders: the hyphen character is usually retained in the text, because it may be regarded as part of the way a compound or other lexical item is spelled. Deciding whether a compound is to be decomposed into its constituent parts, and if so how, is a different question, involving consideration of many other phenomena in addition to the simple presence of a hyphen.\par
When it appears at the end of a printed or written line however, the hard hyphen generally indicates that—contrary to what might be expected—a word is not yet complete, but continues on the next line (or over the next page or column or other boundary). The hyphen character is not, in this case, part of the word, but just a signal that the word continues over the break. Unfortunately, few languages distinguish these two cases visually, which necessarily poses a problem for text encoders. Suppose, for example, that we wish to investigate a diachronic English corpus for occurrences of \textit{tea-pot} and \textit{teapot}, to find evidence for the point at which this compound becomes lexicalized. Any case where the word is hyphenated across a linebreak, like this: \par\hfill\bgroup\exampleFont\vskip 10pt\begin{shaded}
\obeyspaces tea-\newline
pot\end{shaded}
\par\egroup 
 is ambiguous: there is no simple way of deciding which of the two spellings was intended.\par
As elsewhere, therefore, encoders have a range of choices: \begin{itemize}
\item They may decide simply to remove any end-of-line hyphenation from the encoded text, on the grounds that its presence is purely a secondary matter of formatting. This will obviously apply also if line endings are themselves regarded as unimportant.
\item Alternatively, they may decide to record the presence of the hyphen, perhaps on the grounds that it provides useful morphological information; perhaps in order to retain information about the visual appearance of the original source. In either case, they need to decide whether to record it explicitly, by including an appropriate punctuation character in the text data, or implicitly by supplying an appropriate symbolic value for one or more of the attributes on the \hyperref[TEI.lb]{<lb>} or other milestone element used to record the fact of the line division. If the hyphen is included in the character data of the TEI document, it might be marked up using the \hyperref[TEI.pc]{<pc>} (punctuation character) tag, which allows the encoder to express information about its function as a separator, through the {\itshape force} attribute. For example, the example above could be encoded with a {\itshape force} value of "inter" to indicate that the punctuation mark may or may not be a word separator (See also \textit{\hyperref[AIPC]{17.1.2.\ Below the Word Level}}). \par\bgroup\index{pc=<pc>|exampleindex}\index{force=@force!<pc>|exampleindex}\exampleFont \begin{shaded}\noindent\mbox{}tea{<\textbf{pc}\hspace*{1em}{force}="{inter}">}-{</\textbf{pc}>}pot\end{shaded}\egroup\par 
\end{itemize}  A similar range of possibilities applies equally to the representation of other common punctuation marks, notably quotation marks, as discussed in \textit{\hyperref[COHQQ]{3.3.3.\ Quotation}}.\par
The ‘text data’ of which XML documents are composed is decomposable into smaller units, here called \textit{orthographic tokens}, even if those units are not explicitly indicated by the XML markup. The ambiguity of the end-of-line hyphen also causes problems in the way a processor identifies such tokens in the absence of explicit markup. If token boundaries are not explicitly marked (for example using the \hyperref[TEI.seg]{<seg>} or \hyperref[TEI.w]{<w>} elements), for most languages a processor will rely on character class information to determine where they are to be found: some punctuation characters are considered to be word-breaking, while others are not. In XML, the newline character in text data is a kind of whitespace, and is therefore word breaking. However, it is generally unsafe to assume that whitespace adjacent to markup tags will always be preserved, and it is decidedly unsafe to assume that markup tags themselves are equivalent to whitespace.\par
The \hyperref[TEI.lb]{<lb>}, \hyperref[TEI.pb]{<pb>}, and \hyperref[TEI.cb]{<cb>} elements are notable exceptions to this general rule, since their function is precisely to represent (or replace) line, page, or column breaks, which, as noted above, are generally considered to be equivalent to whitespace. These elements provide a more reliable way of preserving the lineation, pagination, etc of a source document, since the encoder should not assume that (untagged) line breaks etc. in an XML source file will necessarily be preserved.\par
To control the intended tokenization, the encoder may use the {\itshape break} attribute on such elements to indicate whether or not the element is to be regarded as equivalent to whitespace. This attribute can take the values yes or no to indicate whether or not the element corresponds with a token boundary. The value maybe is also available, for cases where the encoder does not wish (or is unable) to determine whether the orthographic token concerned is broken by the line ending.\par
As a final complication, it should be noted that in some languages, particularly German and Dutch, the spelling of a word may be altered in the presence of end of line hyphenation. For example, in Dutch, the word \textit{opaatje} ( \textit{granddad}), occurring at the end of a line may be hyphenated as \textit{opa-tje}, with a single letter a. An encoder wishing to preserve the original form of this orthographic token in a printed text while at the same time facilitating its recognition as the word \textit{opaatje} will therefore need to rely on a more sophisticated process than simply removing the hyphen. This is however essentially the same as any other form of normalization accompanying the recognition of variations in spelling or morphology: as such it may be encoded using the \hyperref[TEI.choice]{<choice>} element discussed in \textit{\hyperref[COED]{3.5.\ Simple Editorial Changes}}, or the more sophisticated mechanisms for linguistic analysis discussed in chapter \textit{\hyperref[AI]{17.\ Simple Analytic Mechanisms}}.
\subsection[{Highlighting and Quotation}]{Highlighting and Quotation}\label{COHQ}\par
This section deals with a variety of textual features, all of which have in common that they are frequently realized in conventional printing practice by the use of such features as underlining, italic fonts, or quotation marks, collectively referred to here as \textit{highlighting}. After an initial discussion of this phenomenon and alternate approaches to encoding it, this section describes ways of encoding the following textual features, all of which are conventionally rendered using some kind of highlighting: \begin{itemize}
\item emphasis, foreign words and other linguistically distinct uses of highlighting
\item representation of speech and thought, quotation, etc.
\item technical terms, glosses, etc.
\end{itemize} 
\subsubsection[{What Is Highlighting?}]{What Is Highlighting?}\label{COHQW}\par
By \textit{highlighting} we mean the use of any combination of typographic features (font, size, hue, etc.) in a printed or written text in order to distinguish some passage of a text from its surroundings.\footnote{Although the way in which a spoken text is performed, (for example, the voice quality, loudness, etc.) might be regarded as analogous to ‘highlighting’ in this sense, these Guidelines recommend distinct elements for the encoding of such ‘highlighting’ in spoken texts. See further section \textit{\hyperref[TSSASH]{8.3.6.\ Shifts}}.} The purpose of highlighting is generally to draw the reader's attention to some feature or characteristic of the passage highlighted; this section describes the elements recommended by these Guidelines for the encoding of such textual features.\par
In conventionally printed modern texts, highlighting is often employed to identify words or phrases which are regarded as being one or more of the following: \begin{itemize}
\item distinct in some way—as foreign, dialectal, archaic, technical, etc.
\item emphatic, and which would for example be stressed when spoken
\item not part of the body of the text, for example cross-references, titles, headings, labels, etc.
\item identified with a distinct narrative stream, for example an internal monologue or commentary.
\item attributed by the narrator to some other agency, either within the text or outside it: for example, direct speech or quotation.
\item set apart from the text in some other way: for example, proverbial phrases, words mentioned but not used, names of persons and places in older texts, editorial corrections or additions, etc.
\end{itemize} \par
The textual functions indicated by highlighting may not be rendered consistently in different parts of a text or in different texts. (For example, a foreign word may appear in italics if the surrounding text is in roman, but in roman if the surrounding text is in italics.) For this reason, these Guidelines distinguish between the encoding of rendering itself and the encoding of the underlying feature expressed by it. \par
Highlighting as such may be encoded by using one of the global attributes {\itshape rend}, {\itshape rendition}, or {\itshape style} (see further \textit{\hyperref[STGA]{1.3.1.1.\ Global Attributes}}). This allows the encoder both to specify the function of a highlighted phrase or word, by selecting the appropriate element described here or elsewhere in these Guidelines, and to further describe the way in which it is highlighted, by means of an attribute. If the encoder wishes to offer no interpretation of the feature underlying the use of highlighting in the source text, then the \hyperref[TEI.hi]{<hi>} element may be used, which indicates only that the text so tagged was highlighted in some way. 
\begin{sansreflist}
  
\item [\textbf{<hi>}] (highlighted) marks a word or phrase as graphically distinct from the surrounding text, for reasons concerning which no claim is made.
\end{sansreflist}
 The \hyperref[TEI.hi]{<hi>} element is provided by the \textsf{model.hiLike} class.\par
The possible values carried by the {\itshape rend} attribute are not formally defined in this version of the Guidelines. It may be used to document any peculiarity of the way a given segment of text was rendered in the original source text, and may thus express a very large range of typographic or other features, by no means restricted to typeface, type size, etc. The {\itshape style} attribute, by contrast, defines the way the source text was rendered using a formally defined style language, such as the W3C standard Cascading Stylesheet Language (\cite{CSS1}). The complementary {\itshape rendition} attribute is used to point to one or more fragments expressed using such a language which have been predefined in the TEI header using the \hyperref[TEI.rendition]{<rendition>} element discussed in section \textit{\hyperref[HD57]{2.3.4.\ The Tagging Declaration}}.\par
Where it is both appropriate and feasible, these Guidelines recommend that the textual feature marked by the highlighting should be encoded, rather than just the simple fact of the highlighting. This is for the following reasons: \begin{itemize}
\item the same kind of highlighting may be used for different purposes in different contexts
\item the same textual function may be highlighted in different ways in different contexts
\item for analytic purposes, it is in general more useful to know the intended function of a highlighted phrase than simply that it is distinct.
\end{itemize} \par
In many, if not most, cases the underlying function of a highlighted phrase will be obvious and non-controversial, since the distinctions indicated by a change of highlighting correspond with distinctions discussed elsewhere in these Guidelines. The elements available to record such distinctions are, for the most part, members of the \textsf{model.emphLike} class. This and the \textsf{model.hiLike} class mentioned above constitute the \textsf{model.highlighted} class, which is a phrase level class. Members of this class may appear anywhere within paragraph level elements.\par
The distinction between the two classes is simple, and typified by the two elements \hyperref[TEI.hi]{<hi>} and \hyperref[TEI.emph]{<emph>}: the former marks simply that a passage is typographically distinct in some way, while the latter asserts that a passage is linguistically emphasized for some purpose. These two properties, though often combined, are not identical. It should however be recognized, however, that cases do exist in which it is not economically feasible to mark the underlying function (e.g. in the preparation of large text corpora), as well as cases in which it is not intellectually appropriate (as in the transcription of some older materials, or in the preparation of material for the study of typographic practice). In such cases, the \hyperref[TEI.hi]{<hi>} element or some other element from the \textsf{model.hiLike} class should be used.\par
Elements which are sometimes realized by typographic distinction but which are not discussed in this section include \hyperref[TEI.title]{<title>} (discussed in section \textit{\hyperref[COBI]{3.12.\ Bibliographic Citations and References}}) and \hyperref[TEI.name]{<name>} (discussed in section \textit{\hyperref[CONARS]{3.6.1.\ Referring Strings}}).
\subsubsection[{Emphasis, Foreign Words, and Unusual Language}]{Emphasis, Foreign Words, and Unusual Language}\label{COHQH}\par
This subsection discusses the following elements: 
\begin{sansreflist}
  
\item [\textbf{<foreign>}] (foreign) identifies a word or phrase as belonging to some language other than that of the surrounding text.
\item [\textbf{<emph>}] (emphasized) marks words or phrases which are stressed or emphasized for linguistic or rhetorical effect.
\item [\textbf{<distinct>}] identifies any word or phrase which is regarded as linguistically distinct, for example as archaic, technical, dialectal, non-preferred, etc., or as forming part of a sublanguage.
\end{sansreflist}
 These elements are all members of the \textsf{model.emphLike} class.
\paragraph[{Foreign Words or Expressions}]{Foreign Words or Expressions}\label{COHQHF}\par
Words or phrases which are not in the main language of the text should be tagged as such, at least where the fact is indicated in the text. Where the word or phrase concerned is already distinguished from the rest of the text by virtue of its function (for example, because it is a name, a technical term, a quotation, a mentioned word, etc.) then the global {\itshape xml:lang} attribute should be used to specify additionally that its language distinguishes it from the surrounding text. Any element in the TEI scheme may take an {\itshape xml:lang} attribute, which specifies both the writing system and the language used by its content (see sections \textit{\hyperref[STGAla]{1.3.1.1.2.\ Language Indicators}} and \textit{\hyperref[CHSH]{vi.1\ Language Identification}} for discussion of this attribute and its values respectively). Where there is no other applicable element, the element \hyperref[TEI.foreign]{<foreign>} may be used to provide a peg onto which the {\itshape xml:lang} may be attached. \par\bgroup\index{q=<q>|exampleindex}\index{foreign=<foreign>|exampleindex}\index{foreign=<foreign>|exampleindex}\index{q=<q>|exampleindex}\exampleFont \begin{shaded}\noindent\mbox{}{<\textbf{q}>}Aren't you confusing {<\textbf{foreign}\hspace*{1em}{xml:lang}="{la}">}post hoc{</\textbf{foreign}>} with {<\textbf{foreign}\hspace*{1em}{xml:lang}="{la}">}propter\mbox{}\newline 
\hspace*{1em}\hspace*{1em} hoc{</\textbf{foreign}>}?{</\textbf{q}>} said the Bee Master. \mbox{}\newline 
{<\textbf{q}>}Wax-moth only succeed when\mbox{}\newline 
 weak bees let them in.{</\textbf{q}>}\end{shaded}\egroup\par \par
The \hyperref[TEI.foreign]{<foreign>} element should not be used to represent foreign words which are mentioned or glossed within the text: for these use the appropriate element from section \textit{\hyperref[COHTG]{3.4.1.\ Terms and Glosses}} below. Compare the following example sentences: \par\bgroup\index{foreign=<foreign>|exampleindex}\exampleFont \begin{shaded}\noindent\mbox{}John eats a {<\textbf{foreign}\hspace*{1em}{xml:lang}="{fr}">}croissant{</\textbf{foreign}>} every morning.\end{shaded}\egroup\par \noindent  \par\bgroup\index{mentioned=<mentioned>|exampleindex}\exampleFont \begin{shaded}\noindent\mbox{}{<\textbf{mentioned}\hspace*{1em}{xml:lang}="{fr}">}Croissant{</\textbf{mentioned}>} is difficult to\mbox{}\newline 
 pronounce with your mouth full.\end{shaded}\egroup\par \noindent  \par\bgroup\index{term=<term>|exampleindex}\exampleFont \begin{shaded}\noindent\mbox{}A {<\textbf{term}\hspace*{1em}{xml:lang}="{fr}">}croissant{</\textbf{term}>} is a crescent-shaped\mbox{}\newline 
 piece of light, buttery, pastry that is usually eaten for\mbox{}\newline 
 breakfast, especially in France.\end{shaded}\egroup\par \par
Elements which do not explicitly state the language of their content by means of an {\itshape xml:lang} attribute are understood to inherit a value for it from their parent element. In the general case, therefore, it is recommended practice to supply a default value for {\itshape xml:lang} on the root \hyperref[TEI.TEI]{<TEI>} or \hyperref[TEI.text]{<text>} element, as further discussed in section \textit{\hyperref[STGAla]{1.3.1.1.2.\ Language Indicators}}
\paragraph[{Emphatic Words and Phrases}]{Emphatic Words and Phrases}\label{COHQHE}\par
The \hyperref[TEI.emph]{<emph>} element is provided to mark words or phrases which are \textit{linguistically} emphatic or stressed. Text which is only typographically ‘emphasized’ falls into the class of highlighted text, and may be tagged with the \hyperref[TEI.hi]{<hi>} element. In printed works, emphasis is generally indicated by devices such as the use of an italic font, a large typeface, or extra wide letter spacing; in manuscripts and typescripts, it is usually indicated by the use of underlining. As the following examples demonstrate, an encoder may choose whether or not to make explicit the particular type of rendition associated with the emphasis. If a source text consistently renders a particular feature (e.g. emphasis or words in foreign languages) in a particular way, the rendering associated with that feature may be described in the TEI header using the \hyperref[TEI.rendition]{<rendition>} element. The {\itshape rend}, {\itshape rendition}, or {\itshape style} attributes may then be used to describe examples which deviate from the norm. For example, assuming that the TEI header has defined a default rendering for the \hyperref[TEI.emph]{<emph>} element, the following encoding would use it: \par\bgroup\index{q=<q>|exampleindex}\index{emph=<emph>|exampleindex}\exampleFont \begin{shaded}\noindent\mbox{}{<\textbf{q}>}Sex, sir, is {<\textbf{emph}>}purely{</\textbf{emph}>} a question\mbox{}\newline 
 of appetite!{</\textbf{q}>} Tarr exclaimed.\end{shaded}\egroup\par \noindent  If on the other hand no such default has been defined for the element, the encoder may specify it informally using the {\itshape rend} attribute: \par\bgroup\index{q=<q>|exampleindex}\index{q=<q>|exampleindex}\index{emph=<emph>|exampleindex}\index{rend=@rend!<emph>|exampleindex}\exampleFont \begin{shaded}\noindent\mbox{}{<\textbf{q}>}What it all comes to is this,{</\textbf{q}>} he said.\mbox{}\newline 
\mbox{}\newline 
{<\textbf{q}>}\mbox{}\newline 
\hspace*{1em}{<\textbf{emph}\hspace*{1em}{rend}="{italic}">}What does Christopher Robin do in the morning\mbox{}\newline 
\hspace*{1em}\hspace*{1em} nowadays?{</\textbf{emph}>}\mbox{}\newline 
{</\textbf{q}>}\end{shaded}\egroup\par \noindent  If the encoder wishes to express information about the rendition used in the source using a formal language such as CSS, then the {\itshape style} attribute can be used in a similar way: \par\bgroup\index{q=<q>|exampleindex}\index{q=<q>|exampleindex}\index{emph=<emph>|exampleindex}\index{style=@style!<emph>|exampleindex}\exampleFont \begin{shaded}\noindent\mbox{}{<\textbf{q}>}What it all comes to is this,{</\textbf{q}>} he said.\mbox{}\newline 
\mbox{}\newline 
{<\textbf{q}>}\mbox{}\newline 
\hspace*{1em}{<\textbf{emph}\hspace*{1em}{style}="{font-style: italic}">}What does Christopher Robin do in\mbox{}\newline 
\hspace*{1em}\hspace*{1em} the morning nowadays?{</\textbf{emph}>}\mbox{}\newline 
{</\textbf{q}>}\end{shaded}\egroup\par \par
In cases where the rendition of a source needs to be indicated several times in a document, it may be more convenient to provide a default value using the \hyperref[TEI.rendition]{<rendition>} element in the header. If a small number of distinct values are required, it may also be convenient to define them all by means of a series of \hyperref[TEI.rendition]{<rendition>} elements which can then be referenced from the elements in question by means of the global {\itshape rendition} attribute: \par\bgroup\index{l=<l>|exampleindex}\index{name=<name>|exampleindex}\index{rendition=@rendition!<name>|exampleindex}\index{l=<l>|exampleindex}\index{emph=<emph>|exampleindex}\index{rendition=@rendition!<emph>|exampleindex}\index{rendition=<rendition>|exampleindex}\index{scheme=@scheme!<rendition>|exampleindex}\exampleFont \begin{shaded}\noindent\mbox{}{<\textbf{l}>}Here Thou, great {<\textbf{name}\hspace*{1em}{rendition}="{\#italic}">}Anna{</\textbf{name}>}!\mbox{}\newline 
 whom three Realms obey,{</\textbf{l}>}\mbox{}\newline 
{<\textbf{l}>}Doth sometimes Counsel take —\mbox{}\newline 
 and sometimes {<\textbf{emph}\hspace*{1em}{rendition}="{\#italic}">}Tea{</\textbf{emph}>}.{</\textbf{l}>}\mbox{}\newline 
\textit{<!-- in the header ... -->}\mbox{}\newline 
{<\textbf{rendition}\hspace*{1em}{xml:id}="{italic}"\hspace*{1em}{scheme}="{css}">}font-style: italic{</\textbf{rendition}>}\end{shaded}\egroup\par \noindent  Further information on the use of the \hyperref[TEI.rendition]{<rendition>} element is provided at \textit{\hyperref[HD57]{2.3.4.\ The Tagging Declaration}}.\par
The \hyperref[TEI.hi]{<hi>} element is used to mark words or phrases which are highlighted in some way, but for which identification of the intended distinction is difficult, controversial, or impossible. It enables an encoder simply to record the fact of highlighting, possibly describing it by the use of a {\itshape rend}, {\itshape style}, or {\itshape rendition} attribute, as discussed above, without however taking a position as to the function of the highlighting. This may also be useful if the text is to be processed in two stages: representing simply typographic distinctions during a first pass, and then replacing the \hyperref[TEI.hi]{<hi>} elements with more specific elements in a second pass.\par
Some simple examples: \par\bgroup\index{hi=<hi>|exampleindex}\index{rend=@rend!<hi>|exampleindex}\index{hi=<hi>|exampleindex}\index{rend=@rend!<hi>|exampleindex}\exampleFont \begin{shaded}\noindent\mbox{}{<\textbf{hi}\hspace*{1em}{rend}="{gothic}">}And this Indenture further witnesseth{</\textbf{hi}>}\mbox{}\newline 
 that the said {<\textbf{hi}\hspace*{1em}{rend}="{italic}">}Walter Shandy{</\textbf{hi}>}, merchant,\mbox{}\newline 
 in consideration of the said intended marriage ...\end{shaded}\egroup\par \noindent  In this example, the first highlighted phrase uses black letter or gothic print to mimic the appearance of a legal document, and italic to mark \textit{Walter Shandy} as a name. In a second pass, the elements \hyperref[TEI.head]{<head>} or \hyperref[TEI.label]{<label>} might be appropriate for the first use, and the element \hyperref[TEI.name]{<name>} for the second. \par\bgroup\index{hi=<hi>|exampleindex}\index{rend=@rend!<hi>|exampleindex}\exampleFont \begin{shaded}\noindent\mbox{}The heaviest rain, and snow, and hail, and\mbox{}\newline 
 sleet, could boast of the advantage over him in only one respect. They\mbox{}\newline 
 often {<\textbf{hi}\hspace*{1em}{rend}="{quoted}">}came down{</\textbf{hi}>} handsomely, and Scrooge never\mbox{}\newline 
 did.\end{shaded}\egroup\par \noindent  In this example, the phrase \textit{came down} uses inverted commas to indicate a play on words.\footnote{The Oxford English Dictionary documents the phrase \textit{to come down} in the sense ‘to bring or put down; {\itshape esp.} to lay down money; to make a disbursement’ as being in use, mostly in colloquial or humorous contexts, from at least 1700 to the latter half of the 19th century.} In a second pass, the element \hyperref[TEI.soCalled]{<soCalled>} might be preferred as a means of indicating that the narrator is distancing himself from this usage.
\paragraph[{Other Linguistically Distinct Material}]{Other Linguistically Distinct Material}\label{COHQHD}\par
For some kinds of analysis, it may be desirable to encode the linguistic distinctiveness of words and phrases with more delicacy than is allowed by the \hyperref[TEI.foreign]{<foreign>} element. The \hyperref[TEI.distinct]{<distinct>} element is provided for this purpose. Its attributes allow for additional information characterizing the nature of the linguistic distinction to be made in two distinct ways: the {\itshape type} attribute simply assigns a user-defined code of some kind to the word or phrase which assigns it to some register, sub-language, etc. No recommendations as to the set of values for this attribute are provided at this time, as little consensus exists in the field.\par
Alternatively, the remaining three attributes may be used in combination to place a word or phrase on a three-dimensional scale sometimes used in descriptive linguistics, as for example in \hyperref[CO-BIBL-1]{Mattheier et al, 1988}. The {\itshape time} attribute places a word or phrase \textit{diachronically}, for example as archaic, old-fashioned, contemporary, futuristic, etc.; the {\itshape space} attribute places a word or phrase \textit{diatopically}, that is, with respect to a geographical classification, for example as national, regional, international, etc.; the {\itshape social} attribute places a word or phrase \textit{diastratically}, that is, with respect to a social classification, for example as technical, polite, impolite, restricted, etc. Again, no recommendations are made for the values of these attributes at this time; the encoder should provide a description of the scheme used in the appropriate section of the header (see section \textit{\hyperref[HD5]{2.3.\ The Encoding Description}}).\par
Examples: \par\bgroup\index{distinct=<distinct>|exampleindex}\index{type=@type!<distinct>|exampleindex}\index{distinct=<distinct>|exampleindex}\index{type=@type!<distinct>|exampleindex}\exampleFont \begin{shaded}\noindent\mbox{}Next morning a boy in that dormitory confided to his\mbox{}\newline 
 bosom friend, a {<\textbf{distinct}\hspace*{1em}{type}="{psSlang}">}fag{</\textbf{distinct}>} of\mbox{}\newline 
 Macrea's, that there was trouble in their midst which\mbox{}\newline 
 King {<\textbf{distinct}\hspace*{1em}{type}="{archaic}">}would fain{</\textbf{distinct}>} keep\mbox{}\newline 
 secret.\end{shaded}\egroup\par \noindent  \par\bgroup\index{distinct=<distinct>|exampleindex}\index{time=@time!<distinct>|exampleindex}\index{space=@space!<distinct>|exampleindex}\index{social=@social!<distinct>|exampleindex}\index{distinct=<distinct>|exampleindex}\index{time=@time!<distinct>|exampleindex}\exampleFont \begin{shaded}\noindent\mbox{}Next morning a boy in that dormitory confided to his\mbox{}\newline 
 bosom friend, a\mbox{}\newline 
{<\textbf{distinct}\hspace*{1em}{time}="{1900}"\hspace*{1em}{space}="{GB}"\mbox{}\newline 
\hspace*{1em}{social}="{publicschool}">}fag{</\textbf{distinct}>}\mbox{}\newline 
 of Macrea's, that there was trouble in their midst which\mbox{}\newline 
 King {<\textbf{distinct}\hspace*{1em}{time}="{archaic}">}would fain{</\textbf{distinct}>} keep\mbox{}\newline 
 secret.\end{shaded}\egroup\par \noindent  Where more complex (or more rigorous) interpretive analyses of the associations of a word are required, the more detailed and general mechanisms described in chapter \textit{\hyperref[FS]{18.\ Feature Structures}} should be preferred to these simple characterizations. It may also be preferable to record the kinds of analysis suggested here by means of the simple annotation element \hyperref[TEI.note]{<note>} described in section \textit{\hyperref[CONO]{3.9.\ Notes, Annotation, and Indexing}}, or the \hyperref[TEI.span]{<span>} element described in section \textit{\hyperref[AISP]{17.3.\ Spans and Interpretations}}.
\subsubsection[{Quotation}]{Quotation}\label{COHQQ}\par
One form of presentational variation found particularly frequently in written and printed texts is the use of quotation marks. As with the typographic variations discussed in the preceding section, it is generally helpful to separate the encoding of the underlying textual feature (for example, a quotation or a piece of direct speech) from the encoding of its rendering (for example, the use of a particular style of quotation marks).\par
This section discusses the following elements, all of which are often rendered by the use of quotation marks: 
\begin{sansreflist}
  
\item [\textbf{<q>}] (quoted) contains material which is distinguished from the surrounding text using quotation marks or a similar method, for any one of a variety of reasons including, but not limited to: direct speech or thought, technical terms or jargon, authorial distance, quotations from elsewhere, and passages that are mentioned but not used.
\item [\textbf{<said>}] (speech or thought) indicates passages thought or spoken aloud, whether explicitly indicated in the source or not, whether directly or indirectly reported, whether by real people or fictional characters.\hfil\\[-10pt]\begin{sansreflist}
    \item[@{\itshape direct}]
  may be used to indicate whether the quoted matter is regarded as direct or indirect speech.
    \item[@{\itshape aloud}]
  may be used to indicate whether the quoted matter is regarded as having been vocalized or signed.
\end{sansreflist}  
\item [\textbf{<quote>}] (quotation) contains a phrase or passage attributed by the narrator or author to some agency external to the text.
\item [\textbf{att.global.source}] provides an attribute used by elements to point to an external source.\hfil\\[-10pt]\begin{sansreflist}
    \item[@{\itshape source}]
  specifies the source from which some aspect of this element is drawn.
\end{sansreflist}  
\item [\textbf{<cit>}] (cited quotation) contains a quotation from some other document, together with a bibliographic reference to its source. In a dictionary it may contain an example text with at least one occurrence of the word form, used in the sense being described, or a translation of the headword, or an example.
\item [\textbf{<mentioned>}] marks words or phrases mentioned, not used.
\item [\textbf{<soCalled>}] (so called) contains a word or phrase for which the author or narrator indicates a disclaiming of responsibility, for example by the use of scare quotes or italics.
\end{sansreflist}
 The elements \hyperref[TEI.mentioned]{<mentioned>} and \hyperref[TEI.soCalled]{<soCalled>} are members of the class \textsf{model.emphLike} while \hyperref[TEI.q]{<q>} stems from \textsf{model.hiLike}; the element \hyperref[TEI.said]{<said>} is a member of the class \textsf{model.attributable} in its own right, while \hyperref[TEI.cit]{<cit>} and \hyperref[TEI.quote]{<quote>} are members of \textsf{model.quoteLike}, a subclass of \textsf{model.attributable}. This class is a subclass of \textsf{model.inter}; hence all of these elements are permitted both within and between paragraph-level elements.\par
The most common and important use of quotation marks is, of course, to mark \textit{quotation}, by which we mean simply any part of the text which the author or narrator wishes to attribute to some agency other than the narrative voice. The \hyperref[TEI.q]{<q>} element may be used if no further distinction beyond this is judged necessary. If it is felt necessary to distinguish such passages further, for example to indicate whether they are regarded as speech, writing, or thought, either the {\itshape type} attribute or one of the more specialized elements discussed in this section may be used. For example, the element \hyperref[TEI.quote]{<quote>} may be used for written passages cited from other works, or the element \hyperref[TEI.said]{<said>} for words or phrases represented as being spoken or thought by people or characters within the current work. The \hyperref[TEI.soCalled]{<soCalled>} element is used for cases where the author or narrator distances him or herself from the words in question without however attributing them to any other voice in particular. The \hyperref[TEI.mentioned]{<mentioned>} element is appropriate for a case where a word or phrase is being discussed in the body of a text rather than forming part of the text directly.\par
As noted above, if the distinction among these various reasons why a passage is offset from surrounding text cannot be made reliably, or is not of interest, then any representation of speech, thought, or writing may simply be marked using the \hyperref[TEI.q]{<q>} element.\par
Quotation may be indicated in a printed source by changes in type face, by special punctuation marks (single or double or angled quotes, dashes, etc.) and by layout (indented paragraphs, etc.), or it may not be explicitly represented at all. If these characteristics are of interest, one or other of the global {\itshape rend} or {\itshape rendition} attributes discussed in section \textit{\hyperref[STGA]{1.3.1.1.\ Global Attributes}} may be used to record them.\par
Quotation marks themselves may, like other punctuation marks, be felt for some purposes to be worth retaining within a text, quite independently of their description by the {\itshape rend} attribute. This should generally be done using the appropriate Unicode character, or, if this is not possible, a numeric character reference (see \textit{\hyperref[SG-er]{v.7.1\ Character References}}). If the encoder decides both to retain the quotation marks and to represent their function by means of an explicit tag such as \hyperref[TEI.quote]{<quote>}, the quotation marks should be included within the element, rather than outside it, as in the first example below: \par\bgroup\index{said=<said>|exampleindex}\index{said=<said>|exampleindex}\index{said=<said>|exampleindex}\exampleFont \begin{shaded}\noindent\mbox{}Adolphe se tourna vers lui :\mbox{}\newline 
{<\textbf{said}>}— Alors, Albert, quoi de neuf?{</\textbf{said}>}\mbox{}\newline 
{<\textbf{said}>}— Pas grand-chose.{</\textbf{said}>}\mbox{}\newline 
{<\textbf{said}>}— Il fait beau,{</\textbf{said}>} dit Robert.\end{shaded}\egroup\par \noindent  Alternatively, since this use of the leading mdash is very common typographic practice, it may be considered unnecessary to retain it in the encoding. Its presence in the source might instead be signalled using one of the attributes {\itshape rend}, {\itshape style}, or {\itshape rendition}. This kind of rendering might be predefined using a \hyperref[TEI.rendition]{<rendition>} element, which can then be referenced using the {\itshape rendition} attribute as follows: \par\bgroup\index{said=<said>|exampleindex}\index{rendition=@rendition!<said>|exampleindex}\index{said=<said>|exampleindex}\index{rendition=@rendition!<said>|exampleindex}\index{said=<said>|exampleindex}\index{rendition=@rendition!<said>|exampleindex}\index{rendition=<rendition>|exampleindex}\index{scope=@scope!<rendition>|exampleindex}\index{quotation=<quotation>|exampleindex}\index{marks=@marks!<quotation>|exampleindex}\exampleFont \begin{shaded}\noindent\mbox{}Adolphe se tourna vers lui :\mbox{}\newline 
{<\textbf{said}\hspace*{1em}{rendition}="{\#dashBefore}">}Alors,\mbox{}\newline 
 Albert, quoi de neuf ?{</\textbf{said}>}\mbox{}\newline 
{<\textbf{said}\hspace*{1em}{rendition}="{\#dashBefore}">}Pas grand-chose.{</\textbf{said}>}\mbox{}\newline 
{<\textbf{said}\hspace*{1em}{rendition}="{\#dashBefore}">}Il fait beau,{</\textbf{said}>}\mbox{}\newline 
 dit Robert.\mbox{}\newline 
\mbox{}\newline 
\textit{<!-- ...  within the header  -->}\mbox{}\newline 
{<\textbf{rendition}\hspace*{1em}{xml:id}="{dashBefore}"\mbox{}\newline 
\hspace*{1em}{scope}="{before}">}content: '— '{</\textbf{rendition}>}\mbox{}\newline 
\textit{<!-- ... -->}\mbox{}\newline 
{<\textbf{quotation}\hspace*{1em}{marks}="{none}"/>}\end{shaded}\egroup\par \par
Whatever policy is adopted, the encoder should document the decision in some way, for example by using the \hyperref[TEI.quotation]{<quotation>} element provided in the TEI header (see \textit{\hyperref[HD53]{2.3.3.\ The Editorial Practices Declaration}}) to indicate that quotation marks have not been retained in the encoding; their presence in the source is implied by the {\itshape rendition} attribute values supplied.\par
Whether or not the quotation marks are suppressed, their presence and nature may be described using some appropriate set of conventions in the {\itshape rend} attribute. These conventions may be entirely idiosyncratic, and hence not necessarily useful for interchange, as in the following example: \par\bgroup\index{said=<said>|exampleindex}\index{rend=@rend!<said>|exampleindex}\index{said=<said>|exampleindex}\index{rend=@rend!<said>|exampleindex}\exampleFont \begin{shaded}\noindent\mbox{}{<\textbf{said}\hspace*{1em}{rend}="{pre(‘) post(’)}">}Who-e debel\mbox{}\newline 
 you?{</\textbf{said}>} — he at last said —\mbox{}\newline 
{<\textbf{said}\hspace*{1em}{rend}="{pre(‘) post(’)}">}you no speak-e,\mbox{}\newline 
 damme, I kill-e.{</\textbf{said}>} And so saying,\mbox{}\newline 
 the lighted tomahawk began flourishing\mbox{}\newline 
 about me in the dark.\end{shaded}\egroup\par \par
Such passages might more effectively be encoded without loss of rendering information by using the {\itshape rendition} attribute and its associated \hyperref[TEI.rendition]{<rendition>} element as described in section \textit{\hyperref[HD57-1]{2.3.4.1.\ Rendition}}. If the rendition of passages tagged as \hyperref[TEI.said]{<said>} is generally uniform throughout a text, then the \hyperref[TEI.rendition]{<rendition>} element may used to specify a default rendering, in which case the same section might simply be tagged: \par\bgroup\index{said=<said>|exampleindex}\index{said=<said>|exampleindex}\index{tagsDecl=<tagsDecl>|exampleindex}\index{partial=@partial!<tagsDecl>|exampleindex}\index{rendition=<rendition>|exampleindex}\index{scheme=@scheme!<rendition>|exampleindex}\index{selector=@selector!<rendition>|exampleindex}\index{scope=@scope!<rendition>|exampleindex}\index{rendition=<rendition>|exampleindex}\index{scheme=@scheme!<rendition>|exampleindex}\index{selector=@selector!<rendition>|exampleindex}\index{scope=@scope!<rendition>|exampleindex}\exampleFont \begin{shaded}\noindent\mbox{}{<\textbf{said}>}Who-e debel\mbox{}\newline 
 you?{</\textbf{said}>} — he at last said —\mbox{}\newline 
{<\textbf{said}>}you no speak-e,\mbox{}\newline 
 damme, I kill-e.{</\textbf{said}>} And so saying,\mbox{}\newline 
 the lighted tomahawk began flourishing\mbox{}\newline 
 about me in the dark.\mbox{}\newline 
\mbox{}\newline 
\textit{<!-- in the header: -->}\mbox{}\newline 
{<\textbf{tagsDecl}\hspace*{1em}{partial}="{true}">}\mbox{}\newline 
\hspace*{1em}{<\textbf{rendition}\hspace*{1em}{scheme}="{css}"\hspace*{1em}{selector}="{said}"\mbox{}\newline 
\hspace*{1em}\hspace*{1em}{scope}="{before}">}content:"‘";{</\textbf{rendition}>}\mbox{}\newline 
\hspace*{1em}{<\textbf{rendition}\hspace*{1em}{scheme}="{css}"\hspace*{1em}{selector}="{said}"\mbox{}\newline 
\hspace*{1em}\hspace*{1em}{scope}="{after}">}content:"’";{</\textbf{rendition}>}\mbox{}\newline 
{</\textbf{tagsDecl}>}\end{shaded}\egroup\par \par
As members of the \textsf{att.ascribed.directed} class, elements \hyperref[TEI.said]{<said>} and \hyperref[TEI.q]{<q>} share the following attributes: 
\begin{sansreflist}
  
\item [\textbf{att.ascribed}] provides attributes for elements representing speech or action that can be ascribed to a specific individual.\hfil\\[-10pt]\begin{sansreflist}
    \item[@{\itshape who}]
  indicates the person, or group of people, to whom the element content is ascribed.
\end{sansreflist}  
\item [\textbf{att.ascribed.directed}] provides attributes for elements representing speech or action that can be directed at a group or individual.\hfil\\[-10pt]\begin{sansreflist}
    \item[@{\itshape toWhom}]
  indicates the person, or group of people, to whom a speech act or action is directed.
\end{sansreflist}  
\end{sansreflist}
 These may be used to make explicit who is speaking and who is being spoken to: \par\bgroup\index{said=<said>|exampleindex}\index{who=@who!<said>|exampleindex}\index{toWhom=@toWhom!<said>|exampleindex}\index{said=<said>|exampleindex}\index{who=@who!<said>|exampleindex}\index{toWhom=@toWhom!<said>|exampleindex}\index{said=<said>|exampleindex}\index{who=@who!<said>|exampleindex}\index{standOff=<standOff>|exampleindex}\index{listPerson=<listPerson>|exampleindex}\index{type=@type!<listPerson>|exampleindex}\index{person=<person>|exampleindex}\index{persName=<persName>|exampleindex}\index{person=<person>|exampleindex}\index{persName=<persName>|exampleindex}\index{person=<person>|exampleindex}\index{persName=<persName>|exampleindex}\exampleFont \begin{shaded}\noindent\mbox{}Adolphe se tourna vers lui :\mbox{}\newline 
{<\textbf{said}\hspace*{1em}{who}="{\#Ado}"\hspace*{1em}{toWhom}="{\#Alb}">}— Alors, Albert, quoi de neuf?{</\textbf{said}>}\mbox{}\newline 
{<\textbf{said}\hspace*{1em}{who}="{\#Alb}"\hspace*{1em}{toWhom}="{\#Ado}">}— Pas grand-chose.{</\textbf{said}>}\mbox{}\newline 
{<\textbf{said}\hspace*{1em}{who}="{\#Rob}">}— Il fait beau,{</\textbf{said}>} dit Robert.\mbox{}\newline 
\mbox{}\newline 
\textit{<!-- ... elsewhere in the document -->}\mbox{}\newline 
{<\textbf{standOff}>}\mbox{}\newline 
\hspace*{1em}{<\textbf{listPerson}\hspace*{1em}{type}="{speakers}">}\mbox{}\newline 
\hspace*{1em}\hspace*{1em}{<\textbf{person}\hspace*{1em}{xml:id}="{Ado}">}\mbox{}\newline 
\hspace*{1em}\hspace*{1em}\hspace*{1em}{<\textbf{persName}>}Adolphe{</\textbf{persName}>}\mbox{}\newline 
\hspace*{1em}\hspace*{1em}{</\textbf{person}>}\mbox{}\newline 
\hspace*{1em}\hspace*{1em}{<\textbf{person}\hspace*{1em}{xml:id}="{Alb}">}\mbox{}\newline 
\hspace*{1em}\hspace*{1em}\hspace*{1em}{<\textbf{persName}>}Albert{</\textbf{persName}>}\mbox{}\newline 
\hspace*{1em}\hspace*{1em}{</\textbf{person}>}\mbox{}\newline 
\hspace*{1em}\hspace*{1em}{<\textbf{person}\hspace*{1em}{xml:id}="{Rob}">}\mbox{}\newline 
\hspace*{1em}\hspace*{1em}\hspace*{1em}{<\textbf{persName}>}Robert{</\textbf{persName}>}\mbox{}\newline 
\hspace*{1em}\hspace*{1em}{</\textbf{person}>}\mbox{}\newline 
\hspace*{1em}{</\textbf{listPerson}>}\mbox{}\newline 
{</\textbf{standOff}>}\end{shaded}\egroup\par \noindent  The {\itshape who} attribute may be supplied whether or not an indication of the speaker is given explicitly in the text. Likewise, the {\itshape toWhom} attribute may be supplied to indicate the encoder's interpretation of the intended recipient of the speech. They may take the form (as above) of an abbreviated or normalized form of the speaker's name, but the role of these attributes is to act as a pointer to a location elsewhere in the text, or another document, where data about each party may be supplied. While these attributes could point to any source of information about the parties involved, the most appropriate place for such information is within the \textit{participant description} component of the TEI header, as further discussed in \textit{\hyperref[CCAHPA]{15.2.2.\ The Participant Description}} but for simple cases like the above, a simple list of speakers located in the front or back matter of the text may suffice.\par
It may also be useful to distinguish representations of speech from representations of thought, in modern printed texts often indicated by a change of typeface. The {\itshape aloud} attribute is provided for this purpose, as in this example: \par\bgroup\index{said=<said>|exampleindex}\index{aloud=@aloud!<said>|exampleindex}\index{said=<said>|exampleindex}\index{aloud=@aloud!<said>|exampleindex}\index{said=<said>|exampleindex}\index{aloud=@aloud!<said>|exampleindex}\exampleFont \begin{shaded}\noindent\mbox{}{<\textbf{said}\hspace*{1em}{aloud}="{true}">}Oh yes,{</\textbf{said}>} said Henry, \mbox{}\newline 
{<\textbf{said}\hspace*{1em}{aloud}="{false}">}I mean\mbox{}\newline 
 Gordon Macrae, for example…{</\textbf{said}>}\mbox{}\newline 
{<\textbf{said}\hspace*{1em}{aloud}="{false}">}Jungian\mbox{}\newline 
 Analyst with Winebox! That's what you called him, you callous bastard,\mbox{}\newline 
 didn't you? Eh? Eh?{</\textbf{said}>}\end{shaded}\egroup\par \par
Quoted matter may be embedded within quoted matter, as when one speaker reports the speech of another: \par\bgroup\index{said=<said>|exampleindex}\index{who=@who!<said>|exampleindex}\index{said=<said>|exampleindex}\index{who=@who!<said>|exampleindex}\index{list=<list>|exampleindex}\index{type=@type!<list>|exampleindex}\index{item=<item>|exampleindex}\index{item=<item>|exampleindex}\exampleFont \begin{shaded}\noindent\mbox{}{<\textbf{said}\hspace*{1em}{who}="{\#Wilson}">}Spaulding, he came down into the office just this day\mbox{}\newline 
 eight weeks with this very paper in his hand, and he says:—\mbox{}\newline 
{<\textbf{said}\hspace*{1em}{who}="{\#WilsonSpaulding}">}I wish to the Lord, Mr. Wilson, that I was a\mbox{}\newline 
\hspace*{1em}\hspace*{1em} red-headed man.{</\textbf{said}>}\mbox{}\newline 
{</\textbf{said}>}\mbox{}\newline 
\textit{<!-- ... -->}\mbox{}\newline 
{<\textbf{list}\hspace*{1em}{type}="{speakers}">}\mbox{}\newline 
\hspace*{1em}{<\textbf{item}\hspace*{1em}{xml:id}="{Wilson}">}Wilson{</\textbf{item}>}\mbox{}\newline 
\hspace*{1em}{<\textbf{item}\hspace*{1em}{xml:id}="{WilsonSpaulding}">}Spaulding reported by Wilson{</\textbf{item}>}\mbox{}\newline 
\textit{<!-- ...-->}\mbox{}\newline 
{</\textbf{list}>}\end{shaded}\egroup\par \par
Direct speech nested in this way is treated in the same way as elsewhere: a change of rendition may occur, but the same element should be used. An encoder may however choose to distinguish between direct speech which contains quotations from extra-textual matter and direct speech itself, as in the following example: \par\bgroup\index{p=<p>|exampleindex}\index{said=<said>|exampleindex}\index{said=<said>|exampleindex}\index{quote=<quote>|exampleindex}\index{l=<l>|exampleindex}\index{l=<l>|exampleindex}\index{l=<l>|exampleindex}\exampleFont \begin{shaded}\noindent\mbox{}{<\textbf{p}>}\mbox{}\newline 
\hspace*{1em}{<\textbf{said}>}The Lord! The Lord! It is Sakya Muni himself,{</\textbf{said}>} the lama half\mbox{}\newline 
 sobbed; and under his breath began the wonderful Buddhist\mbox{}\newline 
 invocation:-{<\textbf{said}>}\mbox{}\newline 
\hspace*{1em}\hspace*{1em}{<\textbf{quote}>}\mbox{}\newline 
\hspace*{1em}\hspace*{1em}\hspace*{1em}{<\textbf{l}>}To Him the Way — the Law — Apart —{</\textbf{l}>}\mbox{}\newline 
\hspace*{1em}\hspace*{1em}\hspace*{1em}{<\textbf{l}>}Whom Maya held beneath her heart{</\textbf{l}>}\mbox{}\newline 
\hspace*{1em}\hspace*{1em}\hspace*{1em}{<\textbf{l}>}Ananda's Lord — the Bodhisat{</\textbf{l}>}\mbox{}\newline 
\hspace*{1em}\hspace*{1em}{</\textbf{quote}>}\mbox{}\newline 
\hspace*{1em}\hspace*{1em} And He is here! The Most Excellent Law is here also. My\mbox{}\newline 
\hspace*{1em}\hspace*{1em} pilgrimage is well begun. And what work! What work!{</\textbf{said}>}\mbox{}\newline 
{</\textbf{p}>}\end{shaded}\egroup\par \par
Quotations from other works are often accompanied by a reference to their source. The \hyperref[TEI.cit]{<cit>} element may be used to group together the quotation and its associated bibliographic reference, which should be encoded using the elements for bibliographic references discussed in section \textit{\hyperref[COBI]{3.12.\ Bibliographic Citations and References}}, as in the following example. \par\bgroup\index{div=<div>|exampleindex}\index{type=@type!<div>|exampleindex}\index{head=<head>|exampleindex}\index{epigraph=<epigraph>|exampleindex}\index{cit=<cit>|exampleindex}\index{quote=<quote>|exampleindex}\index{l=<l>|exampleindex}\index{l=<l>|exampleindex}\index{bibl=<bibl>|exampleindex}\index{title=<title>|exampleindex}\index{author=<author>|exampleindex}\index{p=<p>|exampleindex}\exampleFont \begin{shaded}\noindent\mbox{}{<\textbf{div}\hspace*{1em}{xml:id}="{mm01}"\hspace*{1em}{type}="{chapter}">}\mbox{}\newline 
\hspace*{1em}{<\textbf{head}>}Chapter 1{</\textbf{head}>}\mbox{}\newline 
\hspace*{1em}{<\textbf{epigraph}>}\mbox{}\newline 
\hspace*{1em}\hspace*{1em}{<\textbf{cit}>}\mbox{}\newline 
\hspace*{1em}\hspace*{1em}\hspace*{1em}{<\textbf{quote}>}\mbox{}\newline 
\hspace*{1em}\hspace*{1em}\hspace*{1em}\hspace*{1em}{<\textbf{l}>}Since I can do no good because a woman{</\textbf{l}>}\mbox{}\newline 
\hspace*{1em}\hspace*{1em}\hspace*{1em}\hspace*{1em}{<\textbf{l}>}Reach constantly at something that is near it.{</\textbf{l}>}\mbox{}\newline 
\hspace*{1em}\hspace*{1em}\hspace*{1em}{</\textbf{quote}>}\mbox{}\newline 
\hspace*{1em}\hspace*{1em}\hspace*{1em}{<\textbf{bibl}>}\mbox{}\newline 
\hspace*{1em}\hspace*{1em}\hspace*{1em}\hspace*{1em}{<\textbf{title}>}The Maid's Tragedy{</\textbf{title}>}\mbox{}\newline 
\hspace*{1em}\hspace*{1em}\hspace*{1em}\hspace*{1em}{<\textbf{author}>}Beaumont and Fletcher{</\textbf{author}>}\mbox{}\newline 
\hspace*{1em}\hspace*{1em}\hspace*{1em}{</\textbf{bibl}>}\mbox{}\newline 
\hspace*{1em}\hspace*{1em}{</\textbf{cit}>}\mbox{}\newline 
\hspace*{1em}{</\textbf{epigraph}>}\mbox{}\newline 
\hspace*{1em}{<\textbf{p}>}Miss Brooke had that kind of beauty which seems to be thrown into\mbox{}\newline 
\hspace*{1em}\hspace*{1em} relief by poor dress...{</\textbf{p}>}\mbox{}\newline 
{</\textbf{div}>}\end{shaded}\egroup\par \noindent  Like other bibliographic references, the citation associated with a quotation may be represented simply by a cross-reference, as in this example: \par\bgroup\index{cit=<cit>|exampleindex}\index{quote=<quote>|exampleindex}\index{ref=<ref>|exampleindex}\exampleFont \begin{shaded}\noindent\mbox{}Lexicography has shown little sign of being affected by the\mbox{}\newline 
 work of followers of J.R. Firth, probably best summarized\mbox{}\newline 
 in his slogan, {<\textbf{cit}>}\mbox{}\newline 
\hspace*{1em}{<\textbf{quote}>}You shall know a word by the company it keeps.{</\textbf{quote}>}\mbox{}\newline 
\hspace*{1em}{<\textbf{ref}>}(Firth, 1957){</\textbf{ref}>}\mbox{}\newline 
{</\textbf{cit}>}\end{shaded}\egroup\par \noindent  It is also common for quotations to be separated from their bibliographic reference by intervening text, which makes the use of \hyperref[TEI.cit]{<cit>} impractical. In such circumstances, the quotation can be linked to a bibliographical reference using {\itshape source}: \par\bgroup\index{bibl=<bibl>|exampleindex}\index{quote=<quote>|exampleindex}\index{source=@source!<quote>|exampleindex}\index{title=<title>|exampleindex}\exampleFont \begin{shaded}\noindent\mbox{}{<\textbf{bibl}\hspace*{1em}{xml:id}="{tlk\textunderscore 36}">}Tolkien (1936){</\textbf{bibl}>} tells us that\mbox{}\newline 
{<\textbf{quote}\hspace*{1em}{source}="{\#tlk\textunderscore 36}">}\mbox{}\newline 
\hspace*{1em}{<\textbf{title}>}Beowulf{</\textbf{title}>} is in fact so interesting as\mbox{}\newline 
 poetry, in places poetry so powerful, that this quite\mbox{}\newline 
 overshadows the historical content\mbox{}\newline 
{</\textbf{quote}>}.\end{shaded}\egroup\par \noindent  {\itshape source} could also be used to point to a complete external bibliographic reference in a \hyperref[TEI.listBibl]{<listBibl>} elsewhere in the document, or external to it.\par
Unlike most of the other elements discussed in this chapter, direct speech and quotations may frequently contain other high-level elements such as paragraphs or verse lines, as well as being themselves contained by such elements. Three possible solutions exist for this well-known structural problem: \begin{itemize}
\item the quotation is broken into segments, each of which is entirely contained within a paragraph
\item the quotation is marked up using stand-off markup
\item the quotation boundaries are represented by empty segment boundary delimiter elements
\end{itemize}  For further discussion and several examples, see chapter \textit{\hyperref[NH]{20.\ Non-hierarchical Structures}}.\par
Finally, in this section, the element \hyperref[TEI.soCalled]{<soCalled>} is provided for all cases in which quotation marks are used to distance the quoted text from the narrator or speaker. Common examples include the ‘scare’ quotes often found in newspaper headlines and advertising copy, where the effect is to cast doubts on the veracity of an assertion: \par\bgroup\index{head=<head>|exampleindex}\index{soCalled=<soCalled>|exampleindex}\exampleFont \begin{shaded}\noindent\mbox{}{<\textbf{head}>}PM dodges {<\textbf{soCalled}>}election threat{</\textbf{soCalled}>} in interview{</\textbf{head}>}\end{shaded}\egroup\par \par
The same element should be used to mark a variety of special ironic usages. Some further examples follow: \par\bgroup\index{soCalled=<soCalled>|exampleindex}\exampleFont \begin{shaded}\noindent\mbox{}He hated {<\textbf{soCalled}>}good{</\textbf{soCalled}>} books.\end{shaded}\egroup\par \noindent  \par\bgroup\index{soCalled=<soCalled>|exampleindex}\exampleFont \begin{shaded}\noindent\mbox{}{<\textbf{soCalled}>}Croissants{</\textbf{soCalled}>} indeed! toast not good enough for you?\end{shaded}\egroup\par \noindent   \par\bgroup\index{soCalled=<soCalled>|exampleindex}\index{mentioned=<mentioned>|exampleindex}\exampleFont \begin{shaded}\noindent\mbox{}Although Chomsky's decision that all NL\mbox{}\newline 
 sentences are finite objects was never justified by arguments from\mbox{}\newline 
 the attested properties of NLs, it did have a certain\mbox{}\newline 
{<\textbf{soCalled}>}social{</\textbf{soCalled}>} justification. It was commonly assumed in\mbox{}\newline 
 works on logic until fairly recently that the notion\mbox{}\newline 
{<\textbf{mentioned}>}language{</\textbf{mentioned}>} is necessarily restricted to finite\mbox{}\newline 
 strings.\end{shaded}\egroup\par 
\subsection[{Terms and Glosses, Ruby Annotations, and Equivalents and Descriptions}]{Terms and Glosses, Ruby Annotations, and Equivalents and Descriptions}\label{COHTGEDR}\par
This section describes a set of textual elements which are used to provide a gloss, alternate identification, or description of something.
\subsubsection[{Terms and Glosses}]{Terms and Glosses}\label{COHTG}\par
Technical terms are often italicized or emboldened upon first mention in printed texts; an explanation or gloss is sometimes given in quotation marks. Linguistic analyses conventionally cite words in languages under discussion in italics, providing a gloss immediately following marked with single quotation marks. Other texts in which individual words or phrases are \textit{mentioned} (for example, as examples) rather than \textit{used} may mark them either with italics or with quotation marks, and will gloss them less regularly.
\begin{sansreflist}
  
\item [\textbf{<term>}] (term) contains a single-word, multi-word, or symbolic designation which is regarded as a technical term.
\item [\textbf{<gloss>}] (gloss) identifies a phrase or word used to provide a gloss or definition for some other word or phrase.
\end{sansreflist}
 These elements are also members of the class \textsf{model.emphLike}.\par
A \hyperref[TEI.term]{<term>} may appear with or without a gloss, as may a \hyperref[TEI.mentioned]{<mentioned>} element. Where the \hyperref[TEI.gloss]{<gloss>} is present, it may be linked to the term it is glossing by means of its {\itshape target} attribute. To establish such a link, the encoder should give an {\itshape xml:id} value to the \hyperref[TEI.term]{<term>} or \hyperref[TEI.mentioned]{<mentioned>} element and provide that value, prefixed by a sharp sign, as the value of the {\itshape target} attribute on the \hyperref[TEI.gloss]{<gloss>} element. For example: \par\bgroup\index{term=<term>|exampleindex}\index{rend=@rend!<term>|exampleindex}\index{gloss=<gloss>|exampleindex}\index{target=@target!<gloss>|exampleindex}\exampleFont \begin{shaded}\noindent\mbox{}We may define {<\textbf{term}\hspace*{1em}{xml:id}="{TDPv}"\hspace*{1em}{rend}="{sc}">}discoursal point of view{</\textbf{term}>}\mbox{}\newline 
 as \mbox{}\newline 
{<\textbf{gloss}\hspace*{1em}{target}="{\#TDPv}">}the relationship, expressed through discourse\mbox{}\newline 
 structure, between the implied author or some other addresser,\mbox{}\newline 
 and the fiction.{</\textbf{gloss}>}\end{shaded}\egroup\par \noindent  \par\bgroup\index{gloss=<gloss>|exampleindex}\index{rend=@rend!<gloss>|exampleindex}\index{target=@target!<gloss>|exampleindex}\index{term=<term>|exampleindex}\exampleFont \begin{shaded}\noindent\mbox{}{<\textbf{gloss}\hspace*{1em}{rend}="{unmarked}"\hspace*{1em}{target}="{\#PRSR}">}A computational device that infers\mbox{}\newline 
 structure from grammatical strings of words{</\textbf{gloss}>} is known as a\mbox{}\newline 
{<\textbf{term}\hspace*{1em}{xml:id}="{PRSR}">}parser{</\textbf{term}>}, and much of the history of NLP over the\mbox{}\newline 
 last 20 years has been occupied with the design of parsers.\end{shaded}\egroup\par \par
In the absence of a \hyperref[TEI.gloss]{<gloss>} in the source text, a \hyperref[TEI.term]{<term>} can also be associated with a gloss or definition by means of its {\itshape ref} attribute, as in this imaginary example: \par\bgroup\index{term=<term>|exampleindex}\index{ref=@ref!<term>|exampleindex}\index{rend=@rend!<term>|exampleindex}\exampleFont \begin{shaded}\noindent\mbox{}We discuss Leech's concept of {<\textbf{term}\hspace*{1em}{ref}="{\#TDPV2}"\hspace*{1em}{rend}="{sc}">}discoursal point of view{</\textbf{term}>} below. \end{shaded}\egroup\par \noindent  Here, the value \textsf{\#TDPV2} references some other XML element in the same document. This might be, for example, an entry in a list of technical vocabulary given in the document header, or in an appendix. This attribute could also reference an entry in some centrally maintained terminological database, perhaps using one of the pointing mechanisms discussed in section \textit{\hyperref[SAXP]{16.2.\ Pointing Mechanisms}}.\par
Note that the element \hyperref[TEI.term]{<term>} is intended for use with words or phrases identified as terminological in nature; where words or phrases are simply being cited, discussed, or glossed in a text, it will often be more appropriate to use the \hyperref[TEI.mentioned]{<mentioned>} element, as in the following example: \par\bgroup\index{mentioned=<mentioned>|exampleindex}\index{gloss=<gloss>|exampleindex}\index{target=@target!<gloss>|exampleindex}\index{mentioned=<mentioned>|exampleindex}\index{gloss=<gloss>|exampleindex}\index{target=@target!<gloss>|exampleindex}\exampleFont \begin{shaded}\noindent\mbox{}There is thus a striking accentual difference between a verbal\mbox{}\newline 
 form like {<\textbf{mentioned}\hspace*{1em}{xml:id}="{cw234}"\hspace*{1em}{xml:lang}="{grc}">}eluthemen{</\textbf{mentioned}>}\mbox{}\newline 
{<\textbf{gloss}\hspace*{1em}{target}="{\#cw234}">}we were released,{</\textbf{gloss}>} accented on the\mbox{}\newline 
 second syllable of the word, and its participial derivative\mbox{}\newline 
\mbox{}\newline 
{<\textbf{mentioned}\hspace*{1em}{xml:id}="{cw235}"\hspace*{1em}{xml:lang}="{grc}">}lutheis{</\textbf{mentioned}>}\mbox{}\newline 
{<\textbf{gloss}\hspace*{1em}{target}="{\#cw235}">}released,{</\textbf{gloss}>} accented on the last.\end{shaded}\egroup\par \par
For technical terminology in particular, and generally in terminological studies, it may be useful to associate an instance of a term within a text with a canonical definition for it, which is stored either elsewhere in the same text (for example in a glossary of terms) or externally, for example in a database, authority file, or published standard. The attributes {\itshape key} and {\itshape ref} discussed in section \textit{\hyperref[CONARS]{3.6.1.\ Referring Strings}} below are available on the \hyperref[TEI.term]{<term>} element for this purpose.
\paragraph[{Some Further Examples}]{Some Further Examples}\label{COHTGEG}\par
As a simple example of the elements discussed here and in \textit{\hyperref[COHQ]{3.3.\ Highlighting and Quotation}}, consider the following sentence: 
\begin{quote}On the one hand the \textit{Nibelungenlied} is associated with the new rise of romance of twelfth-century France, the \textit{romans d'antiquité}, the romances of Chrétien de Troyes, and the German adaptations of these works by Heinrich van Veldeke, Hartmann von Aue, and Wolfram von Eschenbach.\end{quote}
 A first approximation to the encoding of this sentence might be simply to record the fact that the phrases printed above in italics are highlighted, as follows: \par\bgroup\index{hi=<hi>|exampleindex}\index{rend=@rend!<hi>|exampleindex}\index{hi=<hi>|exampleindex}\index{rend=@rend!<hi>|exampleindex}\exampleFont \begin{shaded}\noindent\mbox{}On the one hand the {<\textbf{hi}\hspace*{1em}{rend}="{italic}">}Nibelungenlied{</\textbf{hi}>} is\mbox{}\newline 
 associated with the new rise of romance of twelfth-century France,\mbox{}\newline 
 the {<\textbf{hi}\hspace*{1em}{xml:lang}="{fr}"\hspace*{1em}{rend}="{italic}">}romans d'antiquité{</\textbf{hi}>},\mbox{}\newline 
 the romances of Chrétien de Troyes, ...\end{shaded}\egroup\par \noindent  This encoding would, however, lose the important distinction between an italicized title and an italicized foreign phrase. Many other phrases might also be italicized in the text, and a retrieval program seeking to identify foreign terms (for example) would not be able to produce reliable results by simply looking for italicized words. Where economic and intellectual constraints permit, therefore, it would be preferable to encode both the function of the highlighted phrases and their appearance, as follows: \par\bgroup\index{title=<title>|exampleindex}\index{rend=@rend!<title>|exampleindex}\index{foreign=<foreign>|exampleindex}\index{rend=@rend!<foreign>|exampleindex}\exampleFont \begin{shaded}\noindent\mbox{}On the one hand the {<\textbf{title}\hspace*{1em}{rend}="{italic}">}Nibelungenlied{</\textbf{title}>}\mbox{}\newline 
 is associated with the new rise of romance of twelfth-century France,\mbox{}\newline 
 the {<\textbf{foreign}\hspace*{1em}{rend}="{italic}">}romans d'antiquité{</\textbf{foreign}>}, the\mbox{}\newline 
 romances of Chrétien de Troyes, ...\end{shaded}\egroup\par \par
In this example, the decision as to which textual features are distinguished by the highlighting is relatively uncontroversial. As a less straightforward example, consider the use of italic font in the following passage: 
\begin{quote}A pretty common case, I believe; in all {\itshape vehement} debatings. She says I am {\itshape too witty}; Anglicé, {\itshape too pert}; I, that she is {\itshape too wise}; that is to say, being likewise put into English, {\itshape not so young as she has been}: in short, she is grown so much into a {\itshape mother}, that she had forgotten she ever was a {\itshape daughter}. ...\end{quote}
\par
Clearly, the word \textit{vehement} is not italicized for the same reason as the phrase \textit{not so young as she has been}; the former is emphasized, while the latter is proverbial. It also provides an ironic gloss for the words \textit{too wise}, in the same way as \textit{too pert} glosses \textit{too witty}. The glossed phrases are not, however, technical terms or cited words, but quoted phrases, as if the writer were putting words into her own and her mother's mouths. Finally, the words \textit{mother} and \textit{daughter} are apparently italicized simply to oppose them in the sentence; certainly they do not fit into any of the categories so far proposed as reasons for italicizing. Note also that the word \textit{Anglicé} is not italicized although it is not generally considered an English word.\par
The following sample encoding for the above passage attempts to take into account all the above points: \par\bgroup\index{emph=<emph>|exampleindex}\index{q=<q>|exampleindex}\index{rend=@rend!<q>|exampleindex}\index{foreign=<foreign>|exampleindex}\index{rend=@rend!<foreign>|exampleindex}\index{gloss=<gloss>|exampleindex}\index{rend=@rend!<gloss>|exampleindex}\index{q=<q>|exampleindex}\index{rend=@rend!<q>|exampleindex}\index{gloss=<gloss>|exampleindex}\index{rend=@rend!<gloss>|exampleindex}\index{hi=<hi>|exampleindex}\index{rend=@rend!<hi>|exampleindex}\index{hi=<hi>|exampleindex}\index{rend=@rend!<hi>|exampleindex}\exampleFont \begin{shaded}\noindent\mbox{}A pretty common case, I believe; in all {<\textbf{emph}>}vehement{</\textbf{emph}>}\mbox{}\newline 
 debatings. She says I am {<\textbf{q}\hspace*{1em}{rend}="{italic}">}too witty{</\textbf{q}>};\mbox{}\newline 
{<\textbf{foreign}\hspace*{1em}{xml:lang}="{la}"\hspace*{1em}{rend}="{roman}">}Anglicé{</\textbf{foreign}>},\mbox{}\newline 
{<\textbf{gloss}\hspace*{1em}{rend}="{italic}">}too pert{</\textbf{gloss}>}; I, that she is\mbox{}\newline 
{<\textbf{q}\hspace*{1em}{rend}="{italic}">} too wise{</\textbf{q}>}; that is to say, being likewise\mbox{}\newline 
 put into English, {<\textbf{gloss}\hspace*{1em}{rend}="{italic}">}not so young as she has\mbox{}\newline 
 been{</\textbf{gloss}>}: in short, she is grown so much into a\mbox{}\newline 
{<\textbf{hi}\hspace*{1em}{rend}="{italic}">}mother{</\textbf{hi}>}, that she had forgotten she ever\mbox{}\newline 
 was a {<\textbf{hi}\hspace*{1em}{rend}="{italic}">}daughter{</\textbf{hi}>}.\end{shaded}\egroup\par \noindent        
\subsubsection[{Ruby Annotations}]{Ruby Annotations}\label{COHTGRB}\par
The word \textit{ruby} (or \textit{rubi}) refers to a particular method of glossing runs of text which is common in East Asian scripts. In horizontally-oriented text, ruby annotations typically appear above the text being glossed, while in vertical runs of text they may appear to the left or right, or both, also oriented vertically. An English example of a ruby annotation might look like this:\begin{figure}[htbp]
\noindent\noindent\includegraphics[width=10empt,]{Images/rbTEI.png}\end{figure}
\par
In Japanese, furigana ({\textJapanese 振り仮名}) ruby annotations are often used to provide pronunciation guidance for readers; characters from the largely phonetic hiragana or katakana syllabaries accompany Chinese characters, like this:\begin{figure}[htbp]
\noindent\noindent\includegraphics[width=35empt,]{Images/rbNhkEasy.png}
\caption{The first line of a news story from NHK News Web Easy intended for Japanese learners, in which every word composed of Chinese characters has a ruby gloss.}\end{figure}
\par
Pinyin ruby annotations are also used in Chinese to provide pronunciation guidance, and Zhuyin ({\textChinese 注音}) phonetic symbols (commonly known as \textit{bopomofo}) are used in Taiwan for the same purpose.\par
The TEI schema provides many different ways of encoding glosses and annotations, from the simple and flexible \hyperref[TEI.note]{<note>} element to a native implementation of the Web Annotation Data Model (\textit{\hyperref[SASOann]{16.11.\ Annotations}}). However, ruby is a particular, distinct, and widely-used form of annotation that appears in script, print, calligraphy, and web pages, and the TEI therefore provides specific elements for it: 
\begin{sansreflist}
  
\item [\textbf{<ruby>}] (ruby container) contains a passage of base text along with its associated ruby gloss(es).
\item [\textbf{<rb>}] (ruby base) contains the base text annotated by a ruby gloss.
\item [\textbf{<rt>}] (ruby text) contains a ruby text, an annotation closely associated with a passage of the main text.
\end{sansreflist}
 The \hyperref[TEI.rt]{<rt>} element is a member of \textsf{att.placement}, and thus the {\itshape place} attribute may be used to indicate where the ruby gloss is with respect to the base text: 
\begin{sansreflist}
  
\item [\textbf{att.placement}] provides attributes for describing where on the source page or object a textual element appears.\hfil\\[-10pt]\begin{sansreflist}
    \item[@{\itshape place}]
  specifies where this item is placed.
\end{sansreflist}  
\end{sansreflist}
 The most relevant suggested values of {\itshape place} for ruby text are above, below, left, and right.\par
In its simplest representation, a glossed form consists of an \hyperref[TEI.rb]{<rb>} (ruby base) element containing the base form, an \hyperref[TEI.rt]{<rt>} (ruby text) element containing the gloss, and a \hyperref[TEI.ruby]{<ruby>} element which wraps them together: \par\bgroup\index{p=<p>|exampleindex}\index{ruby=<ruby>|exampleindex}\index{rb=<rb>|exampleindex}\index{rt=<rt>|exampleindex}\index{place=@place!<rt>|exampleindex}\exampleFont \begin{shaded}\noindent\mbox{}{<\textbf{p}\hspace*{1em}{xml:lang}="{ja}">}\mbox{}\newline 
\textit{<!--{\textJapanese ...}-->}\mbox{}\newline 
\hspace*{1em}{<\textbf{ruby}>}\mbox{}\newline 
\hspace*{1em}\hspace*{1em}{<\textbf{rb}>}{\textJapanese 大学}{</\textbf{rb}>}\mbox{}\newline 
\hspace*{1em}\hspace*{1em}{<\textbf{rt}\hspace*{1em}{place}="{above}">}{\textJapanese だいがく}{</\textbf{rt}>}\mbox{}\newline 
\hspace*{1em}{</\textbf{ruby}>}\mbox{}\newline 
\textit{<!--{\textJapanese ...}-->}\mbox{}\newline 
{</\textbf{p}>}\end{shaded}\egroup\par \noindent  In the example above, the word {\textJapanese 大学} (\textit{daigaku} = university) is provided with a phonation gloss in hiragana. The full gloss is applied to the complete word. However, it might instead be broken down by character: \par\bgroup\index{p=<p>|exampleindex}\index{ruby=<ruby>|exampleindex}\index{rb=<rb>|exampleindex}\index{rt=<rt>|exampleindex}\index{place=@place!<rt>|exampleindex}\index{ruby=<ruby>|exampleindex}\index{rb=<rb>|exampleindex}\index{rt=<rt>|exampleindex}\index{place=@place!<rt>|exampleindex}\exampleFont \begin{shaded}\noindent\mbox{}{<\textbf{p}\hspace*{1em}{xml:lang}="{ja}">}\mbox{}\newline 
\textit{<!--{\textJapanese ...}-->}\mbox{}\newline 
\hspace*{1em}{<\textbf{ruby}>}\mbox{}\newline 
\hspace*{1em}\hspace*{1em}{<\textbf{rb}>}{\textJapanese 大}{</\textbf{rb}>}\mbox{}\newline 
\hspace*{1em}\hspace*{1em}{<\textbf{rt}\hspace*{1em}{place}="{above}">}{\textJapanese だい}{</\textbf{rt}>}\mbox{}\newline 
\hspace*{1em}{</\textbf{ruby}>}\mbox{}\newline 
\hspace*{1em}{<\textbf{ruby}>}\mbox{}\newline 
\hspace*{1em}\hspace*{1em}{<\textbf{rb}>}{\textJapanese 学}{</\textbf{rb}>}\mbox{}\newline 
\hspace*{1em}\hspace*{1em}{<\textbf{rt}\hspace*{1em}{place}="{above}">}{\textJapanese がく}{</\textbf{rt}>}\mbox{}\newline 
\hspace*{1em}{</\textbf{ruby}>}\mbox{}\newline 
\textit{<!--{\textJapanese ...}-->}\mbox{}\newline 
{</\textbf{p}>}\end{shaded}\egroup\par \noindent  Here is a similar example from Taiwan using bopomofo (pinyin \textit{píngzi} = bottle)\footnote{Taken from \xref{https://en.wikipedia.org/wiki/Bopomofo\#Example}{Wikipedia}.}: \par\bgroup\index{p=<p>|exampleindex}\index{ruby=<ruby>|exampleindex}\index{rb=<rb>|exampleindex}\index{rt=<rt>|exampleindex}\index{place=@place!<rt>|exampleindex}\index{ruby=<ruby>|exampleindex}\index{rb=<rb>|exampleindex}\index{rt=<rt>|exampleindex}\index{place=@place!<rt>|exampleindex}\exampleFont \begin{shaded}\noindent\mbox{}{<\textbf{p}\hspace*{1em}{xml:lang}="{zh-TW}">}\mbox{}\newline 
\textit{<!--{\textChinese ...}-->}\mbox{}\newline 
\hspace*{1em}{<\textbf{ruby}>}\mbox{}\newline 
\hspace*{1em}\hspace*{1em}{<\textbf{rb}>}{\textChinese 瓶}{</\textbf{rb}>}\mbox{}\newline 
\hspace*{1em}\hspace*{1em}{<\textbf{rt}\hspace*{1em}{place}="{right}">}{\textChinese ㄆㄧㄥˊ}{</\textbf{rt}>}\mbox{}\newline 
\hspace*{1em}{</\textbf{ruby}>}\mbox{}\newline 
\hspace*{1em}{<\textbf{ruby}>}\mbox{}\newline 
\hspace*{1em}\hspace*{1em}{<\textbf{rb}>}{\textChinese 子}{</\textbf{rb}>}\mbox{}\newline 
\hspace*{1em}\hspace*{1em}{<\textbf{rt}\hspace*{1em}{place}="{right}">}{\textChinese ˙ㄗ}{</\textbf{rt}>}\mbox{}\newline 
\hspace*{1em}{</\textbf{ruby}>}\mbox{}\newline 
\textit{<!--{\textChinese ...}-->}\mbox{}\newline 
{</\textbf{p}>}\end{shaded}\egroup\par \noindent  Where {\itshape place} is not provided, the default assumption is that the ruby gloss is \textit{above} where the text is horizontal, and to the \textit{right} of the text where it is vertical. See \textit{\hyperref[WDWM]{5.6.\ Writing Modes}} for a detailed guide to writing modes and text directionality.\par
The same ruby base may be accompanied by more than one gloss. Here, the Japanese word {\textJapanese 打球場} (dakyūba, or  \textit{billiard hall}) is glossed with two different pronunciations: ‘biriādo’ (its English equivalent) and ‘dakyū’, a phonation guide for the first two characters. \begin{figure}[htbp]
\noindent\includegraphics[width=149pt,height=276pt,]{Images/billiardhall.jpg}
\caption{\textit{Billiard hall} with two ruby glosses. \xref{http://school.nijl.ac.jp/kindai/NIJL/NIJL-01116.html\#37}{{\textJapanese 国文学研究資料館所蔵：：英国／龍動新繁昌記}}.}\end{figure}
 This example is intriguing in that the right-side ruby glosses apply to the first and second characters respectively, but the left-side gloss applies to the whole word as a unit. We use this instance to exemplify multiple approaches to encoding the same phenomena, which may be appropriate for different projects or editorial preferences. First, using the same segmentation approach as demonstrated for {\textJapanese 大学} above, but with nesting: \par\bgroup\index{p=<p>|exampleindex}\index{style=@style!<p>|exampleindex}\index{ruby=<ruby>|exampleindex}\index{rb=<rb>|exampleindex}\index{ruby=<ruby>|exampleindex}\index{rb=<rb>|exampleindex}\index{rt=<rt>|exampleindex}\index{place=@place!<rt>|exampleindex}\index{ruby=<ruby>|exampleindex}\index{rb=<rb>|exampleindex}\index{rt=<rt>|exampleindex}\index{place=@place!<rt>|exampleindex}\index{rt=<rt>|exampleindex}\index{place=@place!<rt>|exampleindex}\exampleFont \begin{shaded}\noindent\mbox{}{<\textbf{p}\hspace*{1em}{style}="{writing-mode: vertical-rl}"\mbox{}\newline 
\hspace*{1em}{xml:lang}="{ja}">}\mbox{}\newline 
\textit{<!--{\textJapanese ...}-->}\mbox{}\newline 
\hspace*{1em}{<\textbf{ruby}>}\mbox{}\newline 
\hspace*{1em}\hspace*{1em}{<\textbf{rb}>}\mbox{}\newline 
\hspace*{1em}\hspace*{1em}\hspace*{1em}{<\textbf{ruby}>}\mbox{}\newline 
\hspace*{1em}\hspace*{1em}\hspace*{1em}\hspace*{1em}{<\textbf{rb}>}{\textJapanese 打}{</\textbf{rb}>}\mbox{}\newline 
\hspace*{1em}\hspace*{1em}\hspace*{1em}\hspace*{1em}{<\textbf{rt}\hspace*{1em}{place}="{right}">}{\textJapanese ダ}{</\textbf{rt}>}\mbox{}\newline 
\hspace*{1em}\hspace*{1em}\hspace*{1em}{</\textbf{ruby}>}\mbox{}\newline 
\hspace*{1em}\hspace*{1em}\hspace*{1em}{<\textbf{ruby}>}\mbox{}\newline 
\hspace*{1em}\hspace*{1em}\hspace*{1em}\hspace*{1em}{<\textbf{rb}>}{\textJapanese 球}{</\textbf{rb}>}\mbox{}\newline 
\hspace*{1em}\hspace*{1em}\hspace*{1em}\hspace*{1em}{<\textbf{rt}\hspace*{1em}{place}="{right}">}{\textJapanese キウ}{</\textbf{rt}>}\mbox{}\newline 
\hspace*{1em}\hspace*{1em}\hspace*{1em}{</\textbf{ruby}>}{\textJapanese }\mbox{}\newline 
\hspace*{1em}\hspace*{1em}\hspace*{1em}\hspace*{1em} {\textJapanese 場}\mbox{}\newline 
\hspace*{1em}\hspace*{1em}{</\textbf{rb}>}\mbox{}\newline 
\hspace*{1em}\hspace*{1em}{<\textbf{rt}\hspace*{1em}{place}="{left}">}{\textJapanese ビリヤード}{</\textbf{rt}>}\mbox{}\newline 
\hspace*{1em}{</\textbf{ruby}>}\mbox{}\newline 
\textit{<!--{\textJapanese ...}-->}\mbox{}\newline 
{</\textbf{p}>}\end{shaded}\egroup\par \noindent  We could also use a standoff approach with \hyperref[TEI.anchor]{<anchor>} elements and pointers: \par\bgroup\index{p=<p>|exampleindex}\index{style=@style!<p>|exampleindex}\index{ruby=<ruby>|exampleindex}\index{rb=<rb>|exampleindex}\index{anchor=<anchor>|exampleindex}\index{anchor=<anchor>|exampleindex}\index{anchor=<anchor>|exampleindex}\index{anchor=<anchor>|exampleindex}\index{rt=<rt>|exampleindex}\index{place=@place!<rt>|exampleindex}\index{from=@from!<rt>|exampleindex}\index{to=@to!<rt>|exampleindex}\index{rt=<rt>|exampleindex}\index{place=@place!<rt>|exampleindex}\index{from=@from!<rt>|exampleindex}\index{to=@to!<rt>|exampleindex}\index{rt=<rt>|exampleindex}\index{place=@place!<rt>|exampleindex}\index{from=@from!<rt>|exampleindex}\index{to=@to!<rt>|exampleindex}\exampleFont \begin{shaded}\noindent\mbox{}{<\textbf{p}\hspace*{1em}{style}="{writing-mode: vertical-rl}"\mbox{}\newline 
\hspace*{1em}{xml:lang}="{ja}">}\mbox{}\newline 
\textit{<!--{\textJapanese ...}-->}\mbox{}\newline 
\hspace*{1em}{<\textbf{ruby}>}\mbox{}\newline 
\hspace*{1em}\hspace*{1em}{<\textbf{rb}>}\mbox{}\newline 
\hspace*{1em}\hspace*{1em}\hspace*{1em}{<\textbf{anchor}\hspace*{1em}{xml:id}="{da}"/>}{\textJapanese 打}\mbox{}\newline 
\hspace*{1em}\hspace*{1em}{<\textbf{anchor}\hspace*{1em}{xml:id}="{kyuu}"/>}{\textJapanese 球}\mbox{}\newline 
\hspace*{1em}\hspace*{1em}{<\textbf{anchor}\hspace*{1em}{xml:id}="{ba}"/>}{\textJapanese 場}\mbox{}\newline 
\hspace*{1em}\hspace*{1em}{<\textbf{anchor}\hspace*{1em}{xml:id}="{owari}"/>}\mbox{}\newline 
\hspace*{1em}\hspace*{1em}{</\textbf{rb}>}\mbox{}\newline 
\hspace*{1em}\hspace*{1em}{<\textbf{rt}\hspace*{1em}{place}="{left}"\hspace*{1em}{from}="{\#da}"\hspace*{1em}{to}="{\#owari}">}{\textJapanese ビリヤード}{</\textbf{rt}>}\mbox{}\newline 
\hspace*{1em}\hspace*{1em}{<\textbf{rt}\hspace*{1em}{place}="{right}"\hspace*{1em}{from}="{\#da}"\hspace*{1em}{to}="{\#kyuu}">}{\textJapanese ダ}{</\textbf{rt}>}\mbox{}\newline 
\hspace*{1em}\hspace*{1em}{<\textbf{rt}\hspace*{1em}{place}="{right}"\hspace*{1em}{from}="{\#kyuu}"\hspace*{1em}{to}="{\#ba}">}{\textJapanese キウ}{</\textbf{rt}>}\mbox{}\newline 
\hspace*{1em}{</\textbf{ruby}>}\mbox{}\newline 
\textit{<!--{\textJapanese ...}-->}\mbox{}\newline 
{</\textbf{p}>}\end{shaded}\egroup\par \noindent  Alternatively, if the encoding itself already includes segmentation below the word level, we can use the existing elements instead of adding \hyperref[TEI.anchor]{<anchor>}s: \par\bgroup\index{p=<p>|exampleindex}\index{style=@style!<p>|exampleindex}\index{ruby=<ruby>|exampleindex}\index{rb=<rb>|exampleindex}\index{c=<c>|exampleindex}\index{c=<c>|exampleindex}\index{c=<c>|exampleindex}\index{rt=<rt>|exampleindex}\index{place=@place!<rt>|exampleindex}\index{target=@target!<rt>|exampleindex}\index{rt=<rt>|exampleindex}\index{place=@place!<rt>|exampleindex}\index{target=@target!<rt>|exampleindex}\index{rt=<rt>|exampleindex}\index{place=@place!<rt>|exampleindex}\index{target=@target!<rt>|exampleindex}\exampleFont \begin{shaded}\noindent\mbox{}{<\textbf{p}\hspace*{1em}{style}="{writing-mode: vertical-rl}"\mbox{}\newline 
\hspace*{1em}{xml:lang}="{ja}">}\mbox{}\newline 
\textit{<!--{\textJapanese ...}-->}\mbox{}\newline 
\hspace*{1em}{<\textbf{ruby}>}\mbox{}\newline 
\hspace*{1em}\hspace*{1em}{<\textbf{rb}\hspace*{1em}{xml:id}="{dakyuuba}">}\mbox{}\newline 
\hspace*{1em}\hspace*{1em}\hspace*{1em}{<\textbf{c}\hspace*{1em}{xml:id}="{chr1}">}{\textJapanese 打}{</\textbf{c}>}\mbox{}\newline 
\hspace*{1em}\hspace*{1em}\hspace*{1em}{<\textbf{c}\hspace*{1em}{xml:id}="{chr2}">}{\textJapanese 球}{</\textbf{c}>}\mbox{}\newline 
\hspace*{1em}\hspace*{1em}\hspace*{1em}{<\textbf{c}>}{\textJapanese 場}{</\textbf{c}>}\mbox{}\newline 
\hspace*{1em}\hspace*{1em}{</\textbf{rb}>}\mbox{}\newline 
\hspace*{1em}\hspace*{1em}{<\textbf{rt}\hspace*{1em}{place}="{left}"\hspace*{1em}{target}="{\#dakyuuba}">}{\textJapanese ビリヤード}{</\textbf{rt}>}\mbox{}\newline 
\hspace*{1em}\hspace*{1em}{<\textbf{rt}\hspace*{1em}{place}="{right}"\hspace*{1em}{target}="{\#chr1}">}{\textJapanese ダ}{</\textbf{rt}>}\mbox{}\newline 
\hspace*{1em}\hspace*{1em}{<\textbf{rt}\hspace*{1em}{place}="{right}"\hspace*{1em}{target}="{\#chr2}">}{\textJapanese キウ}{</\textbf{rt}>}\mbox{}\newline 
\hspace*{1em}{</\textbf{ruby}>}\mbox{}\newline 
\textit{<!--{\textJapanese ...}-->}\mbox{}\newline 
{</\textbf{p}>}\end{shaded}\egroup\par \par
The current support for ruby is rudimentary, and in future releases of the Guidelines we expect to see more development of these features and recommendations.
\subsubsection[{Equivalents and Descriptions}]{Equivalents and Descriptions}\label{COHTGED}\par
Another group of elements is used to supply different kinds of names for objects described by the TEI. Examples of this are documentation of elements, attributes, classes (and also attribute values where appropriate), and description of glyphs. 
\begin{sansreflist}
  
\item [\textbf{<altIdent>}] (alternate identifier) supplies the recommended XML name for an element, class, attribute, etc. in some language.
\item [\textbf{<desc>}] (description) contains a short description of the purpose, function, or use of its parent element, or when the parent is a documentation element, describes or defines the object being documented. 
\item [\textbf{<equiv>}] (equivalent) specifies a component which is considered equivalent to the parent element, either by co-reference, or by external link.\hfil\\[-10pt]\begin{sansreflist}
    \item[@{\itshape uri}]
  (uniform resource identifier) references the underlying concept of which the parent is a representation by means of some external identifier
    \item[@{\itshape filter}]
  references an external script which contains a method to transform instances of this element to canonical TEI
    \item[@{\itshape name}]
  a single word which follows the rules defining a legal XML name (see \url{http://www.w3.org/TR/REC-xml/\#dt-name}), naming the underlying concept of which the parent is a representation.
    \item[@{\itshape predicate [att.predicate]}]
  the condition under which the element bearing this attribute applies, given as an XPath predicate expression.
\end{sansreflist}  
\end{sansreflist}
 Along with the \hyperref[TEI.gloss]{<gloss>} element mentioned above, these elements constitute the \textsf{model.glossLike} class. They are described in more detail in \textit{\hyperref[TDcrystalsCEdc]{22.4.1.\ Description of Components}}.
\subsection[{Simple Editorial Changes}]{Simple Editorial Changes}\label{COED}\par
As in editing a printed text, so in encoding a text in electronic form, it may be necessary to accommodate editorial comment on the text and to render account of any changes made to the text in preparing it. The tags described in this section may be used to record such editorial interventions, whether made by the encoder, by the editor of a printed edition used as a copy text, by earlier editors, or by the copyists of manuscripts.\par
The tags described here handle most common types of editorial intervention and stereotyped comment; where less structured commentary of other types is to be included, it may be marked using the \hyperref[TEI.note]{<note>} element described in section \textit{\hyperref[CONO]{3.9.\ Notes, Annotation, and Indexing}}. Systematic interpretive annotation is also possible using the various methods described in chapter \textit{\hyperref[SA]{16.\ Linking, Segmentation, and Alignment}}. The examples given here illustrate only simple cases of editorial intervention; in particular, they permit economical encoding of a simple set of alternative readings of a short span of text. To encode multiple views of large or heterogeneous spans of text, the mechanisms described in chapter \textit{\hyperref[SA]{16.\ Linking, Segmentation, and Alignment}} should be used. To encode multiple witnesses of a particular text, a similar mechanism designed specifically for critical editions is described in chapter \textit{\hyperref[TC]{12.\ Critical Apparatus}}.\par
For most of the elements discussed here, some encoders may wish to indicate both a \textit{responsibility}, that is, a code indicating the person or agency responsible for making the editorial intervention in question, and also an indication of the degree of \textit{certainty} which the encoder wishes to associate with the intervention. These requirements are served by the \textsf{att.global.responsibility} class, along with \textsf{att.global.source} and \textsf{att.dimensions}. Any of the elements discussed here thus may potentially carry any of the following optional attributes: 
\begin{sansreflist}
  
\item [\textbf{att.global.responsibility}] provides attributes indicating the agent responsible for some aspect of the text, the markup or something asserted by the markup, and the degree of certainty associated with it.\hfil\\[-10pt]\begin{sansreflist}
    \item[@{\itshape cert}]
  (certainty) signifies the degree of certainty associated with the intervention or interpretation.
    \item[@{\itshape resp}]
  (responsible party) indicates the agency responsible for the intervention or interpretation, for example an editor or transcriber.
\end{sansreflist}  
\item [\textbf{att.global.source}] provides an attribute used by elements to point to an external source.\hfil\\[-10pt]\begin{sansreflist}
    \item[@{\itshape source}]
  specifies the source from which some aspect of this element is drawn.
\end{sansreflist}  
\item [\textbf{att.editLike}] provides attributes describing the nature of an encoded scholarly intervention or interpretation of any kind.\hfil\\[-10pt]\begin{sansreflist}
    \item[@{\itshape evidence}]
  indicates the nature of the evidence supporting the reliability or accuracy of the intervention or interpretation.
\end{sansreflist}  
\item [\textbf{att.dimensions}] provides attributes for describing the size of physical objects.\hfil\\[-10pt]\begin{sansreflist}
    \item[@{\itshape unit}]
  names the unit used for the measurement
    \item[@{\itshape quantity}]
  specifies the length in the units specified
    \item[@{\itshape extent}]
  indicates the size of the object concerned using a project-specific vocabulary combining quantity and units in a single string of words.
    \item[@{\itshape precision}]
  characterizes the precision of the values specified by the other attributes.
    \item[@{\itshape scope}]
  where the measurement summarizes more than one observation, specifies the applicability of this measurement.
\end{sansreflist}  
\end{sansreflist}
\par
Many of the elements discussed here can be used in two ways. Their primary purpose is to indicate that the text encoded as the element's content represents an editorial intervention (or non-intervention) of a specific kind, indicated by the element itself. However, pairs or other meaningful groupings of such elements can also be supplied, wrapped within a special purpose \hyperref[TEI.choice]{<choice>} element: 
\begin{sansreflist}
  
\item [\textbf{<choice>}] (choice) groups a number of alternative encodings for the same point in a text.
\end{sansreflist}
 This element enables the encoder to represent for example a text in its ‘original’ uncorrected and unaltered form, alongside the same text in one or more ‘edited’ forms. This usage permits software to switch automatically between one ‘view’ of a text and another, so that (for example) a stylesheet may be set to display either the text in its original form or after the application of editorial interventions of particular kinds.\par
Elements which can be combined in this way constitute the \textsf{model.choicePart} class. The default members of this class are \hyperref[TEI.sic]{<sic>}, \hyperref[TEI.corr]{<corr>}, \hyperref[TEI.reg]{<reg>}, \hyperref[TEI.orig]{<orig>}, \hyperref[TEI.unclear]{<unclear>}, \hyperref[TEI.supplied]{<supplied>}, \hyperref[TEI.abbr]{<abbr>}, \hyperref[TEI.expan]{<expan>}, \hyperref[TEI.ex]{<ex>}, \hyperref[TEI.am]{<am>} and \hyperref[TEI.seg]{<seg>}; some of their functions and usage are described further below.\par
Three categories of editorial intervention are discussed in this section: \begin{itemize}
\item indication or correction of apparent errors 
\item indication or regularization of variant, irregular, non-standard, or eccentric forms
\item editorial additions, suppressions, and omissions
\end{itemize} \par
A more extended treatment of the use of these tags in transcriptional and editorial work is given in chapter \textit{\hyperref[PH]{11.\ Representation of Primary Sources}}.
\subsubsection[{Apparent Errors}]{Apparent Errors}\label{COEDCOR}\par
When the copy text is manifestly faulty, an encoder or transcriber may elect simply to correct it without comment, although for scholarly purposes it will often be more generally useful to record both the correction and the original state of the text. The elements described here enable all three approaches, and allows the last to be done in such a way as make it easy for software to present either the original or the correction. 
\begin{sansreflist}
  
\item [\textbf{<sic>}] (Latin for thus or so) contains text reproduced although apparently incorrect or inaccurate.
\item [\textbf{<corr>}] (correction) contains the correct form of a passage apparently erroneous in the copy text.
\end{sansreflist}
\par
The following examples show alternative treatment of the same material. The copy text reads: 
\begin{quote}Another property of computer-assisted historical research is that data modelling must permit any one textual feature or part of a textual feature to be a part of more than one information model and to allow the researcher to draw on several such models simultaneously, for example, to select from a machine-readable text those marginal comments which indicate that the date's mentioned in the main body of the text are incorrect.\end{quote}
\par
An encoder may choose to correct the typographic error, either silently or with an indication that a correction has been made, as follows: \par\bgroup\index{corr=<corr>|exampleindex}\exampleFont \begin{shaded}\noindent\mbox{}… marginal comments which indicate that the {<\textbf{corr}>}dates{</\textbf{corr}>}\mbox{}\newline 
 mentioned in the main body of the text are incorrect.\end{shaded}\egroup\par \par
Alternatively, the encoder may simply record the typographic error without correcting it, either without comment or with a \hyperref[TEI.sic]{<sic>} element to indicate the error is not a transcription error in the encoding: \par\bgroup\index{sic=<sic>|exampleindex}\exampleFont \begin{shaded}\noindent\mbox{}… marginal comments which indicate that the {<\textbf{sic}>}date's{</\textbf{sic}>}\mbox{}\newline 
 mentioned in the main body of the text are incorrect.\end{shaded}\egroup\par \par
If the encoder elects both to record the original source text and to provide a correction for the sake of word-search and other programs, both \hyperref[TEI.sic]{<sic>} and \hyperref[TEI.corr]{<corr>} are used, wrapped in a \hyperref[TEI.choice]{<choice>}: \par\bgroup\index{choice=<choice>|exampleindex}\index{corr=<corr>|exampleindex}\index{sic=<sic>|exampleindex}\exampleFont \begin{shaded}\noindent\mbox{}… marginal comments which indicate that the\mbox{}\newline 
{<\textbf{choice}>}\mbox{}\newline 
\hspace*{1em}{<\textbf{corr}>}dates{</\textbf{corr}>}\mbox{}\newline 
\hspace*{1em}{<\textbf{sic}>}date's{</\textbf{sic}>}\mbox{}\newline 
{</\textbf{choice}>} mentioned in the main body of the text are\mbox{}\newline 
 incorrect.\end{shaded}\egroup\par \noindent The \hyperref[TEI.sic]{<sic>} and \hyperref[TEI.corr]{<corr>} elements can appear in either order.\par
If it is desired to indicate the person or edition responsible for the emendation, this might be done as follows: \par\bgroup\index{choice=<choice>|exampleindex}\index{corr=<corr>|exampleindex}\index{resp=@resp!<corr>|exampleindex}\index{sic=<sic>|exampleindex}\index{respStmt=<respStmt>|exampleindex}\index{resp=<resp>|exampleindex}\index{name=<name>|exampleindex}\exampleFont \begin{shaded}\noindent\mbox{}… marginal comments which indicate that the\mbox{}\newline 
{<\textbf{choice}>}\mbox{}\newline 
\hspace*{1em}{<\textbf{corr}\hspace*{1em}{resp}="{\#msm}">}dates{</\textbf{corr}>}\mbox{}\newline 
\hspace*{1em}{<\textbf{sic}>}date's{</\textbf{sic}>}\mbox{}\newline 
{</\textbf{choice}>} mentioned in the main body of the text are\mbox{}\newline 
 incorrect.\mbox{}\newline 
\mbox{}\newline 
\textit{<!-- within the header for this document ... -->}\mbox{}\newline 
{<\textbf{respStmt}\hspace*{1em}{xml:id}="{msm}">}\mbox{}\newline 
\hspace*{1em}{<\textbf{resp}>}editor{</\textbf{resp}>}\mbox{}\newline 
\hspace*{1em}{<\textbf{name}>}C.M. Sperberg-McQueen{</\textbf{name}>}\mbox{}\newline 
{</\textbf{respStmt}>}\end{shaded}\egroup\par \noindent Here the {\itshape resp} attribute has been used to indicate responsibility for the correction. Its value (\#msm) is an example of the \textit{pointer} values discussed in section \textit{\hyperref[COXR]{3.7.\ Simple Links and Cross-References}}; in this case, it points to a \hyperref[TEI.respStmt]{<respStmt>} element within the TEI header, but any element might be indicated in this way, including for example a \hyperref[TEI.name]{<name>} element, or (if the module described in \textit{\hyperref[ND]{13.\ Names, Dates, People, and Places}} has been included) a \hyperref[TEI.person]{<person>} element. The {\itshape resp} attribute is available for all elements which are members of the \textsf{att.global.responsibility} class. The same class makes available a {\itshape cert} attribute, which may be used to indicate the degree of editorial confidence in a particular correction, as in the following example: \par\bgroup\index{choice=<choice>|exampleindex}\index{corr=<corr>|exampleindex}\index{cert=@cert!<corr>|exampleindex}\index{sic=<sic>|exampleindex}\exampleFont \begin{shaded}\noindent\mbox{}An {<\textbf{choice}>}\mbox{}\newline 
\hspace*{1em}{<\textbf{corr}\hspace*{1em}{cert}="{high}">}Autumn{</\textbf{corr}>}\mbox{}\newline 
\hspace*{1em}{<\textbf{sic}>}Antony{</\textbf{sic}>}\mbox{}\newline 
{</\textbf{choice}>} it was,\mbox{}\newline 
 That grew the more by reaping\end{shaded}\egroup\par \noindent  See further the discussion in section \textit{\hyperref[PHCC]{11.3.1.3.\ Correction and Conjecture}}.\par
Where, as here, the correction takes the form of adding text not otherwise present in the text being encoded, the encoder should use the \hyperref[TEI.corr]{<corr>} element. Where the correction is present in the text being encoded, and consists of some combination of visible additions and deletions, the elements \hyperref[TEI.add]{<add>} or \hyperref[TEI.del]{<del>} should be used: see further section \textit{\hyperref[COEDADD]{3.5.3.\ Additions, Deletions, and Omissions}} below. Where the correction takes the form of addition of material not present in the original because of physical damage or illegibility, the \hyperref[TEI.supplied]{<supplied>} element may be used. Where the ‘correction’ is simply a matter of expanding an abbreviation the \hyperref[TEI.ex]{<ex>} element may be used. These and other elements to support the detailed encoding of authorial or scribal interventions of this kind are all provided by the module described in chapter \textit{\hyperref[PH]{11.\ Representation of Primary Sources}}.
\subsubsection[{Regularization and Normalization}]{Regularization and Normalization}\label{COEDREG}\par
When the source text makes extensive use of variant forms or non-standard spellings, it may be desirable for a number of reasons to \textit{regularize} it: that is, to provide ‘standard’ or ‘regularized’ forms equivalent to the non-standard forms.\footnote{In some contexts, the term \textit{regularization} has a narrower and more specific significance than that proposed here: the \hyperref[TEI.reg]{<reg>} element may be used for any kind of regularization, including normalization, standardization, and modernization.}\par
As with other such changes to the copy text, the changes may be made silently (in which case the TEI header should specify the types of silent changes made) or may be explicitly marked using the following elements: 
\begin{sansreflist}
  
\item [\textbf{<reg>}] (regularization) contains a reading which has been regularized or normalized in some sense.
\item [\textbf{<orig>}] (original form) contains a reading which is marked as following the original, rather than being normalized or corrected.
\item [\textbf{<choice>}] (choice) groups a number of alternative encodings for the same point in a text.
\end{sansreflist}
\par
Typical applications for these elements include the production of editions intended for student or lay readers, linguistic research in which spelling or usage variation is not the main question at issue, production of spelling dictionaries, etc.\par
Consider this 16th-century text: 
\begin{quote}how godly a dede it is to overthrowe so wicked a race the world may judge: for my part I thinke there canot be a greater sacryfice to God.\end{quote}
\par
An encoder may choose to preserve the original spelling of this text, but simply flag it as nonstandard by using the \hyperref[TEI.orig]{<orig>} element with no attributes specified, as follows: \par\bgroup\index{p=<p>|exampleindex}\index{orig=<orig>|exampleindex}\index{orig=<orig>|exampleindex}\index{orig=<orig>|exampleindex}\index{orig=<orig>|exampleindex}\index{orig=<orig>|exampleindex}\exampleFont \begin{shaded}\noindent\mbox{}{<\textbf{p}>}...how godly a {<\textbf{orig}>}dede{</\textbf{orig}>} it is to\mbox{}\newline 
{<\textbf{orig}>}overthrowe{</\textbf{orig}>} so wicked a race the\mbox{}\newline 
 world may judge: for my part I {<\textbf{orig}>}thinke{</\textbf{orig}>}\mbox{}\newline 
 there {<\textbf{orig}>}canot{</\textbf{orig}>} be a greater\mbox{}\newline 
{<\textbf{orig}>}sacryfice{</\textbf{orig}>} to God{</\textbf{p}>}\end{shaded}\egroup\par \par
Alternatively, the encoder may simply indicate that certain words have been modernized by using the \hyperref[TEI.reg]{<reg>} element with no attributes specified, as follows:\par\bgroup\index{p=<p>|exampleindex}\index{reg=<reg>|exampleindex}\index{reg=<reg>|exampleindex}\index{reg=<reg>|exampleindex}\index{reg=<reg>|exampleindex}\index{reg=<reg>|exampleindex}\exampleFont \begin{shaded}\noindent\mbox{}{<\textbf{p}>}...how godly a\mbox{}\newline 
{<\textbf{reg}>}deed{</\textbf{reg}>} it is to {<\textbf{reg}>}overthrow{</\textbf{reg}>} so wicked a race the\mbox{}\newline 
 world may judge: for my part I {<\textbf{reg}>}think{</\textbf{reg}>}\mbox{}\newline 
 there {<\textbf{reg}>}cannot{</\textbf{reg}>} be a greater\mbox{}\newline 
{<\textbf{reg}>}sacrifice{</\textbf{reg}>} to God.{</\textbf{p}>}\end{shaded}\egroup\par \par
Alternatively, the encoder may elect to record both old and new spellings, so that (for example) the same electronic text may serve as the basis of an old- or new-spelling edition: \par\bgroup\index{p=<p>|exampleindex}\index{choice=<choice>|exampleindex}\index{orig=<orig>|exampleindex}\index{reg=<reg>|exampleindex}\index{choice=<choice>|exampleindex}\index{orig=<orig>|exampleindex}\index{reg=<reg>|exampleindex}\index{choice=<choice>|exampleindex}\index{orig=<orig>|exampleindex}\index{reg=<reg>|exampleindex}\index{choice=<choice>|exampleindex}\index{orig=<orig>|exampleindex}\index{reg=<reg>|exampleindex}\index{choice=<choice>|exampleindex}\index{orig=<orig>|exampleindex}\index{reg=<reg>|exampleindex}\exampleFont \begin{shaded}\noindent\mbox{}{<\textbf{p}>}...how godly a {<\textbf{choice}>}\mbox{}\newline 
\hspace*{1em}\hspace*{1em}{<\textbf{orig}>}dede{</\textbf{orig}>}\mbox{}\newline 
\hspace*{1em}\hspace*{1em}{<\textbf{reg}>}deed{</\textbf{reg}>}\mbox{}\newline 
\hspace*{1em}{</\textbf{choice}>} it is to\mbox{}\newline 
{<\textbf{choice}>}\mbox{}\newline 
\hspace*{1em}\hspace*{1em}{<\textbf{orig}>}overthrowe{</\textbf{orig}>}\mbox{}\newline 
\hspace*{1em}\hspace*{1em}{<\textbf{reg}>}overthrow{</\textbf{reg}>}\mbox{}\newline 
\hspace*{1em}{</\textbf{choice}>} so wicked a race the\mbox{}\newline 
 world may judge: for my part I {<\textbf{choice}>}\mbox{}\newline 
\hspace*{1em}\hspace*{1em}{<\textbf{orig}>}thinke{</\textbf{orig}>}\mbox{}\newline 
\hspace*{1em}\hspace*{1em}{<\textbf{reg}>}think{</\textbf{reg}>}\mbox{}\newline 
\hspace*{1em}{</\textbf{choice}>}\mbox{}\newline 
 there {<\textbf{choice}>}\mbox{}\newline 
\hspace*{1em}\hspace*{1em}{<\textbf{orig}>}canot{</\textbf{orig}>}\mbox{}\newline 
\hspace*{1em}\hspace*{1em}{<\textbf{reg}>}cannot{</\textbf{reg}>}\mbox{}\newline 
\hspace*{1em}{</\textbf{choice}>} be a greater\mbox{}\newline 
{<\textbf{choice}>}\mbox{}\newline 
\hspace*{1em}\hspace*{1em}{<\textbf{orig}>}sacryfice{</\textbf{orig}>}\mbox{}\newline 
\hspace*{1em}\hspace*{1em}{<\textbf{reg}>}sacrifice{</\textbf{reg}>}\mbox{}\newline 
\hspace*{1em}{</\textbf{choice}>} to God.{</\textbf{p}>}\end{shaded}\egroup\par \par
As elsewhere, the {\itshape resp} attribute may be used to specify the agency responsible for the regularization.
\subsubsection[{Additions, Deletions, and Omissions}]{Additions, Deletions, and Omissions}\label{COEDADD}\par
The following elements are used to indicate when words or phrases have been omitted from, added to, or marked for deletion from, a text. Like the other editorial elements, they allow for a wide range of editorial practices: 
\begin{sansreflist}
  
\item [\textbf{<gap>}] (gap) indicates a point where material has been omitted in a transcription, whether for editorial reasons described in the TEI header, as part of sampling practice, or because the material is illegible, invisible, or inaudible.\hfil\\[-10pt]\begin{sansreflist}
    \item[@{\itshape reason}]
  (reason) gives the reason for omission
\end{sansreflist}  
\item [\textbf{<unclear>}] (unclear) contains a word, phrase, or passage which cannot be transcribed with certainty because it is illegible or inaudible in the source.\hfil\\[-10pt]\begin{sansreflist}
    \item[@{\itshape reason}]
  indicates why the material is hard to transcribe.
\end{sansreflist}  
\item [\textbf{<add>}] (addition) contains letters, words, or phrases inserted in the source text by an author, scribe, or a previous annotator or corrector.
\item [\textbf{<del>}] (deletion) contains a letter, word, or passage deleted, marked as deleted, or otherwise indicated as superfluous or spurious in the copy text by an author, scribe, or a previous annotator or corrector.
\end{sansreflist}
\par
Encoders may choose to omit parts of the copy text for reasons ranging from illegibility of the source or impossibility of transcribing it, to editorial policy, e.g. a systematic exclusion of poetry or prose from an encoding. The full details of the policy decisions concerned should be documented in the TEI header (see section \textit{\hyperref[HD5]{2.3.\ The Encoding Description}}). Each place in the text at which omission has taken place should be marked with a \hyperref[TEI.gap]{<gap>} element, with optionally further information about the reason for the omission, its extent, and the person or agency responsible for it, as in the following examples: \par\bgroup\index{gap=<gap>|exampleindex}\index{reason=@reason!<gap>|exampleindex}\index{unit=@unit!<gap>|exampleindex}\index{quantity=@quantity!<gap>|exampleindex}\exampleFont \begin{shaded}\noindent\mbox{}{<\textbf{gap}\hspace*{1em}{reason}="{illegible}"\hspace*{1em}{unit}="{word}"\mbox{}\newline 
\hspace*{1em}{quantity}="{2}"/>}\end{shaded}\egroup\par \noindent \par\bgroup\index{gap=<gap>|exampleindex}\index{reason=@reason!<gap>|exampleindex}\index{extent=@extent!<gap>|exampleindex}\exampleFont \begin{shaded}\noindent\mbox{}{<\textbf{gap}\hspace*{1em}{reason}="{overwriting illegible}"\mbox{}\newline 
\hspace*{1em}{extent}="{several characters}"/>}\end{shaded}\egroup\par \noindent  Note that the extent of the gap may be marked precisely using attributes {\itshape unit} and {\itshape quantity}, or more descriptively using the {\itshape extent} attribute. Other, more detailed, options are also available for representing dimensions of any kind; see further \textit{\hyperref[msdim]{10.3.4.\ Dimensions}}.\par
The \hyperref[TEI.desc]{<desc>} element may be used to supply a description of the material omitted, where that is considered useful: \par\bgroup\index{gap=<gap>|exampleindex}\index{reason=@reason!<gap>|exampleindex}\index{quantity=@quantity!<gap>|exampleindex}\index{unit=@unit!<gap>|exampleindex}\index{desc=<desc>|exampleindex}\exampleFont \begin{shaded}\noindent\mbox{}{<\textbf{gap}\hspace*{1em}{reason}="{sampling}"\hspace*{1em}{quantity}="{120}"\mbox{}\newline 
\hspace*{1em}{unit}="{lines}">}\mbox{}\newline 
\hspace*{1em}{<\textbf{desc}>}irrelevant commentary{</\textbf{desc}>}\mbox{}\newline 
{</\textbf{gap}>}\end{shaded}\egroup\par \noindent  \par\bgroup\index{gap=<gap>|exampleindex}\index{reason=@reason!<gap>|exampleindex}\index{extent=@extent!<gap>|exampleindex}\index{desc=<desc>|exampleindex}\exampleFont \begin{shaded}\noindent\mbox{}… Their arrangement with respect to Jupiter and to each other was as follows:\mbox{}\newline 
{<\textbf{gap}\hspace*{1em}{reason}="{sampling}"\hspace*{1em}{extent}="{restOfPage}">}\mbox{}\newline 
\hspace*{1em}{<\textbf{desc}>}astrological figure{</\textbf{desc}>}\mbox{}\newline 
{</\textbf{gap}>}\mbox{}\newline 
 That is, there were two stars on the easterly side and one to the\mbox{}\newline 
 west; …\end{shaded}\egroup\par \noindent     \par
The \hyperref[TEI.add]{<add>} and \hyperref[TEI.del]{<del>} elements may be used to record where words or phrases have been added or deleted in the copy text. They are not appropriate where longer passages have been added or deleted, which span several elements; for these, the elements \hyperref[TEI.addSpan]{<addSpan>} and \hyperref[TEI.delSpan]{<delSpan>} described in chapter \textit{\hyperref[PHAD]{11.3.1.4.\ Additions and Deletions}} should be used.\par
Additions to a text may be recorded for a number of reasons. Sometimes they are marked in a distinctive way in the source text, for example by brackets or insertion above the line (\textit{supralinear} insertion), as in the following example, taken from a 19th century manuscript: \par\bgroup\index{add=<add>|exampleindex}\index{place=@place!<add>|exampleindex}\exampleFont \begin{shaded}\noindent\mbox{}The story I am going to relate is true as to its main facts,\mbox{}\newline 
 and as to the consequences {<\textbf{add}\hspace*{1em}{place}="{above}">}of\mbox{}\newline 
 these facts{</\textbf{add}>} from which this tale takes its title.\end{shaded}\egroup\par \par
The \hyperref[TEI.add]{<add>} element should not be used to mark editorial changes, such as supplying a word omitted by mistake from the source text or a passage present in another version. In these cases, either the \hyperref[TEI.corr]{<corr>} or \hyperref[TEI.supplied]{<supplied>} tags should be used, as discussed above in section \textit{\hyperref[COEDCOR]{3.5.1.\ Apparent Errors}}, and in section \textit{\hyperref[PHCC]{11.3.1.3.\ Correction and Conjecture}}, respectively.\par
The \hyperref[TEI.unclear]{<unclear>} element is used to mark passages in the original which cannot be read with confidence, or about which the transcriber is uncertain for other reasons, as for example when transcribing a partially inaudible or illegible source. Its {\itshape reason} and {\itshape resp} attributes are used, as with the \hyperref[TEI.gap]{<gap>} element, to indicate the cause of uncertainty and the person responsible for the conjectured reading.\par
For example: \par\bgroup\index{l=<l>|exampleindex}\index{l=<l>|exampleindex}\index{unclear=<unclear>|exampleindex}\index{reason=@reason!<unclear>|exampleindex}\exampleFont \begin{shaded}\noindent\mbox{}{<\textbf{l}>}And where the sandy mountain Fenwick scald{</\textbf{l}>}\mbox{}\newline 
{<\textbf{l}>}\mbox{}\newline 
\hspace*{1em}{<\textbf{unclear}\hspace*{1em}{reason}="{ink blot}">}The{</\textbf{unclear}>} sea between\mbox{}\newline 
 yet hence his pray'r prevail'd\mbox{}\newline 
{</\textbf{l}>}\end{shaded}\egroup\par \noindent  or from a spoken text: \par\bgroup\index{p=<p>|exampleindex}\index{unclear=<unclear>|exampleindex}\index{reason=@reason!<unclear>|exampleindex}\exampleFont \begin{shaded}\noindent\mbox{}{<\textbf{p}>}... and then {<\textbf{unclear}\hspace*{1em}{reason}="{passingTruck}">}marbled queen{</\textbf{unclear}>}...{</\textbf{p}>}\end{shaded}\egroup\par \par
Where the material affected is entirely illegible or inaudible, the \hyperref[TEI.gap]{<gap>} element discussed above should be used in preference.\par
The \hyperref[TEI.del]{<del>} element is used to mark material which is deleted in the source but which can still be read with some degree of confidence, as opposed to material which has been omitted by the encoder or transcriber either because it is entirely illegible or for some other reason. This is of particular importance in transcribing manuscript material, though deletion is also found in printed texts, sometimes for humorous purposes: \par\bgroup\index{l=<l>|exampleindex}\index{l=<l>|exampleindex}\index{l=<l>|exampleindex}\index{l=<l>|exampleindex}\index{del=<del>|exampleindex}\index{rend=@rend!<del>|exampleindex}\index{l=<l>|exampleindex}\exampleFont \begin{shaded}\noindent\mbox{}{<\textbf{l}>}One day I will sojourn to your shores{</\textbf{l}>}\mbox{}\newline 
{<\textbf{l}>}I live in the middle of England{</\textbf{l}>}\mbox{}\newline 
{<\textbf{l}>}But!{</\textbf{l}>}\mbox{}\newline 
{<\textbf{l}>}Norway! My soul resides in your watery\mbox{}\newline 
{<\textbf{del}\hspace*{1em}{rend}="{overstrike}">}fiords fyords fiiords{</\textbf{del}>}\mbox{}\newline 
{</\textbf{l}>}\mbox{}\newline 
{<\textbf{l}>}Inlets.{</\textbf{l}>}\end{shaded}\egroup\par \par
The {\itshape rend} attribute may be used to distinguish different methods of deletion in manuscript or typescript material, as in this line from the typescript of Eliot's \textit{Waste Land}: \par\bgroup\index{l=<l>|exampleindex}\index{del=<del>|exampleindex}\index{rend=@rend!<del>|exampleindex}\index{del=<del>|exampleindex}\index{type=@type!<del>|exampleindex}\exampleFont \begin{shaded}\noindent\mbox{}{<\textbf{l}>}\mbox{}\newline 
\hspace*{1em}{<\textbf{del}\hspace*{1em}{rend}="{overtyped}">}Mein{</\textbf{del}>} Frisch\mbox{}\newline 
{<\textbf{del}\hspace*{1em}{type}="{overstrike}">}schwebt{</\textbf{del}>} weht der Wind\mbox{}\newline 
{</\textbf{l}>}\end{shaded}\egroup\par \par
Deletion in manuscript or typescript is often associated with addition: \par\bgroup\index{l=<l>|exampleindex}\index{del=<del>|exampleindex}\index{rend=@rend!<del>|exampleindex}\index{add=<add>|exampleindex}\index{place=@place!<add>|exampleindex}\exampleFont \begin{shaded}\noindent\mbox{}{<\textbf{l}>}\mbox{}\newline 
\hspace*{1em}{<\textbf{del}\hspace*{1em}{rend}="{overstrike}">}Inviolable{</\textbf{del}>}\mbox{}\newline 
\hspace*{1em}{<\textbf{add}\hspace*{1em}{place}="{below}">}Inexplicable{</\textbf{add}>}\mbox{}\newline 
 splendour of Corinthian white and gold\mbox{}\newline 
{</\textbf{l}>}\end{shaded}\egroup\par \noindent  The \hyperref[TEI.subst]{<subst>} element discussed in \textit{\hyperref[PHSU]{11.3.1.5.\ Substitutions}} provides a way of grouping additions and deletions of this kind.\par
The \hyperref[TEI.del]{<del>} element should not be used where the deletion is such that material cannot be read with confidence, or read at all, or where the material has been omitted by the transcriber or editor for some other reason. Where the material deleted cannot be read with confidence, the \hyperref[TEI.unclear]{<unclear>} tag should be used with the {\itshape reason} attribute indicating that the difficulty of transcription is due to deletion. Where material has been omitted by the transcriber or editor, this may be indicated by use of the \hyperref[TEI.gap]{<gap>} element. A deletion in which some parts may be read but not others may thus be represented by one or more \hyperref[TEI.gap]{<gap>} elements intermingled with text, all contained by a \hyperref[TEI.del]{<del>} element. Text supplied or marked as unneccessary by an editor should be marked with the \hyperref[TEI.supplied]{<supplied>} and \hyperref[TEI.surplus]{<surplus>} elements (discussed in \textit{\hyperref[PHOM]{11.3.1.7.\ Text Omitted from or Supplied in the Transcription}}) rather than \hyperref[TEI.add]{<add>} and \hyperref[TEI.del]{<del>}. These two sets of elements allow the encoder to distinguish editorial changes from those visible in the source text.
\subsection[{Names, Numbers, Dates, Abbreviations, and Addresses}]{Names, Numbers, Dates, Abbreviations, and Addresses}\label{CONA}\par
This section describes a number of textual features which it is often convenient to distinguish from their surrounding text. Names, dates, and numbers are likely to be of particular importance to the scholar treating a text as source for a database; distinguishing such items from the surrounding text is however equally important to the scholar primarily interested in lexis.\par
The treatment of these textual features proposed here is not intended to be exhaustive: fuller treatments for names, numbers, measures, and dates are provided in the names and dates module (see chapter \textit{\hyperref[ND]{13.\ Names, Dates, People, and Places}}); more detailed treatment of abbreviations is provided by the transcription module (see section \textit{\hyperref[PHAB]{11.3.1.2.\ Abbreviation and Expansion}}).
\subsubsection[{Referring Strings}]{Referring Strings}\label{CONARS}\par
A \textit{referring string} is a phrase which refers to some person, place, object, etc. Two elements are provided to mark such strings: 
\begin{sansreflist}
  
\item [\textbf{<rs>}] (referencing string) contains a general purpose name or referring string.
\item [\textbf{<name>}] (name, proper noun) contains a proper noun or noun phrase.
\end{sansreflist}
 Both the \hyperref[TEI.name]{<name>} and \hyperref[TEI.rs]{<rs>} elements are members of the \textsf{att.typed} class, from which they inherit the following attributes: 
\begin{sansreflist}
  
\item [\textbf{att.typed}] provides attributes which can be used to classify or subclassify elements in any way.\hfil\\[-10pt]\begin{sansreflist}
    \item[@{\itshape type}]
  characterizes the element in some sense, using any convenient classification scheme or typology.
    \item[@{\itshape subtype}]
  (subtype) provides a sub-categorization of the element, if needed
\end{sansreflist}  
\end{sansreflist}
 which may be used to further categorize the kind of object referred to.\par
Examples include: \par\bgroup\index{p=<p>|exampleindex}\index{q=<q>|exampleindex}\index{rs=<rs>|exampleindex}\index{type=@type!<rs>|exampleindex}\index{q=<q>|exampleindex}\index{rs=<rs>|exampleindex}\index{type=@type!<rs>|exampleindex}\exampleFont \begin{shaded}\noindent\mbox{}{<\textbf{p}>}\mbox{}\newline 
\hspace*{1em}{<\textbf{q}>}My dear\mbox{}\newline 
\hspace*{1em}{<\textbf{rs}\hspace*{1em}{type}="{person}">}Mr. Bennet{</\textbf{rs}>}\mbox{}\newline 
\hspace*{1em}{</\textbf{q}>}, said his lady to\mbox{}\newline 
 him one day, {<\textbf{q}>}have you heard that {<\textbf{rs}\hspace*{1em}{type}="{place}">}Netherfield Park{</\textbf{rs}>} is let at last?{</\textbf{q}>}\mbox{}\newline 
{</\textbf{p}>}\end{shaded}\egroup\par \noindent  \par\bgroup\index{p=<p>|exampleindex}\index{rs=<rs>|exampleindex}\index{type=@type!<rs>|exampleindex}\exampleFont \begin{shaded}\noindent\mbox{}{<\textbf{p}>}Collectors of water-rents were appointed by the\mbox{}\newline 
{<\textbf{rs}\hspace*{1em}{type}="{org}">}Watering Committee{</\textbf{rs}>}.\mbox{}\newline 
 They were paid a commission not exceeding four per\mbox{}\newline 
 cent, and gave bond.{</\textbf{p}>}\end{shaded}\egroup\par \noindent  \par\bgroup\index{p=<p>|exampleindex}\index{rs=<rs>|exampleindex}\index{type=@type!<rs>|exampleindex}\index{rs=<rs>|exampleindex}\index{type=@type!<rs>|exampleindex}\index{q=<q>|exampleindex}\exampleFont \begin{shaded}\noindent\mbox{}{<\textbf{p}>}It being one of the principles of the\mbox{}\newline 
{<\textbf{rs}\hspace*{1em}{type}="{org}">}Circumlocution Office{</\textbf{rs}>} never, on any\mbox{}\newline 
 account whatsoever, to give a straightforward answer,\mbox{}\newline 
{<\textbf{rs}\hspace*{1em}{type}="{person}">}Mr Barnacle{</\textbf{rs}>} said, {<\textbf{q}>}Possibly.{</\textbf{q}>}\mbox{}\newline 
{</\textbf{p}>}\end{shaded}\egroup\par \par
As the following example shows, the \hyperref[TEI.rs]{<rs>} element may be used for any reference to a person, place, etc., not only to references in the form of a proper noun or noun phrase. \par\bgroup\index{p=<p>|exampleindex}\index{q=<q>|exampleindex}\index{rs=<rs>|exampleindex}\index{type=@type!<rs>|exampleindex}\index{rs=<rs>|exampleindex}\index{type=@type!<rs>|exampleindex}\exampleFont \begin{shaded}\noindent\mbox{}{<\textbf{p}>}\mbox{}\newline 
\hspace*{1em}{<\textbf{q}>}My dear {<\textbf{rs}\hspace*{1em}{type}="{person}">}Mr. Bennet{</\textbf{rs}>}\mbox{}\newline 
\hspace*{1em}{</\textbf{q}>}, said\mbox{}\newline 
{<\textbf{rs}\hspace*{1em}{type}="{person}">}his lady{</\textbf{rs}>} to him one day ... \mbox{}\newline 
{</\textbf{p}>}\end{shaded}\egroup\par \par
The \hyperref[TEI.name]{<name>} element by contrast is provided for the special case of referencing strings which consist only of proper nouns; it may be used synonymously with the \hyperref[TEI.rs]{<rs>} element, or nested within it if a referring string contains a mixture of common and proper nouns. The following example shows an alternative way of encoding the short sentence from \textit{Pride and Prejudice} quoted above: \par\bgroup\index{p=<p>|exampleindex}\index{q=<q>|exampleindex}\index{name=<name>|exampleindex}\index{type=@type!<name>|exampleindex}\index{rs=<rs>|exampleindex}\index{type=@type!<rs>|exampleindex}\index{q=<q>|exampleindex}\index{name=<name>|exampleindex}\index{type=@type!<name>|exampleindex}\exampleFont \begin{shaded}\noindent\mbox{}{<\textbf{p}>}\mbox{}\newline 
\hspace*{1em}{<\textbf{q}>}My dear {<\textbf{name}\hspace*{1em}{type}="{person}">}Mr. Bennet{</\textbf{name}>},{</\textbf{q}>} said {<\textbf{rs}\hspace*{1em}{type}="{person}">}his lady{</\textbf{rs}>} to him one day,\mbox{}\newline 
{<\textbf{q}>}have you heard that {<\textbf{name}\hspace*{1em}{type}="{place}">}Netherfield Park{</\textbf{name}>} is let at last?{</\textbf{q}>}\mbox{}\newline 
{</\textbf{p}>}\end{shaded}\egroup\par \noindent  As the following example shows, a proper name may be nested within a referring string: \par\bgroup\index{rs=<rs>|exampleindex}\index{name=<name>|exampleindex}\exampleFont \begin{shaded}\noindent\mbox{}{<\textbf{rs}>}His Excellency the Life President, {<\textbf{name}>}Ngwazi Dr H. Kamuzu Banda{</\textbf{name}>}\mbox{}\newline 
{</\textbf{rs}>}\end{shaded}\egroup\par \par
Simply tagging something as a name is generally not enough to enable automatic processing of personal names into the canonical forms usually required for reference purposes. The name as it appears in the text may be inconsistently spelled, partial, or vague. Moreover, name prefixes such as \textit{van} or \textit{de la} may or may not be included as part of the reference form of a name, depending on the language and country of origin of the bearer.\par
Two issues arise in this context: firstly, there may be a need to encode a regularized form of a name, distinct from the actual form in the source to hand; secondly, there may be a need to identify the particular person, place, etc. referred to by the name, irrespective of whether the name itself is normalized or not. The element \hyperref[TEI.reg]{<reg>}, introduced in \textit{\hyperref[COEDREG]{3.5.2.\ Regularization and Normalization}} is provided for the former purpose; the attributes {\itshape key} or {\itshape ref} for the latter.\par
The {\itshape key} and {\itshape ref} attributes are common to all members of the \textsf{att.canonical} class and are defined as follows: 
\begin{sansreflist}
  
\item [\textbf{att.canonical}] provides attributes which can be used to associate a representation such as a name or title with canonical information about the object being named or referenced.\hfil\\[-10pt]\begin{sansreflist}
    \item[@{\itshape key}]
  provides an externally-defined means of identifying the entity (or entities) being named, using a coded value of some kind.
    \item[@{\itshape ref}]
  (reference) provides an explicit means of locating a full definition or identity for the entity being named by means of one or more URIs.
\end{sansreflist}  
\end{sansreflist}
\par
A very useful application for them is as a means of gathering together all references to the same individual or location scattered throughout a document: \par\bgroup\index{q=<q>|exampleindex}\index{rs=<rs>|exampleindex}\index{key=@key!<rs>|exampleindex}\index{type=@type!<rs>|exampleindex}\index{rs=<rs>|exampleindex}\index{key=@key!<rs>|exampleindex}\index{type=@type!<rs>|exampleindex}\index{q=<q>|exampleindex}\index{rs=<rs>|exampleindex}\index{key=@key!<rs>|exampleindex}\index{type=@type!<rs>|exampleindex}\exampleFont \begin{shaded}\noindent\mbox{}{<\textbf{q}>}My dear\mbox{}\newline 
{<\textbf{rs}\hspace*{1em}{key}="{BENM1}"\hspace*{1em}{type}="{person}">}Mr. Bennet{</\textbf{rs}>},{</\textbf{q}>} said\mbox{}\newline 
{<\textbf{rs}\hspace*{1em}{key}="{BENM2}"\hspace*{1em}{type}="{person}">}his lady{</\textbf{rs}>} to him one day,\mbox{}\newline 
\mbox{}\newline 
{<\textbf{q}>}have you heard that\mbox{}\newline 
{<\textbf{rs}\hspace*{1em}{key}="{NETP1}"\hspace*{1em}{type}="{place}">}Netherfield Park{</\textbf{rs}>} is let at\mbox{}\newline 
 last?{</\textbf{q}>}\end{shaded}\egroup\par \noindent  \par\bgroup\index{p=<p>|exampleindex}\index{name=<name>|exampleindex}\index{key=@key!<name>|exampleindex}\index{type=@type!<name>|exampleindex}\index{rs=<rs>|exampleindex}\index{key=@key!<rs>|exampleindex}\exampleFont \begin{shaded}\noindent\mbox{}{<\textbf{p}>}\mbox{}\newline 
\hspace*{1em}{<\textbf{name}\hspace*{1em}{key}="{VOM1}"\hspace*{1em}{type}="{person}">}Mme. de Volanges{</\textbf{name}>} \mbox{}\newline 
 marie {<\textbf{rs}\hspace*{1em}{key}="{VOM2}">}sa fille{</\textbf{rs}>}:\mbox{}\newline 
 c'est encore un secret;\mbox{}\newline 
 mais elle m'en a fait part hier.\mbox{}\newline 
{</\textbf{p}>}\end{shaded}\egroup\par \par
The value of the {\itshape key} attribute may be an unexpanded code, as in the examples above, with no particular significance. More usually however, it will be an externally defined code of some kind, as provided by a standard reference source. \par\bgroup\index{p=<p>|exampleindex}\index{name=<name>|exampleindex}\index{key=@key!<name>|exampleindex}\index{type=@type!<name>|exampleindex}\exampleFont \begin{shaded}\noindent\mbox{}{<\textbf{p}>}\mbox{}\newline 
\hspace*{1em}{<\textbf{name}\hspace*{1em}{key}="{LHR}"\hspace*{1em}{type}="{airport}">}Heathrow{</\textbf{name}>}\mbox{}\newline 
{</\textbf{p}>}\end{shaded}\egroup\par \par
The standard reference source should be documented, for example using a \hyperref[TEI.taxonomy]{<taxonomy>} element in the TEI header.\par
The {\itshape ref} attribute can be used to point directly to some other resource providing more information about the entity named by the element, such as an authority record in a database, an encylopaedia entry, another element in the same or a different document etc. \par\bgroup\index{p=<p>|exampleindex}\index{name=<name>|exampleindex}\index{ref=@ref!<name>|exampleindex}\index{type=@type!<name>|exampleindex}\exampleFont \begin{shaded}\noindent\mbox{}{<\textbf{p}>}\mbox{}\newline 
\hspace*{1em}{<\textbf{name}\hspace*{1em}{ref}="{http://en.wikipedia.org/wiki/Heathrow\textunderscore airport}"\mbox{}\newline 
\hspace*{1em}\hspace*{1em}{type}="{airport}">}Heathrow{</\textbf{name}>}\mbox{}\newline 
{</\textbf{p}>}\end{shaded}\egroup\par \par
This use should be distinguished from the use of a nested \hyperref[TEI.reg]{<reg>} (regularization) element to provide the standard form of a referring string, as in this example: \par\bgroup\index{p=<p>|exampleindex}\index{rs=<rs>|exampleindex}\index{key=@key!<rs>|exampleindex}\index{type=@type!<rs>|exampleindex}\index{reg=<reg>|exampleindex}\exampleFont \begin{shaded}\noindent\mbox{}{<\textbf{p}>}My personal life during\mbox{}\newline 
 the administration of {<\textbf{rs}\hspace*{1em}{key}="{POJA1}"\hspace*{1em}{type}="{person}">}Col. Polk\mbox{}\newline 
\hspace*{1em}\hspace*{1em} ({<\textbf{reg}>}Polk, James K.{</\textbf{reg}>}){</\textbf{rs}>} has but poorly compensated me for the\mbox{}\newline 
 suspended enjoyments and pursuits of private and professional\mbox{}\newline 
 spheres{</\textbf{p}>}\end{shaded}\egroup\par \par
No particular syntax is proposed for the values of the {\itshape key} attribute, since its form will depend entirely on practice within a given project. For the same reason, this attribute is not recommended in data interchange, since there is no way of ensuring that the values used by one project are distinct from those used by another. In such a situation, a preferable approach for magic tokens which follows standard practice on the Web is to use a {\itshape ref} attribute whose value is a tag URI as defined in \hyperref[RFC4151]{RFC 4151}. For example: \par\bgroup\index{p=<p>|exampleindex}\index{name=<name>|exampleindex}\index{ref=@ref!<name>|exampleindex}\index{type=@type!<name>|exampleindex}\index{rs=<rs>|exampleindex}\index{ref=@ref!<rs>|exampleindex}\exampleFont \begin{shaded}\noindent\mbox{}{<\textbf{p}>}\mbox{}\newline 
\hspace*{1em}{<\textbf{name}\hspace*{1em}{ref}="{tag:projectname.org,2012:VOM1}"\mbox{}\newline 
\hspace*{1em}\hspace*{1em}{type}="{person}">}Mme. de Volanges{</\textbf{name}>} marie {<\textbf{rs}\hspace*{1em}{ref}="{tag:theworksoflaclos.org,2012:VOM2}">}sa fille{</\textbf{rs}>}: c'est encore un secret;\mbox{}\newline 
 mais elle m'en a fait part hier.\mbox{}\newline 
{</\textbf{p}>}\end{shaded}\egroup\par \noindent  The inclusion of the domain name of the party responsible for tagging (\textsf{theworksoflaclos.org}), as specified in RFC 4151, helps ensure uniqueness of magic token values across TEI encoding projects, allowing for improved interchange of TEI documents.\par
The \hyperref[TEI.choice]{<choice>} element discussed in \textit{\hyperref[COED]{3.5.\ Simple Editorial Changes}} may be used if it is desired to record both a normalized form of a name and the name used in the source being encoded: \par\bgroup\index{p=<p>|exampleindex}\index{name=<name>|exampleindex}\index{ref=@ref!<name>|exampleindex}\index{type=@type!<name>|exampleindex}\index{choice=<choice>|exampleindex}\index{orig=<orig>|exampleindex}\index{reg=<reg>|exampleindex}\index{name=<name>|exampleindex}\index{ref=@ref!<name>|exampleindex}\index{type=@type!<name>|exampleindex}\index{name=<name>|exampleindex}\index{ref=@ref!<name>|exampleindex}\index{type=@type!<name>|exampleindex}\exampleFont \begin{shaded}\noindent\mbox{}{<\textbf{p}>}\mbox{}\newline 
\hspace*{1em}{<\textbf{name}\hspace*{1em}{ref}="{tag:projectname.org,2012:WADLM1}"\mbox{}\newline 
\hspace*{1em}\hspace*{1em}{type}="{person}">}\mbox{}\newline 
\hspace*{1em}\hspace*{1em}{<\textbf{choice}>}\mbox{}\newline 
\hspace*{1em}\hspace*{1em}\hspace*{1em}{<\textbf{orig}>}Walter de la Mare{</\textbf{orig}>}\mbox{}\newline 
\hspace*{1em}\hspace*{1em}\hspace*{1em}{<\textbf{reg}>}de la Mare, Walter{</\textbf{reg}>}\mbox{}\newline 
\hspace*{1em}\hspace*{1em}{</\textbf{choice}>}\mbox{}\newline 
\hspace*{1em}{</\textbf{name}>}\mbox{}\newline 
 was born at {<\textbf{name}\hspace*{1em}{ref}="{tag:projectname.org,2012:Ch1}"\mbox{}\newline 
\hspace*{1em}\hspace*{1em}{type}="{place}">}Charlton{</\textbf{name}>}, in\mbox{}\newline 
{<\textbf{name}\hspace*{1em}{ref}="{tag:projectname.org,2012:KT1}"\mbox{}\newline 
\hspace*{1em}\hspace*{1em}{type}="{county}">}Kent{</\textbf{name}>}, in 1873.\mbox{}\newline 
{</\textbf{p}>}\end{shaded}\egroup\par \par
The \hyperref[TEI.index]{<index>} element discussed in \textit{\hyperref[CONOIX]{3.9.2.\ Index Entries}} may be more appropriate if the function of the regularization is to provide a consistent index: \par\bgroup\index{p=<p>|exampleindex}\index{name=<name>|exampleindex}\index{type=@type!<name>|exampleindex}\index{name=<name>|exampleindex}\index{type=@type!<name>|exampleindex}\index{index=<index>|exampleindex}\index{term=<term>|exampleindex}\exampleFont \begin{shaded}\noindent\mbox{}{<\textbf{p}>}\mbox{}\newline 
\hspace*{1em}{<\textbf{name}\hspace*{1em}{type}="{place}">}Montaillou{</\textbf{name}>} is not a large parish.\mbox{}\newline 
 At the time of the events which led to\mbox{}\newline 
{<\textbf{name}\hspace*{1em}{type}="{person}">}Fournier{<\textbf{index}>}\mbox{}\newline 
\hspace*{1em}\hspace*{1em}\hspace*{1em}{<\textbf{term}>}Benedict XII, Pope of Avignon (Jacques Fournier){</\textbf{term}>}\mbox{}\newline 
\hspace*{1em}\hspace*{1em}{</\textbf{index}>}\mbox{}\newline 
\hspace*{1em}{</\textbf{name}>}'s\mbox{}\newline 
 investigations, the local population consisted of between 200 and 250 inhabitants.\mbox{}\newline 
{</\textbf{p}>}\end{shaded}\egroup\par \noindent  Although adequate for many simple applications, these methods have two inconveniences: if the name occurs many times, then its regularized form would be repeated many times; and the burden of additional XML markup in the body of the text may be inconvenient to maintain and complex to process. For applications such as onomastics, relating to persons or places named rather than the name itself, or wherever a detailed analysis of the component parts of a name is needed, the specialized elements described in chapter \textit{\hyperref[ND]{13.\ Names, Dates, People, and Places}} or the analytical tools described in chapter \textit{\hyperref[FS]{18.\ Feature Structures}} should be used.
\subsubsection[{Addresses}]{Addresses}\label{CONAAD}\par
These Guidelines propose the following elements to distinguish postal and electronic addresses: 
\begin{sansreflist}
  
\item [\textbf{<address>}] (address) contains a postal address, for example of a publisher, an organization, or an individual.
\item [\textbf{<email>}] (electronic mail address) contains an email address identifying a location to which email messages can be delivered.
\end{sansreflist}
 These two elements constitute the class of \textsf{model.addressLike} elements; for other kinds of address this class may be extended by adding new elements if necessary.\par
These Guidelines provide no particular means for encoding the substructure of an email address (for example, distinguishing the local part from the domain part), nor of distinguishing personal email addresses from generic or fictitious ones. \par\bgroup\index{email=<email>|exampleindex}\exampleFont \begin{shaded}\noindent\mbox{}{<\textbf{email}>}info@tei-c.org{</\textbf{email}>}\end{shaded}\egroup\par \par
The simplest way of encoding a postal address is to regard it as a series of distinct lines, just as they might be written on an envelope. The following element supports this view: 
\begin{sansreflist}
  
\item [\textbf{<addrLine>}] (address line) contains one line of a postal address.
\end{sansreflist}
 Here is an example of a postal address encoded using this approach: \par\bgroup\index{address=<address>|exampleindex}\index{addrLine=<addrLine>|exampleindex}\index{addrLine=<addrLine>|exampleindex}\index{addrLine=<addrLine>|exampleindex}\exampleFont \begin{shaded}\noindent\mbox{}{<\textbf{address}>}\mbox{}\newline 
\hspace*{1em}{<\textbf{addrLine}>}110 Southmoor Road,{</\textbf{addrLine}>}\mbox{}\newline 
\hspace*{1em}{<\textbf{addrLine}>}Oxford OX2 6RB,{</\textbf{addrLine}>}\mbox{}\newline 
\hspace*{1em}{<\textbf{addrLine}>}UK{</\textbf{addrLine}>}\mbox{}\newline 
{</\textbf{address}>}\end{shaded}\egroup\par \par
Alternatively, an address may be encoded as a structure of more semantically rich elements. The class \textsf{model.addrPart} element class identifies a number of such possible components: 
\begin{sansreflist}
  
\item [\textbf{<street>}] contains a full street address including any name or number identifying a building as well as the name of the street or route on which it is located.
\item [\textbf{<name>}] (name, proper noun) contains a proper noun or noun phrase.
\item [\textbf{<postCode>}] (postal code) contains a numerical or alphanumeric code used as part of a postal address to simplify sorting or delivery of mail.
\item [\textbf{<postBox>}] (postal box or post office box) contains a number or other identifier for some postal delivery point other than a street address.
\item [\textbf{model.nameLike}] groups elements which name or refer to a person, place, or organization. \par 
\begin{longtable}{P{0.4008733624454148\textwidth}P{0.44912663755458515\textwidth}}
\hyperref[TEI.model.nameLike.agent]{model.nameLike.agent}\tabcellsep groups elements which contain names of individuals or corporate bodies.\\
\hyperref[TEI.model.offsetLike]{model.offsetLike}\tabcellsep groups elements which can appear only as part of a place name.\\
\hyperref[TEI.model.persNamePart]{model.persNamePart}\tabcellsep groups elements which form part of a personal name.\\
\hyperref[TEI.model.placeStateLike]{model.placeStateLike}\tabcellsep groups elements which describe changing states of a place.\end{longtable} \par
  \par 
\begin{longtable}{P{0.17031539888682745\textwidth}P{0.6796846011131725\textwidth}}
\hyperref[TEI.idno]{idno}\tabcellsep (identifier) supplies any form of identifier used to identify some object, such as a bibliographic item, a person, a title, an organization, etc. in a standardized way.\\
\hyperref[TEI.lang]{lang}\tabcellsep (language name) contains the name of a language mentioned in etymological or other linguistic discussion.\\
\hyperref[TEI.objectName]{objectName}\tabcellsep (name of an object) contains a proper noun or noun phrase used to refer to an object.\\
\hyperref[TEI.rs]{rs}\tabcellsep (referencing string) contains a general purpose name or referring string.\end{longtable} \par
 
\item [\textbf{model.persNamePart}] groups elements which form part of a personal name. \par 
\begin{longtable}{P{0.17762331838565024\textwidth}P{0.6723766816143498\textwidth}}
\hyperref[TEI.addName]{addName}\tabcellsep (additional name) contains an additional name component, such as a nickname, epithet, or alias, or any other descriptive phrase used within a personal name.\\
\hyperref[TEI.forename]{forename}\tabcellsep (forename) contains a forename, given or baptismal name.\\
\hyperref[TEI.genName]{genName}\tabcellsep (generational name component) contains a name component used to distinguish otherwise similar names on the basis of the relative ages or generations of the persons named.\\
\hyperref[TEI.nameLink]{nameLink}\tabcellsep (name link) contains a connecting phrase or link used within a name but not regarded as part of it, such as \textit{van der} or \textit{of}.\\
\hyperref[TEI.persPronouns]{persPronouns}\tabcellsep (personal pronouns) indicates the personal pronouns used, or assumed to be used, by the individual being described.\\
\hyperref[TEI.roleName]{roleName}\tabcellsep (role name) contains a name component which indicates that the referent has a particular role or position in society, such as an official title or rank.\\
\hyperref[TEI.surname]{surname}\tabcellsep (surname) contains a family (inherited) name, as opposed to a given, baptismal, or nick name.\end{longtable} \par
 
\item [\textbf{model.placeNamePart}] groups elements which form part of a place name. \par 
\begin{longtable}{P{0.1712285456187895\textwidth}P{0.6787714543812104\textwidth}}
\hyperref[TEI.bloc]{bloc}\tabcellsep (bloc) contains the name of a geo-political unit consisting of two or more nation states or countries.\\
\hyperref[TEI.country]{country}\tabcellsep (country) contains the name of a geo-political unit, such as a nation, country, colony, or commonwealth, larger than or administratively superior to a region and smaller than a bloc.\\
\hyperref[TEI.district]{district}\tabcellsep (district) contains the name of any kind of subdivision of a settlement, such as a parish, ward, or other administrative or geographic unit.\\
\hyperref[TEI.geogName]{geogName}\tabcellsep (geographical name) identifies a name associated with some geographical feature such as Windrush Valley or Mount Sinai.\\
\hyperref[TEI.placeName]{placeName}\tabcellsep (place name) contains an absolute or relative place name.\\
\hyperref[TEI.region]{region}\tabcellsep (region) contains the name of an administrative unit such as a state, province, or county, larger than a settlement, but smaller than a country.\\
\hyperref[TEI.settlement]{settlement}\tabcellsep (settlement) contains the name of a settlement such as a city, town, or village identified as a single geo-political or administrative unit.\end{longtable} \par
 
\end{sansreflist}
 Any number of elements from the \textsf{model.addrPart} class may appear within an address and in any order. None of them is required.\par
Where code letters are commonly used in addresses (for example, to identify regions or countries) a useful practice is to supply the full name of the region or country as the content of the element, but to supply the abbreviatory code as the value of the global {\itshape n} attribute, so that (for example) an application preparing formatted labels can readily find the required information. Other components of addresses may be represented using the general-purpose \hyperref[TEI.name]{<name>} element or (when the additional module for names and dates is included) the more specialized elements provided for that purpose.\par
Using just the elements defined by the core module, the above address could thus be represented as follows: \par\bgroup\index{address=<address>|exampleindex}\index{street=<street>|exampleindex}\index{name=<name>|exampleindex}\index{type=@type!<name>|exampleindex}\index{postCode=<postCode>|exampleindex}\index{name=<name>|exampleindex}\index{type=@type!<name>|exampleindex}\exampleFont \begin{shaded}\noindent\mbox{}{<\textbf{address}>}\mbox{}\newline 
\hspace*{1em}{<\textbf{street}>}110 Southmoor Road{</\textbf{street}>}\mbox{}\newline 
\hspace*{1em}{<\textbf{name}\hspace*{1em}{type}="{city}">}Oxford{</\textbf{name}>}\mbox{}\newline 
\hspace*{1em}{<\textbf{postCode}>}OX2 6RB{</\textbf{postCode}>}\mbox{}\newline 
\hspace*{1em}{<\textbf{name}\hspace*{1em}{type}="{country}">}United Kingdom{</\textbf{name}>}\mbox{}\newline 
{</\textbf{address}>}\end{shaded}\egroup\par \par
The order of elements within an address is highly culture-specific, and is therefore unconstrained: \par\bgroup\index{address=<address>|exampleindex}\index{name=<name>|exampleindex}\index{type=@type!<name>|exampleindex}\index{name=<name>|exampleindex}\index{type=@type!<name>|exampleindex}\index{postCode=<postCode>|exampleindex}\index{name=<name>|exampleindex}\index{type=@type!<name>|exampleindex}\index{street=<street>|exampleindex}\exampleFont \begin{shaded}\noindent\mbox{}{<\textbf{address}>}\mbox{}\newline 
\hspace*{1em}{<\textbf{name}\hspace*{1em}{type}="{org}">}Università di Bologna{</\textbf{name}>}\mbox{}\newline 
\hspace*{1em}{<\textbf{name}\hspace*{1em}{type}="{country}">}Italy{</\textbf{name}>}\mbox{}\newline 
\hspace*{1em}{<\textbf{postCode}>}40126{</\textbf{postCode}>}\mbox{}\newline 
\hspace*{1em}{<\textbf{name}\hspace*{1em}{type}="{city}">}Bologna{</\textbf{name}>}\mbox{}\newline 
\hspace*{1em}{<\textbf{street}>}via Marsala 24{</\textbf{street}>}\mbox{}\newline 
{</\textbf{address}>}\end{shaded}\egroup\par \par
A telephone number (normally outside of the \hyperref[TEI.address]{<address>} element) might be tagged with an \hyperref[TEI.addrLine]{<addrLine>} and \hyperref[TEI.ref]{<ref>}, with the number itself appearing in the \texttt{tel:} namespace: \par\bgroup\index{addrLine=<addrLine>|exampleindex}\index{ref=<ref>|exampleindex}\index{target=@target!<ref>|exampleindex}\exampleFont \begin{shaded}\noindent\mbox{}{<\textbf{addrLine}>}\mbox{}\newline 
\hspace*{1em}{<\textbf{ref}\hspace*{1em}{target}="{tel:+1-201-555-0123}">}(201) 555 0123{</\textbf{ref}>}\mbox{}\newline 
{</\textbf{addrLine}>}\end{shaded}\egroup\par \par
For further discussion of ways of regularizing the names of places, see section \textit{\hyperref[CONA]{3.6.\ Names, Numbers, Dates, Abbreviations, and Addresses}}. A full postal address may also include the name of the addressee, tagged as above using the general purpose \hyperref[TEI.name]{<name>} element.\par
When a schema includes the names and dates module discussed in chapter \textit{\hyperref[ND]{13.\ Names, Dates, People, and Places}}, a large number of more specific elements such as \hyperref[TEI.country]{<country>} or \hyperref[TEI.settlement]{<settlement>} will be available from the class \textsf{model.addrPart}. The above example might then be encoded as follows: \par\bgroup\index{address=<address>|exampleindex}\index{street=<street>|exampleindex}\index{settlement=<settlement>|exampleindex}\index{postCode=<postCode>|exampleindex}\index{country=<country>|exampleindex}\exampleFont \begin{shaded}\noindent\mbox{}{<\textbf{address}>}\mbox{}\newline 
\hspace*{1em}{<\textbf{street}>}110 Southmoor Road{</\textbf{street}>}\mbox{}\newline 
\hspace*{1em}{<\textbf{settlement}>}Oxford{</\textbf{settlement}>}\mbox{}\newline 
\hspace*{1em}{<\textbf{postCode}>}OX2 6RB{</\textbf{postCode}>}\mbox{}\newline 
\hspace*{1em}{<\textbf{country}>}United Kingdom{</\textbf{country}>}\mbox{}\newline 
{</\textbf{address}>}\end{shaded}\egroup\par 
\subsubsection[{Numbers and Measures}]{Numbers and Measures}\label{CONANU}\par
This section describes elements provided for the simple encoding of numbers and measurements and gives some indication of circumstances in which this may usefully be done. The following phrase level elements are provided for this purpose: 
\begin{sansreflist}
  
\item [\textbf{<num>}] (number) contains a number, written in any form.\hfil\\[-10pt]\begin{sansreflist}
    \item[@{\itshape type}]
  indicates the type of numeric value.
    \item[@{\itshape value}]
  supplies the value of the number in standard form.
\end{sansreflist}  
\item [\textbf{<measure>}] (measure) contains a word or phrase referring to some quantity of an object or commodity, usually comprising a number, a unit, and a commodity name.\hfil\\[-10pt]\begin{sansreflist}
    \item[@{\itshape type}]
  specifies the type of measurement in any convenient typology.
\end{sansreflist}  
\item [\textbf{<measureGrp>}] (measure group) contains a group of dimensional specifications which relate to the same object, for example the height and width of a manuscript page.
\end{sansreflist}
\par
Like names or abbreviations, numbers can occur virtually anywhere in a text. Numbers are special in that they can be written with either letters or digits (\textit{twenty-one}, \textit{xxi}, and \textit{21}) and their presentation is language-dependent (e.g. English \textit{5th} becomes Greek \textit{5.}; English \textit{123,456.78} equals French \textit{123.456,78}).\par
For many kinds of application, e.g. natural-language processing or machine translation, numbers are not regarded as ‘lexical’ in the same way as other parts of a text. For these and other applications, the \hyperref[TEI.num]{<num>} element provides a convenient method of distinguishing numbers from the surrounding text. For other kinds of application, numbers are only useful if normalized: here the \hyperref[TEI.num]{<num>} element is useful precisely because it provides a standardized way of representing a numerical value.\par
For example: \par\bgroup\index{num=<num>|exampleindex}\index{value=@value!<num>|exampleindex}\index{num=<num>|exampleindex}\index{type=@type!<num>|exampleindex}\index{value=@value!<num>|exampleindex}\index{num=<num>|exampleindex}\index{type=@type!<num>|exampleindex}\index{value=@value!<num>|exampleindex}\index{num=<num>|exampleindex}\index{type=@type!<num>|exampleindex}\index{value=@value!<num>|exampleindex}\index{num=<num>|exampleindex}\index{type=@type!<num>|exampleindex}\index{value=@value!<num>|exampleindex}\exampleFont \begin{shaded}\noindent\mbox{}{<\textbf{num}\hspace*{1em}{value}="{33}">}xxxiii{</\textbf{num}>}\mbox{}\newline 
{<\textbf{num}\hspace*{1em}{type}="{cardinal}"\hspace*{1em}{value}="{21}">}twenty-one{</\textbf{num}>}\mbox{}\newline 
{<\textbf{num}\hspace*{1em}{type}="{percentage}"\hspace*{1em}{value}="{10}">}ten percent{</\textbf{num}>}\mbox{}\newline 
{<\textbf{num}\hspace*{1em}{type}="{percentage}"\hspace*{1em}{value}="{10}">}10\%{</\textbf{num}>}\mbox{}\newline 
{<\textbf{num}\hspace*{1em}{type}="{ordinal}"\hspace*{1em}{value}="{5}">}5th{</\textbf{num}>}\end{shaded}\egroup\par \noindent  \par\bgroup\index{num=<num>|exampleindex}\index{type=@type!<num>|exampleindex}\index{value=@value!<num>|exampleindex}\index{num=<num>|exampleindex}\index{type=@type!<num>|exampleindex}\index{value=@value!<num>|exampleindex}\exampleFont \begin{shaded}\noindent\mbox{}{<\textbf{num}\hspace*{1em}{type}="{fraction}"\hspace*{1em}{value}="{0.5}">}one half{</\textbf{num}>}\mbox{}\newline 
{<\textbf{num}\hspace*{1em}{type}="{fraction}"\hspace*{1em}{value}="{0.5}">}1/2{</\textbf{num}>}\end{shaded}\egroup\par \par
Sometimes it may be desired to mark something as numerical which cannot be accurately normalized, for example an expression such as \textit{dozens}; less frequently the number may be recognisable linguistically as such but may use a notation with which the encoder is unfamiliar. To help in these situations, the \hyperref[TEI.num]{<num>} element may also bear either or both of the following attributes from the \textsf{att.ranging} class: 
\begin{sansreflist}
  
\item [\textbf{att.ranging}] provides attributes for describing numerical ranges.\hfil\\[-10pt]\begin{sansreflist}
    \item[@{\itshape atLeast}]
  gives a minimum estimated value for the approximate measurement.
    \item[@{\itshape atMost}]
  gives a maximum estimated value for the approximate measurement.
\end{sansreflist}  
\end{sansreflist}
\par
In its fullest form, a \textit{measure} consists of a number, a phrase expressing units of measure and a phrase expressing the commodity being measured, though not all of these components need be present in every case. It may be helpful to distinguish measures from surrounding text for two reasons. Firstly, a measure may be expressed using a particular notation or system of abbreviations which the encoder does not wish to regard as lexical. Secondly, a quantitative application may wish to distinguish and normalize the internal components of a measure, in order to perform calculations on them.\par
Consider, as an example of the first case, the following list of Celia's charms, in which the encoder has chosen to make explicit the measurements: \par\bgroup\index{div=<div>|exampleindex}\index{n=@n!<div>|exampleindex}\index{list=<list>|exampleindex}\index{type=@type!<list>|exampleindex}\index{label=<label>|exampleindex}\index{item=<item>|exampleindex}\index{label=<label>|exampleindex}\index{item=<item>|exampleindex}\index{label=<label>|exampleindex}\index{item=<item>|exampleindex}\index{label=<label>|exampleindex}\index{item=<item>|exampleindex}\index{label=<label>|exampleindex}\index{item=<item>|exampleindex}\index{label=<label>|exampleindex}\index{item=<item>|exampleindex}\index{label=<label>|exampleindex}\index{item=<item>|exampleindex}\index{measure=<measure>|exampleindex}\index{label=<label>|exampleindex}\index{item=<item>|exampleindex}\index{measure=<measure>|exampleindex}\exampleFont \begin{shaded}\noindent\mbox{}{<\textbf{div}\hspace*{1em}{n}="{2}">}\mbox{}\newline 
\hspace*{1em}{<\textbf{list}\hspace*{1em}{type}="{gloss}">}\mbox{}\newline 
\hspace*{1em}\hspace*{1em}{<\textbf{label}>}Age{</\textbf{label}>}\mbox{}\newline 
\hspace*{1em}\hspace*{1em}{<\textbf{item}>}Unimportant{</\textbf{item}>}\mbox{}\newline 
\hspace*{1em}\hspace*{1em}{<\textbf{label}>}Head{</\textbf{label}>}\mbox{}\newline 
\hspace*{1em}\hspace*{1em}{<\textbf{item}>}Small and round{</\textbf{item}>}\mbox{}\newline 
\hspace*{1em}\hspace*{1em}{<\textbf{label}>}Eyes{</\textbf{label}>}\mbox{}\newline 
\hspace*{1em}\hspace*{1em}{<\textbf{item}>}Green{</\textbf{item}>}\mbox{}\newline 
\hspace*{1em}\hspace*{1em}{<\textbf{label}>}Complexion{</\textbf{label}>}\mbox{}\newline 
\hspace*{1em}\hspace*{1em}{<\textbf{item}>}White{</\textbf{item}>}\mbox{}\newline 
\hspace*{1em}\hspace*{1em}{<\textbf{label}>}Hair{</\textbf{label}>}\mbox{}\newline 
\hspace*{1em}\hspace*{1em}{<\textbf{item}>}yellow{</\textbf{item}>}\mbox{}\newline 
\hspace*{1em}\hspace*{1em}{<\textbf{label}>}Features{</\textbf{label}>}\mbox{}\newline 
\hspace*{1em}\hspace*{1em}{<\textbf{item}>}Mobile{</\textbf{item}>}\mbox{}\newline 
\hspace*{1em}\hspace*{1em}{<\textbf{label}>}Neck{</\textbf{label}>}\mbox{}\newline 
\hspace*{1em}\hspace*{1em}{<\textbf{item}>}\mbox{}\newline 
\hspace*{1em}\hspace*{1em}\hspace*{1em}{<\textbf{measure}>}13¾"{</\textbf{measure}>}\mbox{}\newline 
\hspace*{1em}\hspace*{1em}{</\textbf{item}>}\mbox{}\newline 
\hspace*{1em}\hspace*{1em}{<\textbf{label}>}Upper arm{</\textbf{label}>}\mbox{}\newline 
\hspace*{1em}\hspace*{1em}{<\textbf{item}>}\mbox{}\newline 
\hspace*{1em}\hspace*{1em}\hspace*{1em}{<\textbf{measure}>}11"{</\textbf{measure}>}\mbox{}\newline 
\hspace*{1em}\hspace*{1em}{</\textbf{item}>}\mbox{}\newline 
\textit{<!--...-->}\mbox{}\newline 
\hspace*{1em}{</\textbf{list}>}\mbox{}\newline 
\textit{<!-- ... -->}\mbox{}\newline 
{</\textbf{div}>}\end{shaded}\egroup\par \noindent  In the same way, it may be convenient to mark representations of currency which might otherwise be misinterpreted as lexical: \par\bgroup\index{p=<p>|exampleindex}\index{measure=<measure>|exampleindex}\index{type=@type!<measure>|exampleindex}\exampleFont \begin{shaded}\noindent\mbox{}{<\textbf{p}>}...the sum of\mbox{}\newline 
{<\textbf{measure}\hspace*{1em}{type}="{currency}">}12s 6d{</\textbf{measure}>}...{</\textbf{p}>}\end{shaded}\egroup\par \par
In general, normalization of a measure will require specification of one or more of its three parts: the quantity, the units, and possibly also the commodity being measured. This is accomplished by supplying values for the three attributes {\itshape quantity}, {\itshape unit}, and {\itshape commodity}, which are supplied by the \textsf{att.measurement} class: 
\begin{sansreflist}
  
\item [\textbf{att.measurement}] provides attributes to represent a regularized or normalized measurement.\hfil\\[-10pt]\begin{sansreflist}
    \item[@{\itshape quantity}]
  (quantity) specifies the number of the specified units that comprise the measurement
    \item[@{\itshape unit}]
  (unit) indicates the units used for the measurement, usually using the standard symbol for the desired units.
    \item[@{\itshape commodity}]
  (commodity) indicates the substance that is being measured
\end{sansreflist}  
\end{sansreflist}
 With these attributes, the measurement of Celia's neck may be specified in a normalized form: \par\bgroup\index{measure=<measure>|exampleindex}\index{quantity=@quantity!<measure>|exampleindex}\index{unit=@unit!<measure>|exampleindex}\exampleFont \begin{shaded}\noindent\mbox{}{<\textbf{measure}\hspace*{1em}{quantity}="{13.75}"\hspace*{1em}{unit}="{in}">}13¾"{</\textbf{measure}>}\end{shaded}\egroup\par \noindent  Such techniques are particularly useful when representing historical data such as inventories: \par\bgroup\index{list=<list>|exampleindex}\index{item=<item>|exampleindex}\index{measure=<measure>|exampleindex}\index{type=@type!<measure>|exampleindex}\index{quantity=@quantity!<measure>|exampleindex}\index{unit=@unit!<measure>|exampleindex}\index{commodity=@commodity!<measure>|exampleindex}\index{item=<item>|exampleindex}\index{measure=<measure>|exampleindex}\index{type=@type!<measure>|exampleindex}\index{quantity=@quantity!<measure>|exampleindex}\index{unit=@unit!<measure>|exampleindex}\index{commodity=@commodity!<measure>|exampleindex}\index{item=<item>|exampleindex}\index{measure=<measure>|exampleindex}\index{type=@type!<measure>|exampleindex}\index{quantity=@quantity!<measure>|exampleindex}\index{unit=@unit!<measure>|exampleindex}\index{commodity=@commodity!<measure>|exampleindex}\exampleFont \begin{shaded}\noindent\mbox{}{<\textbf{list}>}\mbox{}\newline 
\hspace*{1em}{<\textbf{item}>}\mbox{}\newline 
\hspace*{1em}\hspace*{1em}{<\textbf{measure}\hspace*{1em}{type}="{volume}"\hspace*{1em}{quantity}="{2}"\mbox{}\newline 
\hspace*{1em}\hspace*{1em}\hspace*{1em}{unit}="{bag}"\hspace*{1em}{commodity}="{hops}">}ii bags hops{</\textbf{measure}>}\mbox{}\newline 
\hspace*{1em}{</\textbf{item}>}\mbox{}\newline 
\hspace*{1em}{<\textbf{item}>}\mbox{}\newline 
\hspace*{1em}\hspace*{1em}{<\textbf{measure}\hspace*{1em}{type}="{volume}"\hspace*{1em}{quantity}="{6}"\mbox{}\newline 
\hspace*{1em}\hspace*{1em}\hspace*{1em}{unit}="{truss}"\hspace*{1em}{commodity}="{cloth}">}six trusses Woolen and linen goods{</\textbf{measure}>}\mbox{}\newline 
\hspace*{1em}{</\textbf{item}>}\mbox{}\newline 
\hspace*{1em}{<\textbf{item}>}\mbox{}\newline 
\hspace*{1em}\hspace*{1em}{<\textbf{measure}\hspace*{1em}{type}="{weight}"\hspace*{1em}{quantity}="{5}"\mbox{}\newline 
\hspace*{1em}\hspace*{1em}\hspace*{1em}{unit}="{ton}"\hspace*{1em}{commodity}="{coal}">}5 tonnes coale{</\textbf{measure}>}\mbox{}\newline 
\hspace*{1em}{</\textbf{item}>}\mbox{}\newline 
{</\textbf{list}>}\end{shaded}\egroup\par \par
The \hyperref[TEI.measureGrp]{<measureGrp>} element is provided as a means of grouping several related measurements together, either because the measurement involves several dimensions (for example height and width) or to avoid the need to repeat all the normalizing attributes: \par\bgroup\index{measureGrp=<measureGrp>|exampleindex}\index{type=@type!<measureGrp>|exampleindex}\index{unit=@unit!<measureGrp>|exampleindex}\index{measure=<measure>|exampleindex}\index{type=@type!<measure>|exampleindex}\index{quantity=@quantity!<measure>|exampleindex}\index{measure=<measure>|exampleindex}\index{type=@type!<measure>|exampleindex}\index{quantity=@quantity!<measure>|exampleindex}\index{measure=<measure>|exampleindex}\index{type=@type!<measure>|exampleindex}\index{quantity=@quantity!<measure>|exampleindex}\exampleFont \begin{shaded}\noindent\mbox{}{<\textbf{measureGrp}\hspace*{1em}{type}="{volume}"\hspace*{1em}{unit}="{in}">}\mbox{}\newline 
\hspace*{1em}{<\textbf{measure}\hspace*{1em}{type}="{height}"\hspace*{1em}{quantity}="{14}">}xiv{</\textbf{measure}>}\mbox{}\newline 
\hspace*{1em}{<\textbf{measure}\hspace*{1em}{type}="{width}"\hspace*{1em}{quantity}="{5}">}v{</\textbf{measure}>}\mbox{}\newline 
\hspace*{1em}{<\textbf{measure}\hspace*{1em}{type}="{depth}"\hspace*{1em}{quantity}="{10}">}x{</\textbf{measure}>}\mbox{}\newline 
{</\textbf{measureGrp}>}\end{shaded}\egroup\par \noindent  \par
The \hyperref[TEI.unit]{<unit>} element may be applied when units of measurement require more detailed encoding about the system they belong to or the what kind of phenomenon they designate. The \hyperref[TEI.unit]{<unit>} element may carry the {\itshape unit} attribute to indicate a standard value, as well as other optional attributes for indicating type, language, and other distinguishing characteristics. 
\begin{sansreflist}
  
\item [\textbf{<unit>}] contains a symbol, a word or a phrase referring to a unit of measurement in any kind of formal or informal system.
\item [\textbf{att.measurement}] provides attributes to represent a regularized or normalized measurement.\hfil\\[-10pt]\begin{sansreflist}
    \item[@{\itshape unit}]
  (unit) indicates the units used for the measurement, usually using the standard symbol for the desired units.
\end{sansreflist}  
\item [\textbf{att.typed}] provides attributes which can be used to classify or subclassify elements in any way.
\item [\textbf{att.global}] provides attributes common to all elements in the TEI encoding scheme.
\end{sansreflist}
 A \hyperref[TEI.measure]{<measure>} element may contain a combination of \hyperref[TEI.num]{<num>} and \hyperref[TEI.unit]{<unit>} elements: \par\bgroup\index{measure=<measure>|exampleindex}\index{type=@type!<measure>|exampleindex}\index{num=<num>|exampleindex}\index{num=<num>|exampleindex}\index{num=<num>|exampleindex}\index{num=<num>|exampleindex}\index{unit=<unit>|exampleindex}\index{type=@type!<unit>|exampleindex}\index{unit=@unit!<unit>|exampleindex}\exampleFont \begin{shaded}\noindent\mbox{}{<\textbf{measure}\hspace*{1em}{type}="{list}">}\mbox{}\newline 
\hspace*{1em}{<\textbf{num}>}1{</\textbf{num}>}, {<\textbf{num}>}2{</\textbf{num}>}, {<\textbf{num}>}5{</\textbf{num}>}, {<\textbf{num}>}7{</\textbf{num}>}\mbox{}\newline 
\hspace*{1em}{<\textbf{unit}\hspace*{1em}{type}="{length}"\hspace*{1em}{unit}="{mm}">}millimètres{</\textbf{unit}>}\mbox{}\newline 
{</\textbf{measure}>}\end{shaded}\egroup\par \noindent  The unit element may also be nested to indicate a complex unit and its component parts, for example, to indicate that rate combines space and time: \par\bgroup\index{p=<p>|exampleindex}\index{num=<num>|exampleindex}\index{value=@value!<num>|exampleindex}\index{unit=<unit>|exampleindex}\index{type=@type!<unit>|exampleindex}\index{unit=@unit!<unit>|exampleindex}\index{unit=<unit>|exampleindex}\index{type=@type!<unit>|exampleindex}\index{unit=<unit>|exampleindex}\index{type=@type!<unit>|exampleindex}\exampleFont \begin{shaded}\noindent\mbox{}{<\textbf{p}>}Light travels at {<\textbf{num}\hspace*{1em}{value}="{3E10}">}3×10\textasciicircum 10{</\textbf{num}>}\mbox{}\newline 
\hspace*{1em}{<\textbf{unit}\hspace*{1em}{type}="{rate}"\hspace*{1em}{unit}="{cm/s}">}\mbox{}\newline 
\hspace*{1em}\hspace*{1em}{<\textbf{unit}\hspace*{1em}{type}="{space}">}cm{</\textbf{unit}>} per {<\textbf{unit}\hspace*{1em}{type}="{time}">}second{</\textbf{unit}>}\mbox{}\newline 
\hspace*{1em}{</\textbf{unit}>}.{</\textbf{p}>}\end{shaded}\egroup\par 
\subsubsection[{Dates and Times}]{Dates and Times}\label{CONADA}\par
Dates and times, like numbers, can appear in widely varying culture- and language-dependent forms, and can pose similar problems in automatic language processing. Such elements constitute the \textsf{model.dateLike} class, of which the default members are: 
\begin{sansreflist}
  
\item [\textbf{<date>}] (date) contains a date in any format.
\item [\textbf{<time>}] (time) contains a phrase defining a time of day in any format.
\end{sansreflist}
 These elements have some additional attributes by virtue of being members of the \textsf{att.datable} and \textsf{att.duration} classes which, in turn, are members of the \textsf{att.datable.w3c} and \textsf{att.duration.w3c} classes. In particular, the {\itshape when} and {\itshape calendar} attributes will be discussed here: 
\begin{sansreflist}
  
\item [\textbf{att.datable.w3c}] provides attributes for normalization of elements that contain datable events conforming to the W3C \textit{XML Schema Part 2: Datatypes Second Edition}.\hfil\\[-10pt]\begin{sansreflist}
    \item[@{\itshape when}]
  supplies the value of the date or time in a standard form, e.g. yyyy-mm-dd.
\end{sansreflist}  
\item [\textbf{att.datable}] provides attributes for normalization of elements that contain dates, times, or datable events.\hfil\\[-10pt]\begin{sansreflist}
    \item[@{\itshape calendar}]
  indicates the system or calendar to which the date represented by the content of this element belongs.
\end{sansreflist}  
\end{sansreflist}
\par
Dates can occur virtually anywhere in a text, but in some contexts (e.g. bibliographic citations) their encoding is recommended or required rather than optional. Times can also appear anywhere but encoding these is more generally optional. \par
Partial dates or times (e.g. \textit{1990}, \textit{September 1990}, \textit{twelvish}) can be expressed in the {\itshape when} attribute by simply omitting a part of the value supplied. Imprecise dates or times (for example \textit{early August}, \textit{some time after ten and before twelve}) may be expressed as date or time ranges.\par
These mechanisms are useful primarily for fully specified dates or times known with certainty. If component parts of dates or times are to be marked up, or if a more complex analysis of the meaning of a temporal expression is required, the techniques described in chapter \textit{\hyperref[ND]{13.\ Names, Dates, People, and Places}} should be used in preference to the simple method outlined here.\par
Where the certainty (i.e. reliability) of the date or time is in question, the encoder should record this fact using the mechanisms discussed in chapter \textit{\hyperref[CE]{21.\ Certainty, Precision, and Responsibility}}. The same chapter also discusses various methods of recording the precision of numerical or temporal assertions.\par
The {\itshape when} attribute is a useful way of normalizing or disambiguating dates and times which can appear in many formats, as the following examples show: \par\bgroup\index{date=<date>|exampleindex}\index{when=@when!<date>|exampleindex}\exampleFont \begin{shaded}\noindent\mbox{}{<\textbf{date}\hspace*{1em}{when}="{1980-02-12}">}12/2/1980{</\textbf{date}>}\end{shaded}\egroup\par \noindent  \par\bgroup\index{date=<date>|exampleindex}\index{when=@when!<date>|exampleindex}\exampleFont \begin{shaded}\noindent\mbox{}Given on the {<\textbf{date}\hspace*{1em}{when}="{1977-06-12}">}Twelfth Day of June\mbox{}\newline 
 in the Year of Our Lord One Thousand Nine Hundred and\mbox{}\newline 
 Seventy-seven of the Republic the Two Hundredth and first\mbox{}\newline 
 and of the University the Eighty-Sixth.{</\textbf{date}>}\end{shaded}\egroup\par \par
The {\itshape when} attribute always supplies a normalized representation of the date given as content of the \hyperref[TEI.date]{<date>} element. The format used should be a valid W3C schema datatype.\footnote{The datatypes are taken from the W3C Recommendation \textit{XML Schema Part 2: Datatypes Second Edition}. The permitted datatypes are: \begin{itemize}
\item \xref{http://www.w3.org/TR/xmlschema-2/\#date}{date}
\item \xref{http://www.w3.org/TR/xmlschema-2/\#gYear}{gYear}
\item \xref{http://www.w3.org/TR/xmlschema-2/\#gMonth}{gMonth}
\item \xref{http://www.w3.org/TR/xmlschema-2/\#gDay}{gDay}
\item \xref{http://www.w3.org/TR/xmlschema-2/\#gYearMonth}{gYearMonth}
\item \xref{http://www.w3.org/TR/xmlschema-2/\#gMonthDay}{gMonthDay}
\item \xref{http://www.w3.org/TR/xmlschema-2/\#time}{time}
\item \xref{http://www.w3.org/TR/xmlschema-2/\#dateTime}{dateTime}
\end{itemize}  There is one exception: these Guidelines permit a time to be expressed as only a number of hours, or as a number of hours and minutes, as per ISO 8601:2004 section 4.2.2.3 and 4.3.3. The W3C \textsf{time} and \textsf{dateTime} datatypes require that the minutes and seconds be included in the normalized value if they are to be correctly processed for example when sorting.} Some typical examples follow: \par\bgroup\index{date=<date>|exampleindex}\index{when=@when!<date>|exampleindex}\index{date=<date>|exampleindex}\index{when=@when!<date>|exampleindex}\index{date=<date>|exampleindex}\index{when=@when!<date>|exampleindex}\index{date=<date>|exampleindex}\index{when=@when!<date>|exampleindex}\index{date=<date>|exampleindex}\index{when=@when!<date>|exampleindex}\index{date=<date>|exampleindex}\index{when=@when!<date>|exampleindex}\index{time=<time>|exampleindex}\index{when=@when!<time>|exampleindex}\index{date=<date>|exampleindex}\index{when=@when!<date>|exampleindex}\exampleFont \begin{shaded}\noindent\mbox{}{<\textbf{date}\hspace*{1em}{when}="{2001}">}The year 2001{</\textbf{date}>}\mbox{}\newline 
{<\textbf{date}\hspace*{1em}{when}="{2001-09}">}September 2001{</\textbf{date}>}\mbox{}\newline 
{<\textbf{date}\hspace*{1em}{when}="{2001-09-11}">}11 Sep 01{</\textbf{date}>}\mbox{}\newline 
{<\textbf{date}\hspace*{1em}{when}="{--09-11}">}9/11{</\textbf{date}>}\mbox{}\newline 
{<\textbf{date}\hspace*{1em}{when}="{--09}">}September{</\textbf{date}>}\mbox{}\newline 
{<\textbf{date}\hspace*{1em}{when}="{---11}">}Eleventh of the month{</\textbf{date}>}\mbox{}\newline 
{<\textbf{time}\hspace*{1em}{when}="{08:48:00}">}8:48{</\textbf{time}>}\mbox{}\newline 
{<\textbf{date}\hspace*{1em}{when}="{2001-09-11T12:48:00}">}Sept 11th, 12 minutes before 9 am{</\textbf{date}>}\end{shaded}\egroup\par \noindent Note in the last example the use of a normalized representation for the date string which includes a time: this example could thus equally well be tagged using the \hyperref[TEI.time]{<time>} element.\par
The following examples demonstrate the use of the \hyperref[TEI.date]{<date>} element to mark a period of time:\par\bgroup\index{p=<p>|exampleindex}\index{date=<date>|exampleindex}\index{from=@from!<date>|exampleindex}\index{to=@to!<date>|exampleindex}\exampleFont \begin{shaded}\noindent\mbox{}{<\textbf{p}>}Those five years —\mbox{}\newline 
{<\textbf{date}\hspace*{1em}{from}="{1918}"\hspace*{1em}{to}="{1923}">}1918 to 1923{</\textbf{date}>}\mbox{}\newline 
 — had been, he suspected,\mbox{}\newline 
 somehow very important.{</\textbf{p}>}\end{shaded}\egroup\par \noindent  \par\bgroup\index{p=<p>|exampleindex}\index{date=<date>|exampleindex}\index{notBefore=@notBefore!<date>|exampleindex}\index{notAfter=@notAfter!<date>|exampleindex}\index{title=<title>|exampleindex}\index{date=<date>|exampleindex}\index{when=@when!<date>|exampleindex}\exampleFont \begin{shaded}\noindent\mbox{}{<\textbf{p}>}The Eddic poems are preserved in a unique\mbox{}\newline 
 manuscript (Codex Regius 2365) from\mbox{}\newline 
{<\textbf{date}\hspace*{1em}{notBefore}="{1250}"\hspace*{1em}{notAfter}="{1300}">}the second half of the\mbox{}\newline 
\hspace*{1em}\hspace*{1em} thirteenth century{</\textbf{date}>}, and {<\textbf{title}>}Hervarar\mbox{}\newline 
\hspace*{1em}\hspace*{1em} saga{</\textbf{title}>} dates from {<\textbf{date}\hspace*{1em}{when}="{1300}">}around 1300{</\textbf{date}>}.{</\textbf{p}>}\end{shaded}\egroup\par \par
The {\itshape calendar} attribute may be used to specify a date in any calendar system; if the {\itshape when} attribute is also supplied, it should specify the equivalent date in the Gregorian calendar.
\subsubsection[{Abbreviations and Their Expansions}]{Abbreviations and Their Expansions}\label{CONAAB}\par
It is sometimes desirable to mark abbreviations in the copy text, whether to trigger special processing for them, to provide the full form of the word or phrase abbreviated, or to allow for different possible expansions of the abbreviation. Abbreviations may be transcribed as they stand, or expanded; they may be left unmarked, or marked using these tags: 
\begin{sansreflist}
  
\item [\textbf{<abbr>}] (abbreviation) contains an abbreviation of any sort.
\item [\textbf{<expan>}] (expansion) contains the expansion of an abbreviation.
\end{sansreflist}
\par
The \hyperref[TEI.abbr]{<abbr>} element is useful as a means of distinguishing semi-lexical items such as acronyms or jargon: \par\bgroup\index{abbr=<abbr>|exampleindex}\index{abbr=<abbr>|exampleindex}\exampleFont \begin{shaded}\noindent\mbox{}We can sum up the above discussion as follows:\mbox{}\newline 
 the identity of a {<\textbf{abbr}>}CC{</\textbf{abbr}>} is defined by that calibration of values which\mbox{}\newline 
 motivates the elements of its {<\textbf{abbr}>}GSP{</\textbf{abbr}>}; ...\end{shaded}\egroup\par \noindent  \par\bgroup\index{abbr=<abbr>|exampleindex}\index{abbr=<abbr>|exampleindex}\index{abbr=<abbr>|exampleindex}\exampleFont \begin{shaded}\noindent\mbox{}Every manufacturer of {<\textbf{abbr}>}3GL{</\textbf{abbr}>} or {<\textbf{abbr}>}4GL{</\textbf{abbr}>}\mbox{}\newline 
 languages is currently nailing on {<\textbf{abbr}>}OOP{</\textbf{abbr}>} extensions.\end{shaded}\egroup\par \par
The {\itshape type} attribute may be used to distinguish types of abbreviation by their function:\par\bgroup\index{abbr=<abbr>|exampleindex}\index{type=@type!<abbr>|exampleindex}\index{abbr=<abbr>|exampleindex}\index{type=@type!<abbr>|exampleindex}\index{abbr=<abbr>|exampleindex}\index{type=@type!<abbr>|exampleindex}\exampleFont \begin{shaded}\noindent\mbox{}{<\textbf{abbr}\hspace*{1em}{type}="{title}">}Dr.{</\textbf{abbr}>}\mbox{}\newline 
{<\textbf{abbr}\hspace*{1em}{type}="{initial}">}M.{</\textbf{abbr}>} Deegan is\mbox{}\newline 
 the Director of the {<\textbf{abbr}\hspace*{1em}{type}="{acronym}">}CTI{</\textbf{abbr}>} Centre for Textual Studies.\end{shaded}\egroup\par \par
Abbreviations such as \textit{Dr. M.} above may be treated as two abbreviations, as above, or as one: \par\bgroup\index{abbr=<abbr>|exampleindex}\index{abbr=<abbr>|exampleindex}\exampleFont \begin{shaded}\noindent\mbox{}{<\textbf{abbr}>}Dr. M.{</\textbf{abbr}>} Deegan is\mbox{}\newline 
 the Director of the {<\textbf{abbr}>}CTI{</\textbf{abbr}>} Centre for Textual Studies.\end{shaded}\egroup\par \par
The \hyperref[TEI.expan]{<expan>} element may be used simply to record that an abbreviation has been silently expanded by the encoder, perhaps for reasons of house style or editorial policy. It should always include the whole of an abbreviated phrase or word. More usually however this will be combined with the \hyperref[TEI.abbr]{<abbr>} element inside a \hyperref[TEI.choice]{<choice>} element to record both the abbreviation and its expansion: \par\bgroup\index{choice=<choice>|exampleindex}\index{expan=<expan>|exampleindex}\index{abbr=<abbr>|exampleindex}\exampleFont \begin{shaded}\noindent\mbox{} the\mbox{}\newline 
{<\textbf{choice}>}\mbox{}\newline 
\hspace*{1em}{<\textbf{expan}>}World Wide Web Consortium{</\textbf{expan}>}\mbox{}\newline 
\hspace*{1em}{<\textbf{abbr}>}W3C{</\textbf{abbr}>}\mbox{}\newline 
{</\textbf{choice}>}\end{shaded}\egroup\par \noindent  Nested abbreviations may also be handled in this way: \par\bgroup\index{choice=<choice>|exampleindex}\index{abbr=<abbr>|exampleindex}\index{expan=<expan>|exampleindex}\index{choice=<choice>|exampleindex}\index{abbr=<abbr>|exampleindex}\index{expan=<expan>|exampleindex}\exampleFont \begin{shaded}\noindent\mbox{}{<\textbf{choice}>}\mbox{}\newline 
\hspace*{1em}{<\textbf{abbr}>}RELAX NG{</\textbf{abbr}>}\mbox{}\newline 
\hspace*{1em}{<\textbf{expan}>}regular\mbox{}\newline 
\hspace*{1em}\hspace*{1em} language for {<\textbf{choice}>}\mbox{}\newline 
\hspace*{1em}\hspace*{1em}\hspace*{1em}{<\textbf{abbr}>}XML{</\textbf{abbr}>}\mbox{}\newline 
\hspace*{1em}\hspace*{1em}\hspace*{1em}{<\textbf{expan}>}extensible markup\mbox{}\newline 
\hspace*{1em}\hspace*{1em}\hspace*{1em}\hspace*{1em}\hspace*{1em}\hspace*{1em} language{</\textbf{expan}>}\mbox{}\newline 
\hspace*{1em}\hspace*{1em}{</\textbf{choice}>}, next\mbox{}\newline 
\hspace*{1em}\hspace*{1em} generation{</\textbf{expan}>}\mbox{}\newline 
{</\textbf{choice}>}\end{shaded}\egroup\par \par
Abbreviation is a particularly important feature of manuscript and other source materials, the transcription of which needs more detailed treatment than is possible using these simple elements. A more detailed set of recommendations is discussed in \textit{\hyperref[PHCH]{11.3.1.\ Altered, Corrected, and Erroneous Texts}}, which includes additional elements made available for the purpose by the \textsf{transcr} module.
\subsection[{Simple Links and Cross-References}]{Simple Links and Cross-References}\label{COXR}\par
Cross-references or links between one location in a document and one or more other locations, either in the same or different XML documents, may be encoded using the elements \hyperref[TEI.ptr]{<ptr>} and \hyperref[TEI.ref]{<ref>}, as discussed in this section. These elements both ‘point’ from one location in a document, the place that the element itself appears, to another (or to several), specified by means of a {\itshape target} attribute, supplied by the \textsf{att.pointing} class: 
\begin{sansreflist}
  
\item [\textbf{att.pointing}] provides a set of attributes used by all elements which point to other elements by means of one or more URI references.\hfil\\[-10pt]\begin{sansreflist}
    \item[@{\itshape target}]
  specifies the destination of the reference by supplying one or more URI References
\end{sansreflist}  
\end{sansreflist}
 Linkages of several other kinds are also provided for in these guidelines; see further chapter \textit{\hyperref[SA]{16.\ Linking, Segmentation, and Alignment}}.\par
The value of the {\itshape target} attribute, wherever it appears, provides a way of pointing to some other element using a method standardized by the W3C consortium, and known as the \textit{XPointer} mechanism. This permits a range of complexity, from the very simple (a reference to the value of the target element's {\itshape xml:id} attribute) to the more complex usage of a full URI with embedded XPointers. For example, the source of the following paragraph looks something like this: \par\bgroup\index{p=<p>|exampleindex}\index{note=<note>|exampleindex}\index{place=@place!<note>|exampleindex}\index{ptr=<ptr>|exampleindex}\index{target=@target!<ptr>|exampleindex}\index{ptr=<ptr>|exampleindex}\index{target=@target!<ptr>|exampleindex}\index{ptr=<ptr>|exampleindex}\index{target=@target!<ptr>|exampleindex}\index{ptr=<ptr>|exampleindex}\index{target=@target!<ptr>|exampleindex}\index{ptr=<ptr>|exampleindex}\index{target=@target!<ptr>|exampleindex}\index{div=<div>|exampleindex}\exampleFont \begin{shaded}\noindent\mbox{}{<\textbf{p}>}...\mbox{}\newline 
 The complete XPointer specification is managed by the W3C{<\textbf{note}\hspace*{1em}{place}="{foot}">}\mbox{}\newline 
\hspace*{1em}\hspace*{1em}{<\textbf{ptr}\hspace*{1em}{target}="{http://www.w3.org/TR/xptr-framework/}"/>},\mbox{}\newline 
\hspace*{1em}{<\textbf{ptr}\hspace*{1em}{target}="{http://www.w3.org/TR/xptr-element/}"/>},\mbox{}\newline 
\hspace*{1em}{<\textbf{ptr}\hspace*{1em}{target}="{http://www.w3.org/TR/xptr-xmlns/}"/>}, and\mbox{}\newline 
\hspace*{1em}{<\textbf{ptr}\hspace*{1em}{target}="{http://www.w3.org/TR/xptr-xpointer/}"/>}\mbox{}\newline 
\hspace*{1em}{</\textbf{note}>};\mbox{}\newline 
 for a discussion of TEI schemes for XPointer, see\mbox{}\newline 
{<\textbf{ptr}\hspace*{1em}{target}="{\#eSATS}"/>}.{</\textbf{p}>}\mbox{}\newline 
\textit{<!--... -->}\mbox{}\newline 
{<\textbf{div}\hspace*{1em}{xml:id}="{eSATS}">}\mbox{}\newline 
\textit{<!--... -->}\mbox{}\newline 
{</\textbf{div}>}\end{shaded}\egroup\par \noindent  Alternatively, if no explicit link is to be encoded, but it is simply required to mark the phrase as a cross-reference, the \hyperref[TEI.ref]{<ref>} element may be used without a {\itshape target} attribute.\par
For an introduction to the use of links in general, see \textit{\hyperref[SA]{16.\ Linking, Segmentation, and Alignment}}. The complete XPointer specification is managed by the W3C\footnote{\url{http://www.w3.org/TR/xptr-framework/}, \url{http://www.w3.org/TR/xptr-element/}, \url{http://www.w3.org/TR/xptr-xmlns/}, and \url{http://www.w3.org/TR/xptr-xpointer/}}; for a discussion of TEI schemes for XPointer, see \textit{\hyperref[SATS]{16.2.4.\ TEI XPointer Schemes}}.\par

\begin{sansreflist}
  
\item [\textbf{<ptr>}] (pointer) defines a pointer to another location.
\item [\textbf{<ref>}] (reference) defines a reference to another location, possibly modified by additional text or comment.
\end{sansreflist}
\par
The elements \hyperref[TEI.ptr]{<ptr>} and \hyperref[TEI.ref]{<ref>} are the default members of the phrase-level model class \textsf{model.ptrLike}. As members of the classes \textsf{att.pointing}, \textsf{att.typed}, \textsf{att.cReferencing}, and \textsf{att.internetMedia} they also carry the following attributes: 
\begin{sansreflist}
  
\item [\textbf{att.pointing}] provides a set of attributes used by all elements which point to other elements by means of one or more URI references.\hfil\\[-10pt]\begin{sansreflist}
    \item[@{\itshape target}]
  specifies the destination of the reference by supplying one or more URI References
    \item[@{\itshape evaluate}]
  (evaluate) specifies the intended meaning when the target of a pointer is itself a pointer.
\end{sansreflist}  
\item [\textbf{att.cReferencing}] provides an attribute which may be used to supply a \textit{canonical reference} as a means of identifying the target of a pointer.\hfil\\[-10pt]\begin{sansreflist}
    \item[@{\itshape cRef}]
  (canonical reference) specifies the destination of the pointer by supplying a canonical reference expressed using the scheme defined in a \hyperref[TEI.refsDecl]{<refsDecl>} element in the TEI header
\end{sansreflist}  
\item [\textbf{att.typed}] provides attributes which can be used to classify or subclassify elements in any way.\hfil\\[-10pt]\begin{sansreflist}
    \item[@{\itshape type}]
  characterizes the element in some sense, using any convenient classification scheme or typology.
    \item[@{\itshape subtype}]
  (subtype) provides a sub-categorization of the element, if needed
\end{sansreflist}  
\item [\textbf{att.internetMedia}] provides attributes for specifying the type of a computer resource using a standard taxonomy.\hfil\\[-10pt]\begin{sansreflist}
    \item[@{\itshape mimeType}]
  (MIME media type) specifies the applicable multimedia internet mail extension (MIME) media type
\end{sansreflist}  
\end{sansreflist}
\par
The two elements may be used in the same way; the difference between them is simply that while the \hyperref[TEI.ptr]{<ptr>} element is empty, the \hyperref[TEI.ref]{<ref>} element may contain phrases specifying, or describing more exactly, the target of a cross-reference, which form the content of the element. Since its content thus serves as a human-readable pointer, in the simplest case a \hyperref[TEI.ref]{<ref>} element need not identify its target in any other way. For example: \par\bgroup\index{ref=<ref>|exampleindex}\exampleFont \begin{shaded}\noindent\mbox{}See {<\textbf{ref}>}section 12 on page 34{</\textbf{ref}>}.\end{shaded}\egroup\par \par
More usually, it will be desirable to identify the target of the cross-reference using either the {\itshape target} or the {\itshape cRef} attribute, so that processing software can access it directly, for example to implement a linkage, to generate an appropriate reference, or to give an error message if it cannot be found. Assuming that section 12 in the previous example has been tagged \par\bgroup\index{div1=<div1>|exampleindex}\exampleFont \begin{shaded}\noindent\mbox{}{<\textbf{div1}\hspace*{1em}{xml:id}="{SEC12}">}\mbox{}\newline 
\textit{<!-- ... -->}\mbox{}\newline 
{</\textbf{div1}>}\end{shaded}\egroup\par \noindent  then the same cross-reference might more exactly be encoded as \par\bgroup\index{ref=<ref>|exampleindex}\index{target=@target!<ref>|exampleindex}\exampleFont \begin{shaded}\noindent\mbox{}See especially {<\textbf{ref}\hspace*{1em}{target}="{\#SEC12}">}section 12 on page 34{</\textbf{ref}>}.\end{shaded}\egroup\par \par
If the cross-reference itself is to be generated according to a fixed pattern, or if no text is to appear in the body of the cross-reference, the \hyperref[TEI.ptr]{<ptr>} element would be used as follows: \par\bgroup\index{ptr=<ptr>|exampleindex}\index{target=@target!<ptr>|exampleindex}\exampleFont \begin{shaded}\noindent\mbox{}See in particular {<\textbf{ptr}\hspace*{1em}{target}="{\#SEC12}"/>}.\end{shaded}\egroup\par \par
The {\itshape cRef} attribute may be used to express the target of a cross reference using some canonical referencing scheme, such as those typically used for ancient texts. In this case, the referencing scheme must be defined using the \hyperref[TEI.cRefPattern]{<cRefPattern>} or \hyperref[TEI.citeStructure]{<citeStructure>} elements discussed below (\textit{\hyperref[CORS6]{3.11.4.\ Declaring Reference Systems}}); the definition these provide may be used to translate the value of the {\itshape cRef} attribute into a conventional pointer value, such as one that might be supplied by the {\itshape target} attribute. It is an error to supply both {\itshape cRef} and {\itshape target} values.\par
When the {\itshape target} attribute is used, a cross reference may point to any number of locations simultaneously, simply by giving more than one identifier as the value of its {\itshape target} attribute. This may be particularly useful where an analytic index is to be encoded, as in the following example: \par\bgroup\index{list=<list>|exampleindex}\index{item=<item>|exampleindex}\index{ptr=<ptr>|exampleindex}\index{target=@target!<ptr>|exampleindex}\index{item=<item>|exampleindex}\index{ptr=<ptr>|exampleindex}\index{target=@target!<ptr>|exampleindex}\index{item=<item>|exampleindex}\index{ptr=<ptr>|exampleindex}\index{target=@target!<ptr>|exampleindex}\index{item=<item>|exampleindex}\index{ptr=<ptr>|exampleindex}\index{target=@target!<ptr>|exampleindex}\exampleFont \begin{shaded}\noindent\mbox{}{<\textbf{list}>}\mbox{}\newline 
\hspace*{1em}{<\textbf{item}>}Saints aid rejected in mel. {<\textbf{ptr}\hspace*{1em}{target}="{\#p299}"/>}\mbox{}\newline 
\hspace*{1em}{</\textbf{item}>}\mbox{}\newline 
\hspace*{1em}{<\textbf{item}>}Sallets censured {<\textbf{ptr}\hspace*{1em}{target}="{\#p143 \#p144}"/>}\mbox{}\newline 
\hspace*{1em}{</\textbf{item}>}\mbox{}\newline 
\hspace*{1em}{<\textbf{item}>}Sanguine mel. signs {<\textbf{ptr}\hspace*{1em}{target}="{\#p263}"/>}\mbox{}\newline 
\hspace*{1em}{</\textbf{item}>}\mbox{}\newline 
\hspace*{1em}{<\textbf{item}>}Scilla or sea onyon, a purger of mel. {<\textbf{ptr}\hspace*{1em}{target}="{\#p442}"/>}\mbox{}\newline 
\hspace*{1em}{</\textbf{item}>}\mbox{}\newline 
{</\textbf{list}>}\end{shaded}\egroup\par \noindent  Here the targets of the cross-references are simply page numbers; it is assumed that corresponding elements with identifiers \textsf{p299}, \textsf{p143}, etc. have been provided in the body of the text, for example as page breaks \par\bgroup\index{pb=<pb>|exampleindex}\index{pb=<pb>|exampleindex}\index{pb=<pb>|exampleindex}\index{pb=<pb>|exampleindex}\index{pb=<pb>|exampleindex}\exampleFont \begin{shaded}\noindent\mbox{}{<\textbf{pb}\hspace*{1em}{xml:id}="{p143}"/>}\mbox{}\newline 
 ...\mbox{}\newline 
{<\textbf{pb}\hspace*{1em}{xml:id}="{p144}"/>}\mbox{}\newline 
 ...\mbox{}\newline 
{<\textbf{pb}\hspace*{1em}{xml:id}="{p263}"/>}\mbox{}\newline 
 ...\mbox{}\newline 
{<\textbf{pb}\hspace*{1em}{xml:id}="{p299}"/>}\mbox{}\newline 
 ...\mbox{}\newline 
{<\textbf{pb}\hspace*{1em}{xml:id}="{p442}"/>}\mbox{}\newline 
 ...\mbox{}\newline 
\end{shaded}\egroup\par \par
A similar method may be used to link annotations on a text with the sigla used to encode their points of attachment in a text. For example: \par\bgroup\index{ref=<ref>|exampleindex}\index{target=@target!<ref>|exampleindex}\index{type=@type!<ref>|exampleindex}\index{note=<note>|exampleindex}\index{type=@type!<note>|exampleindex}\exampleFont \begin{shaded}\noindent\mbox{}annotated text {<\textbf{ref}\hspace*{1em}{target}="{\#a51}"\hspace*{1em}{type}="{noteAnchor}">}⁵¹{</\textbf{ref}>}\mbox{}\newline 
\textit{<!-- ... -->}\mbox{}\newline 
{<\textbf{note}\hspace*{1em}{xml:id}="{a51}"\hspace*{1em}{type}="{footnote}">}text of annotation{</\textbf{note}>}\end{shaded}\egroup\par \par
The {\itshape type} attribute may be used, as elsewhere, to categorize the cross-reference according to any system of importance to the encoder. If bibliographic references require special processing (e.g. in order to provide a consistent short-form reference), they might be tagged thus: \par\bgroup\index{term=<term>|exampleindex}\index{rend=@rend!<term>|exampleindex}\index{ptr=<ptr>|exampleindex}\index{target=@target!<ptr>|exampleindex}\index{ptr=<ptr>|exampleindex}\index{type=@type!<ptr>|exampleindex}\index{target=@target!<ptr>|exampleindex}\index{figure=<figure>|exampleindex}\index{graphic=<graphic>|exampleindex}\index{url=@url!<graphic>|exampleindex}\index{bibl=<bibl>|exampleindex}\exampleFont \begin{shaded}\noindent\mbox{}Similar forms, often called\mbox{}\newline 
{<\textbf{term}\hspace*{1em}{rend}="{ldquo rdquo}">}rewriting systems{</\textbf{term}>}, have a long history\mbox{}\newline 
 among mathematicians, but the specific form of {<\textbf{ptr}\hspace*{1em}{target}="{\#fig22}"/>}\mbox{}\newline 
 was first studied extensively by Chomsky {<\textbf{ptr}\hspace*{1em}{type}="{bibliog}"\hspace*{1em}{target}="{\#chom59}"/>}.\mbox{}\newline 
\mbox{}\newline 
\textit{<!-- ... -->}\mbox{}\newline 
{<\textbf{figure}\hspace*{1em}{xml:id}="{fig22}">}\mbox{}\newline 
\hspace*{1em}{<\textbf{graphic}\hspace*{1em}{url}="{fig22.jpg}"/>}\mbox{}\newline 
{</\textbf{figure}>}\mbox{}\newline 
\textit{<!-- elsewhere, in the bibliography -->}\mbox{}\newline 
{<\textbf{bibl}\hspace*{1em}{xml:id}="{chom59}">}\mbox{}\newline 
\textit{<!-- citation for the book referenced above -->}\mbox{}\newline 
{</\textbf{bibl}>}\end{shaded}\egroup\par \noindent  The value bibliog for the {\itshape type} attribute on the second \hyperref[TEI.ptr]{<ptr>} element here might be used to indicate that the object being referenced here is a bibliographic entry rather than a simple cross-reference to an illustration, as is the first \hyperref[TEI.ptr]{<ptr>}. In either case, the value of the {\itshape target} attribute is a pointer to some other element.\par
The \hyperref[TEI.ptr]{<ptr>} and \hyperref[TEI.ref]{<ref>} elements have many applications in addition to the simple cross-referencing facilities illustrated in this section. In conjunction with the analytic tools discussed in chapters \textit{\hyperref[SA]{16.\ Linking, Segmentation, and Alignment}}, \textit{\hyperref[AI]{17.\ Simple Analytic Mechanisms}}, and \textit{\hyperref[FS]{18.\ Feature Structures}}, they may be used to link analyses of a text to their object, to combine corresponding segments of a text, or to align segments of a text with a temporal or other axis or with each other.\par
Where the {\itshape target} attribute of \hyperref[TEI.ptr]{<ptr>} or \hyperref[TEI.ref]{<ref>} points to an external resource available on the network, the {\itshape mimeType} attribute may be used to specify the mime type of that resource; this may be important for to enable appropriate processing. For example: \par\bgroup\index{p=<p>|exampleindex}\index{ref=<ref>|exampleindex}\index{target=@target!<ref>|exampleindex}\index{mimeType=@mimeType!<ref>|exampleindex}\exampleFont \begin{shaded}\noindent\mbox{}{<\textbf{p}>}The current version of the TEI Guidelines source code\mbox{}\newline 
 is available in the TEI GitHub Repository; {<\textbf{ref}\hspace*{1em}{target}="{https://github.com/TEIC/TEI/blob/dev/P5/Source/guidelines-en.xml}"\mbox{}\newline 
\hspace*{1em}\hspace*{1em}{mimeType}="{application/tei+xml}">}guidelines-en.xml{</\textbf{ref}>} \mbox{}\newline 
 is the root document used to create the English version\mbox{}\newline 
 of these Guidelines.{</\textbf{p}>}\end{shaded}\egroup\par 
\subsection[{Lists}]{Lists}\label{COLI}\par
The following elements are provided for the encoding of lists, their constituent items, and the labels or headings associated with them: 
\begin{sansreflist}
  
\item [\textbf{<list>}] (list) contains any sequence of items organized as a list.
\item [\textbf{<item>}] (item) contains one component of a list.
\item [\textbf{<label>}] (label) contains any label or heading used to identify part of a text, typically but not exclusively in a list or glossary.
\item [\textbf{<head>}] (heading) contains any type of heading, for example the title of a section, or the heading of a list, glossary, manuscript description, etc.
\item [\textbf{<headLabel>}] (heading for list labels) contains the heading for the label or term column in a glossary list or similar structured list.
\item [\textbf{<headItem>}] (heading for list items) contains the heading for the item or gloss column in a glossary list or similar structured list.
\end{sansreflist}
\par
The \hyperref[TEI.list]{<list>} element may be used to mark any kind of \textit{list}: numbered, lettered, bulleted, or unmarked. Lists formatted as such in the copy text should in general be encoded using this element, with an appropriate value for the {\itshape rend} attribute. Suggested values for {\itshape rend} include: \begin{itemize}
\item \textit{bulleted} (items preceded by bullets or similar markings)
\item \textit{inline} (items rendered within continuous prose, with no linebreaks)
\item \textit{numbered} (items preceded by numbers or letters)
\item \textit{simple} (items rendered as blocks, but with no bullet or number)
\end{itemize}  Some of these values may of course be combined; a list may be inline, but also be rendered with numbers. An example appears below. For more sophisticated and detailed description of list rendering, consider using the {\itshape style} attribute with Cascading Stylesheet properties and values, as described in the W3C's \xref{http://www.w3.org/TR/css-lists-3/}{CSS Lists and Counters Module Level 3}.\par
Each distinct item in the list should be encoded as a distinct \hyperref[TEI.item]{<item>} element. If the numbering or other identification for the items in a list is unremarkable and may be reconstructed by any processing program, no enumerator need be specified. If however an enumerator is retained in the encoded text, it may be supplied either by using the {\itshape n} attribute on the \hyperref[TEI.item]{<item>} element, or by using a \hyperref[TEI.label]{<label>} element. The following examples are thus equivalent: \par\bgroup\index{list=<list>|exampleindex}\index{rend=@rend!<list>|exampleindex}\index{label=<label>|exampleindex}\index{item=<item>|exampleindex}\index{label=<label>|exampleindex}\index{item=<item>|exampleindex}\exampleFont \begin{shaded}\noindent\mbox{}I will add two facts, which have seldom occurred in\mbox{}\newline 
 the composition of six, or even five quartos.\mbox{}\newline 
{<\textbf{list}\hspace*{1em}{rend}="{inline numbered}">}\mbox{}\newline 
\hspace*{1em}{<\textbf{label}>}(1){</\textbf{label}>}\mbox{}\newline 
\hspace*{1em}{<\textbf{item}>}My first rough manuscript, without any\mbox{}\newline 
\hspace*{1em}\hspace*{1em} intermediate copy, has been sent to the press.{</\textbf{item}>}\mbox{}\newline 
\hspace*{1em}{<\textbf{label}>}(2){</\textbf{label}>}\mbox{}\newline 
\hspace*{1em}{<\textbf{item}>}Not a sheet has been seen by any human\mbox{}\newline 
\hspace*{1em}\hspace*{1em} eyes, excepting those of the author and the printer:\mbox{}\newline 
\hspace*{1em}\hspace*{1em} the faults and the merits are exclusively my own.{</\textbf{item}>}\mbox{}\newline 
{</\textbf{list}>}\end{shaded}\egroup\par \noindent  \par\bgroup\index{list=<list>|exampleindex}\index{rend=@rend!<list>|exampleindex}\index{item=<item>|exampleindex}\index{n=@n!<item>|exampleindex}\index{item=<item>|exampleindex}\index{n=@n!<item>|exampleindex}\exampleFont \begin{shaded}\noindent\mbox{}I will add two facts, which have seldom occurred in\mbox{}\newline 
 the composition of six, or even five quartos.\mbox{}\newline 
{<\textbf{list}\hspace*{1em}{rend}="{inline numbered}">}\mbox{}\newline 
\hspace*{1em}{<\textbf{item}\hspace*{1em}{n}="{1}">}My first rough manuscript, without any\mbox{}\newline 
\hspace*{1em}\hspace*{1em} intermediate copy, has been sent to the press.{</\textbf{item}>}\mbox{}\newline 
\hspace*{1em}{<\textbf{item}\hspace*{1em}{n}="{2}">}Not a sheet has been seen by any human\mbox{}\newline 
\hspace*{1em}\hspace*{1em} eyes, excepting those of the author and the printer:\mbox{}\newline 
\hspace*{1em}\hspace*{1em} the faults and the merits are exclusively my own.{</\textbf{item}>}\mbox{}\newline 
{</\textbf{list}>}\end{shaded}\egroup\par \noindent  The two styles may not be mixed in the same list: if one item is preceded by a label, all must be.\par
A list need not necessarily be displayed in list format. For example, the following is a reasonable encoding of a list which (in the original) is simply printed as a single paragraph: \par\bgroup\index{list=<list>|exampleindex}\index{rend=@rend!<list>|exampleindex}\index{item=<item>|exampleindex}\index{n=@n!<item>|exampleindex}\index{item=<item>|exampleindex}\index{n=@n!<item>|exampleindex}\index{item=<item>|exampleindex}\index{n=@n!<item>|exampleindex}\index{item=<item>|exampleindex}\index{n=@n!<item>|exampleindex}\index{item=<item>|exampleindex}\index{n=@n!<item>|exampleindex}\index{item=<item>|exampleindex}\index{n=@n!<item>|exampleindex}\index{item=<item>|exampleindex}\index{n=@n!<item>|exampleindex}\index{item=<item>|exampleindex}\index{n=@n!<item>|exampleindex}\index{item=<item>|exampleindex}\index{n=@n!<item>|exampleindex}\index{item=<item>|exampleindex}\index{n=@n!<item>|exampleindex}\index{item=<item>|exampleindex}\index{n=@n!<item>|exampleindex}\index{item=<item>|exampleindex}\index{n=@n!<item>|exampleindex}\index{item=<item>|exampleindex}\index{n=@n!<item>|exampleindex}\index{item=<item>|exampleindex}\index{n=@n!<item>|exampleindex}\exampleFont \begin{shaded}\noindent\mbox{}On those remote pages it is written that animals are\mbox{}\newline 
 divided into {<\textbf{list}\hspace*{1em}{rend}="{inline}">}\mbox{}\newline 
\hspace*{1em}{<\textbf{item}\hspace*{1em}{n}="{a}">}those that belong to the Emperor, {</\textbf{item}>}\mbox{}\newline 
\hspace*{1em}{<\textbf{item}\hspace*{1em}{n}="{b}">}embalmed ones, {</\textbf{item}>}\mbox{}\newline 
\hspace*{1em}{<\textbf{item}\hspace*{1em}{n}="{c}">}those that are trained, {</\textbf{item}>}\mbox{}\newline 
\hspace*{1em}{<\textbf{item}\hspace*{1em}{n}="{d}">}suckling pigs, {</\textbf{item}>}\mbox{}\newline 
\hspace*{1em}{<\textbf{item}\hspace*{1em}{n}="{e}">}mermaids, {</\textbf{item}>}\mbox{}\newline 
\hspace*{1em}{<\textbf{item}\hspace*{1em}{n}="{f}">}fabulous ones, {</\textbf{item}>}\mbox{}\newline 
\hspace*{1em}{<\textbf{item}\hspace*{1em}{n}="{g}">}stray dogs, {</\textbf{item}>}\mbox{}\newline 
\hspace*{1em}{<\textbf{item}\hspace*{1em}{n}="{h}">}those that are included in this classification, {</\textbf{item}>}\mbox{}\newline 
\hspace*{1em}{<\textbf{item}\hspace*{1em}{n}="{i}">}those that tremble as if they were mad, {</\textbf{item}>}\mbox{}\newline 
\hspace*{1em}{<\textbf{item}\hspace*{1em}{n}="{j}">}innumerable ones, {</\textbf{item}>}\mbox{}\newline 
\hspace*{1em}{<\textbf{item}\hspace*{1em}{n}="{k}">}those drawn with a very fine camel's-hair brush, {</\textbf{item}>}\mbox{}\newline 
\hspace*{1em}{<\textbf{item}\hspace*{1em}{n}="{l}">}others, {</\textbf{item}>}\mbox{}\newline 
\hspace*{1em}{<\textbf{item}\hspace*{1em}{n}="{m}">}those that have just broken a flower vase, {</\textbf{item}>}\mbox{}\newline 
\hspace*{1em}{<\textbf{item}\hspace*{1em}{n}="{n}">}those that resemble flies from a distance. {</\textbf{item}>}\mbox{}\newline 
{</\textbf{list}>}\end{shaded}\egroup\par \par
A list may be given a heading or title, for which the \hyperref[TEI.head]{<head>} element should be used, as in the next example, which also demonstrates simple use of the \hyperref[TEI.label]{<label>} element to mark a tabular or glossary list in which each item is associated with a word or phrase rather than a numeric or alphabetic enumerator: \par\bgroup\index{list=<list>|exampleindex}\index{type=@type!<list>|exampleindex}\index{head=<head>|exampleindex}\index{label=<label>|exampleindex}\index{item=<item>|exampleindex}\index{label=<label>|exampleindex}\index{item=<item>|exampleindex}\index{label=<label>|exampleindex}\index{item=<item>|exampleindex}\index{label=<label>|exampleindex}\index{item=<item>|exampleindex}\index{label=<label>|exampleindex}\index{item=<item>|exampleindex}\exampleFont \begin{shaded}\noindent\mbox{}{<\textbf{list}\hspace*{1em}{type}="{gloss}">}\mbox{}\newline 
\hspace*{1em}{<\textbf{head}>}Report of the conduct and progress of Ernest Pontifex.\mbox{}\newline 
\hspace*{1em}\hspace*{1em} Upper Vth form — half term ending Midsummer 1851{</\textbf{head}>}\mbox{}\newline 
\hspace*{1em}{<\textbf{label}>}Classics{</\textbf{label}>}\mbox{}\newline 
\hspace*{1em}{<\textbf{item}>}Idle listless and unimproving{</\textbf{item}>}\mbox{}\newline 
\hspace*{1em}{<\textbf{label}>}Mathematics{</\textbf{label}>}\mbox{}\newline 
\hspace*{1em}{<\textbf{item}>}ditto{</\textbf{item}>}\mbox{}\newline 
\hspace*{1em}{<\textbf{label}>}Divinity{</\textbf{label}>}\mbox{}\newline 
\hspace*{1em}{<\textbf{item}>}ditto{</\textbf{item}>}\mbox{}\newline 
\hspace*{1em}{<\textbf{label}>}Conduct in house{</\textbf{label}>}\mbox{}\newline 
\hspace*{1em}{<\textbf{item}>}Orderly{</\textbf{item}>}\mbox{}\newline 
\hspace*{1em}{<\textbf{label}>}General conduct{</\textbf{label}>}\mbox{}\newline 
\hspace*{1em}{<\textbf{item}>}Not satisfactory, on account of his great\mbox{}\newline 
\hspace*{1em}\hspace*{1em} unpunctuality and inattention to duties{</\textbf{item}>}\mbox{}\newline 
{</\textbf{list}>}\end{shaded}\egroup\par \par
In such a list, the individual items have internal structure. In complex cases, where list items contain many components, the list is better treated as a \textit{table},  on which see chapter \textit{\hyperref[FT]{14.\ Tables, Formulæ, Graphics and Notated Music}}. A particularly important instance of the simple two-column table is the ‘glossary list’, which should be marked by the tag <list type="gloss">. In such lists, each \hyperref[TEI.label]{<label>} element contains a term and each \hyperref[TEI.item]{<item>} its gloss; it is a semantic error for a list tagged with \texttt{type="gloss"} not to have labels. For example: \par\bgroup\index{list=<list>|exampleindex}\index{type=@type!<list>|exampleindex}\index{head=<head>|exampleindex}\index{label=<label>|exampleindex}\index{item=<item>|exampleindex}\index{label=<label>|exampleindex}\index{item=<item>|exampleindex}\index{label=<label>|exampleindex}\index{item=<item>|exampleindex}\index{label=<label>|exampleindex}\index{item=<item>|exampleindex}\index{label=<label>|exampleindex}\index{item=<item>|exampleindex}\exampleFont \begin{shaded}\noindent\mbox{}{<\textbf{list}\hspace*{1em}{type}="{gloss}">}\mbox{}\newline 
\hspace*{1em}{<\textbf{head}>}Unit Three — Vocabulary{</\textbf{head}>}\mbox{}\newline 
\hspace*{1em}{<\textbf{label}\hspace*{1em}{xml:lang}="{la}">}acerbus, -a, -um {</\textbf{label}>}\mbox{}\newline 
\hspace*{1em}{<\textbf{item}>}bitter, harsh{</\textbf{item}>}\mbox{}\newline 
\hspace*{1em}{<\textbf{label}\hspace*{1em}{xml:lang}="{la}">}ager, agrī, M. {</\textbf{label}>}\mbox{}\newline 
\hspace*{1em}{<\textbf{item}>}field{</\textbf{item}>}\mbox{}\newline 
\hspace*{1em}{<\textbf{label}\hspace*{1em}{xml:lang}="{la}">}audiō, īre,\mbox{}\newline 
\hspace*{1em}\hspace*{1em} īvī, ītus {</\textbf{label}>}\mbox{}\newline 
\hspace*{1em}{<\textbf{item}>}hear, listen (to){</\textbf{item}>}\mbox{}\newline 
\hspace*{1em}{<\textbf{label}\hspace*{1em}{xml:lang}="{la}">}bellum, -ī, N. {</\textbf{label}>}\mbox{}\newline 
\hspace*{1em}{<\textbf{item}>}war{</\textbf{item}>}\mbox{}\newline 
\hspace*{1em}{<\textbf{label}\hspace*{1em}{xml:lang}="{la}">}bonus, -a, -um {</\textbf{label}>}\mbox{}\newline 
\hspace*{1em}{<\textbf{item}>}good{</\textbf{item}>}\mbox{}\newline 
{</\textbf{list}>}\end{shaded}\egroup\par \noindent  Additionally, the \hyperref[TEI.term]{<term>} and \hyperref[TEI.gloss]{<gloss>} elements discussed in section \textit{\hyperref[COHTG]{3.4.1.\ Terms and Glosses}} might be used to make explicit the role that each column in the glossary list has, as follows: \par\bgroup\index{list=<list>|exampleindex}\index{type=@type!<list>|exampleindex}\index{head=<head>|exampleindex}\index{label=<label>|exampleindex}\index{term=<term>|exampleindex}\index{item=<item>|exampleindex}\index{gloss=<gloss>|exampleindex}\index{label=<label>|exampleindex}\index{term=<term>|exampleindex}\index{item=<item>|exampleindex}\index{gloss=<gloss>|exampleindex}\index{label=<label>|exampleindex}\index{term=<term>|exampleindex}\index{item=<item>|exampleindex}\index{gloss=<gloss>|exampleindex}\index{label=<label>|exampleindex}\index{term=<term>|exampleindex}\index{item=<item>|exampleindex}\index{gloss=<gloss>|exampleindex}\index{label=<label>|exampleindex}\index{term=<term>|exampleindex}\index{item=<item>|exampleindex}\index{gloss=<gloss>|exampleindex}\exampleFont \begin{shaded}\noindent\mbox{}{<\textbf{list}\hspace*{1em}{type}="{gloss}">}\mbox{}\newline 
\hspace*{1em}{<\textbf{head}>}Unit Three — Vocabulary{</\textbf{head}>}\mbox{}\newline 
\hspace*{1em}{<\textbf{label}>}\mbox{}\newline 
\hspace*{1em}\hspace*{1em}{<\textbf{term}\hspace*{1em}{xml:lang}="{la}">}acerbus, -a, -um{</\textbf{term}>}\mbox{}\newline 
\hspace*{1em}{</\textbf{label}>}\mbox{}\newline 
\hspace*{1em}{<\textbf{item}>}\mbox{}\newline 
\hspace*{1em}\hspace*{1em}{<\textbf{gloss}>}bitter, harsh{</\textbf{gloss}>}\mbox{}\newline 
\hspace*{1em}{</\textbf{item}>}\mbox{}\newline 
\hspace*{1em}{<\textbf{label}>}\mbox{}\newline 
\hspace*{1em}\hspace*{1em}{<\textbf{term}\hspace*{1em}{xml:lang}="{la}">}ager, agrī, M. {</\textbf{term}>}\mbox{}\newline 
\hspace*{1em}{</\textbf{label}>}\mbox{}\newline 
\hspace*{1em}{<\textbf{item}>}\mbox{}\newline 
\hspace*{1em}\hspace*{1em}{<\textbf{gloss}>}field{</\textbf{gloss}>}\mbox{}\newline 
\hspace*{1em}{</\textbf{item}>}\mbox{}\newline 
\hspace*{1em}{<\textbf{label}>}\mbox{}\newline 
\hspace*{1em}\hspace*{1em}{<\textbf{term}\hspace*{1em}{xml:lang}="{la}">}audiō, -īre, -īvī, -ītus{</\textbf{term}>}\mbox{}\newline 
\hspace*{1em}{</\textbf{label}>}\mbox{}\newline 
\hspace*{1em}{<\textbf{item}>}\mbox{}\newline 
\hspace*{1em}\hspace*{1em}{<\textbf{gloss}>}hear, listen (to){</\textbf{gloss}>}\mbox{}\newline 
\hspace*{1em}{</\textbf{item}>}\mbox{}\newline 
\hspace*{1em}{<\textbf{label}>}\mbox{}\newline 
\hspace*{1em}\hspace*{1em}{<\textbf{term}\hspace*{1em}{xml:lang}="{la}">}bellum, -ī, N. {</\textbf{term}>}\mbox{}\newline 
\hspace*{1em}{</\textbf{label}>}\mbox{}\newline 
\hspace*{1em}{<\textbf{item}>}\mbox{}\newline 
\hspace*{1em}\hspace*{1em}{<\textbf{gloss}>}war{</\textbf{gloss}>}\mbox{}\newline 
\hspace*{1em}{</\textbf{item}>}\mbox{}\newline 
\hspace*{1em}{<\textbf{label}>}\mbox{}\newline 
\hspace*{1em}\hspace*{1em}{<\textbf{term}\hspace*{1em}{xml:lang}="{la}">}bonus, -a, -um{</\textbf{term}>}\mbox{}\newline 
\hspace*{1em}{</\textbf{label}>}\mbox{}\newline 
\hspace*{1em}{<\textbf{item}>}\mbox{}\newline 
\hspace*{1em}\hspace*{1em}{<\textbf{gloss}>}good{</\textbf{gloss}>}\mbox{}\newline 
\hspace*{1em}{</\textbf{item}>}\mbox{}\newline 
{</\textbf{list}>}\end{shaded}\egroup\par \noindent  Note in the above examples the use of the global {\itshape xml:lang} attribute to specify on the \hyperref[TEI.label]{<label>} (or \hyperref[TEI.term]{<term>}) element what language the term is from. For further discussion of the {\itshape xml:lang} attribute see section \textit{\hyperref[STGA]{1.3.1.1.\ Global Attributes}}, and section \textit{\hyperref[CHSH]{vi.1\ Language Identification}}. A more elaborate markup for this glossary would distinguish the headword forms from the grammatical information (principal parts and gender), perhaps using elements taken from \textit{\hyperref[DI]{9.\ Dictionaries}}.\par
In addition to the \hyperref[TEI.head]{<head>} element used to supply a title or heading for the whole list, headings for the two columns of a glossary-style list may be specified using the two special elements \hyperref[TEI.headLabel]{<headLabel>} and \hyperref[TEI.headItem]{<headItem>}: \par\bgroup\index{list=<list>|exampleindex}\index{type=@type!<list>|exampleindex}\index{headLabel=<headLabel>|exampleindex}\index{headItem=<headItem>|exampleindex}\index{label=<label>|exampleindex}\index{item=<item>|exampleindex}\index{label=<label>|exampleindex}\index{item=<item>|exampleindex}\index{label=<label>|exampleindex}\index{item=<item>|exampleindex}\index{label=<label>|exampleindex}\index{item=<item>|exampleindex}\index{label=<label>|exampleindex}\index{item=<item>|exampleindex}\exampleFont \begin{shaded}\noindent\mbox{}The simple, straightforward statement of an idea is\mbox{}\newline 
 preferable to the use of a worn-out expression.\mbox{}\newline 
{<\textbf{list}\hspace*{1em}{type}="{gloss}">}\mbox{}\newline 
\hspace*{1em}{<\textbf{headLabel}>}TRITE{</\textbf{headLabel}>}\mbox{}\newline 
\hspace*{1em}{<\textbf{headItem}>}SIMPLE, STRAIGHTFORWARD{</\textbf{headItem}>}\mbox{}\newline 
\hspace*{1em}{<\textbf{label}>}bury the hatchet {</\textbf{label}>}\mbox{}\newline 
\hspace*{1em}{<\textbf{item}>}stop fighting, make peace{</\textbf{item}>}\mbox{}\newline 
\hspace*{1em}{<\textbf{label}>}at loose ends {</\textbf{label}>}\mbox{}\newline 
\hspace*{1em}{<\textbf{item}>}disorganized{</\textbf{item}>}\mbox{}\newline 
\hspace*{1em}{<\textbf{label}>}on speaking terms {</\textbf{label}>}\mbox{}\newline 
\hspace*{1em}{<\textbf{item}>}friendly{</\textbf{item}>}\mbox{}\newline 
\hspace*{1em}{<\textbf{label}>}fair and square {</\textbf{label}>}\mbox{}\newline 
\hspace*{1em}{<\textbf{item}>}completely honest{</\textbf{item}>}\mbox{}\newline 
\hspace*{1em}{<\textbf{label}>}at death's door {</\textbf{label}>}\mbox{}\newline 
\hspace*{1em}{<\textbf{item}>}near death{</\textbf{item}>}\mbox{}\newline 
{</\textbf{list}>}\end{shaded}\egroup\par \par
The elements \hyperref[TEI.label]{<label>}, \hyperref[TEI.head]{<head>}, \hyperref[TEI.headLabel]{<headLabel>}, and \hyperref[TEI.headItem]{<headItem>} may contain only phrase-level elements. The \hyperref[TEI.item]{<item>} element however may contain paragraphs or other ‘chunks’, including other lists. In this example, a glossary list contains two items, each of which is itself a simple list: \par\bgroup\index{list=<list>|exampleindex}\index{type=@type!<list>|exampleindex}\index{label=<label>|exampleindex}\index{item=<item>|exampleindex}\index{list=<list>|exampleindex}\index{rend=@rend!<list>|exampleindex}\index{item=<item>|exampleindex}\index{item=<item>|exampleindex}\index{item=<item>|exampleindex}\index{label=<label>|exampleindex}\index{item=<item>|exampleindex}\index{list=<list>|exampleindex}\index{rend=@rend!<list>|exampleindex}\index{item=<item>|exampleindex}\index{item=<item>|exampleindex}\index{item=<item>|exampleindex}\exampleFont \begin{shaded}\noindent\mbox{}{<\textbf{list}\hspace*{1em}{type}="{gloss}">}\mbox{}\newline 
\hspace*{1em}{<\textbf{label}>}EVIL{</\textbf{label}>}\mbox{}\newline 
\hspace*{1em}{<\textbf{item}>}\mbox{}\newline 
\hspace*{1em}\hspace*{1em}{<\textbf{list}\hspace*{1em}{rend}="{bulleted}">}\mbox{}\newline 
\hspace*{1em}\hspace*{1em}\hspace*{1em}{<\textbf{item}>}I am cast upon a horrible desolate island, void\mbox{}\newline 
\hspace*{1em}\hspace*{1em}\hspace*{1em}\hspace*{1em}\hspace*{1em}\hspace*{1em} of all hope of recovery.{</\textbf{item}>}\mbox{}\newline 
\hspace*{1em}\hspace*{1em}\hspace*{1em}{<\textbf{item}>}I am singled out and separated as it were from\mbox{}\newline 
\hspace*{1em}\hspace*{1em}\hspace*{1em}\hspace*{1em}\hspace*{1em}\hspace*{1em} all the world to be miserable.{</\textbf{item}>}\mbox{}\newline 
\hspace*{1em}\hspace*{1em}\hspace*{1em}{<\textbf{item}>}I am divided from mankind — a solitaire; one\mbox{}\newline 
\hspace*{1em}\hspace*{1em}\hspace*{1em}\hspace*{1em}\hspace*{1em}\hspace*{1em} banished from human society.{</\textbf{item}>}\mbox{}\newline 
\hspace*{1em}\hspace*{1em}{</\textbf{list}>}\mbox{}\newline 
\hspace*{1em}{</\textbf{item}>}\mbox{}\newline 
\hspace*{1em}{<\textbf{label}>}GOOD{</\textbf{label}>}\mbox{}\newline 
\hspace*{1em}{<\textbf{item}>}\mbox{}\newline 
\hspace*{1em}\hspace*{1em}{<\textbf{list}\hspace*{1em}{rend}="{bulleted}">}\mbox{}\newline 
\hspace*{1em}\hspace*{1em}\hspace*{1em}{<\textbf{item}>}But I am alive; and not drowned, as all my\mbox{}\newline 
\hspace*{1em}\hspace*{1em}\hspace*{1em}\hspace*{1em}\hspace*{1em}\hspace*{1em} ship's company were.{</\textbf{item}>}\mbox{}\newline 
\hspace*{1em}\hspace*{1em}\hspace*{1em}{<\textbf{item}>}But I am singled out, too, from all the ship's\mbox{}\newline 
\hspace*{1em}\hspace*{1em}\hspace*{1em}\hspace*{1em}\hspace*{1em}\hspace*{1em} crew, to be spared from death...{</\textbf{item}>}\mbox{}\newline 
\hspace*{1em}\hspace*{1em}\hspace*{1em}{<\textbf{item}>}But I am not starved, and perishing on a barren place,\mbox{}\newline 
\hspace*{1em}\hspace*{1em}\hspace*{1em}\hspace*{1em}\hspace*{1em}\hspace*{1em} affording no sustenances....{</\textbf{item}>}\mbox{}\newline 
\hspace*{1em}\hspace*{1em}{</\textbf{list}>}\mbox{}\newline 
\hspace*{1em}{</\textbf{item}>}\mbox{}\newline 
{</\textbf{list}>}\end{shaded}\egroup\par \par
Lists of different types may be nested to arbitrary depths in this way. 
\subsection[{Notes, Annotation, and Indexing}]{Notes, Annotation, and Indexing}\label{CONO}
\subsubsection[{Notes and Simple Annotation}]{Notes and Simple Annotation}\label{CONONO}\par
The following element is provided for the encoding of discursive notes, whether already present in the copy text or supplied by the encoder: 
\begin{sansreflist}
  
\item [\textbf{<note>}] (note) contains a note or annotation.
\end{sansreflist}
\par
A note is any additional comment found in a text, marked in some way as being out of the main textual stream. All notes should be marked using the same tag, \hyperref[TEI.note]{<note>}, whether they appear as block notes in the main text area, at the foot of the page, at the end of the chapter or volume, in the margin, or in some other place.\par
Notes may be in a different hand or typeface, may be authorial or editorial, and may have been added later. Attributes may be used to specify these and other characteristics of notes, as detailed below. \par
A note is usually attached to a specific point or span within a text, which we term here its \textit{point of attachment}. In conventional printed text, the point of attachment is represented by some siglum such as a star or cross, or a superscript digit.\par
When encoding such a text, it is conventional to replace this siglum by the content of the annotation, duly marked up with a \hyperref[TEI.note]{<note>} element. This may not always be possible for example with marginal notes, which may not be anchored to an exact location. For ease of processing, it may be adequate to position marginal notes before the relevant paragraph or other element. In printed texts, it is sometimes conventional to group notes together at the foot of the page on which their points of attachment appear. This practice is not generally recommended for TEI-encoded texts, since the pagination of a particular printed text is unlikely to be of structural significance. In some cases, however, it may be desirable to transcribe notes not at their point of attachment to the text but at their point of appearance, typically at the end of the volume, or the end of the chapter. In such cases, the {\itshape target} attribute of the \hyperref[TEI.note]{<note>} may be used to indicate the point of attachment. It is also possible to encode the point of attachment itself, using the \hyperref[TEI.ptr]{<ptr>} or \hyperref[TEI.ref]{<ref>} element, pointing from that to the body of the \hyperref[TEI.note]{<note>} placed elsewhere.\par
In cases where the note is applied not to a point but to a span of text, not itself represented as a TEI element, the {\itshape target} attribute may use an appropriate pointer expression, for example using the \textsf{range()} function to specify the span of attachment.\par
For further discussion of pointing to points and spans in the text, see section \textit{\hyperref[COXR]{3.7.\ Simple Links and Cross-References}}.\par
In the following example, the {\itshape type} attribute is used to categorise the note as a gloss: \par\bgroup\index{l=<l>|exampleindex}\index{l=<l>|exampleindex}\index{l=<l>|exampleindex}\index{l=<l>|exampleindex}\index{note=<note>|exampleindex}\index{type=@type!<note>|exampleindex}\index{place=@place!<note>|exampleindex}\exampleFont \begin{shaded}\noindent\mbox{}{<\textbf{l}>}The self-same moment I could pray{</\textbf{l}>}\mbox{}\newline 
{<\textbf{l}>}And from my neck so free{</\textbf{l}>}\mbox{}\newline 
{<\textbf{l}>}The albatross fell off, and sank{</\textbf{l}>}\mbox{}\newline 
{<\textbf{l}>}Like lead into the sea.\mbox{}\newline 
{<\textbf{note}\hspace*{1em}{type}="{gloss}"\hspace*{1em}{place}="{margin}">}The spell begins to break{</\textbf{note}>}\mbox{}\newline 
{</\textbf{l}>}\end{shaded}\egroup\par \noindent  As the \hyperref[TEI.note]{<note>} appears within an \hyperref[TEI.l]{<l>} element, we may infer that its point of attachment is in the margin adjacent to the line in question. In the following version of the same text, however, it may be inferred that the note applies to the whole of the stanza: \par\bgroup\index{lg=<lg>|exampleindex}\index{l=<l>|exampleindex}\index{l=<l>|exampleindex}\index{l=<l>|exampleindex}\index{l=<l>|exampleindex}\index{note=<note>|exampleindex}\index{type=@type!<note>|exampleindex}\index{place=@place!<note>|exampleindex}\exampleFont \begin{shaded}\noindent\mbox{}{<\textbf{lg}>}\mbox{}\newline 
\hspace*{1em}{<\textbf{l}>}The self-same moment I could pray{</\textbf{l}>}\mbox{}\newline 
\hspace*{1em}{<\textbf{l}>}And from my neck so free{</\textbf{l}>}\mbox{}\newline 
\hspace*{1em}{<\textbf{l}>}The albatross fell off, and sank{</\textbf{l}>}\mbox{}\newline 
\hspace*{1em}{<\textbf{l}>}Like lead into the sea.{</\textbf{l}>}\mbox{}\newline 
\hspace*{1em}{<\textbf{note}\hspace*{1em}{type}="{gloss}"\hspace*{1em}{place}="{margin}">}The spell begins to break{</\textbf{note}>}\mbox{}\newline 
{</\textbf{lg}>}\end{shaded}\egroup\par \par
This type of annotation, very common in the early printed texts which Coleridge may be presumed to be imitating in this case, may also be regarded as providing a heading or descriptive label for the passage concerned. The encoder may therefore prefer to use the \hyperref[TEI.label]{<label>} element to represent it, as in the following case: \par\bgroup\index{lg=<lg>|exampleindex}\index{l=<l>|exampleindex}\index{l=<l>|exampleindex}\index{l=<l>|exampleindex}\index{l=<l>|exampleindex}\index{label=<label>|exampleindex}\index{place=@place!<label>|exampleindex}\exampleFont \begin{shaded}\noindent\mbox{}{<\textbf{lg}>}\mbox{}\newline 
\hspace*{1em}{<\textbf{l}>}The self-same moment I could pray{</\textbf{l}>}\mbox{}\newline 
\hspace*{1em}{<\textbf{l}>}And from my neck so free{</\textbf{l}>}\mbox{}\newline 
\hspace*{1em}{<\textbf{l}>}The albatross fell off, and sank{</\textbf{l}>}\mbox{}\newline 
\hspace*{1em}{<\textbf{l}>}Like lead into the sea.{</\textbf{l}>}\mbox{}\newline 
\hspace*{1em}{<\textbf{label}\hspace*{1em}{place}="{margin}">}The spell begins to break{</\textbf{label}>}\mbox{}\newline 
{</\textbf{lg}>}\end{shaded}\egroup\par \par
In the following example, a note which appears at the foot of the page in the printed source is given at its point of attachment within the text. The global {\itshape n} attribute is used to indicate the note number: \par\bgroup\index{note=<note>|exampleindex}\index{n=@n!<note>|exampleindex}\index{place=@place!<note>|exampleindex}\index{mentioned=<mentioned>|exampleindex}\index{mentioned=<mentioned>|exampleindex}\index{mentioned=<mentioned>|exampleindex}\index{mentioned=<mentioned>|exampleindex}\exampleFont \begin{shaded}\noindent\mbox{}Collections are ensembles of\mbox{}\newline 
 distinct entities or objects of any sort.{<\textbf{note}\hspace*{1em}{n}="{1}"\hspace*{1em}{place}="{bottom}">}We\mbox{}\newline 
 explain below why we use the uncommon term\mbox{}\newline 
{<\textbf{mentioned}>}collection{</\textbf{mentioned}>} instead of the expected\mbox{}\newline 
{<\textbf{mentioned}>}set{</\textbf{mentioned}>}. Our usage corresponds to the\mbox{}\newline 
{<\textbf{mentioned}>}aggregate{</\textbf{mentioned}>} of many mathematical writings and to\mbox{}\newline 
 the sense of {<\textbf{mentioned}>}class{</\textbf{mentioned}>} found in older logical\mbox{}\newline 
 writings.{</\textbf{note}>} The elements ...\end{shaded}\egroup\par \par
In addition to transcribing notes already present in the copy text, researchers may wish to add their own notes or comments to it. The \hyperref[TEI.note]{<note>} element may be used for either purpose, but it will usually be advisable to distinguish the two categories. One way might be to use the {\itshape type} attribute shown above, categorizing notes as \textit{authorial}, \textit{editorial}, etc. Where notes derive from many sources, or where a more precise attribution is required, the {\itshape resp} attribute may be used to point to a definition of the person or other agency responsible for the content of the note.\par
As a simple example, an edition of the \textit{Ancient Mariner} might include both Coleridge's original glosses and those of a modern commentator: \par\bgroup\index{lg=<lg>|exampleindex}\index{l=<l>|exampleindex}\index{note=<note>|exampleindex}\index{place=@place!<note>|exampleindex}\index{resp=@resp!<note>|exampleindex}\index{type=@type!<note>|exampleindex}\index{note=<note>|exampleindex}\index{place=@place!<note>|exampleindex}\index{resp=@resp!<note>|exampleindex}\exampleFont \begin{shaded}\noindent\mbox{}{<\textbf{lg}>}\mbox{}\newline 
\textit{<!-- ... -->}\mbox{}\newline 
\hspace*{1em}{<\textbf{l}>}The self-same moment I could pray;\mbox{}\newline 
\hspace*{1em}{<\textbf{note}\hspace*{1em}{place}="{margin}"\hspace*{1em}{resp}="{\#STC}"\mbox{}\newline 
\hspace*{1em}\hspace*{1em}\hspace*{1em}{type}="{gloss}">}The spell begins to break{</\textbf{note}>}\mbox{}\newline 
\hspace*{1em}\hspace*{1em}{<\textbf{note}\hspace*{1em}{place}="{bottom}"\hspace*{1em}{resp}="{\#JLL}">}The turning point of the poem...{</\textbf{note}>}\mbox{}\newline 
\hspace*{1em}{</\textbf{l}>}\mbox{}\newline 
{</\textbf{lg}>}\end{shaded}\egroup\par \noindent  For this to be valid, the codes \texttt{\#JLL} and \texttt{\#STC} must point to some more information identifying the agency concerned. The syntax used is identical to that used for other cross-references, as discussed in \textit{\hyperref[COXR]{3.7.\ Simple Links and Cross-References}}; thus in this case, the TEI header for this text might contain a title statement like the following: \par\bgroup\index{titleStmt=<titleStmt>|exampleindex}\index{title=<title>|exampleindex}\index{author=<author>|exampleindex}\index{editor=<editor>|exampleindex}\exampleFont \begin{shaded}\noindent\mbox{}{<\textbf{titleStmt}>}\mbox{}\newline 
\hspace*{1em}{<\textbf{title}>}The Rime of the Ancient Mariner: an annotated edition{</\textbf{title}>}\mbox{}\newline 
\hspace*{1em}{<\textbf{author}\hspace*{1em}{xml:id}="{STC}">}Samuel Taylor Coleridge{</\textbf{author}>}\mbox{}\newline 
\hspace*{1em}{<\textbf{editor}\hspace*{1em}{xml:id}="{JLL}">}John Livingston Lowes{</\textbf{editor}>}\mbox{}\newline 
{</\textbf{titleStmt}>}\end{shaded}\egroup\par \par
When annotating the electronic text by means of analytic notes in some structured vocabulary, e.g. to specify the topics or themes of a text, the \hyperref[TEI.span]{<span>} and \hyperref[TEI.interp]{<interp>} elements may be more effective than the free form \hyperref[TEI.note]{<note>} element; these elements are available when the module for simple analysis is selected (see section \textit{\hyperref[AISP]{17.3.\ Spans and Interpretations}}).
\paragraph[{Encoding Grouped Notes}]{Encoding Grouped Notes}\label{CONONOGR}\par
The following element is provided for the grouping of notes: 
\begin{sansreflist}
  
\item [\textbf{<noteGrp>}] contains a group of notes
\end{sansreflist}
\par
A text may have multiple alternative versions of the same note, such as the same annotation expressed in multiple languages, or both an extensive note and a short form for different audiences. In such cases multiple \hyperref[TEI.note]{<note>} elements may be grouped within a \hyperref[TEI.noteGrp]{<noteGrp>} element.\par
Typically, the \hyperref[TEI.note]{<note>} elements within a \hyperref[TEI.noteGrp]{<noteGrp>} would be differentiated by use of attributes such as {\itshape xml:lang} or {\itshape type}, while sharing the same point of attachment. This differentiation can be made either implicitly in case of inline notes, or explicitly via a {\itshape target} attribute, which may be specified on the \hyperref[TEI.noteGrp]{<noteGrp>} itself.\par
The simple example below demonstrates the grouping of a short and a full version of the same note, where the short version might be intended for use in contexts with space constraints. \par\bgroup\index{p=<p>|exampleindex}\index{noteGrp=<noteGrp>|exampleindex}\index{note=<note>|exampleindex}\index{type=@type!<note>|exampleindex}\index{note=<note>|exampleindex}\index{type=@type!<note>|exampleindex}\exampleFont \begin{shaded}\noindent\mbox{}{<\textbf{p}>}... reuerendos dominos archiepiscopum et canonicos Leopolienses\mbox{}\newline 
 in duplicibus Quatuortemporibus\mbox{}\newline 
{<\textbf{noteGrp}>}\mbox{}\newline 
\hspace*{1em}\hspace*{1em}{<\textbf{note}\hspace*{1em}{type}="{short}">}Quatuor Tempora, so called dry fast days.{</\textbf{note}>}\mbox{}\newline 
\hspace*{1em}\hspace*{1em}{<\textbf{note}\hspace*{1em}{type}="{full}">}Quatuor Tempora, so called dry fast days (Wednesday, Friday, and Saturday)\mbox{}\newline 
\hspace*{1em}\hspace*{1em}\hspace*{1em}\hspace*{1em} falling on each of the quarters of the year. In the first quarter they were called Cinerum\mbox{}\newline 
\hspace*{1em}\hspace*{1em}\hspace*{1em}\hspace*{1em} (following Ash Wednesday), second Spiritus (following Pentecost), third Crucis\mbox{}\newline 
\hspace*{1em}\hspace*{1em}\hspace*{1em}\hspace*{1em} (after the Exaltation of the Holy Cross, September 14th), and Luciae\mbox{}\newline 
\hspace*{1em}\hspace*{1em}\hspace*{1em}\hspace*{1em} in the fourth (after the feast of St. Lucia, December 13th).\mbox{}\newline 
\hspace*{1em}\hspace*{1em}{</\textbf{note}>}\mbox{}\newline 
\hspace*{1em}{</\textbf{noteGrp}>}\mbox{}\newline 
 totaliter expediui.\mbox{}\newline 
{</\textbf{p}>}\end{shaded}\egroup\par 
\subsubsection[{Index Entries}]{Index Entries}\label{CONOIX}\par
The indexing of scholarly texts is a skilled activity, involving substantial amounts of human judgment and analysis. It should not therefore be assumed that simple searching and information retrieval software will be able to meet all the needs addressed by a well-crafted manual index, although it may complement them for example by providing free text search. The role of an index is to provide access via keywords and phrases which are not necessarily present in the text itself, but must be added by the skill of the indexer. 
\paragraph[{Pre-existing Indexes}]{Pre-existing Indexes}\label{CONOIXpre}\par
When encoding a pre-existing text, therefore, if such an index is present it may be advisable to retain it along with the text, rather than attempt to regenerate it automatically. Elements discussed elsewhere in these Guidelines may be used for this purpose. For example, the \hyperref[TEI.div1]{<div1>} element or \hyperref[TEI.div]{<div>} element may be used to mark the section of the text containing the index and the \hyperref[TEI.list]{<list>} element might be used to mark the index itself, each entry being represented by an \hyperref[TEI.item]{<item>} element, possibly containing within it a series of \hyperref[TEI.ptr]{<ptr>} or \hyperref[TEI.ref]{<ref>} elements, as follows: \par\bgroup\index{div=<div>|exampleindex}\index{type=@type!<div>|exampleindex}\index{list=<list>|exampleindex}\index{type=@type!<list>|exampleindex}\index{item=<item>|exampleindex}\index{ref=<ref>|exampleindex}\index{ref=<ref>|exampleindex}\index{ref=<ref>|exampleindex}\index{ref=<ref>|exampleindex}\index{ref=<ref>|exampleindex}\index{item=<item>|exampleindex}\index{ref=<ref>|exampleindex}\index{item=<item>|exampleindex}\index{ref=<ref>|exampleindex}\index{item=<item>|exampleindex}\exampleFont \begin{shaded}\noindent\mbox{}{<\textbf{div}\hspace*{1em}{type}="{index}">}\mbox{}\newline 
\textit{<!--...-->}\mbox{}\newline 
\hspace*{1em}{<\textbf{list}\hspace*{1em}{type}="{index}">}\mbox{}\newline 
\hspace*{1em}\hspace*{1em}{<\textbf{item}>}Women, how cause of mel. {<\textbf{ref}>}193{</\textbf{ref}>}; their vanity in\mbox{}\newline 
\hspace*{1em}\hspace*{1em}\hspace*{1em}\hspace*{1em} apparell taxed, {<\textbf{ref}>}527{</\textbf{ref}>}; their counterfeit tears\mbox{}\newline 
\hspace*{1em}\hspace*{1em}{<\textbf{ref}>}547{</\textbf{ref}>}; their vices {<\textbf{ref}>}601{</\textbf{ref}>}, commended,\mbox{}\newline 
\hspace*{1em}\hspace*{1em}{<\textbf{ref}>}624{</\textbf{ref}>}.{</\textbf{item}>}\mbox{}\newline 
\hspace*{1em}\hspace*{1em}{<\textbf{item}>}Wormwood, good against mel. {<\textbf{ref}>}443{</\textbf{ref}>}\mbox{}\newline 
\hspace*{1em}\hspace*{1em}{</\textbf{item}>}\mbox{}\newline 
\hspace*{1em}\hspace*{1em}{<\textbf{item}>}World taxed, {<\textbf{ref}>}181{</\textbf{ref}>}\mbox{}\newline 
\hspace*{1em}\hspace*{1em}{</\textbf{item}>}\mbox{}\newline 
\hspace*{1em}\hspace*{1em}{<\textbf{item}>}Writers of the cure of mel. 295{</\textbf{item}>}\mbox{}\newline 
\textit{<!--...-->}\mbox{}\newline 
\hspace*{1em}{</\textbf{list}>}\mbox{}\newline 
{</\textbf{div}>}\end{shaded}\egroup\par \par
Note that this simple representation does not capture the nested structure of the first of these index entries. A more accurate representation might entail the use of nested lists like the following: \par\bgroup\index{item=<item>|exampleindex}\index{list=<list>|exampleindex}\index{item=<item>|exampleindex}\index{ref=<ref>|exampleindex}\index{item=<item>|exampleindex}\index{ref=<ref>|exampleindex}\index{item=<item>|exampleindex}\index{ref=<ref>|exampleindex}\index{item=<item>|exampleindex}\index{list=<list>|exampleindex}\index{item=<item>|exampleindex}\index{ref=<ref>|exampleindex}\index{item=<item>|exampleindex}\index{ref=<ref>|exampleindex}\exampleFont \begin{shaded}\noindent\mbox{}{<\textbf{item}>}Women,\mbox{}\newline 
{<\textbf{list}>}\mbox{}\newline 
\hspace*{1em}\hspace*{1em}{<\textbf{item}>}how cause of mel. {<\textbf{ref}>}193{</\textbf{ref}>};{</\textbf{item}>}\mbox{}\newline 
\hspace*{1em}\hspace*{1em}{<\textbf{item}>}their vanity in apparell taxed, {<\textbf{ref}>}527{</\textbf{ref}>};{</\textbf{item}>}\mbox{}\newline 
\hspace*{1em}\hspace*{1em}{<\textbf{item}>}their counterfeit tears {<\textbf{ref}>}547{</\textbf{ref}>};{</\textbf{item}>}\mbox{}\newline 
\hspace*{1em}\hspace*{1em}{<\textbf{item}>}their vices\mbox{}\newline 
\hspace*{1em}\hspace*{1em}{<\textbf{list}>}\mbox{}\newline 
\hspace*{1em}\hspace*{1em}\hspace*{1em}\hspace*{1em}{<\textbf{item}>}\mbox{}\newline 
\hspace*{1em}\hspace*{1em}\hspace*{1em}\hspace*{1em}\hspace*{1em}{<\textbf{ref}>}601{</\textbf{ref}>},{</\textbf{item}>}\mbox{}\newline 
\hspace*{1em}\hspace*{1em}\hspace*{1em}\hspace*{1em}{<\textbf{item}>} commended, {<\textbf{ref}>}624{</\textbf{ref}>}.{</\textbf{item}>}\mbox{}\newline 
\hspace*{1em}\hspace*{1em}\hspace*{1em}{</\textbf{list}>}\mbox{}\newline 
\hspace*{1em}\hspace*{1em}{</\textbf{item}>}\mbox{}\newline 
\hspace*{1em}{</\textbf{list}>}\mbox{}\newline 
{</\textbf{item}>}\end{shaded}\egroup\par \par
The page references, encoded simply as \hyperref[TEI.ref]{<ref>} elements above, might also include direct links to the appropriate location in the encoded text, using (for example) a target attribute to supply the identifier of an associated page break element: \par\bgroup\index{pb=<pb>|exampleindex}\index{ref=<ref>|exampleindex}\index{target=@target!<ref>|exampleindex}\exampleFont \begin{shaded}\noindent\mbox{}\mbox{}\newline 
\textit{<!-- in the text -->}{<\textbf{pb}\hspace*{1em}{xml:id}="{P624}"/>}\mbox{}\newline 
\textit{<!-- start of  page 624 -->}\mbox{}\newline 
\textit{<!-- in the index -->}\mbox{}\newline 
{<\textbf{ref}\hspace*{1em}{target}="{\#P624}">}624{</\textbf{ref}>}\end{shaded}\egroup\par \noindent  For further discussion of this and alternative ways of encoding such links see the discussion in section \textit{\hyperref[SA]{16.\ Linking, Segmentation, and Alignment}}. Note that similar methods may also be used to encode a table of contents, as further exemplified in section \textit{\hyperref[DSFRONT]{4.5.\ Front Matter}}.
\paragraph[{Auto-generated Indexes}]{Auto-generated Indexes}\label{CONOIXgen}\par
It can also be useful, however, to generate a new index from a machine-readable text, whether the text is being written for the first time with the tags here defined, or as an addition to a text transcribed from some other source. Depending on the complexity of the text and its subject matter, such an automatically-generated index may not in itself satisfy all the needs of scholarly users. However it can assist a professional indexer to construct a fully adequate index, which might then be post-edited into the digital text, marked-up along the lines already suggested for preserving pre-existing index material.\par
Indexes generally contain both references to specific pages or sections and references to page ranges or sequences. The same element is used in either case: 
\begin{sansreflist}
  
\item [\textbf{<index>}] (index entry) marks a location to be indexed for whatever purpose.
\end{sansreflist}
\par
Like the \hyperref[TEI.interp]{<interp>} element described in \textit{\hyperref[AISP]{17.3.\ Spans and Interpretations}} this element may be used simply to provide descriptive or interpretive label of some kind for any location within a text, to be processed in any way by analytic software, but its main purpose is to facilitate the generation of an index for a printed version of the text. An \hyperref[TEI.index]{<index>} element may be placed anywhere within a text, between or within other elements. The headwords to be used when making up this index are given by the \hyperref[TEI.term]{<term>} elements within the \hyperref[TEI.index]{<index>} element. The location of the generated index might be specified by means of a processing instruction within the text, such as the following (the exact form of the PI is of course dependent on the application software in use): \par\hfill\bgroup\exampleFont\vskip 10pt\begin{shaded}
\obeyspaces <?tei indexplacement ?>\end{shaded}
\par\egroup 
 Alternatively, the special purpose \hyperref[TEI.divGen]{<divGen>} element might be used.\par
In the simplest case, a single headword is supplied by a \hyperref[TEI.term]{<term>} element contained by an \hyperref[TEI.index]{<index>} element: \par\bgroup\index{p=<p>|exampleindex}\index{index=<index>|exampleindex}\index{term=<term>|exampleindex}\exampleFont \begin{shaded}\noindent\mbox{}{<\textbf{p}>}The students understand procedures for Arabic lemmatisation\mbox{}\newline 
{<\textbf{index}>}\mbox{}\newline 
\hspace*{1em}\hspace*{1em}{<\textbf{term}>}Lemmatization, Arabic{</\textbf{term}>}\mbox{}\newline 
\hspace*{1em}{</\textbf{index}>}and are beginning to build parsers.{</\textbf{p}>}\end{shaded}\egroup\par \par
The effect of this is to document an index entry for the term ‘Lemmatization, Arabic’, which when processed could reference the location of the original \hyperref[TEI.index]{<index>} element.\par
If the subject of Arabic lemmatization is treated at length in a text, then the index entry generated may need to reference a sequence of locations (e.g. page numbers). In such a case it will be necessary to identify the end of the relevant span of text as well as its starting point. This is most conveniently done by supplying an empty \hyperref[TEI.anchor]{<anchor>} element (as discussed in chapter \textit{\hyperref[SA]{16.\ Linking, Segmentation, and Alignment}}) at the appropriate point and pointing to it from the \hyperref[TEI.index]{<index>} element by means of its {\itshape spanTo} attribute, as in this example: \par\bgroup\index{p=<p>|exampleindex}\index{index=<index>|exampleindex}\index{spanTo=@spanTo!<index>|exampleindex}\index{term=<term>|exampleindex}\index{anchor=<anchor>|exampleindex}\exampleFont \begin{shaded}\noindent\mbox{}{<\textbf{p}>}We now turn to the\mbox{}\newline 
 topic of Arabic lemmatisation\mbox{}\newline 
{<\textbf{index}\hspace*{1em}{spanTo}="{\#ALAMEND}">}\mbox{}\newline 
\hspace*{1em}\hspace*{1em}{<\textbf{term}>}Lemmatization, Arabic{</\textbf{term}>}\mbox{}\newline 
\hspace*{1em}{</\textbf{index}>} concerning which it is important to note [...]\mbox{}\newline 
\mbox{}\newline 
\textit{<!-- much learned material omitted here -->}\mbox{}\newline 
 and now we can build our parser.{<\textbf{anchor}\hspace*{1em}{xml:id}="{ALAMEND}"/>}\mbox{}\newline 
{</\textbf{p}>}\end{shaded}\egroup\par \par
This would generate the same index entries as the previous example, but the reference would be to the whole span of text between the location of the \hyperref[TEI.index]{<index>} element and the location of the element identified by the code \textsf{ALAMEND}, rather than a single point, and thus might (for example) include a sequence of page numbers.\par
Although the position of the \hyperref[TEI.index]{<index>} element in the text provides the target location that will be specified in the generated index entry, no part of the text itself is used to construct that entry. Index terms appearing in the entry come solely from the content of \hyperref[TEI.term]{<term>} elements, which consequently may have to repeat words or phrases from the text proper. This need not be done verbatim, thus giving scope for normalization of spelling (as in the example above) or other modifications which may assist generation of an index in a desired form or sequence.\par
Sometimes, for example when index terms are taken from a different language or consist of mathematical formulae or other expressions, even a normalized form of an index term may be insufficient for an application to order it exactly as desired. The {\itshape sortKey} attribute may be used to address this problem, as in the following example: \par\bgroup\index{p=<p>|exampleindex}\index{index=<index>|exampleindex}\index{term=<term>|exampleindex}\index{sortKey=@sortKey!<term>|exampleindex}\exampleFont \begin{shaded}\noindent\mbox{}{<\textbf{p}>}The @ operator\mbox{}\newline 
{<\textbf{index}>}\mbox{}\newline 
\hspace*{1em}\hspace*{1em}{<\textbf{term}\hspace*{1em}{sortKey}="{0000}">}@{</\textbf{term}>}\mbox{}\newline 
\hspace*{1em}{</\textbf{index}>} precedes an\mbox{}\newline 
 attribute name{</\textbf{p}>}\end{shaded}\egroup\par \noindent  Here, an entry for the symbol @ will appear in the index, but will be sorted alphabetically as if it were the string 0000. This technique is also useful when an index entry is to contain some non-Unicode character or glyph represented by the \hyperref[TEI.g]{<g>} element discussed in chapter \textit{\hyperref[WD]{5.\ Characters, Glyphs, and Writing Modes}}. In the following example, we assume that somewhere a definition for this glyph has been provided using the elements described in chapter \textit{\hyperref[WD]{5.\ Characters, Glyphs, and Writing Modes}}, and given the code PrinceGlyph: \par\bgroup\index{char=<char>|exampleindex}\index{p=<p>|exampleindex}\index{index=<index>|exampleindex}\index{term=<term>|exampleindex}\index{sortKey=@sortKey!<term>|exampleindex}\index{g=<g>|exampleindex}\index{ref=@ref!<g>|exampleindex}\exampleFont \begin{shaded}\noindent\mbox{}{<\textbf{char}\hspace*{1em}{xml:id}="{PrinceGlyph}">}\mbox{}\newline 
\textit{<!-- definition of the glyph here -->}\mbox{}\newline 
{</\textbf{char}>}\mbox{}\newline 
{<\textbf{p}>}The Artist formerly known as Prince {<\textbf{index}>}\mbox{}\newline 
\hspace*{1em}\hspace*{1em}{<\textbf{term}\hspace*{1em}{sortKey}="{Prince}">}\mbox{}\newline 
\hspace*{1em}\hspace*{1em}\hspace*{1em}{<\textbf{g}\hspace*{1em}{ref}="{\#PrinceGlyph}"/>}\mbox{}\newline 
\hspace*{1em}\hspace*{1em}{</\textbf{term}>}\mbox{}\newline 
\hspace*{1em}{</\textbf{index}>}...{</\textbf{p}>}\end{shaded}\egroup\par \noindent  Note that if no value is supplied for the sortKey attribute, a sorting application should always use the content of the \hyperref[TEI.term]{<term>} element as a sort key.\par
It is common practice to compile more than one index for a given text. A biography of a poet, for example, may offer an index of references to poems by the subject of the study, another index of works by other writers, an index of places or historical personages etc. The indexName attribute is used to assigning index terms and locations to one or more specific indexes: \par\bgroup\index{p=<p>|exampleindex}\index{index=<index>|exampleindex}\index{indexName=@indexName!<index>|exampleindex}\index{term=<term>|exampleindex}\index{index=<index>|exampleindex}\index{indexName=@indexName!<index>|exampleindex}\index{term=<term>|exampleindex}\exampleFont \begin{shaded}\noindent\mbox{}{<\textbf{p}>}Sir John Ashford\mbox{}\newline 
{<\textbf{index}\hspace*{1em}{indexName}="{INDEX-PERSONS}">}\mbox{}\newline 
\hspace*{1em}\hspace*{1em}{<\textbf{term}>}Ashford, John{</\textbf{term}>}\mbox{}\newline 
\hspace*{1em}{</\textbf{index}>} was,\mbox{}\newline 
 coincidentally, born in\mbox{}\newline 
{<\textbf{index}\hspace*{1em}{indexName}="{INDEX-PLACES}">}\mbox{}\newline 
\hspace*{1em}\hspace*{1em}{<\textbf{term}>}Ashford\mbox{}\newline 
\hspace*{1em}\hspace*{1em}\hspace*{1em}\hspace*{1em} (Kent){</\textbf{term}>}\mbox{}\newline 
\hspace*{1em}{</\textbf{index}>}Ashford...{</\textbf{p}>}\end{shaded}\egroup\par \par
Multi-level indexing is particularly common in scholarly documents. For example, as well as entries such as \textit{TEI}, or \textit{markup}, an index may contain structured entries like \textit{TEI, markup practices, index terms}, where a top level entry \textit{TEI} is followed by a number of second-level subcategories, any or all of which may have a third-level list attached to them and so on. In order to reflect such a hierarchical index listing, \hyperref[TEI.index]{<index>} elements may be nested to the required depth. For example, suppose that we wish to make a structured index entry for ‘lemmatisation’ with subentries for ‘Arabic’, ‘Sanskrit’, etc. The example at the start of this section might then be encoded with nested \hyperref[TEI.index]{<index>} elements: \par\bgroup\index{p=<p>|exampleindex}\index{index=<index>|exampleindex}\index{term=<term>|exampleindex}\index{index=<index>|exampleindex}\index{term=<term>|exampleindex}\exampleFont \begin{shaded}\noindent\mbox{}{<\textbf{p}>}The students understand procedures for Arabic lemmatisation\mbox{}\newline 
{<\textbf{index}>}\mbox{}\newline 
\hspace*{1em}\hspace*{1em}{<\textbf{term}>}lemmatization{</\textbf{term}>}\mbox{}\newline 
\hspace*{1em}\hspace*{1em}{<\textbf{index}>}\mbox{}\newline 
\hspace*{1em}\hspace*{1em}\hspace*{1em}{<\textbf{term}>}arabic{</\textbf{term}>}\mbox{}\newline 
\hspace*{1em}\hspace*{1em}{</\textbf{index}>}\mbox{}\newline 
\hspace*{1em}{</\textbf{index}>}\mbox{}\newline 
 ...{</\textbf{p}>}\end{shaded}\egroup\par \par
The index entry from Burton's \textit{Anatomy of Melancholy} quoted above might be generated in a similar way. To generate such an entry, the body of the text might include, at page 193, an \hyperref[TEI.index]{<index>} element such as \par\bgroup\index{index=<index>|exampleindex}\index{term=<term>|exampleindex}\index{index=<index>|exampleindex}\index{term=<term>|exampleindex}\exampleFont \begin{shaded}\noindent\mbox{}{<\textbf{index}>}\mbox{}\newline 
\hspace*{1em}{<\textbf{term}>}Women{</\textbf{term}>}\mbox{}\newline 
\hspace*{1em}{<\textbf{index}>}\mbox{}\newline 
\hspace*{1em}\hspace*{1em}{<\textbf{term}>}how cause of mel.{</\textbf{term}>}\mbox{}\newline 
\hspace*{1em}{</\textbf{index}>}\mbox{}\newline 
{</\textbf{index}>}\end{shaded}\egroup\par \noindent . Similarly, page 601 of the body text would include an \hyperref[TEI.index]{<index>} element like the following: \par\bgroup\index{index=<index>|exampleindex}\index{term=<term>|exampleindex}\index{index=<index>|exampleindex}\index{term=<term>|exampleindex}\exampleFont \begin{shaded}\noindent\mbox{}{<\textbf{index}>}\mbox{}\newline 
\hspace*{1em}{<\textbf{term}>}Women{</\textbf{term}>}\mbox{}\newline 
\hspace*{1em}{<\textbf{index}>}\mbox{}\newline 
\hspace*{1em}\hspace*{1em}{<\textbf{term}>}their vices{</\textbf{term}>}\mbox{}\newline 
\hspace*{1em}{</\textbf{index}>}\mbox{}\newline 
{</\textbf{index}>}\end{shaded}\egroup\par \noindent  while the \hyperref[TEI.index]{<index>} element at page 624 would have a structure like the following: \par\bgroup\index{index=<index>|exampleindex}\index{term=<term>|exampleindex}\index{index=<index>|exampleindex}\index{term=<term>|exampleindex}\index{index=<index>|exampleindex}\index{term=<term>|exampleindex}\exampleFont \begin{shaded}\noindent\mbox{}{<\textbf{index}>}\mbox{}\newline 
\hspace*{1em}{<\textbf{term}>}Women{</\textbf{term}>}\mbox{}\newline 
\hspace*{1em}{<\textbf{index}>}\mbox{}\newline 
\hspace*{1em}\hspace*{1em}{<\textbf{term}>}their vices{</\textbf{term}>}\mbox{}\newline 
\hspace*{1em}\hspace*{1em}{<\textbf{index}>}\mbox{}\newline 
\hspace*{1em}\hspace*{1em}\hspace*{1em}{<\textbf{term}>}commended{</\textbf{term}>}\mbox{}\newline 
\hspace*{1em}\hspace*{1em}{</\textbf{index}>}\mbox{}\newline 
\hspace*{1em}{</\textbf{index}>}\mbox{}\newline 
{</\textbf{index}>}\end{shaded}\egroup\par \par
When processing such \hyperref[TEI.index]{<index>} elements, the duplication required to make the structure explicit will normally be removed, so as to produce entries like those quoted above. However, this is not required by the encoding recommended here.\par
As noted above, either a processing instruction or a \hyperref[TEI.divGen]{<divGen>} element may be used to mark the place at which an index generated from \hyperref[TEI.index]{<index>} elements should be inserted into the output of a processing program; typically but not necessarily this will be at some point within the back matter of the document. If the \hyperref[TEI.divGen]{<divGen>} element is used, then the {\itshape type} attribute should be used to specify which kind of index is to be generated, and its value should correspond with that of the {\itshape indexName} attribute on the relevant \hyperref[TEI.index]{<index>} elements. \par\bgroup\index{back=<back>|exampleindex}\index{div=<div>|exampleindex}\index{type=@type!<div>|exampleindex}\index{head=<head>|exampleindex}\index{listBibl=<listBibl>|exampleindex}\index{bibl=<bibl>|exampleindex}\index{divGen=<divGen>|exampleindex}\index{n=@n!<divGen>|exampleindex}\index{type=@type!<divGen>|exampleindex}\index{divGen=<divGen>|exampleindex}\index{n=@n!<divGen>|exampleindex}\index{type=@type!<divGen>|exampleindex}\exampleFont \begin{shaded}\noindent\mbox{}{<\textbf{back}>}\mbox{}\newline 
\hspace*{1em}{<\textbf{div}\hspace*{1em}{type}="{appendix}">}\mbox{}\newline 
\hspace*{1em}\hspace*{1em}{<\textbf{head}>}Bibliography{</\textbf{head}>}\mbox{}\newline 
\hspace*{1em}\hspace*{1em}{<\textbf{listBibl}>}\mbox{}\newline 
\hspace*{1em}\hspace*{1em}\hspace*{1em}{<\textbf{bibl}>} ... {</\textbf{bibl}>}\mbox{}\newline 
\hspace*{1em}\hspace*{1em}{</\textbf{listBibl}>}\mbox{}\newline 
\hspace*{1em}{</\textbf{div}>}\mbox{}\newline 
\hspace*{1em}{<\textbf{divGen}\hspace*{1em}{n}="{Index Nominum}"\mbox{}\newline 
\hspace*{1em}\hspace*{1em}{type}="{INDEX-NAMES}"/>}\mbox{}\newline 
\hspace*{1em}{<\textbf{divGen}\hspace*{1em}{n}="{Index Loci}"\hspace*{1em}{type}="{INDEX-PLACES}"/>}\mbox{}\newline 
{</\textbf{back}>}\end{shaded}\egroup\par \par
As this example shows, the global {\itshape n} attribute may also be used to specify a name or identifier for the generated index itself in the usual way. Any additional headings etc. required for the generated index must be specified as content of the \hyperref[TEI.divGen]{<divGen>} element. \par\bgroup\index{back=<back>|exampleindex}\index{divGen=<divGen>|exampleindex}\index{n=@n!<divGen>|exampleindex}\index{type=@type!<divGen>|exampleindex}\index{head=<head>|exampleindex}\exampleFont \begin{shaded}\noindent\mbox{}{<\textbf{back}>}\mbox{}\newline 
\hspace*{1em}{<\textbf{divGen}\hspace*{1em}{n}="{A1}"\hspace*{1em}{type}="{INDEX-NAMES}">}\mbox{}\newline 
\hspace*{1em}\hspace*{1em}{<\textbf{head}>}An Index of Names{</\textbf{head}>}\mbox{}\newline 
\hspace*{1em}{</\textbf{divGen}>}\mbox{}\newline 
{</\textbf{back}>}\end{shaded}\egroup\par \par
If a processing instruction is used, then these parameters for the generated index may be supplied in some other way.\par
One final feature frequently found in manually-created indexes to printed works cannot readily be encoded by the means provided here, namely cross-references internal to the index term listing. For example, if all references to the TEI in a text have been indexed using the index term \textit{Text Encoding Initiative}, it may also be helpful to include an entry under the term \textit{TEI} containing some text such as ‘see Text Encoding Initiative’. Such internal cross-references must be added as part of the post-editing phase for an auto-generated index.
\subsection[{Graphics and Other Non-textual Components}]{Graphics and Other Non-textual Components}\label{COGR}\par
Graphics, such as illustrations or diagrams, appear in many different kinds of text, and often with different purposes. Audio or video clips may also appear. In some cases, such media form an integral part of a text (indeed, some texts—comic books for example—may be almost entirely graphic); in others the graphic or video may be a kind of optional extra. In some cases, the text may be incomprehensible unless the media is included; in others, the presence of the media adds little to the sense of the work. It will therefore be a matter of encoding policy as to whether or how media found in a source text are transferred to a new encoded version of the same. In documents which are ‘born digital’, media such as graphics and other non-textual components may be particularly salient, but their inclusion in an archival form of the document concerned remains an editorial decision.\par
Considered as structural components, media may be anchored to a particular point in the text, or they may \textit{float} either completely freely, or within some defined scope, such as a chapter or section. Time-based media such as audio or video may need to be synchronized with particular parts of a written text. Media of all kinds often contain associated text such as a heading or label. These Guidelines provide the following different elements to indicate their appearance within a text: 
\begin{sansreflist}
  
\item [\textbf{<figure>}] (figure) groups elements representing or containing graphic information such as an illustration, formula, or figure.
\item [\textbf{<media>}] indicates the location of any form of external media such as an audio or video clip etc.
\item [\textbf{<graphic>}] (graphic) indicates the location of a graphic or illustration, either forming part of a text, or providing an image of it.
\item [\textbf{<binaryObject>}] provides encoded binary data representing an inline graphic, audio, video or other object.
\end{sansreflist}
\par
Media files may be encoded in a number of different ways: \begin{itemize}
\item in some non-XML or binary format such as PNG, JPEG, MP3, MP4 etc.
\item in an XML format such as SVG
\item in a TEI XML format such as the notation for graphs and trees described in \textit{\hyperref[GD]{19.\ Graphs, Networks, and Trees}}
\end{itemize}  In the last two cases, the presence of the graphic will be indicated by an appropriate XML element, drawn from the SVG namespace in the second case, and its content will fully define the graphic to be produced. In the first case, however, one of the elements \hyperref[TEI.graphic]{<graphic>} or \hyperref[TEI.media]{<media>} is used to mark the presence of the graphic only and the visual content itself is stored outside the XML document at a location referenced by means of a {\itshape url} attribute. This attribute is provided by membership of these elements in the \textsf{att.resourced} class. Alternatively, if it is small, the media information may be embedded directly within the document using some suitable binary format such as Base64; in this case the \hyperref[TEI.binaryObject]{<binaryObject>} element may be used to contain it.\par
The elements \hyperref[TEI.graphic]{<graphic>}, \hyperref[TEI.media]{<media>}, and \hyperref[TEI.binaryObject]{<binaryObject>} are made available as members of the class \textsf{model.graphicLike} when this module is included in a schema. These elements are also members of the class \textsf{att.media}, from which they inherit the following attributes: 
\begin{sansreflist}
  
\item [\textbf{att.internetMedia}] provides attributes for specifying the type of a computer resource using a standard taxonomy.\hfil\\[-10pt]\begin{sansreflist}
    \item[@{\itshape mimeType}]
  (MIME media type) specifies the applicable multimedia internet mail extension (MIME) media type
\end{sansreflist}  
\item [\textbf{att.media}] provides attributes for specifying display and related properties of external media.\hfil\\[-10pt]\begin{sansreflist}
    \item[@{\itshape width}]
  Where the media are displayed, indicates the display width
    \item[@{\itshape height}]
  Where the media are displayed, indicates the display height
    \item[@{\itshape scale}]
  Where the media are displayed, indicates a scale factor to be applied when generating the desired display size
\end{sansreflist}  
\end{sansreflist}
\par
For example, the following passage indicates that a copy of the image found in the source text may be recovered from the URL \textsf{zigzag2.png} and that this image is in PNG format: \par\bgroup\index{p=<p>|exampleindex}\index{graphic=<graphic>|exampleindex}\index{url=@url!<graphic>|exampleindex}\index{mimeType=@mimeType!<graphic>|exampleindex}\exampleFont \begin{shaded}\noindent\mbox{}{<\textbf{p}>}These were the four lines I moved in\mbox{}\newline 
 through my first, second, third, and\mbox{}\newline 
 fourth volumes. -- In the fifth volume\mbox{}\newline 
 I have been very good, -- the precise\mbox{}\newline 
 line I have described in it being this:\mbox{}\newline 
{<\textbf{graphic}\hspace*{1em}{url}="{zigzag2.png}"\mbox{}\newline 
\hspace*{1em}\hspace*{1em}{mimeType}="{image/png}"/>} \mbox{}\newline 
 By which it appears, that except at the\mbox{}\newline 
 curve, marked A. where I took a trip\mbox{}\newline 
 to Navarre, -- and the indented curve B.\mbox{}\newline 
 which is the short airing when I was\mbox{}\newline 
 there with the Lady Baussiere and her\mbox{}\newline 
 page, -- I have not taken the least frisk\mbox{}\newline 
 ...{</\textbf{p}>}\end{shaded}\egroup\par \par
The media elements are phrase level elements which may be used anywhere that textual content is permitted, within but not between paragraphs or headings. In the following example, the encoder has decided to treat a specific printer's ornament as a heading: \par\bgroup\index{head=<head>|exampleindex}\index{graphic=<graphic>|exampleindex}\index{url=@url!<graphic>|exampleindex}\exampleFont \begin{shaded}\noindent\mbox{}{<\textbf{head}>}\mbox{}\newline 
\hspace*{1em}{<\textbf{graphic}\hspace*{1em}{url}="{http://www.iath.virginia.edu/gants/Ornaments/Heads/hp-ral02.gif}"/>}\mbox{}\newline 
{</\textbf{head}>}\end{shaded}\egroup\par \noindent .\par
The \hyperref[TEI.figure]{<figure>} element discussed in \textit{\hyperref[FTGRA]{14.4.\ Specific Elements for Graphic Images}} provides additional capabilities, for example the ability to combine a number of images into a hierarchically organized structure or a block of images. The \hyperref[TEI.figure]{<figure>} element carries a {\itshape type} attribute, which can be used to distinguish different kinds of graphic component within a single work, for example, maps as opposed to illustrations. It also provides the ability to associate an image with additional information such as a heading or a description.
\subsection[{Reference Systems}]{Reference Systems}\label{CORS}\par
By \textit{reference system} we mean the system by which names or references are associated with particular passages of a text (e.g. \textit{Ps. 23:3} for the third verse of Psalm 23 or \textit{Amores 2.10.7} for Ovid's \textit{Amores}, book 2, poem 10, line 7). Such names make it possible to mark a place within a text and enable other readers to find it again. A reference system may be based on structural units (chapters, paragraphs, sentences; stanza and verse), typographic units (page and line numbers), or divisions created specifically for reference purposes (chapter and verse in Biblical texts). Where one exists, the traditional reference system for a text should be preserved in an electronic transcript of it, if only to make it easier to compare electronic and non-electronic versions of the text.\par
Reference systems may be recorded in TEI-encoded texts in any of the following ways: \begin{itemize}
\item where a reference system exists, and is based on the same logical structure as that of the text's markup, the reference for a passage may be recorded as the value of the global {\itshape xml:id} or {\itshape n} attribute on an appropriate tag, or may be constructed by combining attribute values from several levels of tags, as described below in section \textit{\hyperref[CORS1]{3.11.1.\ Using the xml:id and n Attributes}}.
\item where there is no pre-existing reference system, the global {\itshape xml:id} or {\itshape n} attributes may be used to construct one (e.g. collections and corpora created in electronic form), as described below in section \textit{\hyperref[CORS2]{3.11.2.\ Creating New Reference Systems}}.
\item where a reference system exists which is not based on the same logical structure as that of the text's markup (for example, one based on the page and line numbers of particular editions of the text rather than on the structural divisions of it), any of a variety of methods for encoding the logical structure representing the reference system may be employed, as described in chapter \textit{\hyperref[NH]{20.\ Non-hierarchical Structures}}.
\item where a reference system exists which does not correspond to any particular logical structure, or where the logical structure concerned is of no interest to the encoder except as a means of supporting the referencing system, then references may be encoded by means of \hyperref[TEI.milestone]{<milestone>} elements, which simply mark points in the text at which values in the reference system change, as described below in section \textit{\hyperref[CORS5]{3.11.3.\ Milestone Elements}}.
\end{itemize}  The specific method used to record traditional or new reference systems for a text should be declared in the TEI header, as further described in section \textit{\hyperref[CORS6]{3.11.4.\ Declaring Reference Systems}} and in section \textit{\hyperref[SACR]{16.2.5.\ Canonical References}}.\par
When a text has no pre-existing associated reference system of any kind, these Guidelines recommend as a minimum that at least the page boundaries of the source text be marked using one of the methods outlined in this section. Retaining page breaks in the markup is also recommended for texts which have a detailed reference system of their own. Line breaks in prose texts may be, but need not be, tagged.\footnote{Many encoders find it convenient to retain the line breaks of the original during data entry, to simplify proofreading, but this may be done without inserting a tag for each line break of the original.}
\subsubsection[{Using the xml:id and n Attributes}]{Using the {\itshape xml:id} and {\itshape n} Attributes}\label{CORS1}\par
When traditional reference schemes represent a hierarchical structuring of the text which mirrors that of the marked-up document, the {\itshape n} attribute defined for all elements may be used to indicate the traditional identifier of the relevant structural units. The {\itshape n} attribute may also be used to record the numbering of sections or list items in the copy text if the copy-text numbering is important for some reason, for example because the numbers are out of sequence.\par
For example, a traditional reference to Ovid's \textit{Amores} might be \textit{Amores 2.10.7}—book 2, poem 10, line 7. Book, poem, and line are structural units of the work and will therefore be tagged in any case. (See chapter \textit{\hyperref[VE]{6.\ Verse}} for a discussion of structural units in verse collections.) In such cases, it is convenient to record traditional reference numbers of the structural units using the {\itshape n} attribute. The relevant tags for our example would be: \par\bgroup\index{div1=<div1>|exampleindex}\index{n=@n!<div1>|exampleindex}\index{type=@type!<div1>|exampleindex}\index{div2=<div2>|exampleindex}\index{n=@n!<div2>|exampleindex}\index{type=@type!<div2>|exampleindex}\index{div2=<div2>|exampleindex}\index{n=@n!<div2>|exampleindex}\index{type=@type!<div2>|exampleindex}\index{div3=<div3>|exampleindex}\index{n=@n!<div3>|exampleindex}\index{type=@type!<div3>|exampleindex}\index{div3=<div3>|exampleindex}\index{n=@n!<div3>|exampleindex}\index{type=@type!<div3>|exampleindex}\index{div3=<div3>|exampleindex}\index{n=@n!<div3>|exampleindex}\index{type=@type!<div3>|exampleindex}\index{l=<l>|exampleindex}\index{n=@n!<l>|exampleindex}\index{l=<l>|exampleindex}\index{n=@n!<l>|exampleindex}\index{l=<l>|exampleindex}\index{n=@n!<l>|exampleindex}\exampleFont \begin{shaded}\noindent\mbox{}{<\textbf{div1}\hspace*{1em}{n}="{Amores}"\hspace*{1em}{type}="{volume}">}\mbox{}\newline 
\hspace*{1em}{<\textbf{div2}\hspace*{1em}{n}="{1}"\hspace*{1em}{type}="{book}">}\mbox{}\newline 
\textit{<!-- ... -->}\mbox{}\newline 
\hspace*{1em}{</\textbf{div2}>}\mbox{}\newline 
\hspace*{1em}{<\textbf{div2}\hspace*{1em}{n}="{2}"\hspace*{1em}{type}="{book}">}\mbox{}\newline 
\hspace*{1em}\hspace*{1em}{<\textbf{div3}\hspace*{1em}{n}="{1}"\hspace*{1em}{type}="{poem}">}\mbox{}\newline 
\textit{<!-- ... -->}\mbox{}\newline 
\hspace*{1em}\hspace*{1em}{</\textbf{div3}>}\mbox{}\newline 
\hspace*{1em}\hspace*{1em}{<\textbf{div3}\hspace*{1em}{n}="{2}"\hspace*{1em}{type}="{poem}">}\mbox{}\newline 
\textit{<!-- ... -->}\mbox{}\newline 
\hspace*{1em}\hspace*{1em}{</\textbf{div3}>}\mbox{}\newline 
\textit{<!-- ... -->}\mbox{}\newline 
\hspace*{1em}\hspace*{1em}{<\textbf{div3}\hspace*{1em}{n}="{10}"\hspace*{1em}{type}="{poem}">}\mbox{}\newline 
\hspace*{1em}\hspace*{1em}\hspace*{1em}{<\textbf{l}\hspace*{1em}{n}="{1}">} ... {</\textbf{l}>}\mbox{}\newline 
\hspace*{1em}\hspace*{1em}\hspace*{1em}{<\textbf{l}\hspace*{1em}{n}="{2}">} ... {</\textbf{l}>}\mbox{}\newline 
\textit{<!-- ... -->}\mbox{}\newline 
\hspace*{1em}\hspace*{1em}\hspace*{1em}{<\textbf{l}\hspace*{1em}{n}="{7}">} ... {</\textbf{l}>}\mbox{}\newline 
\hspace*{1em}\hspace*{1em}{</\textbf{div3}>}\mbox{}\newline 
\textit{<!-- ... -->}\mbox{}\newline 
\hspace*{1em}{</\textbf{div2}>}\mbox{}\newline 
\textit{<!-- ... -->}\mbox{}\newline 
{</\textbf{div1}>}\end{shaded}\egroup\par \par
One may also place the entire standard reference for each portion of the text into the appropriate value for the {\itshape n} attribute, though for obvious reasons this takes more space in the file: \par\bgroup\index{div1=<div1>|exampleindex}\index{n=@n!<div1>|exampleindex}\index{type=@type!<div1>|exampleindex}\index{div2=<div2>|exampleindex}\index{n=@n!<div2>|exampleindex}\index{type=@type!<div2>|exampleindex}\index{div2=<div2>|exampleindex}\index{n=@n!<div2>|exampleindex}\index{type=@type!<div2>|exampleindex}\index{div3=<div3>|exampleindex}\index{n=@n!<div3>|exampleindex}\index{type=@type!<div3>|exampleindex}\index{div3=<div3>|exampleindex}\index{n=@n!<div3>|exampleindex}\index{type=@type!<div3>|exampleindex}\index{l=<l>|exampleindex}\index{n=@n!<l>|exampleindex}\exampleFont \begin{shaded}\noindent\mbox{}{<\textbf{div1}\hspace*{1em}{n}="{Amores}"\hspace*{1em}{type}="{volume}">}\mbox{}\newline 
\hspace*{1em}{<\textbf{div2}\hspace*{1em}{n}="{Amores 1}"\hspace*{1em}{type}="{book}">}\mbox{}\newline 
\textit{<!-- ... -->}\mbox{}\newline 
\hspace*{1em}{</\textbf{div2}>}\mbox{}\newline 
\hspace*{1em}{<\textbf{div2}\hspace*{1em}{n}="{Amores 2}"\hspace*{1em}{type}="{book}">}\mbox{}\newline 
\hspace*{1em}\hspace*{1em}{<\textbf{div3}\hspace*{1em}{n}="{Amores 2.1}"\hspace*{1em}{type}="{poem}">}\mbox{}\newline 
\textit{<!-- ... -->}\mbox{}\newline 
\hspace*{1em}\hspace*{1em}{</\textbf{div3}>}\mbox{}\newline 
\textit{<!-- ... -->}\mbox{}\newline 
\hspace*{1em}\hspace*{1em}{<\textbf{div3}\hspace*{1em}{n}="{Amores 2.10}"\hspace*{1em}{type}="{poem}">}\mbox{}\newline 
\textit{<!-- ... -->}\mbox{}\newline 
\hspace*{1em}\hspace*{1em}\hspace*{1em}{<\textbf{l}\hspace*{1em}{n}="{Amores 2.10.7}">} ... {</\textbf{l}>}\mbox{}\newline 
\textit{<!-- ... -->}\mbox{}\newline 
\hspace*{1em}\hspace*{1em}{</\textbf{div3}>}\mbox{}\newline 
\textit{<!-- ... -->}\mbox{}\newline 
\hspace*{1em}{</\textbf{div2}>}\mbox{}\newline 
\textit{<!-- ... -->}\mbox{}\newline 
{</\textbf{div1}>}\end{shaded}\egroup\par \par
If the names used by the traditional reference system can be formulated as identifiers, then the references can be given as values for the {\itshape xml:id} attribute; this requires that the reference be given without internal spaces, begin with a letter or underscore, and contain no characters other than letters, digits, hyphens, underscores, full stops, and the various combining and extender characters, as defined by the XML specification. Unlike values for the {\itshape n} attribute, values for the {\itshape xml:id} attribute must be unique throughout the document. Our example then looks like this: \par\bgroup\index{div1=<div1>|exampleindex}\index{n=@n!<div1>|exampleindex}\index{type=@type!<div1>|exampleindex}\index{div2=<div2>|exampleindex}\index{type=@type!<div2>|exampleindex}\index{div2=<div2>|exampleindex}\index{type=@type!<div2>|exampleindex}\index{div3=<div3>|exampleindex}\index{type=@type!<div3>|exampleindex}\index{div3=<div3>|exampleindex}\index{type=@type!<div3>|exampleindex}\index{l=<l>|exampleindex}\exampleFont \begin{shaded}\noindent\mbox{}{<\textbf{div1}\hspace*{1em}{n}="{Amores}"\hspace*{1em}{type}="{volume}">}\mbox{}\newline 
\hspace*{1em}{<\textbf{div2}\hspace*{1em}{xml:id}="{am.1}"\hspace*{1em}{type}="{book}">}\mbox{}\newline 
\textit{<!-- ... -->}\mbox{}\newline 
\hspace*{1em}{</\textbf{div2}>}\mbox{}\newline 
\hspace*{1em}{<\textbf{div2}\hspace*{1em}{xml:id}="{am.2}"\hspace*{1em}{type}="{book}">}\mbox{}\newline 
\hspace*{1em}\hspace*{1em}{<\textbf{div3}\hspace*{1em}{xml:id}="{am.2.1}"\hspace*{1em}{type}="{poem}">}\mbox{}\newline 
\textit{<!-- ... -->}\mbox{}\newline 
\hspace*{1em}\hspace*{1em}{</\textbf{div3}>}\mbox{}\newline 
\textit{<!-- ... -->}\mbox{}\newline 
\hspace*{1em}\hspace*{1em}{<\textbf{div3}\hspace*{1em}{xml:id}="{am.2.10}"\hspace*{1em}{type}="{poem}">}\mbox{}\newline 
\textit{<!-- ... -->}\mbox{}\newline 
\hspace*{1em}\hspace*{1em}\hspace*{1em}{<\textbf{l}\hspace*{1em}{xml:id}="{am.2.10.7}">} ... {</\textbf{l}>}\mbox{}\newline 
\textit{<!-- ... -->}\mbox{}\newline 
\hspace*{1em}\hspace*{1em}{</\textbf{div3}>}\mbox{}\newline 
\textit{<!-- ... -->}\mbox{}\newline 
\hspace*{1em}{</\textbf{div2}>}\mbox{}\newline 
\textit{<!-- ... -->}\mbox{}\newline 
{</\textbf{div1}>}\end{shaded}\egroup\par \par
To document the usage and to allow automatic processing of these standard references, it is recommended that the TEI header be used to declare whether standard references are recorded in the {\itshape n} or {\itshape xml:id} attributes and which elements may carry standard references or portions of them. For examples of declarations for the reference systems just shown, see section \textit{\hyperref[CORS6]{3.11.4.\ Declaring Reference Systems}}.\par
Using the {\itshape n} attribute one can specify only a single standard referencing system, a limitation not without problems, since some editions may define structural units differently and thus create alternative reference systems. For example, another edition of the \textit{Amores} considers poem 10 a continuation of poem 9, and therefore would specify the same line as \textit{Amores 2.9.31}. In order to record both of these reference systems one could employ any of a variety of methods discussed in chapter \textit{\hyperref[NH]{20.\ Non-hierarchical Structures}}.
\subsubsection[{Creating New Reference Systems}]{Creating New Reference Systems}\label{CORS2}\par
If a text has no canonical reference system of its own, a new custom reference system may be used.\par
The global attributes {\itshape n} and {\itshape xml:id} may be used to assign reference identifiers to segments of the text. Identifiers specified by either attribute apply to the entire element for which they are given. {\itshape xml:id} attributes must be unique within a single document, and {\itshape xml:id} values must begin with a letter. No such restrictions are made on the values of {\itshape n} attributes.\par
Determining a referencing system for a TEI encoding depends on many factors that may either be derived from textual structure, or influenced by extra-textual contingencies such as project and file management concerns. It is important, therefore, that the attribute used, the elements which can bear standard reference identifiers, and the method for constructing standard reference identifiers, should all be declared in the header as described in section \textit{\hyperref[HD54]{2.3.6.\ The Reference System Declaration}}.\par
The Guidelines do not recommend one specific method for creating new referencing systems; however, the rest of this section lists some possibly useful strategies.
\paragraph[{Referencing system derived from markup}]{Referencing system derived from markup}\label{CORS2-1}\par
A new referencing system may be derived from the structure of the electronic text, specifically from the markup of the text. As with any reference system intended for long-term use, it is important to see the reference as an established, unchanging point in the text. Should the text be revised or rearranged, the reference-system identifiers associated with any section of text must stay with that section of text, even if it means the reference numbers fall out of sequence. (A new reference system may always be created beside the old one if out-of-sequence numbers must be avoided.) \par
A convenient method of mechanically generating unique values for {\itshape xml:id} or {\itshape n} attributes based on the structure of the document is to construct, for each element, a \textit{domain-style address} comprising a series of components separated by full stops, with one component for each level of the document hierarchy. Two methods may be used. In the \textit{typed path} form of identifier, each component in the identifier takes the form of an element identifier, a hyphen, and a number, for example \texttt{p-2}. The element name specifies what type of element is to be sought, and the number specifies which occurrence of that element type is to be selected. (The hyphen and number may be omitted if there is only one element of the given type.) In the \textit{untyped path} form of identifier, each component consists of a number, indicating which element in the sequence of nodes at each level is to be selected. To make the resulting identifier a valid XML identifier, it may need to be prefixed with an unchanging alphabetic letter.\par
Identifiers generated with these methods should use the \hyperref[TEI.text]{<text>} element as their starting point, rather than the \hyperref[TEI.TEI]{<TEI>} or \hyperref[TEI.body]{<body>} elements. The \hyperref[TEI.TEI]{<TEI>} element may be taken as a starting point only if identifiers need to be generated for the \hyperref[TEI.teiHeader]{<teiHeader>}, which is not usually the case; using the \hyperref[TEI.body]{<body>} element as a root would prevent assignment of identifiers for the front and back matter. The component corresponding to the root element can be omitted from identifiers, if no confusion will result. In collections and corpora, the component corresponding to the root may be replaced by the unique identifier assigned to the text or sample.\par
In the following example, each element within the \hyperref[TEI.text]{<text>} element has been given a typed-path identifier as its {\itshape xml:id} value, and an untyped-path identifier as its {\itshape n} value; the latter are prefixed with the string \textit{AB}, which may be imagined to be the general identifier for this text. \par\bgroup\index{text=<text>|exampleindex}\index{n=@n!<text>|exampleindex}\index{front=<front>|exampleindex}\index{n=@n!<front>|exampleindex}\index{div=<div>|exampleindex}\index{n=@n!<div>|exampleindex}\index{p=<p>|exampleindex}\index{titlePage=<titlePage>|exampleindex}\index{n=@n!<titlePage>|exampleindex}\index{titlePart=<titlePart>|exampleindex}\index{div=<div>|exampleindex}\index{n=@n!<div>|exampleindex}\index{p=<p>|exampleindex}\index{body=<body>|exampleindex}\index{n=@n!<body>|exampleindex}\index{p=<p>|exampleindex}\index{n=@n!<p>|exampleindex}\index{p=<p>|exampleindex}\index{n=@n!<p>|exampleindex}\index{div=<div>|exampleindex}\index{n=@n!<div>|exampleindex}\index{head=<head>|exampleindex}\index{n=@n!<head>|exampleindex}\index{p=<p>|exampleindex}\index{n=@n!<p>|exampleindex}\index{p=<p>|exampleindex}\index{n=@n!<p>|exampleindex}\index{div=<div>|exampleindex}\index{n=@n!<div>|exampleindex}\index{head=<head>|exampleindex}\index{n=@n!<head>|exampleindex}\index{p=<p>|exampleindex}\index{n=@n!<p>|exampleindex}\index{p=<p>|exampleindex}\index{n=@n!<p>|exampleindex}\exampleFont \begin{shaded}\noindent\mbox{}{<\textbf{text}\hspace*{1em}{xml:id}="{Text-1}"\hspace*{1em}{n}="{AB}">}\mbox{}\newline 
\hspace*{1em}{<\textbf{front}\hspace*{1em}{xml:id}="{Front}"\hspace*{1em}{n}="{AB.1}">}\mbox{}\newline 
\hspace*{1em}\hspace*{1em}{<\textbf{div}\hspace*{1em}{xml:id}="{Front.div-1}"\hspace*{1em}{n}="{AB.1.1}">}\mbox{}\newline 
\hspace*{1em}\hspace*{1em}\hspace*{1em}{<\textbf{p}>} ... {</\textbf{p}>}\mbox{}\newline 
\hspace*{1em}\hspace*{1em}{</\textbf{div}>}\mbox{}\newline 
\hspace*{1em}\hspace*{1em}{<\textbf{titlePage}\hspace*{1em}{xml:id}="{Front.titlePage}"\mbox{}\newline 
\hspace*{1em}\hspace*{1em}\hspace*{1em}{n}="{AB.1.2}">}\mbox{}\newline 
\hspace*{1em}\hspace*{1em}\hspace*{1em}{<\textbf{titlePart}>} ... {</\textbf{titlePart}>}\mbox{}\newline 
\hspace*{1em}\hspace*{1em}{</\textbf{titlePage}>}\mbox{}\newline 
\hspace*{1em}\hspace*{1em}{<\textbf{div}\hspace*{1em}{xml:id}="{Front.div-2}"\hspace*{1em}{n}="{AB.1.3}">}\mbox{}\newline 
\hspace*{1em}\hspace*{1em}\hspace*{1em}{<\textbf{p}>} ... {</\textbf{p}>}\mbox{}\newline 
\hspace*{1em}\hspace*{1em}{</\textbf{div}>}\mbox{}\newline 
\hspace*{1em}{</\textbf{front}>}\mbox{}\newline 
\hspace*{1em}{<\textbf{body}\hspace*{1em}{xml:id}="{Body}"\hspace*{1em}{n}="{AB.2}">}\mbox{}\newline 
\hspace*{1em}\hspace*{1em}{<\textbf{p}\hspace*{1em}{xml:id}="{Body.p-1}"\hspace*{1em}{n}="{AB.2.1}">} ... {</\textbf{p}>}\mbox{}\newline 
\hspace*{1em}\hspace*{1em}{<\textbf{p}\hspace*{1em}{xml:id}="{Body.p-2}"\hspace*{1em}{n}="{AB.2.2}">} ... {</\textbf{p}>}\mbox{}\newline 
\hspace*{1em}\hspace*{1em}{<\textbf{div}\hspace*{1em}{xml:id}="{Body.div-1}"\hspace*{1em}{n}="{AB.2.3}">}\mbox{}\newline 
\hspace*{1em}\hspace*{1em}\hspace*{1em}{<\textbf{head}\hspace*{1em}{xml:id}="{Body.div-1.head}"\mbox{}\newline 
\hspace*{1em}\hspace*{1em}\hspace*{1em}\hspace*{1em}{n}="{AB.2.3.1}">} ... {</\textbf{head}>}\mbox{}\newline 
\hspace*{1em}\hspace*{1em}\hspace*{1em}{<\textbf{p}\hspace*{1em}{xml:id}="{Body.div-1.p-1}"\hspace*{1em}{n}="{AB.2.3.2}">} ... {</\textbf{p}>}\mbox{}\newline 
\hspace*{1em}\hspace*{1em}\hspace*{1em}{<\textbf{p}\hspace*{1em}{xml:id}="{Body.div-1.p-2}"\hspace*{1em}{n}="{AB.2.3.3}">} ... {</\textbf{p}>}\mbox{}\newline 
\hspace*{1em}\hspace*{1em}{</\textbf{div}>}\mbox{}\newline 
\hspace*{1em}\hspace*{1em}{<\textbf{div}\hspace*{1em}{xml:id}="{Body.div-2}"\hspace*{1em}{n}="{AB.2.4}">}\mbox{}\newline 
\hspace*{1em}\hspace*{1em}\hspace*{1em}{<\textbf{head}\hspace*{1em}{xml:id}="{Body.div-2.head}"\mbox{}\newline 
\hspace*{1em}\hspace*{1em}\hspace*{1em}\hspace*{1em}{n}="{AB.2.4.1}">} ... {</\textbf{head}>}\mbox{}\newline 
\hspace*{1em}\hspace*{1em}\hspace*{1em}{<\textbf{p}\hspace*{1em}{xml:id}="{Body.div-2.p-1}"\hspace*{1em}{n}="{AB.2.4.2}">} ... {</\textbf{p}>}\mbox{}\newline 
\hspace*{1em}\hspace*{1em}\hspace*{1em}{<\textbf{p}\hspace*{1em}{xml:id}="{Body.div-2.p-2}"\hspace*{1em}{n}="{AB.2.4.3}">} ... {</\textbf{p}>}\mbox{}\newline 
\hspace*{1em}\hspace*{1em}{</\textbf{div}>}\mbox{}\newline 
\hspace*{1em}{</\textbf{body}>}\mbox{}\newline 
{</\textbf{text}>}\end{shaded}\egroup\par \noindent  The typed and untyped path methods are convenient, but are in no way required for anyone creating a reference system.\par
If the {\itshape xml:id} attribute is used to record the reference identifiers generated, each value should record the entire path. If the {\itshape n} attribute is used, each value may record either the entire path or only the subpath from the parent element. The attribute used, the elements which can bear standard reference identifiers, and the method for constructing standard reference identifiers, should all be declared in the header as described in section \textit{\hyperref[HD54]{2.3.6.\ The Reference System Declaration}}.
\paragraph[{Referencing systems based on project conventions}]{Referencing systems based on project conventions}\label{CORS2-2}\par
A reference system may be based on an agreed project-specific convention for {\itshape xml:id} attributes. Every convention will have strengths and weaknesses and it is left to encoders to make a decision that enables them to locate information in their TEI document.\par
Here are some examples of referencing systems that have been used in TEI project: \begin{itemize}
\item \textbf{Title-based identifiers:} identifiers constructed with a number of characters from the main document title, followed by an incremental number. E.g. HOL001, HOL002, etc. using a fixed number of digits; or without fixed digits: HOL1, HOL2, etc.
\item \textbf{Based on markup, with prefix:} identifiers constructed on the markup itself, as described in the previous section. To facilitate uniqueness in a corpus, each identifier may be prefixed with the identifier of the root \hyperref[TEI.TEI]{<TEI>} element. E.g. RootID-Body-p-1.
\item \textbf{Opaque identifiers:} computed identifiers using either a randomized algorithm or a universally unique identifier (UUID) algorithm. Note that XSLT's function generate-id() only guarantees identifier unique to the document being processed.
\end{itemize} \par
XML well-formedness requires only that xml:id attributes be unique within a single document. However, it is also worth keeping in mind that for operating with referencing systems across a corpus of TEI files it is helpful (or even necessary in some circumstances) to have unique identifiers across the whole corpus.\par
Values of {\itshape xml:id} may be either populated computationally or manually. In the latter case, it is advisable to put measures in place to avoid human error. Custom data types and Schematron rules may be defined in a customization ODD, and a check digit may be added to prevent unwanted changes. \footnote{A check digit is computed from the value of an identifier and appended to the value itself. If the identifier is changed, the check digit would therefore invalidate it.}
\subsubsection[{Milestone Elements}]{Milestone Elements}\label{CORS5}\par
Where the desired reference system does not correspond to any particular structural hierarchy, or the document combines multiple structural hierarchies (as further discussed in \textit{\hyperref[NH]{20.\ Non-hierarchical Structures}}), simpler though less expressive methods may be necessary. In such cases the simplest solution may be just to mark up changes in the reference system where they occur, by using one or more of the following \textit{milestone} elements: 
\begin{sansreflist}
  
\item [\textbf{<milestone>}] (milestone) marks a boundary point separating any kind of section of a text, typically but not necessarily indicating a point at which some part of a standard reference system changes, where the change is not represented by a structural element.
\item [\textbf{<gb>}] (gathering beginning) marks the beginning of a new gathering or quire in a transcribed codex.
\item [\textbf{<pb>}] (page beginning) marks the beginning of a new page in a paginated document.
\item [\textbf{<lb>}] (line beginning) marks the beginning of a new (typographic) line in some edition or version of a text.
\item [\textbf{<cb>}] (column beginning) marks the beginning of a new column of a text on a multi-column page.
\end{sansreflist}
\par
These elements simply mark the points in a text at which some category in a reference system changes. They have no content but subdivide the text into regions, rather in the same way as milestones mark points along a road, thus implicitly dividing it into segments. The elements \hyperref[TEI.gb]{<gb>}, \hyperref[TEI.pb]{<pb>}, \hyperref[TEI.cb]{<cb>}, and \hyperref[TEI.lb]{<lb>} are specialized types of milestone, marking gathering, page, column, and line boundaries respectively. The global {\itshape n} attribute is used in each case to provide a value for the particular unit associated with this milestone (for example, the page or line number). Since it is not structural, validation of a reference system based on \hyperref[TEI.milestone]{<milestone>}s cannot readily be checked by an XML parser, so it will be the responsibility of the encoder or the application software to ensure that they are given in the correct order.\par
Milestone elements are often used as a simple means of capturing the original appearance of an early printed text, which will rarely coincide exactly with structural units, but they are generally useful wherever a text has two or more competing structures. For example, many English novels were first published as serial works, individual parts of which do not always contain a whole number of chapters. An encoder might decide to represent the chapter-based structure using \hyperref[TEI.div1]{<div1>} elements, with \hyperref[TEI.milestone]{<milestone>} elements to mark the points at which individual parts end; or the reverse. Thus, an encoding in which chapters are regarded as more important than parts might encode some work in which chapter three begins in part one and is concluded in part two as follows: \par\bgroup\index{text=<text>|exampleindex}\index{body=<body>|exampleindex}\index{milestone=<milestone>|exampleindex}\index{unit=@unit!<milestone>|exampleindex}\index{n=@n!<milestone>|exampleindex}\index{div1=<div1>|exampleindex}\index{n=@n!<div1>|exampleindex}\index{type=@type!<div1>|exampleindex}\index{p=<p>|exampleindex}\index{div1=<div1>|exampleindex}\index{n=@n!<div1>|exampleindex}\index{type=@type!<div1>|exampleindex}\index{p=<p>|exampleindex}\index{div1=<div1>|exampleindex}\index{n=@n!<div1>|exampleindex}\index{type=@type!<div1>|exampleindex}\index{p=<p>|exampleindex}\index{milestone=<milestone>|exampleindex}\index{unit=@unit!<milestone>|exampleindex}\index{n=@n!<milestone>|exampleindex}\index{p=<p>|exampleindex}\exampleFont \begin{shaded}\noindent\mbox{}{<\textbf{text}>}\mbox{}\newline 
\hspace*{1em}{<\textbf{body}>}\mbox{}\newline 
\hspace*{1em}\hspace*{1em}{<\textbf{milestone}\hspace*{1em}{unit}="{part}"\hspace*{1em}{n}="{1}"/>}\mbox{}\newline 
\hspace*{1em}\hspace*{1em}{<\textbf{div1}\hspace*{1em}{n}="{1}"\hspace*{1em}{type}="{chapter}">}\mbox{}\newline 
\hspace*{1em}\hspace*{1em}\hspace*{1em}{<\textbf{p}>}\mbox{}\newline 
\textit{<!-- ... -->}\mbox{}\newline 
\hspace*{1em}\hspace*{1em}\hspace*{1em}{</\textbf{p}>}\mbox{}\newline 
\hspace*{1em}\hspace*{1em}{</\textbf{div1}>}\mbox{}\newline 
\hspace*{1em}\hspace*{1em}{<\textbf{div1}\hspace*{1em}{n}="{2}"\hspace*{1em}{type}="{chapter}">}\mbox{}\newline 
\hspace*{1em}\hspace*{1em}\hspace*{1em}{<\textbf{p}>}\mbox{}\newline 
\textit{<!-- ... -->}\mbox{}\newline 
\hspace*{1em}\hspace*{1em}\hspace*{1em}{</\textbf{p}>}\mbox{}\newline 
\hspace*{1em}\hspace*{1em}{</\textbf{div1}>}\mbox{}\newline 
\hspace*{1em}\hspace*{1em}{<\textbf{div1}\hspace*{1em}{n}="{3}"\hspace*{1em}{type}="{chapter}">}\mbox{}\newline 
\hspace*{1em}\hspace*{1em}\hspace*{1em}{<\textbf{p}>}\mbox{}\newline 
\textit{<!-- ... -->}\mbox{}\newline 
\hspace*{1em}\hspace*{1em}\hspace*{1em}{</\textbf{p}>}\mbox{}\newline 
\hspace*{1em}\hspace*{1em}\hspace*{1em}{<\textbf{milestone}\hspace*{1em}{unit}="{part}"\hspace*{1em}{n}="{2}"/>}\mbox{}\newline 
\hspace*{1em}\hspace*{1em}\hspace*{1em}{<\textbf{p}>}\mbox{}\newline 
\textit{<!-- ... -->}\mbox{}\newline 
\hspace*{1em}\hspace*{1em}\hspace*{1em}{</\textbf{p}>}\mbox{}\newline 
\hspace*{1em}\hspace*{1em}{</\textbf{div1}>}\mbox{}\newline 
\hspace*{1em}{</\textbf{body}>}\mbox{}\newline 
{</\textbf{text}>}\end{shaded}\egroup\par \noindent  An encoding of the same work in which parts are regarded as more important than chapters might begin as follows: \par\bgroup\index{text=<text>|exampleindex}\index{body=<body>|exampleindex}\index{div1=<div1>|exampleindex}\index{n=@n!<div1>|exampleindex}\index{type=@type!<div1>|exampleindex}\index{milestone=<milestone>|exampleindex}\index{unit=@unit!<milestone>|exampleindex}\index{n=@n!<milestone>|exampleindex}\index{p=<p>|exampleindex}\index{milestone=<milestone>|exampleindex}\index{unit=@unit!<milestone>|exampleindex}\index{n=@n!<milestone>|exampleindex}\index{p=<p>|exampleindex}\index{milestone=<milestone>|exampleindex}\index{unit=@unit!<milestone>|exampleindex}\index{n=@n!<milestone>|exampleindex}\index{p=<p>|exampleindex}\index{div1=<div1>|exampleindex}\index{n=@n!<div1>|exampleindex}\index{type=@type!<div1>|exampleindex}\index{p=<p>|exampleindex}\index{milestone=<milestone>|exampleindex}\index{unit=@unit!<milestone>|exampleindex}\index{n=@n!<milestone>|exampleindex}\index{p=<p>|exampleindex}\exampleFont \begin{shaded}\noindent\mbox{}{<\textbf{text}>}\mbox{}\newline 
\hspace*{1em}{<\textbf{body}>}\mbox{}\newline 
\hspace*{1em}\hspace*{1em}{<\textbf{div1}\hspace*{1em}{n}="{1}"\hspace*{1em}{type}="{part}">}\mbox{}\newline 
\hspace*{1em}\hspace*{1em}\hspace*{1em}{<\textbf{milestone}\hspace*{1em}{unit}="{chapter}"\hspace*{1em}{n}="{1}"/>}\mbox{}\newline 
\hspace*{1em}\hspace*{1em}\hspace*{1em}{<\textbf{p}>}\mbox{}\newline 
\textit{<!-- ... -->}\mbox{}\newline 
\hspace*{1em}\hspace*{1em}\hspace*{1em}{</\textbf{p}>}\mbox{}\newline 
\hspace*{1em}\hspace*{1em}\hspace*{1em}{<\textbf{milestone}\hspace*{1em}{unit}="{chapter}"\hspace*{1em}{n}="{2}"/>}\mbox{}\newline 
\hspace*{1em}\hspace*{1em}\hspace*{1em}{<\textbf{p}>}\mbox{}\newline 
\textit{<!-- ... -->}\mbox{}\newline 
\hspace*{1em}\hspace*{1em}\hspace*{1em}{</\textbf{p}>}\mbox{}\newline 
\hspace*{1em}\hspace*{1em}\hspace*{1em}{<\textbf{milestone}\hspace*{1em}{unit}="{chapter}"\hspace*{1em}{n}="{3}"/>}\mbox{}\newline 
\hspace*{1em}\hspace*{1em}\hspace*{1em}{<\textbf{p}>}\mbox{}\newline 
\textit{<!-- ... -->}\mbox{}\newline 
\hspace*{1em}\hspace*{1em}\hspace*{1em}{</\textbf{p}>}\mbox{}\newline 
\hspace*{1em}\hspace*{1em}{</\textbf{div1}>}\mbox{}\newline 
\hspace*{1em}\hspace*{1em}{<\textbf{div1}\hspace*{1em}{n}="{2}"\hspace*{1em}{type}="{part}">}\mbox{}\newline 
\hspace*{1em}\hspace*{1em}\hspace*{1em}{<\textbf{p}>}\mbox{}\newline 
\textit{<!-- ... -->}\mbox{}\newline 
\hspace*{1em}\hspace*{1em}\hspace*{1em}{</\textbf{p}>}\mbox{}\newline 
\hspace*{1em}\hspace*{1em}\hspace*{1em}{<\textbf{milestone}\hspace*{1em}{unit}="{chapter}"\hspace*{1em}{n}="{4}"/>}\mbox{}\newline 
\hspace*{1em}\hspace*{1em}\hspace*{1em}{<\textbf{p}>}\mbox{}\newline 
\textit{<!-- ... -->}\mbox{}\newline 
\hspace*{1em}\hspace*{1em}\hspace*{1em}{</\textbf{p}>}\mbox{}\newline 
\hspace*{1em}\hspace*{1em}{</\textbf{div1}>}\mbox{}\newline 
\hspace*{1em}{</\textbf{body}>}\mbox{}\newline 
{</\textbf{text}>}\end{shaded}\egroup\par \par
Similarly, when tagging dramatic verse one may wish to privilege stanzas and lines over speeches and speakers, particularly where speeches cross line and line group boundaries. One might also wish to mark changes in narrative voice in a prose text. In either case, a milestone tag may be used to indicate change of speaker: \par\bgroup\index{lg=<lg>|exampleindex}\index{milestone=<milestone>|exampleindex}\index{unit=@unit!<milestone>|exampleindex}\index{n=@n!<milestone>|exampleindex}\index{l=<l>|exampleindex}\index{l=<l>|exampleindex}\index{milestone=<milestone>|exampleindex}\index{unit=@unit!<milestone>|exampleindex}\index{n=@n!<milestone>|exampleindex}\index{l=<l>|exampleindex}\index{l=<l>|exampleindex}\exampleFont \begin{shaded}\noindent\mbox{}{<\textbf{lg}>}\mbox{}\newline 
\hspace*{1em}{<\textbf{milestone}\hspace*{1em}{unit}="{speaker}"\hspace*{1em}{n}="{Man}"/>}\mbox{}\newline 
\hspace*{1em}{<\textbf{l}>}Oh what is this I cannot see{</\textbf{l}>}\mbox{}\newline 
\hspace*{1em}{<\textbf{l}>}With icy hands gets a hold on me{</\textbf{l}>}\mbox{}\newline 
\hspace*{1em}{<\textbf{milestone}\hspace*{1em}{unit}="{speaker}"\hspace*{1em}{n}="{Death}"/>}\mbox{}\newline 
\hspace*{1em}{<\textbf{l}>}Oh I am Death, none can excel{</\textbf{l}>}\mbox{}\newline 
\hspace*{1em}{<\textbf{l}>}I open the doors of heaven and hell{</\textbf{l}>}\mbox{}\newline 
{</\textbf{lg}>}\end{shaded}\egroup\par \par
Milestone tags also make it possible to record the reference systems used in a number of different editions of the same work. The reference system of any one edition can be recreated from a text in which all are marked by simply ignoring all elements that do not specify that edition on their {\itshape ed} attribute.\par
As a simple example, assuming that edition E1 of some collection of poems regards the first two poems as constituting the first book, while edition E2 regards the first poem as prefatory, a markup scheme like the following might be adopted: \par\bgroup\index{milestone=<milestone>|exampleindex}\index{ed=@ed!<milestone>|exampleindex}\index{unit=@unit!<milestone>|exampleindex}\index{milestone=<milestone>|exampleindex}\index{ed=@ed!<milestone>|exampleindex}\index{unit=@unit!<milestone>|exampleindex}\index{milestone=<milestone>|exampleindex}\index{ed=@ed!<milestone>|exampleindex}\index{unit=@unit!<milestone>|exampleindex}\index{milestone=<milestone>|exampleindex}\index{ed=@ed!<milestone>|exampleindex}\index{unit=@unit!<milestone>|exampleindex}\index{milestone=<milestone>|exampleindex}\index{ed=@ed!<milestone>|exampleindex}\index{unit=@unit!<milestone>|exampleindex}\index{milestone=<milestone>|exampleindex}\index{ed=@ed!<milestone>|exampleindex}\index{unit=@unit!<milestone>|exampleindex}\index{milestone=<milestone>|exampleindex}\index{ed=@ed!<milestone>|exampleindex}\index{unit=@unit!<milestone>|exampleindex}\index{milestone=<milestone>|exampleindex}\index{ed=@ed!<milestone>|exampleindex}\index{unit=@unit!<milestone>|exampleindex}\exampleFont \begin{shaded}\noindent\mbox{}{<\textbf{milestone}\hspace*{1em}{ed}="{E1}"\hspace*{1em}{unit}="{work}"/>}\mbox{}\newline 
{<\textbf{milestone}\hspace*{1em}{ed}="{E2}"\hspace*{1em}{unit}="{work}"/>}\mbox{}\newline 
{<\textbf{milestone}\hspace*{1em}{ed}="{E1}"\hspace*{1em}{unit}="{book}"/>}\mbox{}\newline 
{<\textbf{milestone}\hspace*{1em}{ed}="{E1}"\hspace*{1em}{unit}="{poem}"/>}\mbox{}\newline 
{<\textbf{milestone}\hspace*{1em}{ed}="{E2}"\hspace*{1em}{unit}="{poem}"/>}\mbox{}\newline 
{<\textbf{milestone}\hspace*{1em}{ed}="{E2}"\hspace*{1em}{unit}="{book}"/>}\mbox{}\newline 
{<\textbf{milestone}\hspace*{1em}{ed}="{E1}"\hspace*{1em}{unit}="{poem}"/>}\mbox{}\newline 
{<\textbf{milestone}\hspace*{1em}{ed}="{E2}"\hspace*{1em}{unit}="{poem}"/>}\end{shaded}\egroup\par \par
In this case no {\itshape n} value is specified, since the numbers rise predictably and the application can keep a count from the start of the document, if desired.\par
The value of the {\itshape n} attribute may but need not include the identifiers used for any larger sections. That is, either of the following styles is legitimate: \par\bgroup\index{milestone=<milestone>|exampleindex}\index{ed=@ed!<milestone>|exampleindex}\index{unit=@unit!<milestone>|exampleindex}\index{n=@n!<milestone>|exampleindex}\index{milestone=<milestone>|exampleindex}\index{ed=@ed!<milestone>|exampleindex}\index{unit=@unit!<milestone>|exampleindex}\index{n=@n!<milestone>|exampleindex}\index{milestone=<milestone>|exampleindex}\index{ed=@ed!<milestone>|exampleindex}\index{unit=@unit!<milestone>|exampleindex}\index{n=@n!<milestone>|exampleindex}\index{milestone=<milestone>|exampleindex}\index{ed=@ed!<milestone>|exampleindex}\index{unit=@unit!<milestone>|exampleindex}\index{n=@n!<milestone>|exampleindex}\index{milestone=<milestone>|exampleindex}\index{ed=@ed!<milestone>|exampleindex}\index{unit=@unit!<milestone>|exampleindex}\index{n=@n!<milestone>|exampleindex}\exampleFont \begin{shaded}\noindent\mbox{}{<\textbf{milestone}\hspace*{1em}{ed}="{E1}"\hspace*{1em}{unit}="{work}"\hspace*{1em}{n}="{Amores}"/>}\mbox{}\newline 
{<\textbf{milestone}\hspace*{1em}{ed}="{E1}"\hspace*{1em}{unit}="{book}"\hspace*{1em}{n}="{1}"/>}\mbox{}\newline 
{<\textbf{milestone}\hspace*{1em}{ed}="{E1}"\hspace*{1em}{unit}="{poem}"\hspace*{1em}{n}="{1}"/>}\mbox{}\newline 
{<\textbf{milestone}\hspace*{1em}{ed}="{E1}"\hspace*{1em}{unit}="{poem}"\hspace*{1em}{n}="{2}"/>}\mbox{}\newline 
{<\textbf{milestone}\hspace*{1em}{ed}="{E1}"\hspace*{1em}{unit}="{book}"\hspace*{1em}{n}="{2}"/>}\end{shaded}\egroup\par \noindent  or \par\bgroup\index{milestone=<milestone>|exampleindex}\index{ed=@ed!<milestone>|exampleindex}\index{unit=@unit!<milestone>|exampleindex}\index{n=@n!<milestone>|exampleindex}\index{milestone=<milestone>|exampleindex}\index{ed=@ed!<milestone>|exampleindex}\index{unit=@unit!<milestone>|exampleindex}\index{n=@n!<milestone>|exampleindex}\index{milestone=<milestone>|exampleindex}\index{ed=@ed!<milestone>|exampleindex}\index{unit=@unit!<milestone>|exampleindex}\index{n=@n!<milestone>|exampleindex}\index{milestone=<milestone>|exampleindex}\index{ed=@ed!<milestone>|exampleindex}\index{unit=@unit!<milestone>|exampleindex}\index{n=@n!<milestone>|exampleindex}\index{milestone=<milestone>|exampleindex}\index{ed=@ed!<milestone>|exampleindex}\index{unit=@unit!<milestone>|exampleindex}\index{n=@n!<milestone>|exampleindex}\exampleFont \begin{shaded}\noindent\mbox{}{<\textbf{milestone}\hspace*{1em}{ed}="{E1}"\hspace*{1em}{unit}="{work}"\hspace*{1em}{n}="{Amores}"/>}\mbox{}\newline 
{<\textbf{milestone}\hspace*{1em}{ed}="{E1}"\hspace*{1em}{unit}="{book}"\hspace*{1em}{n}="{1}"/>}\mbox{}\newline 
{<\textbf{milestone}\hspace*{1em}{ed}="{E1}"\hspace*{1em}{unit}="{poem}"\hspace*{1em}{n}="{1.1}"/>}\mbox{}\newline 
{<\textbf{milestone}\hspace*{1em}{ed}="{E1}"\hspace*{1em}{unit}="{poem}"\hspace*{1em}{n}="{1.2}"/>}\mbox{}\newline 
{<\textbf{milestone}\hspace*{1em}{ed}="{E1}"\hspace*{1em}{unit}="{book}"\hspace*{1em}{n}="{2}"/>}\end{shaded}\egroup\par \par
When using \hyperref[TEI.milestone]{<milestone>} tags, line numbers may be supplied for every line or only periodically (every fifth, every tenth line). The latter may be simpler; the former is more reliable.\par
The style of numbering used in the values of {\itshape n} is unrestricted: for the example above, I.i, I.ii, and I.iii could have been used equally well if preferred. The special value unnumbered should be reserved for marking sections of text which fall outside the normal numbering system (e.g. chapter heads, poem numbers, titles, or speaker attributions in a verse drama).\par
By default, there are no constraints on the values supplied for the {\itshape ed} attribute. If it is felt appropriate to enforce such a restriction, the techniques described in \textit{\hyperref[MD]{23.3.\ Customization}} may be used, for example to specify that the attribute must specify one of a predefined set of values.\par
See below, section \textit{\hyperref[CORS6]{3.11.4.\ Declaring Reference Systems}}, for examples of declarations for the reference systems just shown.\par
Milestone elements may be used to mark any kind of shift in the properties associated with a piece of text, whether or not would normally be considered a reference system. For example, they may be used to mark changes in narrative voice in a prose text, or changes of speaker in a dramatic text, where these are not marked using structural elements such as \hyperref[TEI.sp]{<sp>}, perhaps in order to avoid a clash of hierarchies.\par
As noted in \textit{\hyperref[COPU-2]{3.2.2.\ Hyphenation}} above, milestone elements such as \hyperref[TEI.lb]{<lb>} or \hyperref[TEI.pb]{<pb>} represent whitespace and are therefore by default assumed to occur between orthographic tokens in the text, where these are not otherwise indicated. By default it is reasonable to assume that words are not broken across page or line boundaries, and that therefore a sequence such as \par\bgroup\index{lb=<lb>|exampleindex}\exampleFont \begin{shaded}\noindent\mbox{}...sed imp{<\textbf{lb}/>}erator dixit...\mbox{}\newline 
\end{shaded}\egroup\par \noindent  should be tokenized as four words (\textit{sed}, \textit{imp}, \textit{erator}, and \textit{dixit}). The {\itshape break} attribute is provided to change the default assumption. To make explicit that \textit{imperator} in the above example should be treated as a single word, a tagging such as the following is recommended: \par\bgroup\index{lb=<lb>|exampleindex}\index{break=@break!<lb>|exampleindex}\exampleFont \begin{shaded}\noindent\mbox{}...sed imp{<\textbf{lb}\hspace*{1em}{break}="{no}"/>}erator dixit...\mbox{}\newline 
\end{shaded}\egroup\par \noindent  Where hyphenation appears before a line or page break, the encoder may or may not choose to record the fact, either explicitly using an appropriate Unicode character, or descriptively for example by means of the {\itshape rend} attribute; see further \textit{\hyperref[COPU-2]{3.2.2.\ Hyphenation}}.
\subsubsection[{Declaring Reference Systems}]{Declaring Reference Systems}\label{CORS6}\par
Whatever kind of reference system is used in an electronic text, it is recommended that the TEI header contain a description of its construction in the \hyperref[TEI.refsDecl]{<refsDecl>} element described in section \textit{\hyperref[HD54]{2.3.6.\ The Reference System Declaration}}. As described there, the declaration may consist either of a formal declaration using the \hyperref[TEI.cRefPattern]{<cRefPattern>} or \hyperref[TEI.citeStructure]{<citeStructure>} elements, or an informal description in prose. One of the former is recommended because unlike prose they can be processed by software.\par
The three examples given in section \textit{\hyperref[CORS1]{3.11.1.\ Using the xml:id and n Attributes}} would be declared as follows. The first example encodes the standard references for Ovid's \textit{Amores} one level at a time, using the {\itshape n} attribute on the \hyperref[TEI.div1]{<div1>}, \hyperref[TEI.div2]{<div2>}, \hyperref[TEI.div3]{<div3>}, and \hyperref[TEI.l]{<l>} tags. The header section for such an encoding should look something like this: \par\bgroup\index{encodingDesc=<encodingDesc>|exampleindex}\index{refsDecl=<refsDecl>|exampleindex}\index{cRefPattern=<cRefPattern>|exampleindex}\index{matchPattern=@matchPattern!<cRefPattern>|exampleindex}\index{replacementPattern=@replacementPattern!<cRefPattern>|exampleindex}\index{p=<p>|exampleindex}\index{list=<list>|exampleindex}\index{item=<item>|exampleindex}\index{label=<label>|exampleindex}\index{att=<att>|exampleindex}\index{gi=<gi>|exampleindex}\index{item=<item>|exampleindex}\index{item=<item>|exampleindex}\index{label=<label>|exampleindex}\index{att=<att>|exampleindex}\index{gi=<gi>|exampleindex}\index{item=<item>|exampleindex}\index{item=<item>|exampleindex}\index{label=<label>|exampleindex}\index{att=<att>|exampleindex}\index{gi=<gi>|exampleindex}\index{item=<item>|exampleindex}\index{att=<att>|exampleindex}\index{gi=<gi>|exampleindex}\index{cRefPattern=<cRefPattern>|exampleindex}\index{matchPattern=@matchPattern!<cRefPattern>|exampleindex}\index{replacementPattern=@replacementPattern!<cRefPattern>|exampleindex}\index{p=<p>|exampleindex}\index{gi=<gi>|exampleindex}\index{att=<att>|exampleindex}\index{cRefPattern=<cRefPattern>|exampleindex}\index{matchPattern=@matchPattern!<cRefPattern>|exampleindex}\index{replacementPattern=@replacementPattern!<cRefPattern>|exampleindex}\index{p=<p>|exampleindex}\index{gi=<gi>|exampleindex}\index{att=<att>|exampleindex}\exampleFont \begin{shaded}\noindent\mbox{}{<\textbf{encodingDesc}>}\mbox{}\newline 
\hspace*{1em}{<\textbf{refsDecl}>}\mbox{}\newline 
\hspace*{1em}\hspace*{1em}{<\textbf{cRefPattern}\hspace*{1em}{matchPattern}="{([\textasciicircum  ]+) ([0-9]+)⃥.([0-9]+)⃥.([0-9]+)}"\mbox{}\newline 
\hspace*{1em}\hspace*{1em}\hspace*{1em}{replacementPattern}="{\#xpath(//div1[@n='\$1']/div2[@n='\$2']/div3[@n='\$3']/l[@n='\$4']}">}\mbox{}\newline 
\hspace*{1em}\hspace*{1em}\hspace*{1em}{<\textbf{p}>}A canonical reference is assembled with\mbox{}\newline 
\hspace*{1em}\hspace*{1em}\hspace*{1em}{<\textbf{list}>}\mbox{}\newline 
\hspace*{1em}\hspace*{1em}\hspace*{1em}\hspace*{1em}\hspace*{1em}{<\textbf{item}>}the name of the {<\textbf{label}>}work{</\textbf{label}>}: the\mbox{}\newline 
\hspace*{1em}\hspace*{1em}\hspace*{1em}\hspace*{1em}\hspace*{1em}{<\textbf{att}>}n{</\textbf{att}>} of a {<\textbf{gi}>}div1{</\textbf{gi}>},{</\textbf{item}>}\mbox{}\newline 
\hspace*{1em}\hspace*{1em}\hspace*{1em}\hspace*{1em}\hspace*{1em}{<\textbf{item}>}a space,{</\textbf{item}>}\mbox{}\newline 
\hspace*{1em}\hspace*{1em}\hspace*{1em}\hspace*{1em}\hspace*{1em}{<\textbf{item}>}the number of the {<\textbf{label}>}book{</\textbf{label}>}: the\mbox{}\newline 
\hspace*{1em}\hspace*{1em}\hspace*{1em}\hspace*{1em}\hspace*{1em}{<\textbf{att}>}n{</\textbf{att}>} of a child {<\textbf{gi}>}div2{</\textbf{gi}>},{</\textbf{item}>}\mbox{}\newline 
\hspace*{1em}\hspace*{1em}\hspace*{1em}\hspace*{1em}\hspace*{1em}{<\textbf{item}>}a full stop{</\textbf{item}>}\mbox{}\newline 
\hspace*{1em}\hspace*{1em}\hspace*{1em}\hspace*{1em}\hspace*{1em}{<\textbf{item}>}the number of the {<\textbf{label}>}poem{</\textbf{label}>}: the\mbox{}\newline 
\hspace*{1em}\hspace*{1em}\hspace*{1em}\hspace*{1em}\hspace*{1em}{<\textbf{att}>}n{</\textbf{att}>} of a child {<\textbf{gi}>}div3{</\textbf{gi}>},{</\textbf{item}>}\mbox{}\newline 
\hspace*{1em}\hspace*{1em}\hspace*{1em}\hspace*{1em}\hspace*{1em}{<\textbf{item}>}the line number: the {<\textbf{att}>}n{</\textbf{att}>} value of a\mbox{}\newline 
\hspace*{1em}\hspace*{1em}\hspace*{1em}\hspace*{1em}\hspace*{1em}\hspace*{1em}\hspace*{1em}\hspace*{1em}\hspace*{1em}\hspace*{1em} child {<\textbf{gi}>}l{</\textbf{gi}>}\mbox{}\newline 
\hspace*{1em}\hspace*{1em}\hspace*{1em}\hspace*{1em}\hspace*{1em}{</\textbf{item}>}\mbox{}\newline 
\hspace*{1em}\hspace*{1em}\hspace*{1em}\hspace*{1em}{</\textbf{list}>}\mbox{}\newline 
\hspace*{1em}\hspace*{1em}\hspace*{1em}{</\textbf{p}>}\mbox{}\newline 
\hspace*{1em}\hspace*{1em}{</\textbf{cRefPattern}>}\mbox{}\newline 
\hspace*{1em}\hspace*{1em}{<\textbf{cRefPattern}\hspace*{1em}{matchPattern}="{([\textasciicircum  ]+) ([0-9]+)⃥.([0-9]+)}"\mbox{}\newline 
\hspace*{1em}\hspace*{1em}\hspace*{1em}{replacementPattern}="{\#xpath(//div1[@n='\$1']/div2[@n='\$2']/div3[@n='\$3']}">}\mbox{}\newline 
\hspace*{1em}\hspace*{1em}\hspace*{1em}{<\textbf{p}>}Same as above, but without the last component (full\mbox{}\newline 
\hspace*{1em}\hspace*{1em}\hspace*{1em}\hspace*{1em}\hspace*{1em}\hspace*{1em} stop followed by the {<\textbf{gi}>}l{</\textbf{gi}>}'s {<\textbf{att}>}n{</\textbf{att}>}.{</\textbf{p}>}\mbox{}\newline 
\hspace*{1em}\hspace*{1em}{</\textbf{cRefPattern}>}\mbox{}\newline 
\hspace*{1em}\hspace*{1em}{<\textbf{cRefPattern}\hspace*{1em}{matchPattern}="{([\textasciicircum  ]+) ([0-9]+)}"\mbox{}\newline 
\hspace*{1em}\hspace*{1em}\hspace*{1em}{replacementPattern}="{\#xpath(//div1[@n='\$1']/div2[@n='\$2']}">}\mbox{}\newline 
\hspace*{1em}\hspace*{1em}\hspace*{1em}{<\textbf{p}>}Same as above, but without the poem component (full\mbox{}\newline 
\hspace*{1em}\hspace*{1em}\hspace*{1em}\hspace*{1em}\hspace*{1em}\hspace*{1em} stop followed by the {<\textbf{gi}>}div3{</\textbf{gi}>}'s {<\textbf{att}>}n{</\textbf{att}>}.{</\textbf{p}>}\mbox{}\newline 
\hspace*{1em}\hspace*{1em}{</\textbf{cRefPattern}>}\mbox{}\newline 
\hspace*{1em}{</\textbf{refsDecl}>}\mbox{}\newline 
{</\textbf{encodingDesc}>}\end{shaded}\egroup\par \par
The second example encodes the same reference system, again using the {\itshape n} attribute on the \hyperref[TEI.div1]{<div1>}, \hyperref[TEI.div2]{<div2>}, \hyperref[TEI.div3]{<div3>}, and \hyperref[TEI.l]{<l>} tags, but giving the reference string in full on each tag. If canonical references are made only to lines, the reference system could be declared as follows: \par\bgroup\index{refsDecl=<refsDecl>|exampleindex}\index{cRefPattern=<cRefPattern>|exampleindex}\index{matchPattern=@matchPattern!<cRefPattern>|exampleindex}\index{replacementPattern=@replacementPattern!<cRefPattern>|exampleindex}\exampleFont \begin{shaded}\noindent\mbox{}{<\textbf{refsDecl}>}\mbox{}\newline 
\hspace*{1em}{<\textbf{cRefPattern}\hspace*{1em}{matchPattern}="{([\textasciicircum  ]+ [0-9]+⃥.[0-9]+⃥.[0-9]+)}"\mbox{}\newline 
\hspace*{1em}\hspace*{1em}{replacementPattern}="{\#xpath(//l[@n='\$1')}"/>}\mbox{}\newline 
{</\textbf{refsDecl}>}\end{shaded}\egroup\par \noindent  Since the entire regular expression is enclosed as a parenthetical subgroup, the entire canonical reference string is sought as the value of the {\itshape n} attribute on an \hyperref[TEI.l]{<l>} element.\par
In order to handle references to poems as well as to individual lines, the declaration for the reference system must be more complicated: \par\bgroup\index{refsDecl=<refsDecl>|exampleindex}\index{cRefPattern=<cRefPattern>|exampleindex}\index{matchPattern=@matchPattern!<cRefPattern>|exampleindex}\index{replacementPattern=@replacementPattern!<cRefPattern>|exampleindex}\index{cRefPattern=<cRefPattern>|exampleindex}\index{matchPattern=@matchPattern!<cRefPattern>|exampleindex}\index{replacementPattern=@replacementPattern!<cRefPattern>|exampleindex}\exampleFont \begin{shaded}\noindent\mbox{}{<\textbf{refsDecl}>}\mbox{}\newline 
\hspace*{1em}{<\textbf{cRefPattern}\hspace*{1em}{matchPattern}="{([\textasciicircum  ]+ [0-9]+⃥.[0-9]+⃥.[0-9]+)}"\mbox{}\newline 
\hspace*{1em}\hspace*{1em}{replacementPattern}="{\#xpath(//l[@n='\$1')}"/>}\mbox{}\newline 
\hspace*{1em}{<\textbf{cRefPattern}\hspace*{1em}{matchPattern}="{([\textasciicircum  ]+ [0-9]+⃥.[0-9]+)}"\mbox{}\newline 
\hspace*{1em}\hspace*{1em}{replacementPattern}="{\#xpath(//div2[@n='\$1')}"/>}\mbox{}\newline 
{</\textbf{refsDecl}>}\end{shaded}\egroup\par \noindent  This declaration indicates that the entire reference string must be sought as the value of the {\itshape n} attribute on a \hyperref[TEI.div1]{<div1>}, \hyperref[TEI.div2]{<div2>}, \hyperref[TEI.div3]{<div3>}, or \hyperref[TEI.l]{<l>} element.\par
The third example encodes the same reference system, this time giving the entire reference string as the value of the {\itshape xml:id} attribute on the relevant tags. The reference system declaration for such an encoding could be: \par\bgroup\index{refsDecl=<refsDecl>|exampleindex}\index{cRefPattern=<cRefPattern>|exampleindex}\index{matchPattern=@matchPattern!<cRefPattern>|exampleindex}\index{replacementPattern=@replacementPattern!<cRefPattern>|exampleindex}\exampleFont \begin{shaded}\noindent\mbox{}{<\textbf{refsDecl}>}\mbox{}\newline 
\hspace*{1em}{<\textbf{cRefPattern}\hspace*{1em}{matchPattern}="{(.*)}"\mbox{}\newline 
\hspace*{1em}\hspace*{1em}{replacementPattern}="{\#\$1}"/>}\mbox{}\newline 
{</\textbf{refsDecl}>}\end{shaded}\egroup\par \noindent  although in general there seems to be little advantage in this case: it is no more difficult to use a standard relative URI reference as the value of {\itshape target}.\par
In cases where a more complete formal declaration of text structure is desirable, for example in systems that will present the contents of a large TEI file in smaller chunks, the \hyperref[TEI.citeStructure]{<citeStructure>} element may be used. This method permits canonical references to be resolved and also allows them to be extracted, so that, for example, a list of resolvable citations may be generated from the document. The example from the \textit{Amores} above could be implemented using \hyperref[TEI.citeStructure]{<citeStructure>} thus: \par\bgroup\index{refsDecl=<refsDecl>|exampleindex}\index{citeStructure=<citeStructure>|exampleindex}\index{unit=@unit!<citeStructure>|exampleindex}\index{match=@match!<citeStructure>|exampleindex}\index{use=@use!<citeStructure>|exampleindex}\index{citeStructure=<citeStructure>|exampleindex}\index{unit=@unit!<citeStructure>|exampleindex}\index{match=@match!<citeStructure>|exampleindex}\index{use=@use!<citeStructure>|exampleindex}\index{delim=@delim!<citeStructure>|exampleindex}\index{citeStructure=<citeStructure>|exampleindex}\index{unit=@unit!<citeStructure>|exampleindex}\index{match=@match!<citeStructure>|exampleindex}\index{use=@use!<citeStructure>|exampleindex}\index{delim=@delim!<citeStructure>|exampleindex}\index{citeStructure=<citeStructure>|exampleindex}\index{unit=@unit!<citeStructure>|exampleindex}\index{match=@match!<citeStructure>|exampleindex}\index{use=@use!<citeStructure>|exampleindex}\index{delim=@delim!<citeStructure>|exampleindex}\exampleFont \begin{shaded}\noindent\mbox{}{<\textbf{refsDecl}>}\mbox{}\newline 
\hspace*{1em}{<\textbf{citeStructure}\hspace*{1em}{unit}="{work}"\hspace*{1em}{match}="{//div1}"\mbox{}\newline 
\hspace*{1em}\hspace*{1em}{use}="{@n}">}\mbox{}\newline 
\hspace*{1em}\hspace*{1em}{<\textbf{citeStructure}\hspace*{1em}{unit}="{book}"\hspace*{1em}{match}="{div2}"\mbox{}\newline 
\hspace*{1em}\hspace*{1em}\hspace*{1em}{use}="{@n}"\hspace*{1em}{delim}="{ }">}\mbox{}\newline 
\hspace*{1em}\hspace*{1em}\hspace*{1em}{<\textbf{citeStructure}\hspace*{1em}{unit}="{poem}"\hspace*{1em}{match}="{div3}"\mbox{}\newline 
\hspace*{1em}\hspace*{1em}\hspace*{1em}\hspace*{1em}{use}="{@n}"\hspace*{1em}{delim}="{.}">}\mbox{}\newline 
\hspace*{1em}\hspace*{1em}\hspace*{1em}\hspace*{1em}{<\textbf{citeStructure}\hspace*{1em}{unit}="{line}"\hspace*{1em}{match}="{l}"\mbox{}\newline 
\hspace*{1em}\hspace*{1em}\hspace*{1em}\hspace*{1em}\hspace*{1em}{use}="{@n}"\hspace*{1em}{delim}="{.}"/>}\mbox{}\newline 
\hspace*{1em}\hspace*{1em}\hspace*{1em}{</\textbf{citeStructure}>}\mbox{}\newline 
\hspace*{1em}\hspace*{1em}{</\textbf{citeStructure}>}\mbox{}\newline 
\hspace*{1em}{</\textbf{citeStructure}>}\mbox{}\newline 
{</\textbf{refsDecl}>}\end{shaded}\egroup\par \par
\hyperref[TEI.citeStructure]{<citeStructure>} also provides a method for attaching informational properties to units of structure, by means of the \hyperref[TEI.citeData]{<citeData>} element. The work, book, and poem divisions above might all have \hyperref[TEI.head]{<head>} elements which provide a title for the section. If we wish that information to be extractable, we can use citeData to specify where it is to be found: \par\bgroup\index{refsDecl=<refsDecl>|exampleindex}\index{citeStructure=<citeStructure>|exampleindex}\index{unit=@unit!<citeStructure>|exampleindex}\index{match=@match!<citeStructure>|exampleindex}\index{use=@use!<citeStructure>|exampleindex}\index{citeData=<citeData>|exampleindex}\index{property=@property!<citeData>|exampleindex}\index{use=@use!<citeData>|exampleindex}\index{citeStructure=<citeStructure>|exampleindex}\index{unit=@unit!<citeStructure>|exampleindex}\index{match=@match!<citeStructure>|exampleindex}\index{use=@use!<citeStructure>|exampleindex}\index{delim=@delim!<citeStructure>|exampleindex}\index{citeData=<citeData>|exampleindex}\index{property=@property!<citeData>|exampleindex}\index{use=@use!<citeData>|exampleindex}\index{citeStructure=<citeStructure>|exampleindex}\index{unit=@unit!<citeStructure>|exampleindex}\index{match=@match!<citeStructure>|exampleindex}\index{use=@use!<citeStructure>|exampleindex}\index{delim=@delim!<citeStructure>|exampleindex}\index{citeData=<citeData>|exampleindex}\index{property=@property!<citeData>|exampleindex}\index{use=@use!<citeData>|exampleindex}\index{citeStructure=<citeStructure>|exampleindex}\index{unit=@unit!<citeStructure>|exampleindex}\index{match=@match!<citeStructure>|exampleindex}\index{use=@use!<citeStructure>|exampleindex}\index{delim=@delim!<citeStructure>|exampleindex}\exampleFont \begin{shaded}\noindent\mbox{}{<\textbf{refsDecl}>}\mbox{}\newline 
\hspace*{1em}{<\textbf{citeStructure}\hspace*{1em}{unit}="{work}"\hspace*{1em}{match}="{//div1}"\mbox{}\newline 
\hspace*{1em}\hspace*{1em}{use}="{@n}">}\mbox{}\newline 
\hspace*{1em}\hspace*{1em}{<\textbf{citeData}\hspace*{1em}{property}="{http://purl.org/dc/terms/title}"\mbox{}\newline 
\hspace*{1em}\hspace*{1em}\hspace*{1em}{use}="{head}"/>}\mbox{}\newline 
\hspace*{1em}\hspace*{1em}{<\textbf{citeStructure}\hspace*{1em}{unit}="{book}"\hspace*{1em}{match}="{div2}"\mbox{}\newline 
\hspace*{1em}\hspace*{1em}\hspace*{1em}{use}="{@n}"\hspace*{1em}{delim}="{ }">}\mbox{}\newline 
\hspace*{1em}\hspace*{1em}\hspace*{1em}{<\textbf{citeData}\hspace*{1em}{property}="{http://purl.org/dc/terms/title}"\mbox{}\newline 
\hspace*{1em}\hspace*{1em}\hspace*{1em}\hspace*{1em}{use}="{head}"/>}\mbox{}\newline 
\hspace*{1em}\hspace*{1em}\hspace*{1em}{<\textbf{citeStructure}\hspace*{1em}{unit}="{poem}"\hspace*{1em}{match}="{div3}"\mbox{}\newline 
\hspace*{1em}\hspace*{1em}\hspace*{1em}\hspace*{1em}{use}="{@n}"\hspace*{1em}{delim}="{.}">}\mbox{}\newline 
\hspace*{1em}\hspace*{1em}\hspace*{1em}\hspace*{1em}{<\textbf{citeData}\hspace*{1em}{property}="{http://purl.org/dc/terms/title}"\mbox{}\newline 
\hspace*{1em}\hspace*{1em}\hspace*{1em}\hspace*{1em}\hspace*{1em}{use}="{head}"/>}\mbox{}\newline 
\hspace*{1em}\hspace*{1em}\hspace*{1em}\hspace*{1em}{<\textbf{citeStructure}\hspace*{1em}{unit}="{line}"\hspace*{1em}{match}="{l}"\mbox{}\newline 
\hspace*{1em}\hspace*{1em}\hspace*{1em}\hspace*{1em}\hspace*{1em}{use}="{@n}"\hspace*{1em}{delim}="{.}"/>}\mbox{}\newline 
\hspace*{1em}\hspace*{1em}\hspace*{1em}{</\textbf{citeStructure}>}\mbox{}\newline 
\hspace*{1em}\hspace*{1em}{</\textbf{citeStructure}>}\mbox{}\newline 
\hspace*{1em}{</\textbf{citeStructure}>}\mbox{}\newline 
{</\textbf{refsDecl}>}\end{shaded}\egroup\par \noindent  The example above maps the head element to the \hyperref[CO-BIBL-2]{Dublin Core} property \texttt{title}. For convenience, property URIs may be abbreviated using \hyperref[TEI.prefixDef]{<prefixDef>}.\par
Reference systems recorded by means of milestone tags can also be declared; the following prose description could be used to declare the example given in section \textit{\hyperref[CORS5]{3.11.3.\ Milestone Elements}}. \par\bgroup\index{refsDecl=<refsDecl>|exampleindex}\index{p=<p>|exampleindex}\exampleFont \begin{shaded}\noindent\mbox{}{<\textbf{refsDecl}>}\mbox{}\newline 
\hspace*{1em}{<\textbf{p}>}Standard references to work, book, poem, and line may be\mbox{}\newline 
\hspace*{1em}\hspace*{1em} constructed from the milestone tags in the text.{</\textbf{p}>}\mbox{}\newline 
{</\textbf{refsDecl}>}\end{shaded}\egroup\par \noindent  Or in this way, using a formal declaration for this reference scheme derived from edition E1. \par\bgroup\index{refsDecl=<refsDecl>|exampleindex}\index{refState=<refState>|exampleindex}\index{ed=@ed!<refState>|exampleindex}\index{unit=@unit!<refState>|exampleindex}\index{delim=@delim!<refState>|exampleindex}\index{refState=<refState>|exampleindex}\index{ed=@ed!<refState>|exampleindex}\index{unit=@unit!<refState>|exampleindex}\index{delim=@delim!<refState>|exampleindex}\index{refState=<refState>|exampleindex}\index{ed=@ed!<refState>|exampleindex}\index{unit=@unit!<refState>|exampleindex}\index{delim=@delim!<refState>|exampleindex}\index{refState=<refState>|exampleindex}\index{ed=@ed!<refState>|exampleindex}\index{unit=@unit!<refState>|exampleindex}\exampleFont \begin{shaded}\noindent\mbox{}{<\textbf{refsDecl}>}\mbox{}\newline 
\hspace*{1em}{<\textbf{refState}\hspace*{1em}{ed}="{E1}"\hspace*{1em}{unit}="{work}"\hspace*{1em}{delim}="{ }"/>}\mbox{}\newline 
\hspace*{1em}{<\textbf{refState}\hspace*{1em}{ed}="{E1}"\hspace*{1em}{unit}="{book}"\hspace*{1em}{delim}="{.}"/>}\mbox{}\newline 
\hspace*{1em}{<\textbf{refState}\hspace*{1em}{ed}="{E1}"\hspace*{1em}{unit}="{poem}"\hspace*{1em}{delim}="{:}"/>}\mbox{}\newline 
\hspace*{1em}{<\textbf{refState}\hspace*{1em}{ed}="{E1}"\hspace*{1em}{unit}="{line}"/>}\mbox{}\newline 
{</\textbf{refsDecl}>}\end{shaded}\egroup\par 
\subsection[{Bibliographic Citations and References}]{Bibliographic Citations and References}\label{COBI}\par
Bibliographic references (that is, full descriptions of bibliographic items such as books, articles, films, broadcasts, songs, etc.) or pointers to them may appear at various places in a TEI text. They are required at several points within the TEI header's source description, as discussed in section \textit{\hyperref[HD3]{2.2.7.\ The Source Description}}; they may also appear within the body of a text, either singly (for example within a footnote), or collected together in a list as a distinct part of a text; detailed bibliographic descriptions of manuscript or other source materials may also be required. These Guidelines propose a number of specialized elements to encode such descriptions, which together constitute the \textsf{model.biblLike} class. 
\begin{sansreflist}
  
\item [\textbf{model.biblLike}] groups elements containing a bibliographic description. \par 
\begin{longtable}{P{0.1619538834951456\textwidth}P{0.6880461165048544\textwidth}}
\hyperref[TEI.bibl]{bibl}\tabcellsep (bibliographic citation) contains a loosely-structured bibliographic citation of which the sub-components may or may not be explicitly tagged.\\
\hyperref[TEI.biblFull]{biblFull}\tabcellsep (fully-structured bibliographic citation) contains a fully-structured bibliographic citation, in which all components of the TEI file description are present.\\
\hyperref[TEI.biblStruct]{biblStruct}\tabcellsep (structured bibliographic citation) contains a structured bibliographic citation, in which only bibliographic sub-elements appear and in a specified order.\\
\hyperref[TEI.listBibl]{listBibl}\tabcellsep (citation list) contains a list of bibliographic citations of any kind.\\
\hyperref[TEI.msDesc]{msDesc}\tabcellsep (manuscript description) contains a description of a single identifiable manuscript or other text-bearing object such as early printed books.\end{longtable} \par
 
\end{sansreflist}
 Lists of such elements may also be encoded using the following element: 
\begin{sansreflist}
  
\item [\textbf{<listBibl>}] (citation list) contains a list of bibliographic citations of any kind.
\end{sansreflist}
\par
In printed texts, the individual constituents of a bibliographic reference are conventionally marked off from each other and from the flow of text by such features as bracketing, italics, special punctuation conventions, underlining, etc. In electronic texts, such distinctions are also important, whether in order to produce acceptably formatted output or to facilitate intelligent retrieval processing,\footnote{For example, to distinguish \textit{London} as an author's name from \textit{London} as a place of publication or as a component of a title.} quite apart from the need to distinguish the reference itself as a textual object with particular linguistic properties.\par
It should be emphasized that for references as for other textual features, the primary or sole consideration is not how the text should be formatted when it is printed or displayed. The distinctions permitted by the scheme outlined here may not necessarily be all that particular formatters or bibliographic styles require, although they should prove adequate to the needs of many such commonly used software systems.\footnote{Among the bibliographic software systems and subsystems consulted in the design of the \hyperref[TEI.biblStruct]{<biblStruct>} structure were BibTeX, Scribe, and ProCite. The distinctions made by all three may be preserved in \hyperref[TEI.biblStruct]{<biblStruct>} structures, though the nature of their design prevents a simple one-to-one mapping from their data elements to TEI elements. For further information, see section \textit{\hyperref[COBIOT]{3.12.4.\ Relationship to Other Bibliographic Schemes}}.} The features distinguished and described below (in section \textit{\hyperref[COBICO]{3.12.2.\ Components of Bibliographic References}}) constitute a set which has been useful for a wide range of bibliographic purposes and in many applications, and which moreover corresponds to a great extent with existing bibliographic and library cataloguing practice. For a fuller account of that practice as applied to electronic texts see section \textit{\hyperref[HD3]{2.2.7.\ The Source Description}}; for a brief mention of related library standards see section \textit{\hyperref[HD8]{2.8.\ Note for Library Cataloguers}}.\par
The most commonly used elements in the \textsf{model.biblLike} class are \hyperref[TEI.biblStruct]{<biblStruct>} and \hyperref[TEI.bibl]{<bibl>}. \hyperref[TEI.biblStruct]{<biblStruct>} will usually be easier to process mechanically than \hyperref[TEI.bibl]{<bibl>} because its structure is more constrained and predictable. It is suited to situations in which the objective is to represent bibliographic information for machine processing directly by other systems or after conversion to some other bibliographic markup formats such as BibTeXML or MODS. Punctuation delimiting the components of a print citation is not permitted directly within a \hyperref[TEI.biblStruct]{<biblStruct>} element; instead, the presence and order of child elements must be used to reconstruct the punctuation required by a particular style.\par
By contrast, \hyperref[TEI.bibl]{<bibl>} allows for considerable flexibility in that it can include both delimiting punctuation and unmarked-up text; and its constituents can also be ordered in any way. This makes it suitable for marking up bibliographies in existing documents, where it is considered important to preserve the form of references in the original document, while also distinguishing important pieces of information such as authors, dates, publishers, and so on. \hyperref[TEI.bibl]{<bibl>} may also be useful when encoding ‘born digital’ documents which require use of a specific style guide when rendering the content; its flexibility makes it easier to provide all the information for a reference in the exact sequence required by the target rendering, including any necessary punctuation and linking words, rather than using an XSLT stylesheet or similar to reorder and punctuate the data.\par
The third element in the \textsf{model.biblLike} class, \hyperref[TEI.biblFull]{<biblFull>}, has a content model based on the \hyperref[TEI.fileDesc]{<fileDesc>} element of the TEI header. Both are based on the International Standard for Bibliographic Description (ISBD), which forms the basis of several national standards for bibliographic citations. The order of child elements in both \hyperref[TEI.biblFull]{<biblFull>} and \hyperref[TEI.fileDesc]{<fileDesc>} corresponds to the order of bibliographic description ‘areas’ in ISBD with two minor exceptions. First, the \hyperref[TEI.extent]{<extent>} element, corresponding to the \textit{physical description area} in ISBD, appears just after the \textit{publication, production, distribution, etc. area} in ISBD, not before it as in TEI. Second, \hyperref[TEI.biblFull]{<biblFull>} and \hyperref[TEI.fileDesc]{<fileDesc>} use the child element \hyperref[TEI.publicationStmt]{<publicationStmt>} to cover not only the \textit{publication, production, distribution, etc. area} but also the \textit{resource identifier and terms of availability area} associated with that publication. Despite these inconsistencies, users encoding citations and attempting to format them according to a standard that closely adheres to ISBD may find that \hyperref[TEI.biblFull]{<biblFull>}, used with its child elements and without delimiting punctuation, provides an appropriate granularity of encoding with elements that can easily be rendered for the reader. However, it is important to note that some ISBD-derived citation formats (such as ANSI/NISO Z39.29 and ГОСТ 7.1) are not entirely conformant to ISBD either, since they may begin with a statement of authorship that does not map to the ISBD statement of responsibility.
\subsubsection[{Methods of Encoding Bibliographic References and Lists of References}]{Methods of Encoding Bibliographic References and Lists of References}\label{COBITY}\par
The members of the \textsf{model.biblLike} class all share a number of possible component sub-elements. For the \hyperref[TEI.bibl]{<bibl>} and \hyperref[TEI.biblStruct]{<biblStruct>} elements, exactly the same sub-elements are concerned, and they are described together in section \textit{\hyperref[COBICO]{3.12.2.\ Components of Bibliographic References}}; for the \hyperref[TEI.biblFull]{<biblFull>} element, the sub-elements concerned are fully described in section \textit{\hyperref[HD2]{2.2.\ The File Description}}.\par
Different levels of specific tagging may be appropriate in different situations. In some cases, it may be felt necessary to mark just the extent of the reference itself, with perhaps a few distinctions being made within it (for example, between the part of the reference which identifies a title or author and the rest). Such references, containing a mixture of text with specialized bibliographic elements, are regarded as \hyperref[TEI.bibl]{<bibl>} elements, and tagged accordingly. For example: \par\bgroup\index{p=<p>|exampleindex}\index{bibl=<bibl>|exampleindex}\index{title=<title>|exampleindex}\exampleFont \begin{shaded}\noindent\mbox{}{<\textbf{p}>}A book which had a great influence on him\mbox{}\newline 
 was {<\textbf{bibl}>}Tufte's {<\textbf{title}>}Envisioning\mbox{}\newline 
\hspace*{1em}\hspace*{1em}\hspace*{1em}\hspace*{1em} Information{</\textbf{title}>}\mbox{}\newline 
\hspace*{1em}{</\textbf{bibl}>}, although he may\mbox{}\newline 
 never have actually read it.{</\textbf{p}>}\end{shaded}\egroup\par \noindent  Indeed, some encoders may find it unnecessary to mark the bibliographic reference at all: \par\bgroup\index{p=<p>|exampleindex}\index{title=<title>|exampleindex}\exampleFont \begin{shaded}\noindent\mbox{}{<\textbf{p}>}A book which had a great influence on him\mbox{}\newline 
 was Tufte's {<\textbf{title}>}Envisioning Information{</\textbf{title}>},\mbox{}\newline 
 although he may never have actually read it.{</\textbf{p}>}\end{shaded}\egroup\par \par
Some bibliographic references are extremely elliptical, often only a string of the form \textit{Baxter, 1983}. If no further details of Baxter's book are given in the source text and none is supplied by the encoder, then the reference thus given should be tagged as a \hyperref[TEI.bibl]{<bibl>}: \par\bgroup\index{bibl=<bibl>|exampleindex}\exampleFont \begin{shaded}\noindent\mbox{}All of this is of course much more fully treated\mbox{}\newline 
 in {<\textbf{bibl}>}Baxter, 1983{</\textbf{bibl}>}.\end{shaded}\egroup\par \noindent  In general, however, normal modern bibliographic practice, and these Guidelines, distinguish between a bibliographic \textit{reference}, which is a self-sufficient description of a bibliographic item, and a bibliographic \textit{pointer}, which is a short-form citation (e.g. \textit{Baxter, 1983}) which serves usually as a place-holder or pointer to a full long-form reference found elsewhere in the text. The usual encoding of short-form references such as \textit{Baxter, 1983} is not as \hyperref[TEI.bibl]{<bibl>} elements but as cross-references to such elements; see section \textit{\hyperref[COBIXR]{3.12.3.\ Bibliographic Pointers }} below.\par
In cases where the encoder wishes to impose more structure on the bibliographic information, for example to make sure it conforms to a particular stylesheet or retrieval processor, the \hyperref[TEI.biblStruct]{<biblStruct>} element should be used. Note that several of the features in this and later examples are explained later in the current section. \par\bgroup\index{biblStruct=<biblStruct>|exampleindex}\index{monogr=<monogr>|exampleindex}\index{author=<author>|exampleindex}\index{persName=<persName>|exampleindex}\index{forename=<forename>|exampleindex}\index{forename=<forename>|exampleindex}\index{full=@full!<forename>|exampleindex}\index{surname=<surname>|exampleindex}\index{idno=<idno>|exampleindex}\index{type=@type!<idno>|exampleindex}\index{idno=<idno>|exampleindex}\index{type=@type!<idno>|exampleindex}\index{title=<title>|exampleindex}\index{level=@level!<title>|exampleindex}\index{imprint=<imprint>|exampleindex}\index{pubPlace=<pubPlace>|exampleindex}\index{publisher=<publisher>|exampleindex}\index{date=<date>|exampleindex}\index{when=@when!<date>|exampleindex}\exampleFont \begin{shaded}\noindent\mbox{}{<\textbf{biblStruct}>}\mbox{}\newline 
\hspace*{1em}{<\textbf{monogr}>}\mbox{}\newline 
\hspace*{1em}\hspace*{1em}{<\textbf{author}>}\mbox{}\newline 
\hspace*{1em}\hspace*{1em}\hspace*{1em}{<\textbf{persName}>}\mbox{}\newline 
\hspace*{1em}\hspace*{1em}\hspace*{1em}\hspace*{1em}{<\textbf{forename}>}Edward{</\textbf{forename}>}\mbox{}\newline 
\hspace*{1em}\hspace*{1em}\hspace*{1em}\hspace*{1em}{<\textbf{forename}\hspace*{1em}{full}="{init}">}R.{</\textbf{forename}>}\mbox{}\newline 
\hspace*{1em}\hspace*{1em}\hspace*{1em}\hspace*{1em}{<\textbf{surname}>}Tufte{</\textbf{surname}>}\mbox{}\newline 
\hspace*{1em}\hspace*{1em}\hspace*{1em}{</\textbf{persName}>}\mbox{}\newline 
\hspace*{1em}\hspace*{1em}\hspace*{1em}{<\textbf{idno}\hspace*{1em}{type}="{scopus}">}6506403994{</\textbf{idno}>}\mbox{}\newline 
\hspace*{1em}\hspace*{1em}\hspace*{1em}{<\textbf{idno}\hspace*{1em}{type}="{lcaf}">}http://id.loc.gov/authorites/names/n50012763.html{</\textbf{idno}>}\mbox{}\newline 
\hspace*{1em}\hspace*{1em}{</\textbf{author}>}\mbox{}\newline 
\hspace*{1em}\hspace*{1em}{<\textbf{title}\hspace*{1em}{level}="{m}">}Envisioning Information{</\textbf{title}>}\mbox{}\newline 
\hspace*{1em}\hspace*{1em}{<\textbf{imprint}>}\mbox{}\newline 
\hspace*{1em}\hspace*{1em}\hspace*{1em}{<\textbf{pubPlace}>}Cheshire, Conn.{</\textbf{pubPlace}>}\mbox{}\newline 
\hspace*{1em}\hspace*{1em}\hspace*{1em}{<\textbf{publisher}>}Graphics Press{</\textbf{publisher}>}\mbox{}\newline 
\hspace*{1em}\hspace*{1em}\hspace*{1em}{<\textbf{date}\hspace*{1em}{when}="{1990}"/>}\mbox{}\newline 
\hspace*{1em}\hspace*{1em}{</\textbf{imprint}>}\mbox{}\newline 
\hspace*{1em}{</\textbf{monogr}>}\mbox{}\newline 
{</\textbf{biblStruct}>}\end{shaded}\egroup\par \par
A more complex and detailed bibliographic structure is provided by the \hyperref[TEI.biblFull]{<biblFull>} element defined in the TEI header module. This element is provided as a means of embedding the file description of one existing digital text within that of another (see further section \textit{\hyperref[HD2]{2.2.\ The File Description}}); however, its use is not confined to digital texts, and it may be used in the same way as any other bibliographic element, as in this example: \par\bgroup\index{biblFull=<biblFull>|exampleindex}\index{titleStmt=<titleStmt>|exampleindex}\index{title=<title>|exampleindex}\index{author=<author>|exampleindex}\index{extent=<extent>|exampleindex}\index{publicationStmt=<publicationStmt>|exampleindex}\index{publisher=<publisher>|exampleindex}\index{pubPlace=<pubPlace>|exampleindex}\index{date=<date>|exampleindex}\exampleFont \begin{shaded}\noindent\mbox{}{<\textbf{biblFull}>}\mbox{}\newline 
\hspace*{1em}{<\textbf{titleStmt}>}\mbox{}\newline 
\hspace*{1em}\hspace*{1em}{<\textbf{title}>}Envisioning Information{</\textbf{title}>}\mbox{}\newline 
\hspace*{1em}\hspace*{1em}{<\textbf{author}>}Tufte, Edward R[olf]{</\textbf{author}>}\mbox{}\newline 
\hspace*{1em}{</\textbf{titleStmt}>}\mbox{}\newline 
\hspace*{1em}{<\textbf{extent}>}126 pp.{</\textbf{extent}>}\mbox{}\newline 
\hspace*{1em}{<\textbf{publicationStmt}>}\mbox{}\newline 
\hspace*{1em}\hspace*{1em}{<\textbf{publisher}>}Graphics Press{</\textbf{publisher}>}\mbox{}\newline 
\hspace*{1em}\hspace*{1em}{<\textbf{pubPlace}>}Cheshire, Conn. USA{</\textbf{pubPlace}>}\mbox{}\newline 
\hspace*{1em}\hspace*{1em}{<\textbf{date}>}1990{</\textbf{date}>}\mbox{}\newline 
\hspace*{1em}{</\textbf{publicationStmt}>}\mbox{}\newline 
{</\textbf{biblFull}>}\end{shaded}\egroup\par \par
A list of bibliographic items, of whatever kind, may be treated in the same way as any other list (see section \textit{\hyperref[COLI]{3.8.\ Lists}}). Alternatively, the specialized \hyperref[TEI.listBibl]{<listBibl>} element may be used. The difference between the two is that a \hyperref[TEI.list]{<list>} contains \hyperref[TEI.item]{<item>} elements, within which bibliographic elements (\hyperref[TEI.bibl]{<bibl>}, \hyperref[TEI.biblStruct]{<biblStruct>}, or \hyperref[TEI.biblFull]{<biblFull>}) may appear, as well as other phrase- and paragraph-level elements, whereas the \hyperref[TEI.listBibl]{<listBibl>} may contain only bibliographic elements, optionally preceded by a heading and a series of introductory paragraphs. For most purposes, good practice would usually require that a \hyperref[TEI.listBibl]{<listBibl>} contain only one kind of bibliographic element, though the following example combines both fully structured \hyperref[TEI.biblStruct]{<biblStruct>} and informal \hyperref[TEI.bibl]{<bibl>} elements: \par\bgroup\index{listBibl=<listBibl>|exampleindex}\index{head=<head>|exampleindex}\index{biblStruct=<biblStruct>|exampleindex}\index{analytic=<analytic>|exampleindex}\index{author=<author>|exampleindex}\index{persName=<persName>|exampleindex}\index{surname=<surname>|exampleindex}\index{forename=<forename>|exampleindex}\index{forename=<forename>|exampleindex}\index{title=<title>|exampleindex}\index{level=@level!<title>|exampleindex}\index{monogr=<monogr>|exampleindex}\index{title=<title>|exampleindex}\index{level=@level!<title>|exampleindex}\index{editor=<editor>|exampleindex}\index{persName=<persName>|exampleindex}\index{forename=<forename>|exampleindex}\index{forename=<forename>|exampleindex}\index{surname=<surname>|exampleindex}\index{imprint=<imprint>|exampleindex}\index{publisher=<publisher>|exampleindex}\index{pubPlace=<pubPlace>|exampleindex}\index{date=<date>|exampleindex}\index{when=@when!<date>|exampleindex}\index{biblScope=<biblScope>|exampleindex}\index{unit=@unit!<biblScope>|exampleindex}\index{from=@from!<biblScope>|exampleindex}\index{to=@to!<biblScope>|exampleindex}\index{note=<note>|exampleindex}\index{title=<title>|exampleindex}\index{level=@level!<title>|exampleindex}\index{bibl=<bibl>|exampleindex}\index{author=<author>|exampleindex}\index{persName=<persName>|exampleindex}\index{forename=<forename>|exampleindex}\index{surname=<surname>|exampleindex}\index{title=<title>|exampleindex}\index{level=@level!<title>|exampleindex}\index{date=<date>|exampleindex}\index{when=@when!<date>|exampleindex}\index{bibl=<bibl>|exampleindex}\index{author=<author>|exampleindex}\index{persName=<persName>|exampleindex}\index{surname=<surname>|exampleindex}\index{forename=<forename>|exampleindex}\index{title=<title>|exampleindex}\index{level=@level!<title>|exampleindex}\index{pubPlace=<pubPlace>|exampleindex}\index{publisher=<publisher>|exampleindex}\index{date=<date>|exampleindex}\index{when=@when!<date>|exampleindex}\exampleFont \begin{shaded}\noindent\mbox{}{<\textbf{listBibl}>}\newline
   {<\textbf{head}>}Bibliography{</\textbf{head}>}\newline
   {<\textbf{biblStruct}\hspace*{1em}{xml:id}="{NELSON80}">}\newline
      {<\textbf{analytic}>}\newline
        {<\textbf{author}>}\newline
          {<\textbf{persName}>}\newline
            {<\textbf{surname}>}Nelson{</\textbf{surname}>}\newline
            {<\textbf{forename}>}Theodore{</\textbf{forename}>}\newline
            {<\textbf{forename}>}Holm{</\textbf{forename}>}\newline
          {</\textbf{persName}>}\newline
        {</\textbf{author}>}\newline
        {<\textbf{title}\hspace*{1em}{level}="{a}">}Replacing the printed word:\newline
             a complete literary system{</\textbf{title}>}\newline
      {</\textbf{analytic}>}\newline
      {<\textbf{monogr}>}\newline
         {<\textbf{title}\hspace*{1em}{level}="{m}">}Information Processing '80:  Proceedings of the IFIPS\newline
             Congress, October 1980{</\textbf{title}>}\newline
         {<\textbf{editor}>}\newline
           {<\textbf{persName}>}\newline
             {<\textbf{forename}>}Simon{</\textbf{forename}>}\newline
             {<\textbf{forename}>}H.{</\textbf{forename}>}\newline
             {<\textbf{surname}>}Lavington{</\textbf{surname}>}\newline
           {</\textbf{persName}>}\newline
         {</\textbf{editor}>}\newline
         {<\textbf{imprint}>}\newline
            {<\textbf{publisher}>}North-Holland{</\textbf{publisher}>}\newline
            {<\textbf{pubPlace}>}Amsterdam{</\textbf{pubPlace}>}\newline
            {<\textbf{date}\hspace*{1em}{when}="{1980}"/>}\newline
         {</\textbf{imprint}>}\newline
         {<\textbf{biblScope}\hspace*{1em}{unit}="{page}"\hspace*{1em}{from}="{1013}"\mbox{}\newline 
\hspace*{1em}{to}="{1023}">}1013–23{</\textbf{biblScope}>}\newline
      {</\textbf{monogr}>}\newline
      {<\textbf{note}>}Apparently a draft of section 4 of \newline
          {<\textbf{title}\hspace*{1em}{level}="{m}">}Literary Machines{</\textbf{title}>}.{</\textbf{note}>}\newline
   {</\textbf{biblStruct}>}   \newline
\newline
  {<\textbf{bibl}\hspace*{1em}{xml:id}="{NELSON88}">}{<\textbf{author}>}{<\textbf{persName}>}{<\textbf{forename}>}Ted{</\textbf{forename}>}\newline
{<\textbf{surname}>}Nelson{</\textbf{surname}>}{</\textbf{persName}>}{</\textbf{author}>}: \newline
{<\textbf{title}\hspace*{1em}{level}="{u}">}Literary Machines{</\textbf{title}>} (privately published, \newline
{<\textbf{date}\hspace*{1em}{when}="{1987}">}1987{</\textbf{date}>}).{</\textbf{bibl}>}\newline
  \newline
  {<\textbf{bibl}\hspace*{1em}{xml:id}="{BAXTER88}">}{<\textbf{author}>}{<\textbf{persName}>}{<\textbf{surname}>}Baxter{</\textbf{surname}>}, \newline
{<\textbf{forename}>}Glen{</\textbf{forename}>}{</\textbf{persName}>}{</\textbf{author}>}:\newline
{<\textbf{title}\hspace*{1em}{level}="{m}">}Glen Baxter His Life: the years of struggle{</\textbf{title}>} \newline
{<\textbf{pubPlace}>}London{</\textbf{pubPlace}>}: {<\textbf{publisher}>}Thames and Hudson{</\textbf{publisher}>}, \newline
{<\textbf{date}\hspace*{1em}{when}="{1988}">}1988{</\textbf{date}>}.{</\textbf{bibl}>}\newline
{</\textbf{listBibl}>}\end{shaded}\egroup\par \noindent  This example also demonstrates the way that bibliographical markup of authors, titles, dates etc. can be handled differently in \hyperref[TEI.biblStruct]{<biblStruct>}s and \hyperref[TEI.bibl]{<bibl>}s. In the two \hyperref[TEI.bibl]{<bibl>} items, the key information is marked up, but it is presented in an order which makes it suitable for direct rendering, with the punctuation included.\par
The \hyperref[TEI.listBibl]{<listBibl>} element is most appropriate for a more formal bibliography. The same \hyperref[TEI.bibl]{<bibl>} or \hyperref[TEI.biblStruct]{<biblStruct>} elements may however be embedded within an ordinary list, thus allowing them to be mixed with running prose or presented informally, as in the following version of the same example: \par\bgroup\index{list=<list>|exampleindex}\index{head=<head>|exampleindex}\index{item=<item>|exampleindex}\index{bibl=<bibl>|exampleindex}\index{author=<author>|exampleindex}\index{title=<title>|exampleindex}\index{level=@level!<title>|exampleindex}\index{title=<title>|exampleindex}\index{level=@level!<title>|exampleindex}\index{editor=<editor>|exampleindex}\index{publisher=<publisher>|exampleindex}\index{pubPlace=<pubPlace>|exampleindex}\index{date=<date>|exampleindex}\index{biblScope=<biblScope>|exampleindex}\index{note=<note>|exampleindex}\index{title=<title>|exampleindex}\index{item=<item>|exampleindex}\index{bibl=<bibl>|exampleindex}\index{title=<title>|exampleindex}\index{item=<item>|exampleindex}\index{bibl=<bibl>|exampleindex}\index{author=<author>|exampleindex}\index{title=<title>|exampleindex}\exampleFont \begin{shaded}\noindent\mbox{}{<\textbf{list}>}\mbox{}\newline 
\hspace*{1em}{<\textbf{head}>}Bibliography{</\textbf{head}>}\mbox{}\newline 
\hspace*{1em}{<\textbf{item}>}\mbox{}\newline 
\hspace*{1em}\hspace*{1em}{<\textbf{bibl}\hspace*{1em}{xml:id}="{NEL80}">}\mbox{}\newline 
\hspace*{1em}\hspace*{1em}\hspace*{1em}{<\textbf{author}>}Nelson, T. H.{</\textbf{author}>}\mbox{}\newline 
\hspace*{1em}\hspace*{1em}\hspace*{1em}{<\textbf{title}\hspace*{1em}{level}="{a}">}Replacing the printed word:\mbox{}\newline 
\hspace*{1em}\hspace*{1em}\hspace*{1em}\hspace*{1em}\hspace*{1em}\hspace*{1em} a complete literary system{</\textbf{title}>}.\mbox{}\newline 
\hspace*{1em}\hspace*{1em}{<\textbf{title}\hspace*{1em}{level}="{m}">}Information Processing '80:\mbox{}\newline 
\hspace*{1em}\hspace*{1em}\hspace*{1em}\hspace*{1em}\hspace*{1em}\hspace*{1em} Proceedings of the IFIPS Congress, October 1980{</\textbf{title}>}.\mbox{}\newline 
\hspace*{1em}\hspace*{1em}{<\textbf{editor}>}Simon H. Lavington{</\textbf{editor}>}\mbox{}\newline 
\hspace*{1em}\hspace*{1em}\hspace*{1em}{<\textbf{publisher}>}North-Holland{</\textbf{publisher}>}:\mbox{}\newline 
\hspace*{1em}\hspace*{1em}{<\textbf{pubPlace}>}Amsterdam{</\textbf{pubPlace}>},\mbox{}\newline 
\hspace*{1em}\hspace*{1em}{<\textbf{date}>}1980{</\textbf{date}>}.\mbox{}\newline 
\hspace*{1em}\hspace*{1em}{<\textbf{biblScope}>}pp 1013–23\mbox{}\newline 
\hspace*{1em}\hspace*{1em}\hspace*{1em}{</\textbf{biblScope}>}\mbox{}\newline 
\hspace*{1em}\hspace*{1em}\hspace*{1em}{<\textbf{note}>}Apparently a draft of section 4 of\mbox{}\newline 
\hspace*{1em}\hspace*{1em}\hspace*{1em}{<\textbf{title}>}Literary Machines{</\textbf{title}>}.{</\textbf{note}>}\mbox{}\newline 
\hspace*{1em}\hspace*{1em}{</\textbf{bibl}>}\mbox{}\newline 
\hspace*{1em}{</\textbf{item}>}\mbox{}\newline 
\hspace*{1em}{<\textbf{item}>}\mbox{}\newline 
\hspace*{1em}\hspace*{1em}{<\textbf{bibl}\hspace*{1em}{xml:id}="{NEL88}">}Ted Nelson: {<\textbf{title}>}Literary Machines{</\textbf{title}>}\mbox{}\newline 
\hspace*{1em}\hspace*{1em}\hspace*{1em}\hspace*{1em} (privately published, 1987){</\textbf{bibl}>}\mbox{}\newline 
\hspace*{1em}{</\textbf{item}>}\mbox{}\newline 
\hspace*{1em}{<\textbf{item}>}\mbox{}\newline 
\hspace*{1em}\hspace*{1em}{<\textbf{bibl}\hspace*{1em}{xml:id}="{BAX88}">}\mbox{}\newline 
\hspace*{1em}\hspace*{1em}\hspace*{1em}{<\textbf{author}>}Baxter, Glen{</\textbf{author}>}\mbox{}\newline 
\hspace*{1em}\hspace*{1em}\hspace*{1em}{<\textbf{title}>}Glen Baxter His Life: the years of struggle{</\textbf{title}>}\mbox{}\newline 
\hspace*{1em}\hspace*{1em}\hspace*{1em}\hspace*{1em} London: Thames and Hudson, 1988.\mbox{}\newline 
\hspace*{1em}\hspace*{1em}{</\textbf{bibl}>}\mbox{}\newline 
\hspace*{1em}{</\textbf{item}>}\mbox{}\newline 
{</\textbf{list}>}\end{shaded}\egroup\par 
\subsubsection[{Components of Bibliographic References}]{Components of Bibliographic References}\label{COBICO}\par
This section discusses commonly occurring components of bibliographic references and elements used for encoding them. They fall into four groups: \begin{itemize}
\item elements for grouping components of the \textit{analytic}, \textit{monographic}, and \textit{series} levels in a structured bibliographic reference
\item titles of various kinds, and statements of intellectual responsibility (authorship, etc.)
\item information relating to the publication, pagination, etc. of an item (most of these constitute the default members of the \textsf{model.biblPart} class)
\item annotation, commentary, and further detail
\end{itemize}  The following sections describe the elements which may be used to represent such information within a \hyperref[TEI.bibl]{<bibl>} or \hyperref[TEI.biblStruct]{<biblStruct>} element. Within the former, elements from the \textsf{model.biblPart} class, other phrase-level elements, and plain text may be combined without other constraint; within the latter, such of these elements as exist for a given reference must be distinguished, and must also be presented in a specific order, discussed further below (section \textit{\hyperref[COBICOO]{3.12.2.9.\ Order of Components within References}}).
\paragraph[{Analytic, Monographic, and Series Levels}]{Analytic, Monographic, and Series Levels}\label{COBICOL}\par
In common library practice a clear distinction is made between an individual item within a larger collection and a free-standing book, journal, or collection. Similarly a book in a series is distinguished sharply from the series within which it appears. An article forming part of a collection which itself appears in a series thus has a bibliographic description with three quite distinct levels of information: \begin{enumerate}
\item the \textit{analytic} level, giving the title, author, etc., of the article;
\item the \textit{monographic} level, giving the title, editor, etc., of the collection;
\item the \textit{series} level, giving the title of the series, possibly the names of its editors, etc., and the number of the volume within that series.
\end{enumerate} In the same way, an article in a journal requires at least two levels of information: the analytic level describing the article itself, and the monographic level describing the journal.\par
A different identifying number may be supplied for any of these three items, that is, for the analytic item, the monographic item, or the series. \par
Within \hyperref[TEI.bibl]{<bibl>}, these three levels may be distinguished simply by the use of the {\itshape level} attribute on \hyperref[TEI.title]{<title>}. They may also be distinguished through the practice of employing nested \hyperref[TEI.bibl]{<bibl>} elements. In this example, for instance, the monograph-level component of the reference is encapsulated in its own \hyperref[TEI.bibl]{<bibl>} within the main \hyperref[TEI.bibl]{<bibl>} for the article: \par\bgroup\index{bibl=<bibl>|exampleindex}\index{type=@type!<bibl>|exampleindex}\index{subtype=@subtype!<bibl>|exampleindex}\index{author=<author>|exampleindex}\index{name=<name>|exampleindex}\index{surname=<surname>|exampleindex}\index{forename=<forename>|exampleindex}\index{title=<title>|exampleindex}\index{level=@level!<title>|exampleindex}\index{bibl=<bibl>|exampleindex}\index{type=@type!<bibl>|exampleindex}\index{title=<title>|exampleindex}\index{level=@level!<title>|exampleindex}\index{date=<date>|exampleindex}\index{when=@when!<date>|exampleindex}\index{biblScope=<biblScope>|exampleindex}\index{unit=@unit!<biblScope>|exampleindex}\index{from=@from!<biblScope>|exampleindex}\index{to=@to!<biblScope>|exampleindex}\exampleFont \begin{shaded}\noindent\mbox{}{<\textbf{bibl}\hspace*{1em}{type}="{article}"\mbox{}\newline 
\hspace*{1em}{subtype}="{magazine\textunderscore article}"\hspace*{1em}{xml:id}="{beaupaire\textunderscore 1911}">}\mbox{}\newline 
\hspace*{1em}{<\textbf{author}>}\mbox{}\newline 
\hspace*{1em}\hspace*{1em}{<\textbf{name}>}\mbox{}\newline 
\hspace*{1em}\hspace*{1em}\hspace*{1em}{<\textbf{surname}>}Beaupaire{</\textbf{surname}>}\mbox{}\newline 
\hspace*{1em}\hspace*{1em}\hspace*{1em}\hspace*{1em} ({<\textbf{forename}>}Edmond{</\textbf{forename}>}){</\textbf{name}>}\mbox{}\newline 
\hspace*{1em}{</\textbf{author}>},\mbox{}\newline 
{<\textbf{title}\hspace*{1em}{level}="{a}">}A propos de la rue de la Femme-sans-Tête{</\textbf{title}>},\mbox{}\newline 
{<\textbf{bibl}\hspace*{1em}{type}="{monogr}">}\mbox{}\newline 
\hspace*{1em}\hspace*{1em}{<\textbf{title}\hspace*{1em}{level}="{j}">}La Cité{</\textbf{title}>},\mbox{}\newline 
\hspace*{1em}{<\textbf{date}\hspace*{1em}{when}="{1911-01}">}janvier 1911{</\textbf{date}>}, pp. {<\textbf{biblScope}\hspace*{1em}{unit}="{page}"\hspace*{1em}{from}="{5}"\hspace*{1em}{to}="{17}">}5-17{</\textbf{biblScope}>}.\mbox{}\newline 
\hspace*{1em}{</\textbf{bibl}>}\mbox{}\newline 
{</\textbf{bibl}>}\end{shaded}\egroup\par \par
Within \hyperref[TEI.biblStruct]{<biblStruct>}, the levels are distinguished by the use of the following distinct elements: 
\begin{sansreflist}
  
\item [\textbf{<analytic>}] (analytic level) contains bibliographic elements describing an item (e.g. an article or poem) published within a monograph or journal and not as an independent publication.
\item [\textbf{<monogr>}] (monographic level) contains bibliographic elements describing an item (e.g. a book or journal) published as an independent item (i.e. as a separate physical object).
\item [\textbf{<series>}] (series information) contains information about the series in which a book or other bibliographic item has appeared.
\end{sansreflist}
\par
For purposes of TEI encoding, journals and anthologies are both treated as monographs; a journal title should thus be tagged as a <title level="j"> element within a \hyperref[TEI.monogr]{<monogr>} element. Individual articles in the journal or collected texts should be treated at the ‘analytic’ level. When an article has been printed in more than one journal or collection, the bibliographic reference may have more than one \hyperref[TEI.monogr]{<monogr>} element, each possibly followed by one or more \hyperref[TEI.series]{<series>} elements. A \hyperref[TEI.series]{<series>} element always relates to the most recently preceding \hyperref[TEI.monogr]{<monogr>} element. (Whether reprints of an article are treated in the same bibliographic reference or a separate one varies among different styles. Library lists typically use a different entry for each publication, while academic footnoting practice typically treats all publications of the same article in a single entry.)\par
The \hyperref[TEI.biblScope]{<biblScope>} element is used to supply further information about the location of some part of a bibliographic reference. It specifies where to find the component in which it appears within the immediately preceding component of a different level.\par
In the following example, Schacter's article \textit{Iolaos} appeared on pages 64 to 70 of a volume entitled \textit{Herakles to Poseidon}, which was itself the second of a four volumes published together under the title \textit{Cults of Boitia}; this last title constituted the 38th volume in the series of \textit{Bulletin of the Institute of Classical Studies Supplements}: \par\bgroup\index{biblStruct=<biblStruct>|exampleindex}\index{analytic=<analytic>|exampleindex}\index{author=<author>|exampleindex}\index{title=<title>|exampleindex}\index{level=@level!<title>|exampleindex}\index{monogr=<monogr>|exampleindex}\index{title=<title>|exampleindex}\index{level=@level!<title>|exampleindex}\index{imprint=<imprint>|exampleindex}\index{date=<date>|exampleindex}\index{biblScope=<biblScope>|exampleindex}\index{unit=@unit!<biblScope>|exampleindex}\index{monogr=<monogr>|exampleindex}\index{title=<title>|exampleindex}\index{level=@level!<title>|exampleindex}\index{imprint=<imprint>|exampleindex}\index{pubPlace=<pubPlace>|exampleindex}\index{extent=<extent>|exampleindex}\index{biblScope=<biblScope>|exampleindex}\index{unit=@unit!<biblScope>|exampleindex}\index{series=<series>|exampleindex}\index{title=<title>|exampleindex}\index{level=@level!<title>|exampleindex}\index{biblScope=<biblScope>|exampleindex}\index{unit=@unit!<biblScope>|exampleindex}\exampleFont \begin{shaded}\noindent\mbox{}{<\textbf{biblStruct}>}\mbox{}\newline 
\hspace*{1em}{<\textbf{analytic}>}\mbox{}\newline 
\hspace*{1em}\hspace*{1em}{<\textbf{author}>}Albert Schachter{</\textbf{author}>}\mbox{}\newline 
\hspace*{1em}\hspace*{1em}{<\textbf{title}\hspace*{1em}{level}="{a}">}Iolaos{</\textbf{title}>}\mbox{}\newline 
\hspace*{1em}{</\textbf{analytic}>}\mbox{}\newline 
\hspace*{1em}{<\textbf{monogr}>}\mbox{}\newline 
\hspace*{1em}\hspace*{1em}{<\textbf{title}\hspace*{1em}{level}="{m}">}Herakles to Poseidon{</\textbf{title}>}\mbox{}\newline 
\hspace*{1em}\hspace*{1em}{<\textbf{imprint}>}\mbox{}\newline 
\hspace*{1em}\hspace*{1em}\hspace*{1em}{<\textbf{date}>}1986{</\textbf{date}>}\mbox{}\newline 
\hspace*{1em}\hspace*{1em}{</\textbf{imprint}>}\mbox{}\newline 
\hspace*{1em}\hspace*{1em}{<\textbf{biblScope}\hspace*{1em}{unit}="{page}">}64-70{</\textbf{biblScope}>}\mbox{}\newline 
\hspace*{1em}{</\textbf{monogr}>}\mbox{}\newline 
\hspace*{1em}{<\textbf{monogr}>}\mbox{}\newline 
\hspace*{1em}\hspace*{1em}{<\textbf{title}\hspace*{1em}{level}="{m}">}Cults of Boiotia{</\textbf{title}>}\mbox{}\newline 
\hspace*{1em}\hspace*{1em}{<\textbf{imprint}>}\mbox{}\newline 
\hspace*{1em}\hspace*{1em}\hspace*{1em}{<\textbf{pubPlace}>}London{</\textbf{pubPlace}>}\mbox{}\newline 
\hspace*{1em}\hspace*{1em}{</\textbf{imprint}>}\mbox{}\newline 
\hspace*{1em}\hspace*{1em}{<\textbf{extent}>}4 vols.{</\textbf{extent}>}\mbox{}\newline 
\hspace*{1em}\hspace*{1em}{<\textbf{biblScope}\hspace*{1em}{unit}="{part}">}2{</\textbf{biblScope}>}\mbox{}\newline 
\hspace*{1em}{</\textbf{monogr}>}\mbox{}\newline 
\hspace*{1em}{<\textbf{series}>}\mbox{}\newline 
\hspace*{1em}\hspace*{1em}{<\textbf{title}\hspace*{1em}{level}="{s}">}Bulletin of the Institute of Classical Studies\mbox{}\newline 
\hspace*{1em}\hspace*{1em}\hspace*{1em}\hspace*{1em} Supplements{</\textbf{title}>}\mbox{}\newline 
\hspace*{1em}\hspace*{1em}{<\textbf{biblScope}\hspace*{1em}{unit}="{volume}">}38{</\textbf{biblScope}>}\mbox{}\newline 
\hspace*{1em}{</\textbf{series}>}\mbox{}\newline 
{</\textbf{biblStruct}>}\end{shaded}\egroup\par \par
In the following example, the article cited has been published twice, once in a journal (where it appeared in volume 40, on pages 3 -46 of the issue of October 1986) and once as a free-standing item, which appeared as number 11 of a German language series. \par\bgroup\index{biblStruct=<biblStruct>|exampleindex}\index{analytic=<analytic>|exampleindex}\index{author=<author>|exampleindex}\index{persName=<persName>|exampleindex}\index{surname=<surname>|exampleindex}\index{forename=<forename>|exampleindex}\index{title=<title>|exampleindex}\index{level=@level!<title>|exampleindex}\index{monogr=<monogr>|exampleindex}\index{title=<title>|exampleindex}\index{level=@level!<title>|exampleindex}\index{imprint=<imprint>|exampleindex}\index{date=<date>|exampleindex}\index{when=@when!<date>|exampleindex}\index{biblScope=<biblScope>|exampleindex}\index{unit=@unit!<biblScope>|exampleindex}\index{biblScope=<biblScope>|exampleindex}\index{unit=@unit!<biblScope>|exampleindex}\index{from=@from!<biblScope>|exampleindex}\index{to=@to!<biblScope>|exampleindex}\index{monogr=<monogr>|exampleindex}\index{title=<title>|exampleindex}\index{level=@level!<title>|exampleindex}\index{editor=<editor>|exampleindex}\index{persName=<persName>|exampleindex}\index{forename=<forename>|exampleindex}\index{forename=<forename>|exampleindex}\index{full=@full!<forename>|exampleindex}\index{surname=<surname>|exampleindex}\index{imprint=<imprint>|exampleindex}\index{pubPlace=<pubPlace>|exampleindex}\index{publisher=<publisher>|exampleindex}\index{date=<date>|exampleindex}\index{when=@when!<date>|exampleindex}\index{series=<series>|exampleindex}\index{title=<title>|exampleindex}\index{level=@level!<title>|exampleindex}\index{respStmt=<respStmt>|exampleindex}\index{resp=<resp>|exampleindex}\index{name=<name>|exampleindex}\index{type=@type!<name>|exampleindex}\index{name=<name>|exampleindex}\index{type=@type!<name>|exampleindex}\index{title=<title>|exampleindex}\index{level=@level!<title>|exampleindex}\index{biblScope=<biblScope>|exampleindex}\index{unit=@unit!<biblScope>|exampleindex}\exampleFont \begin{shaded}\noindent\mbox{}{<\textbf{biblStruct}>}\mbox{}\newline 
\hspace*{1em}{<\textbf{analytic}>}\mbox{}\newline 
\hspace*{1em}\hspace*{1em}{<\textbf{author}>}\mbox{}\newline 
\hspace*{1em}\hspace*{1em}\hspace*{1em}{<\textbf{persName}>}\mbox{}\newline 
\hspace*{1em}\hspace*{1em}\hspace*{1em}\hspace*{1em}{<\textbf{surname}>}Thaller{</\textbf{surname}>}\mbox{}\newline 
\hspace*{1em}\hspace*{1em}\hspace*{1em}\hspace*{1em}{<\textbf{forename}>}Manfred{</\textbf{forename}>}\mbox{}\newline 
\hspace*{1em}\hspace*{1em}\hspace*{1em}{</\textbf{persName}>}\mbox{}\newline 
\hspace*{1em}\hspace*{1em}{</\textbf{author}>}\mbox{}\newline 
\hspace*{1em}\hspace*{1em}{<\textbf{title}\hspace*{1em}{level}="{a}">}A Draft Proposal for a Standard for the\mbox{}\newline 
\hspace*{1em}\hspace*{1em}\hspace*{1em}\hspace*{1em} Coding of Machine Readable Sources{</\textbf{title}>}\mbox{}\newline 
\hspace*{1em}{</\textbf{analytic}>}\mbox{}\newline 
\hspace*{1em}{<\textbf{monogr}>}\mbox{}\newline 
\hspace*{1em}\hspace*{1em}{<\textbf{title}\hspace*{1em}{level}="{j}">}Historical Social Research{</\textbf{title}>}\mbox{}\newline 
\hspace*{1em}\hspace*{1em}{<\textbf{imprint}>}\mbox{}\newline 
\hspace*{1em}\hspace*{1em}\hspace*{1em}{<\textbf{date}\hspace*{1em}{when}="{1986-10}">}October 1986{</\textbf{date}>}\mbox{}\newline 
\hspace*{1em}\hspace*{1em}{</\textbf{imprint}>}\mbox{}\newline 
\hspace*{1em}\hspace*{1em}{<\textbf{biblScope}\hspace*{1em}{unit}="{volume}">}40{</\textbf{biblScope}>}\mbox{}\newline 
\hspace*{1em}\hspace*{1em}{<\textbf{biblScope}\hspace*{1em}{unit}="{page}"\hspace*{1em}{from}="{3}"\hspace*{1em}{to}="{46}">}3-46{</\textbf{biblScope}>}\mbox{}\newline 
\hspace*{1em}{</\textbf{monogr}>}\mbox{}\newline 
\hspace*{1em}{<\textbf{monogr}>}\mbox{}\newline 
\hspace*{1em}\hspace*{1em}{<\textbf{title}\hspace*{1em}{level}="{m}">}Modelling Historical Data:\mbox{}\newline 
\hspace*{1em}\hspace*{1em}\hspace*{1em}\hspace*{1em} Towards a Standard for Encoding and\mbox{}\newline 
\hspace*{1em}\hspace*{1em}\hspace*{1em}\hspace*{1em} Exchanging Machine-Readable Texts{</\textbf{title}>}\mbox{}\newline 
\hspace*{1em}\hspace*{1em}{<\textbf{editor}>}\mbox{}\newline 
\hspace*{1em}\hspace*{1em}\hspace*{1em}{<\textbf{persName}>}\mbox{}\newline 
\hspace*{1em}\hspace*{1em}\hspace*{1em}\hspace*{1em}{<\textbf{forename}>}Daniel{</\textbf{forename}>}\mbox{}\newline 
\hspace*{1em}\hspace*{1em}\hspace*{1em}\hspace*{1em}{<\textbf{forename}\hspace*{1em}{full}="{init}">}I.{</\textbf{forename}>}\mbox{}\newline 
\hspace*{1em}\hspace*{1em}\hspace*{1em}\hspace*{1em}{<\textbf{surname}>}Greenstein{</\textbf{surname}>}\mbox{}\newline 
\hspace*{1em}\hspace*{1em}\hspace*{1em}{</\textbf{persName}>}\mbox{}\newline 
\hspace*{1em}\hspace*{1em}{</\textbf{editor}>}\mbox{}\newline 
\hspace*{1em}\hspace*{1em}{<\textbf{imprint}\hspace*{1em}{xml:lang}="{de}">}\mbox{}\newline 
\hspace*{1em}\hspace*{1em}\hspace*{1em}{<\textbf{pubPlace}>}St. Katharinen{</\textbf{pubPlace}>}\mbox{}\newline 
\hspace*{1em}\hspace*{1em}\hspace*{1em}{<\textbf{publisher}>}Max-Planck-Institut für Geschichte\mbox{}\newline 
\hspace*{1em}\hspace*{1em}\hspace*{1em}\hspace*{1em}\hspace*{1em}\hspace*{1em} In Kommission bei\mbox{}\newline 
\hspace*{1em}\hspace*{1em}\hspace*{1em}\hspace*{1em}\hspace*{1em}\hspace*{1em} Scripta Mercaturae Verlag{</\textbf{publisher}>}\mbox{}\newline 
\hspace*{1em}\hspace*{1em}\hspace*{1em}{<\textbf{date}\hspace*{1em}{when}="{1991}"/>}\mbox{}\newline 
\hspace*{1em}\hspace*{1em}{</\textbf{imprint}>}\mbox{}\newline 
\hspace*{1em}{</\textbf{monogr}>}\mbox{}\newline 
\hspace*{1em}{<\textbf{series}\hspace*{1em}{xml:lang}="{de}">}\mbox{}\newline 
\hspace*{1em}\hspace*{1em}{<\textbf{title}\hspace*{1em}{level}="{s}">}Halbgraue Reihe\mbox{}\newline 
\hspace*{1em}\hspace*{1em}\hspace*{1em}\hspace*{1em} zur Historischen Fachinformatik{</\textbf{title}>}\mbox{}\newline 
\hspace*{1em}\hspace*{1em}{<\textbf{respStmt}>}\mbox{}\newline 
\hspace*{1em}\hspace*{1em}\hspace*{1em}{<\textbf{resp}>}Herausgegeben von{</\textbf{resp}>}\mbox{}\newline 
\hspace*{1em}\hspace*{1em}\hspace*{1em}{<\textbf{name}\hspace*{1em}{type}="{person}">}Manfred Thaller{</\textbf{name}>}\mbox{}\newline 
\hspace*{1em}\hspace*{1em}\hspace*{1em}{<\textbf{name}\hspace*{1em}{type}="{org}">}Max-Planck-Institut für Geschichte{</\textbf{name}>}\mbox{}\newline 
\hspace*{1em}\hspace*{1em}{</\textbf{respStmt}>}\mbox{}\newline 
\hspace*{1em}\hspace*{1em}{<\textbf{title}\hspace*{1em}{level}="{s}">}Serie A: Historische Quellenkunden{</\textbf{title}>}\mbox{}\newline 
\hspace*{1em}\hspace*{1em}{<\textbf{biblScope}\hspace*{1em}{unit}="{volume}">}11{</\textbf{biblScope}>}\mbox{}\newline 
\hspace*{1em}{</\textbf{series}>}\mbox{}\newline 
{</\textbf{biblStruct}>}\end{shaded}\egroup\par \par
The practice of analytic vs. monographic citation, as described here, should be distinguished from the practice of including within one citation a reference to another work, which the encoder considers to be related to in some way: see further \textit{\hyperref[COBIRI]{3.12.2.7.\ Related Items}} below.\par
If an identifier is available for the analytic item, it should be represented by means of an \hyperref[TEI.idno]{<idno>} element placed within the \hyperref[TEI.analytic]{<analytic>} element, as in the following example where a DOI (Digital Object identifier) is supplied for the article in question.\par\bgroup\index{biblStruct=<biblStruct>|exampleindex}\index{analytic=<analytic>|exampleindex}\index{author=<author>|exampleindex}\index{forename=<forename>|exampleindex}\index{forename=<forename>|exampleindex}\index{surname=<surname>|exampleindex}\index{author=<author>|exampleindex}\index{forename=<forename>|exampleindex}\index{surname=<surname>|exampleindex}\index{author=<author>|exampleindex}\index{forename=<forename>|exampleindex}\index{forename=<forename>|exampleindex}\index{surname=<surname>|exampleindex}\index{title=<title>|exampleindex}\index{level=@level!<title>|exampleindex}\index{idno=<idno>|exampleindex}\index{type=@type!<idno>|exampleindex}\index{ref=<ref>|exampleindex}\index{type=@type!<ref>|exampleindex}\index{monogr=<monogr>|exampleindex}\index{title=<title>|exampleindex}\index{level=@level!<title>|exampleindex}\index{imprint=<imprint>|exampleindex}\index{date=<date>|exampleindex}\index{biblScope=<biblScope>|exampleindex}\index{unit=@unit!<biblScope>|exampleindex}\index{biblScope=<biblScope>|exampleindex}\index{unit=@unit!<biblScope>|exampleindex}\index{biblScope=<biblScope>|exampleindex}\index{unit=@unit!<biblScope>|exampleindex}\exampleFont \begin{shaded}\noindent\mbox{}{<\textbf{biblStruct}>}\mbox{}\newline 
\hspace*{1em}{<\textbf{analytic}>}\mbox{}\newline 
\hspace*{1em}\hspace*{1em}{<\textbf{author}>}\mbox{}\newline 
\hspace*{1em}\hspace*{1em}\hspace*{1em}{<\textbf{forename}>}James{</\textbf{forename}>}\mbox{}\newline 
\hspace*{1em}\hspace*{1em}\hspace*{1em}{<\textbf{forename}>}H.{</\textbf{forename}>}\mbox{}\newline 
\hspace*{1em}\hspace*{1em}\hspace*{1em}{<\textbf{surname}>}Coombs{</\textbf{surname}>}\mbox{}\newline 
\hspace*{1em}\hspace*{1em}{</\textbf{author}>}\mbox{}\newline 
\hspace*{1em}\hspace*{1em}{<\textbf{author}>}\mbox{}\newline 
\hspace*{1em}\hspace*{1em}\hspace*{1em}{<\textbf{forename}>}Allen{</\textbf{forename}>}\mbox{}\newline 
\hspace*{1em}\hspace*{1em}\hspace*{1em}{<\textbf{surname}>}Renear{</\textbf{surname}>}\mbox{}\newline 
\hspace*{1em}\hspace*{1em}{</\textbf{author}>}\mbox{}\newline 
\hspace*{1em}\hspace*{1em}{<\textbf{author}>}\mbox{}\newline 
\hspace*{1em}\hspace*{1em}\hspace*{1em}{<\textbf{forename}>}Steven{</\textbf{forename}>}\mbox{}\newline 
\hspace*{1em}\hspace*{1em}\hspace*{1em}{<\textbf{forename}>}J.{</\textbf{forename}>}\mbox{}\newline 
\hspace*{1em}\hspace*{1em}\hspace*{1em}{<\textbf{surname}>}DeRose{</\textbf{surname}>}\mbox{}\newline 
\hspace*{1em}\hspace*{1em}{</\textbf{author}>}\mbox{}\newline 
\hspace*{1em}\hspace*{1em}{<\textbf{title}\hspace*{1em}{level}="{a}">}Markup Systems and The Future of Scholarly Text\mbox{}\newline 
\hspace*{1em}\hspace*{1em}\hspace*{1em}\hspace*{1em} Processing{</\textbf{title}>}\mbox{}\newline 
\hspace*{1em}\hspace*{1em}{<\textbf{idno}\hspace*{1em}{type}="{DOI}">}10.1145/32206.32209{</\textbf{idno}>}\mbox{}\newline 
\hspace*{1em}\hspace*{1em}{<\textbf{ref}\hspace*{1em}{type}="{url}">}http://xml.coverpages.org/coombs.html{</\textbf{ref}>}\mbox{}\newline 
\hspace*{1em}{</\textbf{analytic}>}\mbox{}\newline 
\hspace*{1em}{<\textbf{monogr}>}\mbox{}\newline 
\hspace*{1em}\hspace*{1em}{<\textbf{title}\hspace*{1em}{level}="{j}">}Communications of the ACM{</\textbf{title}>}\mbox{}\newline 
\hspace*{1em}\hspace*{1em}{<\textbf{imprint}>}\mbox{}\newline 
\hspace*{1em}\hspace*{1em}\hspace*{1em}{<\textbf{date}>}1987{</\textbf{date}>}\mbox{}\newline 
\hspace*{1em}\hspace*{1em}{</\textbf{imprint}>}\mbox{}\newline 
\hspace*{1em}\hspace*{1em}{<\textbf{biblScope}\hspace*{1em}{unit}="{volume}">}30{</\textbf{biblScope}>}\mbox{}\newline 
\hspace*{1em}\hspace*{1em}{<\textbf{biblScope}\hspace*{1em}{unit}="{issue}">}11{</\textbf{biblScope}>}\mbox{}\newline 
\hspace*{1em}\hspace*{1em}{<\textbf{biblScope}\hspace*{1em}{unit}="{page}">}933–947{</\textbf{biblScope}>}\mbox{}\newline 
\hspace*{1em}{</\textbf{monogr}>}\mbox{}\newline 
{</\textbf{biblStruct}>}\end{shaded}\egroup\par \par
Punctuation must not appear between the elements within a structured bibliographic entry encoded with \hyperref[TEI.biblStruct]{<biblStruct>} or \hyperref[TEI.biblFull]{<biblFull>}, unless it is contained within the elements it delimits. When (as in most of the examples in this chapter) entries are encoded without any inter-element punctuation, they can be usually be processed more easily by rendering systems able to output bibliographic references in any of several styles.\par
Within a \hyperref[TEI.bibl]{<bibl>} however, it is possible and often convenient to include punctuation. \par\bgroup\index{bibl=<bibl>|exampleindex}\index{author=<author>|exampleindex}\index{persName=<persName>|exampleindex}\index{surname=<surname>|exampleindex}\index{forename=<forename>|exampleindex}\index{forename=<forename>|exampleindex}\index{date=<date>|exampleindex}\index{when=@when!<date>|exampleindex}\index{title=<title>|exampleindex}\index{level=@level!<title>|exampleindex}\index{title=<title>|exampleindex}\index{level=@level!<title>|exampleindex}\index{editor=<editor>|exampleindex}\index{persName=<persName>|exampleindex}\index{forename=<forename>|exampleindex}\index{forename=<forename>|exampleindex}\index{surname=<surname>|exampleindex}\index{biblScope=<biblScope>|exampleindex}\index{unit=@unit!<biblScope>|exampleindex}\index{pubPlace=<pubPlace>|exampleindex}\index{publisher=<publisher>|exampleindex}\index{note=<note>|exampleindex}\index{ref=<ref>|exampleindex}\index{target=@target!<ref>|exampleindex}\index{title=<title>|exampleindex}\index{level=@level!<title>|exampleindex}\exampleFont \begin{shaded}\noindent\mbox{}{<\textbf{bibl}\hspace*{1em}{xml:id}="{NELSON\textunderscore 80}">}\mbox{}\newline 
\hspace*{1em}{<\textbf{author}>}\mbox{}\newline 
\hspace*{1em}\hspace*{1em}{<\textbf{persName}>}\mbox{}\newline 
\hspace*{1em}\hspace*{1em}\hspace*{1em}{<\textbf{surname}>}Nelson{</\textbf{surname}>},\mbox{}\newline 
\hspace*{1em}\hspace*{1em}{<\textbf{forename}>}T.{</\textbf{forename}>}\mbox{}\newline 
\hspace*{1em}\hspace*{1em}\hspace*{1em}{<\textbf{forename}>}H.{</\textbf{forename}>}\mbox{}\newline 
\hspace*{1em}\hspace*{1em}{</\textbf{persName}>}\mbox{}\newline 
\hspace*{1em}{</\textbf{author}>}\mbox{}\newline 
\hspace*{1em}{<\textbf{date}\hspace*{1em}{when}="{1980}">}1980{</\textbf{date}>}.\mbox{}\newline 
{<\textbf{title}\hspace*{1em}{level}="{a}">}Replacing the printed word: a complete literary\mbox{}\newline 
\hspace*{1em}\hspace*{1em} system{</\textbf{title}>}. In {<\textbf{title}\hspace*{1em}{level}="{m}">}Information Processing '80: Proceedings of the\mbox{}\newline 
\hspace*{1em}\hspace*{1em} IFIPS Congress, October 1980{</\textbf{title}>},\mbox{}\newline 
 ed.\mbox{}\newline 
{<\textbf{editor}>}\mbox{}\newline 
\hspace*{1em}\hspace*{1em}{<\textbf{persName}>}\mbox{}\newline 
\hspace*{1em}\hspace*{1em}\hspace*{1em}{<\textbf{forename}>}Simon{</\textbf{forename}>}\mbox{}\newline 
\hspace*{1em}\hspace*{1em}\hspace*{1em}{<\textbf{forename}>}H.{</\textbf{forename}>}\mbox{}\newline 
\hspace*{1em}\hspace*{1em}\hspace*{1em}{<\textbf{surname}>}Lavington{</\textbf{surname}>}\mbox{}\newline 
\hspace*{1em}\hspace*{1em}{</\textbf{persName}>}\mbox{}\newline 
\hspace*{1em}{</\textbf{editor}>},\mbox{}\newline 
{<\textbf{biblScope}\hspace*{1em}{unit}="{page}">}1013-23{</\textbf{biblScope}>}.\mbox{}\newline 
{<\textbf{pubPlace}>}Amsterdam{</\textbf{pubPlace}>}: {<\textbf{publisher}>}North-\mbox{}\newline 
\hspace*{1em}\hspace*{1em} Holland{</\textbf{publisher}>}. ({<\textbf{note}>}Apparently a draft of section 4 of\mbox{}\newline 
\hspace*{1em}{<\textbf{ref}\hspace*{1em}{target}="{\#NELSON\textunderscore 88}">}\mbox{}\newline 
\hspace*{1em}\hspace*{1em}\hspace*{1em}{<\textbf{title}\hspace*{1em}{level}="{m}">}Literary\mbox{}\newline 
\hspace*{1em}\hspace*{1em}\hspace*{1em}\hspace*{1em}\hspace*{1em}\hspace*{1em} Machines{</\textbf{title}>}\mbox{}\newline 
\hspace*{1em}\hspace*{1em}{</\textbf{ref}>}.{</\textbf{note}>})\mbox{}\newline 
\mbox{}\newline 
{</\textbf{bibl}>}\end{shaded}\egroup\par \noindent  This example shows the components sequenced and punctuated according to the Chicago style, with all the relevant data items marked up appropriately. This markup approach can provide easy rendering, if only one styleguide is targeted, or an original source document uses a specific styleguide, while still allowing for automated recovery of key data items such as names of authors, titles etc.
\paragraph[{Titles, Authors, and Editors}]{Titles, Authors, and Editors}\label{COBICOR}\par
Bibliographic references typically include the title of the work being cited and the names of those intellectually responsible for it. For articles in journals or collections, such statements should appear both for the analytic and for the monographic level. The following elements are provided for tagging such elements: 
\begin{sansreflist}
  
\item [\textbf{<title>}] (title) contains a title for any kind of work.
\item [\textbf{<author>}] (author) in a bibliographic reference, contains the name(s) of an author, personal or corporate, of a work; for example in the same form as that provided by a recognized bibliographic name authority.
\item [\textbf{<editor>}] contains a secondary statement of responsibility for a bibliographic item, for example the name of an individual, institution or organization, (or of several such) acting as editor, compiler, translator, etc.
\item [\textbf{<respStmt>}] (statement of responsibility) supplies a statement of responsibility for the intellectual content of a text, edition, recording, or series, where the specialized elements for authors, editors, etc. do not suffice or do not apply. May also be used to encode information about individuals or organizations which have played a role in the production or distribution of a bibliographic work.
\item [\textbf{<resp>}] (responsibility) contains a phrase describing the nature of a person's intellectual responsibility, or an organization's role in the production or distribution of a work.
\item [\textbf{<name>}] (name, proper noun) contains a proper noun or noun phrase.
\item [\textbf{<meeting>}] contains the formalized descriptive title for a meeting or conference, for use in a bibliographic description for an item derived from such a meeting, or as a heading or preamble to publications emanating from it.
\item [\textbf{<sponsor>}] (sponsor) specifies the name of a sponsoring organization or institution.
\item [\textbf{<funder>}] (funding body) specifies the name of an individual, institution, or organization responsible for the funding of a project or text.
\item [\textbf{<distributor>}] (distributor) supplies the name of a person or other agency responsible for the distribution of a text.
\item [\textbf{<principal>}] (principal researcher) supplies the name of the principal researcher responsible for the creation of an electronic text.
\end{sansreflist}
 The elements \hyperref[TEI.author]{<author>}, \hyperref[TEI.editor]{<editor>}, \hyperref[TEI.respStmt]{<respStmt>}, \hyperref[TEI.meeting]{<meeting>}, \hyperref[TEI.sponsor]{<sponsor>}, \hyperref[TEI.funder]{<funder>}, and \hyperref[TEI.principal]{<principal>} are the default members of the \textsf{model.respLike} class, a subclass of the \textsf{model.biblPart} class to which the constituents of the \hyperref[TEI.bibl]{<bibl>} element belong.\par
In bibliographic references, all titles should be tagged as such, whether analytic, monographic, or series titles. The single element \hyperref[TEI.title]{<title>} is used for all these cases. When it appears directly within an \hyperref[TEI.analytic]{<analytic>}, \hyperref[TEI.monogr]{<monogr>}, or \hyperref[TEI.series]{<series>} element, \hyperref[TEI.title]{<title>} is interpreted as belonging to the appropriate level. However, it is recommended that the {\itshape level} attribute be used to signal this explicitly.\par
It is a semantic error to give a value for the {\itshape level} attribute which is inconsistent with the context. The {\itshape level} value a implies the analytic level; the values m, j, and u imply the monographic level; the value s implies the series level. Note, however, that the semantic error occurs only if the nested title is directly enclosed by the \hyperref[TEI.analytic]{<analytic>}, \hyperref[TEI.monogr]{<monogr>}, or \hyperref[TEI.series]{<series>} element; if it is enclosed only indirectly (i.e., nested more deeply), no semantic error need be present. For example, the analytic title may contain a monographic title, as in the following example: \par\bgroup\index{biblStruct=<biblStruct>|exampleindex}\index{analytic=<analytic>|exampleindex}\index{author=<author>|exampleindex}\index{ref=@ref!<author>|exampleindex}\index{persName=<persName>|exampleindex}\index{forename=<forename>|exampleindex}\index{forename=<forename>|exampleindex}\index{surname=<surname>|exampleindex}\index{title=<title>|exampleindex}\index{title=<title>|exampleindex}\index{level=@level!<title>|exampleindex}\index{monogr=<monogr>|exampleindex}\index{title=<title>|exampleindex}\index{level=@level!<title>|exampleindex}\index{imprint=<imprint>|exampleindex}\index{date=<date>|exampleindex}\index{biblScope=<biblScope>|exampleindex}\index{unit=@unit!<biblScope>|exampleindex}\index{biblScope=<biblScope>|exampleindex}\index{unit=@unit!<biblScope>|exampleindex}\exampleFont \begin{shaded}\noindent\mbox{}{<\textbf{biblStruct}>}\mbox{}\newline 
\hspace*{1em}{<\textbf{analytic}>}\mbox{}\newline 
\hspace*{1em}\hspace*{1em}{<\textbf{author}\hspace*{1em}{ref}="{http://id.loc.gov/authorities/names/no2001067434}">}\mbox{}\newline 
\hspace*{1em}\hspace*{1em}\hspace*{1em}{<\textbf{persName}>}\mbox{}\newline 
\hspace*{1em}\hspace*{1em}\hspace*{1em}\hspace*{1em}{<\textbf{forename}>}Lucy{</\textbf{forename}>}\mbox{}\newline 
\hspace*{1em}\hspace*{1em}\hspace*{1em}\hspace*{1em}{<\textbf{forename}>}Allen{</\textbf{forename}>}\mbox{}\newline 
\hspace*{1em}\hspace*{1em}\hspace*{1em}\hspace*{1em}{<\textbf{surname}>}Paton{</\textbf{surname}>}\mbox{}\newline 
\hspace*{1em}\hspace*{1em}\hspace*{1em}{</\textbf{persName}>}\mbox{}\newline 
\hspace*{1em}\hspace*{1em}{</\textbf{author}>}\mbox{}\newline 
\hspace*{1em}\hspace*{1em}{<\textbf{title}>}Notes on Manuscripts of the\mbox{}\newline 
\hspace*{1em}\hspace*{1em}{<\textbf{title}\hspace*{1em}{level}="{m}"\hspace*{1em}{xml:lang}="{fr}">}Prophécies de Merlin{</\textbf{title}>}\mbox{}\newline 
\hspace*{1em}\hspace*{1em}{</\textbf{title}>}\mbox{}\newline 
\hspace*{1em}{</\textbf{analytic}>}\mbox{}\newline 
\hspace*{1em}{<\textbf{monogr}>}\mbox{}\newline 
\hspace*{1em}\hspace*{1em}{<\textbf{title}\hspace*{1em}{level}="{j}">}PMLA{</\textbf{title}>}\mbox{}\newline 
\hspace*{1em}\hspace*{1em}{<\textbf{imprint}>}\mbox{}\newline 
\hspace*{1em}\hspace*{1em}\hspace*{1em}{<\textbf{date}>}1913{</\textbf{date}>}\mbox{}\newline 
\hspace*{1em}\hspace*{1em}{</\textbf{imprint}>}\mbox{}\newline 
\hspace*{1em}\hspace*{1em}{<\textbf{biblScope}\hspace*{1em}{unit}="{volume}">}8{</\textbf{biblScope}>}\mbox{}\newline 
\hspace*{1em}\hspace*{1em}{<\textbf{biblScope}\hspace*{1em}{unit}="{page}">}122{</\textbf{biblScope}>}\mbox{}\newline 
\hspace*{1em}{</\textbf{monogr}>}\mbox{}\newline 
{</\textbf{biblStruct}>}\end{shaded}\egroup\par \noindent  In this case, the analytic title ‘Notes on Manuscripts of the \textit{Prophécies de Merlin}’ needs no {\itshape level} attribute because it is directly contained by an \hyperref[TEI.analytic]{<analytic>} element; the monographic title contained within it, ‘Prophécies de Merlin’, is not semantically erroneous because it is not directly contained by the \hyperref[TEI.analytic]{<analytic>} element.\par
In some bibliographic applications, it may prove useful to distinguish main titles from subordinate titles, parallel titles, etc. The {\itshape type} attribute is provided to allow this distinction to be recorded.\par
The following reference, from a national standard for bibliographic references, illustrates this type of analysis with its distinction between main and subordinate titles. Note that this uses the more flexible \hyperref[TEI.bibl]{<bibl>}, rather than the structured \hyperref[TEI.biblStruct]{<biblStruct>} element: consequently, there is no requirement to tag all the components of the reference (notably the authors). \par\bgroup\index{bibl=<bibl>|exampleindex}\index{title=<title>|exampleindex}\index{level=@level!<title>|exampleindex}\index{type=@type!<title>|exampleindex}\index{title=<title>|exampleindex}\index{level=@level!<title>|exampleindex}\index{type=@type!<title>|exampleindex}\index{title=<title>|exampleindex}\index{level=@level!<title>|exampleindex}\index{date=<date>|exampleindex}\index{biblScope=<biblScope>|exampleindex}\index{unit=@unit!<biblScope>|exampleindex}\index{biblScope=<biblScope>|exampleindex}\index{unit=@unit!<biblScope>|exampleindex}\index{to=@to!<biblScope>|exampleindex}\exampleFont \begin{shaded}\noindent\mbox{}{<\textbf{bibl}>}Saarikoski, Pirkko-Liisa, and Paavo Suomalainen,\mbox{}\newline 
{<\textbf{title}\hspace*{1em}{level}="{a}"\hspace*{1em}{type}="{main}">}Studies on the physiology of\mbox{}\newline 
\hspace*{1em}\hspace*{1em} the hibernating hedgehog, 15{</\textbf{title}>}\mbox{}\newline 
\hspace*{1em}{<\textbf{title}\hspace*{1em}{level}="{a}"\hspace*{1em}{type}="{sub}">}Effects of seasonal\mbox{}\newline 
\hspace*{1em}\hspace*{1em} and temperature changes on the in vitro glycerol release from\mbox{}\newline 
\hspace*{1em}\hspace*{1em} brown adipose tissue{</\textbf{title}>}\mbox{}\newline 
\hspace*{1em}{<\textbf{title}\hspace*{1em}{level}="{j}">}Ann. Acad. Sci. Fenn., Ser. A4{</\textbf{title}>}\mbox{}\newline 
\hspace*{1em}{<\textbf{date}>}1972{</\textbf{date}>}\mbox{}\newline 
\hspace*{1em}{<\textbf{biblScope}\hspace*{1em}{unit}="{volume}">}187{</\textbf{biblScope}>}\mbox{}\newline 
\hspace*{1em}{<\textbf{biblScope}\hspace*{1em}{unit}="{page}"\hspace*{1em}{to}="{4}">}1-4{</\textbf{biblScope}>}\mbox{}\newline 
{</\textbf{bibl}>}\end{shaded}\egroup\par \noindent  \par
Slightly more complex is the distinction made below among main, subordinate, and parallel titles, in an example from the same source (p. 63). The punctuation and the bibliographic analysis are those given in ANSI Z39.29-1977; the punctuation is in the style prescribed by the International Standard Bibliographic Description (ISBD).\footnote{The analysis is not wholly unproblematic: as the text of the standard points out, the first subordinate title is subordinate only to the parallel title in French, while the second is subordinate to both the English main title and the French parallel title, without this relationship being made clear, either in the markup given in the example or in the reference structure offered by the standard.} Again, it is only because this example uses \hyperref[TEI.bibl]{<bibl>} rather than \hyperref[TEI.biblStruct]{<biblStruct>}, that specific punctuation may be included between the component elements of the reference. \par\bgroup\index{bibl=<bibl>|exampleindex}\index{title=<title>|exampleindex}\index{level=@level!<title>|exampleindex}\index{type=@type!<title>|exampleindex}\index{title=<title>|exampleindex}\index{level=@level!<title>|exampleindex}\index{type=@type!<title>|exampleindex}\index{title=<title>|exampleindex}\index{level=@level!<title>|exampleindex}\index{type=@type!<title>|exampleindex}\index{title=<title>|exampleindex}\index{level=@level!<title>|exampleindex}\index{type=@type!<title>|exampleindex}\exampleFont \begin{shaded}\noindent\mbox{}{<\textbf{bibl}>}Tchaikovsky, Peter Ilich.\mbox{}\newline 
{<\textbf{title}\hspace*{1em}{level}="{m}"\hspace*{1em}{type}="{main}">}The swan lake ballet{</\textbf{title}>}\mbox{}\newline 
 = {<\textbf{title}\hspace*{1em}{level}="{m}"\hspace*{1em}{type}="{parallel}"\mbox{}\newline 
\hspace*{1em}\hspace*{1em}{xml:lang}="{fr}">}Le lac des cygnes{</\textbf{title}>}\mbox{}\newline 
 : {<\textbf{title}\hspace*{1em}{level}="{m}"\hspace*{1em}{type}="{sub}"\hspace*{1em}{xml:lang}="{fr}">}grand ballet en 4 actes{</\textbf{title}>}\mbox{}\newline 
 : {<\textbf{title}\hspace*{1em}{level}="{m}"\hspace*{1em}{type}="{sub}">}op. 20{</\textbf{title}>}\mbox{}\newline 
 [Score].\mbox{}\newline 
 New York: Broude Brothers; [1951] (B.B. 59). vi, 685 p.{</\textbf{bibl}>}\end{shaded}\egroup\par \noindent  \par
The elements \hyperref[TEI.author]{<author>} and \hyperref[TEI.editor]{<editor>} have fairly obvious significance for printed books and articles; for other kinds of bibliographic items their proper usage may be less obvious. The \hyperref[TEI.author]{<author>} element should be used for the person or agency with primary responsibility for a work's intellectual content, and the element \hyperref[TEI.editor]{<editor>} for other people or agencies with some responsibility for that content, whether or not they are called ‘editor’. An organization such as a radio or television station is usually accounted ‘author’ of a broadcast, for example, while the author of a government report will usually be the agency which produced it. A translator, illustrator, or compiler, may however be marked by means of the \hyperref[TEI.editor]{<editor>} element, optionally using the {\itshape role} attribute to specify the nature of their responsibility more exactly.\par
Many bibliographic and Linked Data applications require disambiguation of author names using unique identifiers. Both the \hyperref[TEI.author]{<author>} and \hyperref[TEI.editor]{<editor>} elements may contain one or more \hyperref[TEI.idno]{<idno>} elements, to supply such identifiers. Alternatively, if only a single identifier is to be recorded, the {\itshape key} or {\itshape ref} attribute may be used, as further discussed in \textit{\hyperref[CONARS]{3.6.1.\ Referring Strings}}.\par\bgroup\index{bibl=<bibl>|exampleindex}\index{author=<author>|exampleindex}\index{ref=@ref!<author>|exampleindex}\index{editor=<editor>|exampleindex}\index{ref=@ref!<editor>|exampleindex}\index{editor=<editor>|exampleindex}\index{ref=@ref!<editor>|exampleindex}\exampleFont \begin{shaded}\noindent\mbox{}{<\textbf{bibl}>}\mbox{}\newline 
\hspace*{1em}{<\textbf{author}\hspace*{1em}{ref}="{http://viaf.org/viaf/95301405}">}John Warrack{</\textbf{author}>}. „Es waren seine letzten Töne!“\mbox{}\newline 
 In {<\textbf{editor}\hspace*{1em}{ref}="{http://viaf.org/viaf/263865979}">}Joachim Veit{</\textbf{editor}>} \mbox{}\newline 
 and {<\textbf{editor}\hspace*{1em}{ref}="{http://viaf.org/viaf/268371810}">}Frank Ziegler{</\textbf{editor}>} eds. Weber-Studien Bd. 3, Mainz (1996), pp.300–317\mbox{}\newline 
\mbox{}\newline 
{</\textbf{bibl}>}\end{shaded}\egroup\par \par
For anyone else with responsibility for the work, the \hyperref[TEI.respStmt]{<respStmt>} element should be used. The nature of the responsibility is indicated by means of a \hyperref[TEI.resp]{<resp>} element, and the person, organization, etc. responsible by a \hyperref[TEI.name]{<name>}, \hyperref[TEI.persName]{<persName>}, or \hyperref[TEI.orgName]{<orgName>} element. Strings such as ‘unknown’ may be encoded using the \hyperref[TEI.rs]{<rs>} element. A \hyperref[TEI.respStmt]{<respStmt>} should comprise either at least one of the four naming elements (\hyperref[TEI.name]{<name>}, \hyperref[TEI.persName]{<persName>}, \hyperref[TEI.orgName]{<orgName>}, or \hyperref[TEI.rs]{<rs>}) followed by one or more \hyperref[TEI.resp]{<resp>} elements, or at least one \hyperref[TEI.resp]{<resp>} element followed by one or more of the four naming elements.\par
Examples of secondary responsibility of this kind include the roles of illustrator, translator, encoder, and annotator. The \hyperref[TEI.respStmt]{<respStmt>} element may also be used for editors, if it is desired to record the specific terms in which their role is described.\par
Examples of \hyperref[TEI.author]{<author>} and \hyperref[TEI.editor]{<editor>} may be found in sections \textit{\hyperref[COBITY]{3.12.1.\ Methods of Encoding Bibliographic References and Lists of References}}, and \textit{\hyperref[COBICOL]{3.12.2.1.\ Analytic, Monographic, and Series Levels}}; wherever \hyperref[TEI.author]{<author>} and \hyperref[TEI.editor]{<editor>} may occur, the \hyperref[TEI.respStmt]{<respStmt>} element may also occur. When one of these elements precedes or immediately follows a title, it applies to that title; when it follows an \hyperref[TEI.edition]{<edition>} element or occurs within an edition statement, it applies to the edition in question.\par
In this example, the \hyperref[TEI.respStmt]{<respStmt>} elements apply to the work as a whole, not merely to the first edition:  \par\bgroup\index{bibl=<bibl>|exampleindex}\index{author=<author>|exampleindex}\index{title=<title>|exampleindex}\index{level=@level!<title>|exampleindex}\index{respStmt=<respStmt>|exampleindex}\index{resp=<resp>|exampleindex}\index{name=<name>|exampleindex}\index{respStmt=<respStmt>|exampleindex}\index{resp=<resp>|exampleindex}\index{name=<name>|exampleindex}\index{edition=<edition>|exampleindex}\index{pubPlace=<pubPlace>|exampleindex}\index{publisher=<publisher>|exampleindex}\index{date=<date>|exampleindex}\index{extent=<extent>|exampleindex}\index{title=<title>|exampleindex}\index{level=@level!<title>|exampleindex}\index{note=<note>|exampleindex}\index{place=@place!<note>|exampleindex}\index{title=<title>|exampleindex}\index{level=@level!<title>|exampleindex}\index{idno=<idno>|exampleindex}\index{type=@type!<idno>|exampleindex}\exampleFont \begin{shaded}\noindent\mbox{}{<\textbf{bibl}>}\mbox{}\newline 
\hspace*{1em}{<\textbf{author}>}Lominandze, DG{</\textbf{author}>}.\mbox{}\newline 
{<\textbf{title}\hspace*{1em}{level}="{m}">}Cyclotron waves in plasma{</\textbf{title}>}.\mbox{}\newline 
{<\textbf{respStmt}>}\mbox{}\newline 
\hspace*{1em}\hspace*{1em}{<\textbf{resp}>}Translated by{</\textbf{resp}>}\mbox{}\newline 
\hspace*{1em}\hspace*{1em}{<\textbf{name}>}AN. Dellis{</\textbf{name}>}\mbox{}\newline 
\hspace*{1em}{</\textbf{respStmt}>};\mbox{}\newline 
{<\textbf{respStmt}>}\mbox{}\newline 
\hspace*{1em}\hspace*{1em}{<\textbf{resp}>}edited by{</\textbf{resp}>}\mbox{}\newline 
\hspace*{1em}\hspace*{1em}{<\textbf{name}>}SM. Hamberger{</\textbf{name}>}\mbox{}\newline 
\hspace*{1em}{</\textbf{respStmt}>}.\mbox{}\newline 
{<\textbf{edition}>}1st ed.{</\textbf{edition}>}\mbox{}\newline 
\hspace*{1em}{<\textbf{pubPlace}>}Oxford{</\textbf{pubPlace}>}:\mbox{}\newline 
{<\textbf{publisher}>}Pergamon Press{</\textbf{publisher}>},\mbox{}\newline 
{<\textbf{date}>}1981{</\textbf{date}>}.\mbox{}\newline 
{<\textbf{extent}>}206 p.{</\textbf{extent}>}\mbox{}\newline 
\hspace*{1em}{<\textbf{title}\hspace*{1em}{level}="{s}">}International series in natural philosophy{</\textbf{title}>}.\mbox{}\newline 
{<\textbf{note}\hspace*{1em}{place}="{inline}">}Translation of:\mbox{}\newline 
\hspace*{1em}{<\textbf{title}\hspace*{1em}{xml:lang}="{ru-Latn}"\hspace*{1em}{level}="{m}">}Ciklotronnye volny v\mbox{}\newline 
\hspace*{1em}\hspace*{1em}\hspace*{1em}\hspace*{1em} plazme{</\textbf{title}>}.\mbox{}\newline 
\hspace*{1em}{<\textbf{idno}\hspace*{1em}{type}="{ISBN}">}0-08-021680-3{</\textbf{idno}>}.\mbox{}\newline 
\hspace*{1em}{</\textbf{note}>}\mbox{}\newline 
{</\textbf{bibl}>}\end{shaded}\egroup\par \noindent    \par
This example retains the original punctuation and editorial conventions of the source (ISO 690: 1987) and is therefore encoded using the \hyperref[TEI.bibl]{<bibl>} element.\par
In the following example, by contrast, the \hyperref[TEI.respStmt]{<respStmt>} element applies to the edition, and not to the collection \textit{per se} (Moser and Tervooren were not responsible for the first thirty-five printings). As is permissible within a \hyperref[TEI.biblStruct]{<biblStruct>} element, the component elements have been reordered from their appearance on the title page of the volume in order to ensure the correct relationship of the collection title, the edition statement, and the statement of responsibility. \par\bgroup\index{biblStruct=<biblStruct>|exampleindex}\index{monogr=<monogr>|exampleindex}\index{title=<title>|exampleindex}\index{level=@level!<title>|exampleindex}\index{note=<note>|exampleindex}\index{place=@place!<note>|exampleindex}\index{edition=<edition>|exampleindex}\index{respStmt=<respStmt>|exampleindex}\index{resp=<resp>|exampleindex}\index{name=<name>|exampleindex}\index{name=<name>|exampleindex}\index{name=<name>|exampleindex}\index{name=<name>|exampleindex}\index{name=<name>|exampleindex}\index{name=<name>|exampleindex}\index{imprint=<imprint>|exampleindex}\index{pubPlace=<pubPlace>|exampleindex}\index{publisher=<publisher>|exampleindex}\index{date=<date>|exampleindex}\index{biblScope=<biblScope>|exampleindex}\index{unit=@unit!<biblScope>|exampleindex}\exampleFont \begin{shaded}\noindent\mbox{}{<\textbf{biblStruct}>}\mbox{}\newline 
\hspace*{1em}{<\textbf{monogr}\hspace*{1em}{xml:lang}="{de}">}\mbox{}\newline 
\hspace*{1em}\hspace*{1em}{<\textbf{title}\hspace*{1em}{level}="{m}">}Des Minnesangs Frühling{</\textbf{title}>}\mbox{}\newline 
\hspace*{1em}\hspace*{1em}{<\textbf{note}\hspace*{1em}{place}="{inline}">}Mit 1 Faksimile{</\textbf{note}>}\mbox{}\newline 
\hspace*{1em}\hspace*{1em}{<\textbf{edition}>}36., neugestaltete und erweiterte Auflage{</\textbf{edition}>}\mbox{}\newline 
\hspace*{1em}\hspace*{1em}{<\textbf{respStmt}>}\mbox{}\newline 
\hspace*{1em}\hspace*{1em}\hspace*{1em}{<\textbf{resp}>}Unter Benutzung der Ausgaben von {<\textbf{name}>}Karl\mbox{}\newline 
\hspace*{1em}\hspace*{1em}\hspace*{1em}\hspace*{1em}\hspace*{1em}\hspace*{1em}\hspace*{1em}\hspace*{1em} Lachmann{</\textbf{name}>} und {<\textbf{name}>}Moriz Haupt{</\textbf{name}>}, {<\textbf{name}>}Friedrich\mbox{}\newline 
\hspace*{1em}\hspace*{1em}\hspace*{1em}\hspace*{1em}\hspace*{1em}\hspace*{1em}\hspace*{1em}\hspace*{1em} Vogt{</\textbf{name}>} und {<\textbf{name}>}Carl von Kraus{</\textbf{name}>} bearbeitet von{</\textbf{resp}>}\mbox{}\newline 
\hspace*{1em}\hspace*{1em}\hspace*{1em}{<\textbf{name}>}Hugo Moser{</\textbf{name}>}\mbox{}\newline 
\hspace*{1em}\hspace*{1em}\hspace*{1em}{<\textbf{name}>}Helmut Tervooren{</\textbf{name}>}\mbox{}\newline 
\hspace*{1em}\hspace*{1em}{</\textbf{respStmt}>}\mbox{}\newline 
\hspace*{1em}\hspace*{1em}{<\textbf{imprint}>}\mbox{}\newline 
\hspace*{1em}\hspace*{1em}\hspace*{1em}{<\textbf{pubPlace}>}Stuttgart{</\textbf{pubPlace}>}\mbox{}\newline 
\hspace*{1em}\hspace*{1em}\hspace*{1em}{<\textbf{publisher}>}S. Hirzel Verlag{</\textbf{publisher}>}\mbox{}\newline 
\hspace*{1em}\hspace*{1em}\hspace*{1em}{<\textbf{date}>}1977{</\textbf{date}>}\mbox{}\newline 
\hspace*{1em}\hspace*{1em}{</\textbf{imprint}>}\mbox{}\newline 
\hspace*{1em}\hspace*{1em}{<\textbf{biblScope}\hspace*{1em}{unit}="{volume}">}I Texte{</\textbf{biblScope}>}\mbox{}\newline 
\hspace*{1em}{</\textbf{monogr}>}\mbox{}\newline 
{</\textbf{biblStruct}>}\end{shaded}\egroup\par \noindent    \par
The party with a particular responsibility for the intellectual content may vary over time. Likewise, a given individal's responsibility or role may change over time. These situations may be recorded with the \hyperref[TEI.respStmt]{<respStmt>} element. For example, the following could be used when one proofreader took over for another. \par\bgroup\index{respStmt=<respStmt>|exampleindex}\index{resp=<resp>|exampleindex}\index{persName=<persName>|exampleindex}\index{from=@from!<persName>|exampleindex}\index{to=@to!<persName>|exampleindex}\index{persName=<persName>|exampleindex}\index{from=@from!<persName>|exampleindex}\index{to=@to!<persName>|exampleindex}\exampleFont \begin{shaded}\noindent\mbox{}{<\textbf{respStmt}>}\mbox{}\newline 
\hspace*{1em}{<\textbf{resp}>}proofreading{</\textbf{resp}>}\mbox{}\newline 
\hspace*{1em}{<\textbf{persName}\hspace*{1em}{from}="{1994-02}"\hspace*{1em}{to}="{1994-05}">}Ashley Cross{</\textbf{persName}>}\mbox{}\newline 
\hspace*{1em}{<\textbf{persName}\hspace*{1em}{from}="{1994-06}"\hspace*{1em}{to}="{1994-10}">}Loren Noveck{</\textbf{persName}>}\mbox{}\newline 
{</\textbf{respStmt}>}\end{shaded}\egroup\par \noindent  The following example records the fact that one individual had two distinctly different intellectual responsibilities at different times. \par\bgroup\index{respStmt=<respStmt>|exampleindex}\index{persName=<persName>|exampleindex}\index{resp=<resp>|exampleindex}\index{when=@when!<resp>|exampleindex}\index{resp=<resp>|exampleindex}\index{when=@when!<resp>|exampleindex}\exampleFont \begin{shaded}\noindent\mbox{}{<\textbf{respStmt}>}\mbox{}\newline 
\hspace*{1em}{<\textbf{persName}>}Erica Dillon{</\textbf{persName}>}\mbox{}\newline 
\hspace*{1em}{<\textbf{resp}\hspace*{1em}{when}="{2000-08}">}annotated uncredited citations{</\textbf{resp}>}\mbox{}\newline 
\hspace*{1em}{<\textbf{resp}\hspace*{1em}{when}="{2001-03}">}encoded named entities{</\textbf{resp}>}\mbox{}\newline 
{</\textbf{respStmt}>}\end{shaded}\egroup\par \par
Another form of ‘responsibility’ arises when a work is published as the outcome of a conference, workshop or similar meeting. The \hyperref[TEI.meeting]{<meeting>} element may be used to supply this information, as in the following example: \par\bgroup\index{biblStruct=<biblStruct>|exampleindex}\index{monogr=<monogr>|exampleindex}\index{title=<title>|exampleindex}\index{level=@level!<title>|exampleindex}\index{respStmt=<respStmt>|exampleindex}\index{resp=<resp>|exampleindex}\index{name=<name>|exampleindex}\index{meeting=<meeting>|exampleindex}\index{imprint=<imprint>|exampleindex}\index{pubPlace=<pubPlace>|exampleindex}\exampleFont \begin{shaded}\noindent\mbox{}{<\textbf{biblStruct}>}\mbox{}\newline 
\hspace*{1em}{<\textbf{monogr}>}\mbox{}\newline 
\hspace*{1em}\hspace*{1em}{<\textbf{title}\hspace*{1em}{level}="{m}">}Proceedings of a workshop on corpus resources{</\textbf{title}>}\mbox{}\newline 
\hspace*{1em}\hspace*{1em}{<\textbf{respStmt}>}\mbox{}\newline 
\hspace*{1em}\hspace*{1em}\hspace*{1em}{<\textbf{resp}>}Programme Organizer{</\textbf{resp}>}\mbox{}\newline 
\hspace*{1em}\hspace*{1em}\hspace*{1em}{<\textbf{name}>}Geoffrey Leech{</\textbf{name}>}\mbox{}\newline 
\hspace*{1em}\hspace*{1em}{</\textbf{respStmt}>}\mbox{}\newline 
\hspace*{1em}\hspace*{1em}{<\textbf{meeting}>}DTI Speech and Language Technology Club meeting, 3-4\mbox{}\newline 
\hspace*{1em}\hspace*{1em}\hspace*{1em}\hspace*{1em} January 1990, Wadham College, Oxford{</\textbf{meeting}>}\mbox{}\newline 
\hspace*{1em}\hspace*{1em}{<\textbf{imprint}>}\mbox{}\newline 
\hspace*{1em}\hspace*{1em}\hspace*{1em}{<\textbf{pubPlace}>}Oxford{</\textbf{pubPlace}>}\mbox{}\newline 
\hspace*{1em}\hspace*{1em}{</\textbf{imprint}>}\mbox{}\newline 
\hspace*{1em}{</\textbf{monogr}>}\mbox{}\newline 
{</\textbf{biblStruct}>}\end{shaded}\egroup\par 
\paragraph[{Document Identifiers}]{Document Identifiers}\label{COBICOD}\par
Many bibliographic references include identifiers for a work to help with precise identification of an appropriate document. For example, a book in the \textit{Short Title Catalogue} could be referenced with its STC number: \par\bgroup\index{biblStruct=<biblStruct>|exampleindex}\index{monogr=<monogr>|exampleindex}\index{author=<author>|exampleindex}\index{forename=<forename>|exampleindex}\index{surname=<surname>|exampleindex}\index{title=<title>|exampleindex}\index{type=@type!<title>|exampleindex}\index{idno=<idno>|exampleindex}\index{type=@type!<idno>|exampleindex}\index{imprint=<imprint>|exampleindex}\index{pubPlace=<pubPlace>|exampleindex}\index{publisher=<publisher>|exampleindex}\index{date=<date>|exampleindex}\index{when=@when!<date>|exampleindex}\exampleFont \begin{shaded}\noindent\mbox{}{<\textbf{biblStruct}>}\mbox{}\newline 
\hspace*{1em}{<\textbf{monogr}>}\mbox{}\newline 
\hspace*{1em}\hspace*{1em}{<\textbf{author}>}\mbox{}\newline 
\hspace*{1em}\hspace*{1em}\hspace*{1em}{<\textbf{forename}>}John{</\textbf{forename}>}\mbox{}\newline 
\hspace*{1em}\hspace*{1em}\hspace*{1em}{<\textbf{surname}>}Downame{</\textbf{surname}>}\mbox{}\newline 
\hspace*{1em}\hspace*{1em}{</\textbf{author}>}\mbox{}\newline 
\hspace*{1em}\hspace*{1em}{<\textbf{title}\hspace*{1em}{type}="{short}">}Foure treatises tending to disswade all Christians from foure no lesse hainous then common sinnes{</\textbf{title}>}\mbox{}\newline 
\hspace*{1em}\hspace*{1em}{<\textbf{idno}\hspace*{1em}{type}="{stc2ndEd}">}7141{</\textbf{idno}>}\mbox{}\newline 
\hspace*{1em}\hspace*{1em}{<\textbf{imprint}>}\mbox{}\newline 
\hspace*{1em}\hspace*{1em}\hspace*{1em}{<\textbf{pubPlace}>}At London{</\textbf{pubPlace}>}\mbox{}\newline 
\hspace*{1em}\hspace*{1em}\hspace*{1em}{<\textbf{publisher}>}Imprinted by Felix Kyngston, for William Welby, and are to be sold at his shop in Pauls Church-yard at the signe of the Greyhound{</\textbf{publisher}>}\mbox{}\newline 
\hspace*{1em}\hspace*{1em}\hspace*{1em}{<\textbf{date}\hspace*{1em}{when}="{1609}">}1609{</\textbf{date}>}\mbox{}\newline 
\hspace*{1em}\hspace*{1em}{</\textbf{imprint}>}\mbox{}\newline 
\hspace*{1em}{</\textbf{monogr}>}\mbox{}\newline 
{</\textbf{biblStruct}>}\end{shaded}\egroup\par \par
However, some bibliographic references actually \textit{require} identifiers of various types because they do not include a statement of the title and the names of those intellectually responsible for it. The following elements may be used for such purposes: 
\begin{sansreflist}
  
\item [\textbf{<orgName>}] (organization name) contains an organizational name.
\item [\textbf{<idno>}] (identifier) supplies any form of identifier used to identify some object, such as a bibliographic item, a person, a title, an organization, etc. in a standardized way.
\item [\textbf{<classCode>}] (classification code) contains the classification code used for this text in some standard classification system.
\item [\textbf{<date>}] (date) contains a date in any format.
\end{sansreflist}
\par
For example, a citation to a patent typically includes a country or organization code (a two-character code identifying a patent authority) and a serial number for the patent (whose structure varies by patent authority). The citation might also contain a \textit{kind code} (which characterizes a particular publication for the patent and which corresponds to a specific stage in the patent procedure) and the date when the patent was filed with or published by the issuing authority. For bibliographic references to patents, the above elements may be used as follows: \begin{itemize}
\item \hyperref[TEI.orgName]{<orgName>}, within \hyperref[TEI.authority]{<authority>}, may be used to contain the code of the patent authority. The {\itshape type} attribute may be used to specify the type of patent authority (such as a national patent office or a supra-national patent organization).
\item \hyperref[TEI.idno]{<idno>} may be used to contain the serial number assigned by the corresponding patent authority.
\item \hyperref[TEI.classCode]{<classCode>} may be used to contain the kind code of the patent document.
\item \hyperref[TEI.date]{<date>} may be used to contain the date of the patent document. The {\itshape type} attribute may be used to specify whether this corresponds to the filing date of a patent application or the publication date of a patent publication.
\end{itemize} \par
The following reference illustrates an encoding for a patent publication which might be cited in print as ‘United States patent US 6,885,550 B1, issued April 26, 2005’: \par\bgroup\index{biblStruct=<biblStruct>|exampleindex}\index{type=@type!<biblStruct>|exampleindex}\index{status=@status!<biblStruct>|exampleindex}\index{monogr=<monogr>|exampleindex}\index{authority=<authority>|exampleindex}\index{orgName=<orgName>|exampleindex}\index{type=@type!<orgName>|exampleindex}\index{idno=<idno>|exampleindex}\index{type=@type!<idno>|exampleindex}\index{imprint=<imprint>|exampleindex}\index{classCode=<classCode>|exampleindex}\index{scheme=@scheme!<classCode>|exampleindex}\index{date=<date>|exampleindex}\index{type=@type!<date>|exampleindex}\index{when=@when!<date>|exampleindex}\exampleFont \begin{shaded}\noindent\mbox{}{<\textbf{biblStruct}\hspace*{1em}{type}="{patent}"\mbox{}\newline 
\hspace*{1em}{status}="{publication}">}\mbox{}\newline 
\hspace*{1em}{<\textbf{monogr}>}\mbox{}\newline 
\hspace*{1em}\hspace*{1em}{<\textbf{authority}>}\mbox{}\newline 
\hspace*{1em}\hspace*{1em}\hspace*{1em}{<\textbf{orgName}\hspace*{1em}{type}="{national}">}US{</\textbf{orgName}>}\mbox{}\newline 
\hspace*{1em}\hspace*{1em}{</\textbf{authority}>}\mbox{}\newline 
\hspace*{1em}\hspace*{1em}{<\textbf{idno}\hspace*{1em}{type}="{docNumber}">}6885550{</\textbf{idno}>}\mbox{}\newline 
\hspace*{1em}\hspace*{1em}{<\textbf{imprint}>}\mbox{}\newline 
\hspace*{1em}\hspace*{1em}\hspace*{1em}{<\textbf{classCode}\hspace*{1em}{scheme}="{http://www.uspto.gov/}">}B1{</\textbf{classCode}>}\mbox{}\newline 
\hspace*{1em}\hspace*{1em}\hspace*{1em}{<\textbf{date}\hspace*{1em}{type}="{publicationDate}"\mbox{}\newline 
\hspace*{1em}\hspace*{1em}\hspace*{1em}\hspace*{1em}{when}="{2005-04-26}">}April 26, 2005{</\textbf{date}>}\mbox{}\newline 
\hspace*{1em}\hspace*{1em}{</\textbf{imprint}>}\mbox{}\newline 
\hspace*{1em}{</\textbf{monogr}>}\mbox{}\newline 
{</\textbf{biblStruct}>}\end{shaded}\egroup\par 
\paragraph[{Imprint, Size of a Document, and Reprint Information}]{Imprint, Size of a Document, and Reprint Information}\label{COBICOI}\par
By \textit{imprint} is meant all the information relating to the publication of a work: the person or organization by whose authority and in whose name a bibliographic entity such as a book is made public or distributed (whether a commercial publisher or some other organization), the place and the date of publication. It may also include a full address for the publisher or organization. A full bibliographic references will usually also specify the number of pages in a print publication (or equivalent information for non-print materials), and possibly also the specific location of the material being cited within its containing publication. The following elements are provided to hold this information: 
\begin{sansreflist}
  
\item [\textbf{<imprint>}] groups information relating to the publication or distribution of a bibliographic item.
\item [\textbf{<address>}] (address) contains a postal address, for example of a publisher, an organization, or an individual.
\item [\textbf{<pubPlace>}] (publication place) contains the name of the place where a bibliographic item was published.
\item [\textbf{<publisher>}] (publisher) provides the name of the organization responsible for the publication or distribution of a bibliographic item.
\item [\textbf{<date>}] (date) contains a date in any format.
\item [\textbf{<extent>}] (extent) describes the approximate size of a text stored on some carrier medium or of some other object, digital or non-digital, specified in any convenient units.
\item [\textbf{<idno>}] (identifier) supplies any form of identifier used to identify some object, such as a bibliographic item, a person, a title, an organization, etc. in a standardized way.
\end{sansreflist}
 Members of the model classes \textsf{model.imprintPart} and \textsf{model.dateLike} may appear inside an \hyperref[TEI.imprint]{<imprint>} element in a specific location within a \hyperref[TEI.biblStruct]{<biblStruct>}, or alternatively, they may appear alongside any other bibliographic component inside a \hyperref[TEI.bibl]{<bibl>}. 
\begin{sansreflist}
  
\item [\textbf{model.imprintPart}] groups the bibliographic elements which occur inside imprints. \par 
\begin{longtable}{P{0.19436392914653783\textwidth}P{0.6556360708534621\textwidth}}
\hyperref[TEI.biblScope]{biblScope}\tabcellsep (scope of bibliographic reference) defines the scope of a bibliographic reference, for example as a list of page numbers, or a named subdivision of a larger work.\\
\hyperref[TEI.distributor]{distributor}\tabcellsep (distributor) supplies the name of a person or other agency responsible for the distribution of a text.\\
\hyperref[TEI.publisher]{publisher}\tabcellsep (publisher) provides the name of the organization responsible for the publication or distribution of a bibliographic item.\\
\hyperref[TEI.pubPlace]{pubPlace}\tabcellsep (publication place) contains the name of the place where a bibliographic item was published.\end{longtable} \par
 
\item [\textbf{model.dateLike}] groups elements containing temporal expressions. \par 
\begin{longtable}{P{0.2852348993288591\textwidth}P{0.5647651006711409\textwidth}}
\hyperref[TEI.date]{date}\tabcellsep (date) contains a date in any format.\\
\hyperref[TEI.time]{time}\tabcellsep (time) contains a phrase defining a time of day in any format.\end{longtable} \par
 
\end{sansreflist}
\par
For bibliographic purposes, usually only the place (or places) of publication are required, possibly including the name of the country, rather than a full address; the element \hyperref[TEI.pubPlace]{<pubPlace>} is provided for this purpose. Where however the full postal address is likely to be of importance in identifying or locating the bibliographic item concerned, it may be supplied and tagged using the \hyperref[TEI.address]{<address>} element described in section \textit{\hyperref[CONAAD]{3.6.2.\ Addresses}}. Alternatively, if desired, the \hyperref[TEI.rs]{<rs>} or \hyperref[TEI.name]{<name>} elements described in section \textit{\hyperref[CONARS]{3.6.1.\ Referring Strings}} may be used; this involves no claim that the information given is either a full address or the name of a city.\par
The name of the publisher of an item should be marked using the \hyperref[TEI.publisher]{<publisher>} element even if the item is made public (‘published’) by an organization other than a conventional publisher, as is frequently the case with technical reports: \par\bgroup\index{biblStruct=<biblStruct>|exampleindex}\index{monogr=<monogr>|exampleindex}\index{author=<author>|exampleindex}\index{author=<author>|exampleindex}\index{title=<title>|exampleindex}\index{level=@level!<title>|exampleindex}\index{idno=<idno>|exampleindex}\index{type=@type!<idno>|exampleindex}\index{imprint=<imprint>|exampleindex}\index{pubPlace=<pubPlace>|exampleindex}\index{publisher=<publisher>|exampleindex}\index{date=<date>|exampleindex}\index{when=@when!<date>|exampleindex}\index{extent=<extent>|exampleindex}\exampleFont \begin{shaded}\noindent\mbox{}{<\textbf{biblStruct}>}\mbox{}\newline 
\hspace*{1em}{<\textbf{monogr}>}\mbox{}\newline 
\hspace*{1em}\hspace*{1em}{<\textbf{author}>}Nicholas, Charles K.{</\textbf{author}>}\mbox{}\newline 
\hspace*{1em}\hspace*{1em}{<\textbf{author}>}Welsch, Lawrence A.{</\textbf{author}>}\mbox{}\newline 
\hspace*{1em}\hspace*{1em}{<\textbf{title}\hspace*{1em}{level}="{m}">}On the interchangeability of SGML and ODA{</\textbf{title}>}\mbox{}\newline 
\hspace*{1em}\hspace*{1em}{<\textbf{idno}\hspace*{1em}{type}="{NIST}">}NISTIR 4681{</\textbf{idno}>}\mbox{}\newline 
\hspace*{1em}\hspace*{1em}{<\textbf{imprint}>}\mbox{}\newline 
\hspace*{1em}\hspace*{1em}\hspace*{1em}{<\textbf{pubPlace}>}Gaithersburg, MD{</\textbf{pubPlace}>}\mbox{}\newline 
\hspace*{1em}\hspace*{1em}\hspace*{1em}{<\textbf{publisher}>}National Institute of Standards and Technology\mbox{}\newline 
\hspace*{1em}\hspace*{1em}\hspace*{1em}{</\textbf{publisher}>}\mbox{}\newline 
\hspace*{1em}\hspace*{1em}\hspace*{1em}{<\textbf{date}\hspace*{1em}{when}="{1992-01}">}January 1992{</\textbf{date}>}\mbox{}\newline 
\hspace*{1em}\hspace*{1em}{</\textbf{imprint}>}\mbox{}\newline 
\hspace*{1em}\hspace*{1em}{<\textbf{extent}>}19 pp.{</\textbf{extent}>}\mbox{}\newline 
\hspace*{1em}{</\textbf{monogr}>}\mbox{}\newline 
{</\textbf{biblStruct}>}\end{shaded}\egroup\par \noindent  and with dissertations: \par\bgroup\index{biblStruct=<biblStruct>|exampleindex}\index{monogr=<monogr>|exampleindex}\index{author=<author>|exampleindex}\index{title=<title>|exampleindex}\index{level=@level!<title>|exampleindex}\index{idno=<idno>|exampleindex}\index{type=@type!<idno>|exampleindex}\index{note=<note>|exampleindex}\index{place=@place!<note>|exampleindex}\index{imprint=<imprint>|exampleindex}\index{publisher=<publisher>|exampleindex}\index{pubPlace=<pubPlace>|exampleindex}\index{date=<date>|exampleindex}\index{when=@when!<date>|exampleindex}\exampleFont \begin{shaded}\noindent\mbox{}{<\textbf{biblStruct}>}\mbox{}\newline 
\hspace*{1em}{<\textbf{monogr}>}\mbox{}\newline 
\hspace*{1em}\hspace*{1em}{<\textbf{author}>}Hansen, W.{</\textbf{author}>}\mbox{}\newline 
\hspace*{1em}\hspace*{1em}{<\textbf{title}\hspace*{1em}{level}="{u}">}Creation of hierarchic text\mbox{}\newline 
\hspace*{1em}\hspace*{1em}\hspace*{1em}\hspace*{1em} with a computer display{</\textbf{title}>}\mbox{}\newline 
\hspace*{1em}\hspace*{1em}{<\textbf{idno}\hspace*{1em}{type}="{ANL}">}ANL-7818{</\textbf{idno}>}\mbox{}\newline 
\hspace*{1em}\hspace*{1em}{<\textbf{note}\hspace*{1em}{place}="{inline}">}Ph.D. dissertation{</\textbf{note}>}\mbox{}\newline 
\hspace*{1em}\hspace*{1em}{<\textbf{imprint}>}\mbox{}\newline 
\hspace*{1em}\hspace*{1em}\hspace*{1em}{<\textbf{publisher}>}Dept. of Computer Science, Stanford Univ.{</\textbf{publisher}>}\mbox{}\newline 
\hspace*{1em}\hspace*{1em}\hspace*{1em}{<\textbf{pubPlace}>}Stanford, CA{</\textbf{pubPlace}>}\mbox{}\newline 
\hspace*{1em}\hspace*{1em}\hspace*{1em}{<\textbf{date}\hspace*{1em}{when}="{1971-06}">}June 1971{</\textbf{date}>}\mbox{}\newline 
\hspace*{1em}\hspace*{1em}{</\textbf{imprint}>}\mbox{}\newline 
\hspace*{1em}{</\textbf{monogr}>}\mbox{}\newline 
{</\textbf{biblStruct}>}\end{shaded}\egroup\par \par
In this second example, the \hyperref[TEI.idno]{<idno>} element is used to provide the identifier allocated to the thesis by the Argonne National Laboratory. Since it applies to the monographic element, the \hyperref[TEI.idno]{<idno>} should be provided as a direct child of the \hyperref[TEI.monogr]{<monogr>} element, rather than elsewhere in the \hyperref[TEI.biblStruct]{<biblStruct>} element.\par
The specialist elements \hyperref[TEI.publisher]{<publisher>} and \hyperref[TEI.distributor]{<distributor>} are provided to cover the most common roles related to the production and distribution of a bibliographical item, but other roles such as printer and bookseller may also need to be encoded, and \hyperref[TEI.respStmt]{<respStmt>} is available inside \hyperref[TEI.imprint]{<imprint>} for this purpose.\par
When an item has been reprinted, especially reprinted without change from a specific earlier edition, the reprint may appear in a \hyperref[TEI.monogr]{<monogr>} element with only the \hyperref[TEI.imprint]{<imprint>} and other details of the reprint. In the following example, a microform reprint has been issued without any change in the title or authorship. The series statement here applies only to the second \hyperref[TEI.monogr]{<monogr>} element. \par\bgroup\index{biblStruct=<biblStruct>|exampleindex}\index{monogr=<monogr>|exampleindex}\index{author=<author>|exampleindex}\index{title=<title>|exampleindex}\index{type=@type!<title>|exampleindex}\index{title=<title>|exampleindex}\index{type=@type!<title>|exampleindex}\index{note=<note>|exampleindex}\index{place=@place!<note>|exampleindex}\index{imprint=<imprint>|exampleindex}\index{pubPlace=<pubPlace>|exampleindex}\index{publisher=<publisher>|exampleindex}\index{date=<date>|exampleindex}\index{extent=<extent>|exampleindex}\index{monogr=<monogr>|exampleindex}\index{imprint=<imprint>|exampleindex}\index{pubPlace=<pubPlace>|exampleindex}\index{publisher=<publisher>|exampleindex}\index{date=<date>|exampleindex}\index{extent=<extent>|exampleindex}\index{series=<series>|exampleindex}\index{title=<title>|exampleindex}\index{level=@level!<title>|exampleindex}\exampleFont \begin{shaded}\noindent\mbox{}{<\textbf{biblStruct}>}\mbox{}\newline 
\hspace*{1em}{<\textbf{monogr}>}\mbox{}\newline 
\hspace*{1em}\hspace*{1em}{<\textbf{author}>}Shirley, James{</\textbf{author}>}\mbox{}\newline 
\hspace*{1em}\hspace*{1em}{<\textbf{title}\hspace*{1em}{type}="{main}">}The gentlemen of Venice{</\textbf{title}>}\mbox{}\newline 
\hspace*{1em}\hspace*{1em}{<\textbf{title}\hspace*{1em}{type}="{sub}">}a tragi-comedie presented at the private\mbox{}\newline 
\hspace*{1em}\hspace*{1em}\hspace*{1em}\hspace*{1em} house in Salisbury Court by Her Majesties servants{</\textbf{title}>}\mbox{}\newline 
\hspace*{1em}\hspace*{1em}{<\textbf{note}\hspace*{1em}{place}="{inline}">}[Microform]{</\textbf{note}>}\mbox{}\newline 
\hspace*{1em}\hspace*{1em}{<\textbf{imprint}>}\mbox{}\newline 
\hspace*{1em}\hspace*{1em}\hspace*{1em}{<\textbf{pubPlace}>}London{</\textbf{pubPlace}>}\mbox{}\newline 
\hspace*{1em}\hspace*{1em}\hspace*{1em}{<\textbf{publisher}>}H. Moseley{</\textbf{publisher}>}\mbox{}\newline 
\hspace*{1em}\hspace*{1em}\hspace*{1em}{<\textbf{date}>}1655{</\textbf{date}>}\mbox{}\newline 
\hspace*{1em}\hspace*{1em}{</\textbf{imprint}>}\mbox{}\newline 
\hspace*{1em}\hspace*{1em}{<\textbf{extent}>}78 p.{</\textbf{extent}>}\mbox{}\newline 
\hspace*{1em}{</\textbf{monogr}>}\mbox{}\newline 
\hspace*{1em}{<\textbf{monogr}>}\mbox{}\newline 
\hspace*{1em}\hspace*{1em}{<\textbf{imprint}>}\mbox{}\newline 
\hspace*{1em}\hspace*{1em}\hspace*{1em}{<\textbf{pubPlace}>}New York{</\textbf{pubPlace}>}\mbox{}\newline 
\hspace*{1em}\hspace*{1em}\hspace*{1em}{<\textbf{publisher}>}Readex Microprint{</\textbf{publisher}>}\mbox{}\newline 
\hspace*{1em}\hspace*{1em}\hspace*{1em}{<\textbf{date}>}1953{</\textbf{date}>}\mbox{}\newline 
\hspace*{1em}\hspace*{1em}{</\textbf{imprint}>}\mbox{}\newline 
\hspace*{1em}\hspace*{1em}{<\textbf{extent}>}1 microprint card, 23 x 15 cm.{</\textbf{extent}>}\mbox{}\newline 
\hspace*{1em}{</\textbf{monogr}>}\mbox{}\newline 
\hspace*{1em}{<\textbf{series}>}\mbox{}\newline 
\hspace*{1em}\hspace*{1em}{<\textbf{title}\hspace*{1em}{level}="{s}">}Three centuries of drama: English, 1642–1700{</\textbf{title}>}\mbox{}\newline 
\hspace*{1em}{</\textbf{series}>}\mbox{}\newline 
{</\textbf{biblStruct}>}\end{shaded}\egroup\par \noindent  \par
This encoding can be extended to the case of patent documents, where the same patent application is published, with or without changes, at different stages of the patenting procedure. In this case, the kind code and, optionally, the publication date characterize different publications of the same patent application during the procedure. For example: \par\bgroup\index{biblStruct=<biblStruct>|exampleindex}\index{type=@type!<biblStruct>|exampleindex}\index{status=@status!<biblStruct>|exampleindex}\index{monogr=<monogr>|exampleindex}\index{authority=<authority>|exampleindex}\index{orgName=<orgName>|exampleindex}\index{type=@type!<orgName>|exampleindex}\index{idno=<idno>|exampleindex}\index{type=@type!<idno>|exampleindex}\index{imprint=<imprint>|exampleindex}\index{classCode=<classCode>|exampleindex}\index{scheme=@scheme!<classCode>|exampleindex}\index{date=<date>|exampleindex}\index{type=@type!<date>|exampleindex}\index{when=@when!<date>|exampleindex}\index{monogr=<monogr>|exampleindex}\index{imprint=<imprint>|exampleindex}\index{classCode=<classCode>|exampleindex}\index{scheme=@scheme!<classCode>|exampleindex}\index{date=<date>|exampleindex}\index{type=@type!<date>|exampleindex}\index{when=@when!<date>|exampleindex}\exampleFont \begin{shaded}\noindent\mbox{}{<\textbf{biblStruct}\hspace*{1em}{type}="{patent}"\mbox{}\newline 
\hspace*{1em}{status}="{publication}">}\mbox{}\newline 
\hspace*{1em}{<\textbf{monogr}>}\mbox{}\newline 
\hspace*{1em}\hspace*{1em}{<\textbf{authority}>}\mbox{}\newline 
\hspace*{1em}\hspace*{1em}\hspace*{1em}{<\textbf{orgName}\hspace*{1em}{type}="{national}">}EP{</\textbf{orgName}>}\mbox{}\newline 
\hspace*{1em}\hspace*{1em}{</\textbf{authority}>}\mbox{}\newline 
\hspace*{1em}\hspace*{1em}{<\textbf{idno}\hspace*{1em}{type}="{docNumber}">}1558513{</\textbf{idno}>}\mbox{}\newline 
\hspace*{1em}\hspace*{1em}{<\textbf{imprint}>}\mbox{}\newline 
\hspace*{1em}\hspace*{1em}\hspace*{1em}{<\textbf{classCode}\hspace*{1em}{scheme}="{http://www.epo.org/}">}A1{</\textbf{classCode}>}\mbox{}\newline 
\hspace*{1em}\hspace*{1em}\hspace*{1em}{<\textbf{date}\hspace*{1em}{type}="{publicationDate}"\mbox{}\newline 
\hspace*{1em}\hspace*{1em}\hspace*{1em}\hspace*{1em}{when}="{2005-08-03}"/>}\mbox{}\newline 
\hspace*{1em}\hspace*{1em}{</\textbf{imprint}>}\mbox{}\newline 
\hspace*{1em}{</\textbf{monogr}>}\mbox{}\newline 
\hspace*{1em}{<\textbf{monogr}>}\mbox{}\newline 
\hspace*{1em}\hspace*{1em}{<\textbf{imprint}>}\mbox{}\newline 
\hspace*{1em}\hspace*{1em}\hspace*{1em}{<\textbf{classCode}\hspace*{1em}{scheme}="{http://www.epo.org/}">}B1{</\textbf{classCode}>}\mbox{}\newline 
\hspace*{1em}\hspace*{1em}\hspace*{1em}{<\textbf{date}\hspace*{1em}{type}="{publicationDate}"\mbox{}\newline 
\hspace*{1em}\hspace*{1em}\hspace*{1em}\hspace*{1em}{when}="{2009-09-09}"/>}\mbox{}\newline 
\hspace*{1em}\hspace*{1em}{</\textbf{imprint}>}\mbox{}\newline 
\hspace*{1em}{</\textbf{monogr}>}\mbox{}\newline 
{</\textbf{biblStruct}>}\end{shaded}\egroup\par \par
The above bibliographic reference discloses different publications of the patent EP1558513 during the patenting procedure. The first publication from 3 August 2005 has the kind code "A1" indicating that it is a published patent application comprising the European search report issued after carrying out the search at the European Patent Office, whereas the second publication from 9 September 2009 has the kind code "B1" indicating that it was published after the patent application has been granted.\par
An alternative way of handling the above situations would be to use the \hyperref[TEI.relatedItem]{<relatedItem>} element described in section \textit{\hyperref[COBIRI]{3.12.2.7.\ Related Items}} below.
\paragraph[{Scopes and Ranges in Bibliographic Citations}]{Scopes and Ranges in Bibliographic Citations}\label{COBICOB}\par
Many bibliographic citations contain data limiting the citation to one or more volumes, issues, or pages, or to a name or number of a subdivison of the host work. These come in two varieties: \begin{itemize}
\item the scope of a bibliographic reference (encoded using \hyperref[TEI.biblScope]{<biblScope>})
\item the range of a work cited (encoded using \hyperref[TEI.citedRange]{<citedRange>})
\end{itemize}  Where it is desired to distinguish different classes of such information (volume number, page number, chapter number, etc.), the {\itshape unit} attribute may be used with any convenient typology (see the element definitions for \hyperref[TEI.biblScope]{<biblScope>} and \hyperref[TEI.citedRange]{<citedRange>} for some suggested values).\par
A scope of a bibliographic reference defines that the \textit{entire work cited} may be found in particular volumes, issues, pages, etc. For example: \par\bgroup\index{biblStruct=<biblStruct>|exampleindex}\index{analytic=<analytic>|exampleindex}\index{author=<author>|exampleindex}\index{persName=<persName>|exampleindex}\index{surname=<surname>|exampleindex}\index{forename=<forename>|exampleindex}\index{full=@full!<forename>|exampleindex}\index{forename=<forename>|exampleindex}\index{full=@full!<forename>|exampleindex}\index{title=<title>|exampleindex}\index{level=@level!<title>|exampleindex}\index{monogr=<monogr>|exampleindex}\index{author=<author>|exampleindex}\index{persName=<persName>|exampleindex}\index{surname=<surname>|exampleindex}\index{forename=<forename>|exampleindex}\index{full=@full!<forename>|exampleindex}\index{forename=<forename>|exampleindex}\index{full=@full!<forename>|exampleindex}\index{author=<author>|exampleindex}\index{persName=<persName>|exampleindex}\index{surname=<surname>|exampleindex}\index{forename=<forename>|exampleindex}\index{full=@full!<forename>|exampleindex}\index{forename=<forename>|exampleindex}\index{full=@full!<forename>|exampleindex}\index{forename=<forename>|exampleindex}\index{full=@full!<forename>|exampleindex}\index{title=<title>|exampleindex}\index{level=@level!<title>|exampleindex}\index{imprint=<imprint>|exampleindex}\index{pubPlace=<pubPlace>|exampleindex}\index{publisher=<publisher>|exampleindex}\index{date=<date>|exampleindex}\index{when=@when!<date>|exampleindex}\index{biblScope=<biblScope>|exampleindex}\index{unit=@unit!<biblScope>|exampleindex}\index{from=@from!<biblScope>|exampleindex}\index{to=@to!<biblScope>|exampleindex}\index{series=<series>|exampleindex}\index{title=<title>|exampleindex}\index{level=@level!<title>|exampleindex}\index{biblScope=<biblScope>|exampleindex}\index{unit=@unit!<biblScope>|exampleindex}\exampleFont \begin{shaded}\noindent\mbox{}{<\textbf{biblStruct}>}\mbox{}\newline 
\hspace*{1em}{<\textbf{analytic}>}\mbox{}\newline 
\hspace*{1em}\hspace*{1em}{<\textbf{author}>}\mbox{}\newline 
\hspace*{1em}\hspace*{1em}\hspace*{1em}{<\textbf{persName}>}\mbox{}\newline 
\hspace*{1em}\hspace*{1em}\hspace*{1em}\hspace*{1em}{<\textbf{surname}>}Wrigley{</\textbf{surname}>}\mbox{}\newline 
\hspace*{1em}\hspace*{1em}\hspace*{1em}\hspace*{1em}{<\textbf{forename}\hspace*{1em}{full}="{init}">}E.{</\textbf{forename}>}\mbox{}\newline 
\hspace*{1em}\hspace*{1em}\hspace*{1em}\hspace*{1em}{<\textbf{forename}\hspace*{1em}{full}="{init}">}A.{</\textbf{forename}>}\mbox{}\newline 
\hspace*{1em}\hspace*{1em}\hspace*{1em}{</\textbf{persName}>}\mbox{}\newline 
\hspace*{1em}\hspace*{1em}{</\textbf{author}>}\mbox{}\newline 
\hspace*{1em}\hspace*{1em}{<\textbf{title}\hspace*{1em}{level}="{a}">}Parish registers and the historian{</\textbf{title}>}\mbox{}\newline 
\hspace*{1em}{</\textbf{analytic}>}\mbox{}\newline 
\hspace*{1em}{<\textbf{monogr}>}\mbox{}\newline 
\hspace*{1em}\hspace*{1em}{<\textbf{author}>}\mbox{}\newline 
\hspace*{1em}\hspace*{1em}\hspace*{1em}{<\textbf{persName}>}\mbox{}\newline 
\hspace*{1em}\hspace*{1em}\hspace*{1em}\hspace*{1em}{<\textbf{surname}>}Steel{</\textbf{surname}>}\mbox{}\newline 
\hspace*{1em}\hspace*{1em}\hspace*{1em}\hspace*{1em}{<\textbf{forename}\hspace*{1em}{full}="{init}">}D.{</\textbf{forename}>}\mbox{}\newline 
\hspace*{1em}\hspace*{1em}\hspace*{1em}\hspace*{1em}{<\textbf{forename}\hspace*{1em}{full}="{init}">}J.{</\textbf{forename}>}\mbox{}\newline 
\hspace*{1em}\hspace*{1em}\hspace*{1em}{</\textbf{persName}>}\mbox{}\newline 
\hspace*{1em}\hspace*{1em}{</\textbf{author}>}\mbox{}\newline 
\hspace*{1em}\hspace*{1em}{<\textbf{author}>}\mbox{}\newline 
\hspace*{1em}\hspace*{1em}\hspace*{1em}{<\textbf{persName}>}\mbox{}\newline 
\hspace*{1em}\hspace*{1em}\hspace*{1em}\hspace*{1em}{<\textbf{surname}>}Steel{</\textbf{surname}>}\mbox{}\newline 
\hspace*{1em}\hspace*{1em}\hspace*{1em}\hspace*{1em}{<\textbf{forename}\hspace*{1em}{full}="{init}">}A.{</\textbf{forename}>}\mbox{}\newline 
\hspace*{1em}\hspace*{1em}\hspace*{1em}\hspace*{1em}{<\textbf{forename}\hspace*{1em}{full}="{init}">}E.{</\textbf{forename}>}\mbox{}\newline 
\hspace*{1em}\hspace*{1em}\hspace*{1em}\hspace*{1em}{<\textbf{forename}\hspace*{1em}{full}="{init}">}F.{</\textbf{forename}>}\mbox{}\newline 
\hspace*{1em}\hspace*{1em}\hspace*{1em}{</\textbf{persName}>}\mbox{}\newline 
\hspace*{1em}\hspace*{1em}{</\textbf{author}>}\mbox{}\newline 
\hspace*{1em}\hspace*{1em}{<\textbf{title}\hspace*{1em}{level}="{m}">}General sources of births, marriages and deaths before 1837{</\textbf{title}>}\mbox{}\newline 
\hspace*{1em}\hspace*{1em}{<\textbf{imprint}>}\mbox{}\newline 
\hspace*{1em}\hspace*{1em}\hspace*{1em}{<\textbf{pubPlace}>}London{</\textbf{pubPlace}>}\mbox{}\newline 
\hspace*{1em}\hspace*{1em}\hspace*{1em}{<\textbf{publisher}>}Society of Genealogists{</\textbf{publisher}>}\mbox{}\newline 
\hspace*{1em}\hspace*{1em}\hspace*{1em}{<\textbf{date}\hspace*{1em}{when}="{1968}"/>}\mbox{}\newline 
\hspace*{1em}\hspace*{1em}{</\textbf{imprint}>}\mbox{}\newline 
\hspace*{1em}\hspace*{1em}{<\textbf{biblScope}\hspace*{1em}{unit}="{page}"\hspace*{1em}{from}="{155}"\hspace*{1em}{to}="{167}">}155–167{</\textbf{biblScope}>}\mbox{}\newline 
\hspace*{1em}{</\textbf{monogr}>}\mbox{}\newline 
\hspace*{1em}{<\textbf{series}>}\mbox{}\newline 
\hspace*{1em}\hspace*{1em}{<\textbf{title}\hspace*{1em}{level}="{s}">}National index of parish registers{</\textbf{title}>}\mbox{}\newline 
\hspace*{1em}\hspace*{1em}{<\textbf{biblScope}\hspace*{1em}{unit}="{volume}">}1{</\textbf{biblScope}>}\mbox{}\newline 
\hspace*{1em}{</\textbf{series}>}\mbox{}\newline 
{</\textbf{biblStruct}>}\end{shaded}\egroup\par \par
The {\itshape unit} attribute on \hyperref[TEI.biblScope]{<biblScope>} is optional: both the following are legal examples: \par\bgroup\index{biblStruct=<biblStruct>|exampleindex}\index{analytic=<analytic>|exampleindex}\index{author=<author>|exampleindex}\index{author=<author>|exampleindex}\index{title=<title>|exampleindex}\index{level=@level!<title>|exampleindex}\index{monogr=<monogr>|exampleindex}\index{title=<title>|exampleindex}\index{level=@level!<title>|exampleindex}\index{imprint=<imprint>|exampleindex}\index{date=<date>|exampleindex}\index{biblScope=<biblScope>|exampleindex}\index{unit=@unit!<biblScope>|exampleindex}\index{biblScope=<biblScope>|exampleindex}\index{unit=@unit!<biblScope>|exampleindex}\index{biblScope=<biblScope>|exampleindex}\index{unit=@unit!<biblScope>|exampleindex}\exampleFont \begin{shaded}\noindent\mbox{}{<\textbf{biblStruct}>}\mbox{}\newline 
\hspace*{1em}{<\textbf{analytic}>}\mbox{}\newline 
\hspace*{1em}\hspace*{1em}{<\textbf{author}>}Boguraev, Branimir{</\textbf{author}>}\mbox{}\newline 
\hspace*{1em}\hspace*{1em}{<\textbf{author}>}Neff, Mary{</\textbf{author}>}\mbox{}\newline 
\hspace*{1em}\hspace*{1em}{<\textbf{title}\hspace*{1em}{level}="{a}">}Text Representation, Dictionary Structure,\mbox{}\newline 
\hspace*{1em}\hspace*{1em}\hspace*{1em}\hspace*{1em} and Lexical Knowledge{</\textbf{title}>}\mbox{}\newline 
\hspace*{1em}{</\textbf{analytic}>}\mbox{}\newline 
\hspace*{1em}{<\textbf{monogr}>}\mbox{}\newline 
\hspace*{1em}\hspace*{1em}{<\textbf{title}\hspace*{1em}{level}="{j}">}Literary \& Linguistic Computing{</\textbf{title}>}\mbox{}\newline 
\hspace*{1em}\hspace*{1em}{<\textbf{imprint}>}\mbox{}\newline 
\hspace*{1em}\hspace*{1em}\hspace*{1em}{<\textbf{date}>}1992{</\textbf{date}>}\mbox{}\newline 
\hspace*{1em}\hspace*{1em}{</\textbf{imprint}>}\mbox{}\newline 
\hspace*{1em}\hspace*{1em}{<\textbf{biblScope}\hspace*{1em}{unit}="{volume}">}7{</\textbf{biblScope}>}\mbox{}\newline 
\hspace*{1em}\hspace*{1em}{<\textbf{biblScope}\hspace*{1em}{unit}="{issue}">}2{</\textbf{biblScope}>}\mbox{}\newline 
\hspace*{1em}\hspace*{1em}{<\textbf{biblScope}\hspace*{1em}{unit}="{page}">}110-112{</\textbf{biblScope}>}\mbox{}\newline 
\hspace*{1em}{</\textbf{monogr}>}\mbox{}\newline 
{</\textbf{biblStruct}>}\end{shaded}\egroup\par \noindent  \par\bgroup\index{biblStruct=<biblStruct>|exampleindex}\index{analytic=<analytic>|exampleindex}\index{author=<author>|exampleindex}\index{title=<title>|exampleindex}\index{level=@level!<title>|exampleindex}\index{monogr=<monogr>|exampleindex}\index{title=<title>|exampleindex}\index{level=@level!<title>|exampleindex}\index{imprint=<imprint>|exampleindex}\index{date=<date>|exampleindex}\index{when=@when!<date>|exampleindex}\index{biblScope=<biblScope>|exampleindex}\index{biblScope=<biblScope>|exampleindex}\index{from=@from!<biblScope>|exampleindex}\index{to=@to!<biblScope>|exampleindex}\exampleFont \begin{shaded}\noindent\mbox{}{<\textbf{biblStruct}>}\mbox{}\newline 
\hspace*{1em}{<\textbf{analytic}>}\mbox{}\newline 
\hspace*{1em}\hspace*{1em}{<\textbf{author}>}Chesnutt, David{</\textbf{author}>}\mbox{}\newline 
\hspace*{1em}\hspace*{1em}{<\textbf{title}\hspace*{1em}{level}="{a}">}Historical Editions in the States{</\textbf{title}>}\mbox{}\newline 
\hspace*{1em}{</\textbf{analytic}>}\mbox{}\newline 
\hspace*{1em}{<\textbf{monogr}>}\mbox{}\newline 
\hspace*{1em}\hspace*{1em}{<\textbf{title}\hspace*{1em}{level}="{j}">}Computers and the Humanities{</\textbf{title}>}\mbox{}\newline 
\hspace*{1em}\hspace*{1em}{<\textbf{imprint}>}\mbox{}\newline 
\hspace*{1em}\hspace*{1em}\hspace*{1em}{<\textbf{date}\hspace*{1em}{when}="{1991-12}">}(December, 1991):{</\textbf{date}>}\mbox{}\newline 
\hspace*{1em}\hspace*{1em}{</\textbf{imprint}>}\mbox{}\newline 
\hspace*{1em}\hspace*{1em}{<\textbf{biblScope}>}25.6{</\textbf{biblScope}>}\mbox{}\newline 
\hspace*{1em}\hspace*{1em}{<\textbf{biblScope}\hspace*{1em}{from}="{377}"\hspace*{1em}{to}="{380}">}377–380{</\textbf{biblScope}>}\mbox{}\newline 
\hspace*{1em}{</\textbf{monogr}>}\mbox{}\newline 
{</\textbf{biblStruct}>}\end{shaded}\egroup\par \noindent    \par
On the other hand, a cited range encodes that the author \textit{cited only the portion} defined by this range. For example, a footnote following a quotation from page 378 of ‘Historical Editions in the States’ that includes a full bibliographic reference would be encoded using \hyperref[TEI.biblStruct]{<biblStruct>} as follows: \par\bgroup\index{biblStruct=<biblStruct>|exampleindex}\index{analytic=<analytic>|exampleindex}\index{author=<author>|exampleindex}\index{title=<title>|exampleindex}\index{level=@level!<title>|exampleindex}\index{monogr=<monogr>|exampleindex}\index{title=<title>|exampleindex}\index{level=@level!<title>|exampleindex}\index{imprint=<imprint>|exampleindex}\index{date=<date>|exampleindex}\index{when=@when!<date>|exampleindex}\index{biblScope=<biblScope>|exampleindex}\index{biblScope=<biblScope>|exampleindex}\index{unit=@unit!<biblScope>|exampleindex}\index{from=@from!<biblScope>|exampleindex}\index{to=@to!<biblScope>|exampleindex}\index{citedRange=<citedRange>|exampleindex}\exampleFont \begin{shaded}\noindent\mbox{}{<\textbf{biblStruct}>}\mbox{}\newline 
\hspace*{1em}{<\textbf{analytic}>}\mbox{}\newline 
\hspace*{1em}\hspace*{1em}{<\textbf{author}>}Chesnutt, David{</\textbf{author}>}\mbox{}\newline 
\hspace*{1em}\hspace*{1em}{<\textbf{title}\hspace*{1em}{level}="{a}">}Historical Editions in the States{</\textbf{title}>}\mbox{}\newline 
\hspace*{1em}{</\textbf{analytic}>}\mbox{}\newline 
\hspace*{1em}{<\textbf{monogr}>}\mbox{}\newline 
\hspace*{1em}\hspace*{1em}{<\textbf{title}\hspace*{1em}{level}="{j}">}Computers and the Humanities{</\textbf{title}>}\mbox{}\newline 
\hspace*{1em}\hspace*{1em}{<\textbf{imprint}>}\mbox{}\newline 
\hspace*{1em}\hspace*{1em}\hspace*{1em}{<\textbf{date}\hspace*{1em}{when}="{1991-12}">}(December, 1991):{</\textbf{date}>}\mbox{}\newline 
\hspace*{1em}\hspace*{1em}{</\textbf{imprint}>}\mbox{}\newline 
\hspace*{1em}\hspace*{1em}{<\textbf{biblScope}>}25.6{</\textbf{biblScope}>}\mbox{}\newline 
\hspace*{1em}\hspace*{1em}{<\textbf{biblScope}\hspace*{1em}{unit}="{page}"\hspace*{1em}{from}="{377}"\hspace*{1em}{to}="{380}">}377–380{</\textbf{biblScope}>}\mbox{}\newline 
\hspace*{1em}{</\textbf{monogr}>}\mbox{}\newline 
\hspace*{1em}{<\textbf{citedRange}>}378{</\textbf{citedRange}>}\mbox{}\newline 
{</\textbf{biblStruct}>}\end{shaded}\egroup\par 
\paragraph[{Series Information}]{Series Information}\label{COBICOS}\par
Series information may (in \hyperref[TEI.bibl]{<bibl>} elements) or must (in \hyperref[TEI.biblStruct]{<biblStruct>} elements) be enclosed in a \hyperref[TEI.series]{<series>} element or (in a \hyperref[TEI.biblFull]{<biblFull>} element) a \hyperref[TEI.seriesStmt]{<seriesStmt>} element. The title of the series may be tagged <title level="s">, the volume number <biblScope unit="volume">, and responsibility statements for the series (e.g. the name and affiliation of the editor, as in the example in section \textit{\hyperref[COBICOL]{3.12.2.1.\ Analytic, Monographic, and Series Levels}}) may be tagged \hyperref[TEI.editor]{<editor>} or \hyperref[TEI.respStmt]{<respStmt>}. Any identifier associated with the series itself should be marked using the \hyperref[TEI.idno]{<idno>} element.
\paragraph[{Related Items}]{Related Items}\label{COBIRI}\par
In bibliographic parlance, a \textit{related item} is any bibliographic item which, though related to that being defined, is distinct from it. The distinction between analytic and monographic items made above may be thought of as a special case of this kind of ‘related’ item. More usually however, the term is applied to such items as translations, continuations, different versions, parts, etc.\par
The element \hyperref[TEI.relatedItem]{<relatedItem>} is provided as a means of documenting such associated items: 
\begin{sansreflist}
  
\item [\textbf{<relatedItem>}] contains or references some other bibliographic item which is related to the present one in some specified manner, for example as a constituent or alternative version of it.
\end{sansreflist}
\par
In the following example, the first \hyperref[TEI.biblStruct]{<biblStruct>} describes a facsimile edition, and the second describes the work of which it is a facsimile. The relation between the facsimile and its source is represented by means of a \hyperref[TEI.relatedItem]{<relatedItem>} within the first description, which points to the description of the source. \par\bgroup\index{biblStruct=<biblStruct>|exampleindex}\index{monogr=<monogr>|exampleindex}\index{author=<author>|exampleindex}\index{title=<title>|exampleindex}\index{level=@level!<title>|exampleindex}\index{title=<title>|exampleindex}\index{level=@level!<title>|exampleindex}\index{editor=<editor>|exampleindex}\index{imprint=<imprint>|exampleindex}\index{pubPlace=<pubPlace>|exampleindex}\index{publisher=<publisher>|exampleindex}\index{date=<date>|exampleindex}\index{relatedItem=<relatedItem>|exampleindex}\index{type=@type!<relatedItem>|exampleindex}\index{ref=<ref>|exampleindex}\index{target=@target!<ref>|exampleindex}\index{biblStruct=<biblStruct>|exampleindex}\index{monogr=<monogr>|exampleindex}\index{author=<author>|exampleindex}\index{title=<title>|exampleindex}\index{level=@level!<title>|exampleindex}\index{imprint=<imprint>|exampleindex}\index{pubPlace=<pubPlace>|exampleindex}\index{publisher=<publisher>|exampleindex}\index{date=<date>|exampleindex}\exampleFont \begin{shaded}\noindent\mbox{}{<\textbf{biblStruct}\hspace*{1em}{xml:id}="{bibl03}">}\mbox{}\newline 
\hspace*{1em}{<\textbf{monogr}>}\mbox{}\newline 
\hspace*{1em}\hspace*{1em}{<\textbf{author}>}Swinburne, Algernon Charles{</\textbf{author}>}\mbox{}\newline 
\hspace*{1em}\hspace*{1em}{<\textbf{title}\hspace*{1em}{level}="{m}">}Swinburne's {<\textbf{title}\hspace*{1em}{level}="{m}">}Atalanta in Calydon{</\textbf{title}>}: A Facsimile of the\mbox{}\newline 
\hspace*{1em}\hspace*{1em}\hspace*{1em}\hspace*{1em} First Edition{</\textbf{title}>}\mbox{}\newline 
\hspace*{1em}\hspace*{1em}{<\textbf{editor}>}Georges Lafourcade{</\textbf{editor}>}\mbox{}\newline 
\hspace*{1em}\hspace*{1em}{<\textbf{imprint}>}\mbox{}\newline 
\hspace*{1em}\hspace*{1em}\hspace*{1em}{<\textbf{pubPlace}>}London{</\textbf{pubPlace}>}\mbox{}\newline 
\hspace*{1em}\hspace*{1em}\hspace*{1em}{<\textbf{publisher}>}Oxford UP{</\textbf{publisher}>}\mbox{}\newline 
\hspace*{1em}\hspace*{1em}\hspace*{1em}{<\textbf{date}>}1930{</\textbf{date}>}\mbox{}\newline 
\hspace*{1em}\hspace*{1em}{</\textbf{imprint}>}\mbox{}\newline 
\hspace*{1em}{</\textbf{monogr}>}\mbox{}\newline 
\hspace*{1em}{<\textbf{relatedItem}\hspace*{1em}{type}="{otherEdition}">}\mbox{}\newline 
\hspace*{1em}\hspace*{1em}{<\textbf{ref}\hspace*{1em}{target}="{\#bibl04}"/>}\mbox{}\newline 
\hspace*{1em}{</\textbf{relatedItem}>}\mbox{}\newline 
{</\textbf{biblStruct}>}\mbox{}\newline 
{<\textbf{biblStruct}\hspace*{1em}{xml:id}="{bibl04}">}\mbox{}\newline 
\hspace*{1em}{<\textbf{monogr}>}\mbox{}\newline 
\hspace*{1em}\hspace*{1em}{<\textbf{author}>} Swinburne, Algernon Charles{</\textbf{author}>}\mbox{}\newline 
\hspace*{1em}\hspace*{1em}{<\textbf{title}\hspace*{1em}{level}="{m}">}Atalanta in Calydon{</\textbf{title}>}\mbox{}\newline 
\hspace*{1em}\hspace*{1em}{<\textbf{imprint}>}\mbox{}\newline 
\hspace*{1em}\hspace*{1em}\hspace*{1em}{<\textbf{pubPlace}>}London{</\textbf{pubPlace}>}\mbox{}\newline 
\hspace*{1em}\hspace*{1em}\hspace*{1em}{<\textbf{publisher}>}Edward Moxon{</\textbf{publisher}>}\mbox{}\newline 
\hspace*{1em}\hspace*{1em}\hspace*{1em}{<\textbf{date}>}1865{</\textbf{date}>}\mbox{}\newline 
\hspace*{1em}\hspace*{1em}{</\textbf{imprint}>}\mbox{}\newline 
\hspace*{1em}{</\textbf{monogr}>}\mbox{}\newline 
{</\textbf{biblStruct}>}\end{shaded}\egroup\par \par
The \hyperref[TEI.ref]{<ref>} element in the above example could be replaced by the referenced \hyperref[TEI.biblStruct]{<biblStruct>} itself since a \hyperref[TEI.relatedItem]{<relatedItem>} may contain any form of bibliographic reference. For example, one of the examples quoted above might also be encoded as follows: \par\bgroup\index{biblStruct=<biblStruct>|exampleindex}\index{monogr=<monogr>|exampleindex}\index{author=<author>|exampleindex}\index{title=<title>|exampleindex}\index{type=@type!<title>|exampleindex}\index{imprint=<imprint>|exampleindex}\index{pubPlace=<pubPlace>|exampleindex}\index{publisher=<publisher>|exampleindex}\index{date=<date>|exampleindex}\index{extent=<extent>|exampleindex}\index{series=<series>|exampleindex}\index{title=<title>|exampleindex}\index{level=@level!<title>|exampleindex}\index{relatedItem=<relatedItem>|exampleindex}\index{type=@type!<relatedItem>|exampleindex}\index{biblStruct=<biblStruct>|exampleindex}\index{monogr=<monogr>|exampleindex}\index{author=<author>|exampleindex}\index{title=<title>|exampleindex}\index{type=@type!<title>|exampleindex}\index{level=@level!<title>|exampleindex}\index{title=<title>|exampleindex}\index{type=@type!<title>|exampleindex}\index{level=@level!<title>|exampleindex}\index{imprint=<imprint>|exampleindex}\index{pubPlace=<pubPlace>|exampleindex}\index{publisher=<publisher>|exampleindex}\index{date=<date>|exampleindex}\index{when=@when!<date>|exampleindex}\index{extent=<extent>|exampleindex}\exampleFont \begin{shaded}\noindent\mbox{}{<\textbf{biblStruct}>}\mbox{}\newline 
\hspace*{1em}{<\textbf{monogr}>}\mbox{}\newline 
\hspace*{1em}\hspace*{1em}{<\textbf{author}>}Shirley, James{</\textbf{author}>}\mbox{}\newline 
\hspace*{1em}\hspace*{1em}{<\textbf{title}\hspace*{1em}{type}="{main}">}The gentlemen of Venice{</\textbf{title}>}\mbox{}\newline 
\hspace*{1em}\hspace*{1em}{<\textbf{imprint}>}\mbox{}\newline 
\hspace*{1em}\hspace*{1em}\hspace*{1em}{<\textbf{pubPlace}>}New York{</\textbf{pubPlace}>}\mbox{}\newline 
\hspace*{1em}\hspace*{1em}\hspace*{1em}{<\textbf{publisher}>}Readex Microprint{</\textbf{publisher}>}\mbox{}\newline 
\hspace*{1em}\hspace*{1em}\hspace*{1em}{<\textbf{date}>}1953{</\textbf{date}>}\mbox{}\newline 
\hspace*{1em}\hspace*{1em}{</\textbf{imprint}>}\mbox{}\newline 
\hspace*{1em}\hspace*{1em}{<\textbf{extent}>}1 microprint card, 23 x 15 cm.{</\textbf{extent}>}\mbox{}\newline 
\hspace*{1em}{</\textbf{monogr}>}\mbox{}\newline 
\hspace*{1em}{<\textbf{series}>}\mbox{}\newline 
\hspace*{1em}\hspace*{1em}{<\textbf{title}\hspace*{1em}{level}="{s}">}Three centuries of drama: English, 1642–1700{</\textbf{title}>}\mbox{}\newline 
\hspace*{1em}{</\textbf{series}>}\mbox{}\newline 
\hspace*{1em}{<\textbf{relatedItem}\hspace*{1em}{type}="{otherEdition}">}\mbox{}\newline 
\hspace*{1em}\hspace*{1em}{<\textbf{biblStruct}>}\mbox{}\newline 
\hspace*{1em}\hspace*{1em}\hspace*{1em}{<\textbf{monogr}>}\mbox{}\newline 
\hspace*{1em}\hspace*{1em}\hspace*{1em}\hspace*{1em}{<\textbf{author}>}Shirley, James{</\textbf{author}>}\mbox{}\newline 
\hspace*{1em}\hspace*{1em}\hspace*{1em}\hspace*{1em}{<\textbf{title}\hspace*{1em}{type}="{main}"\hspace*{1em}{level}="{m}">}The gentlemen of Venice{</\textbf{title}>}\mbox{}\newline 
\hspace*{1em}\hspace*{1em}\hspace*{1em}\hspace*{1em}{<\textbf{title}\hspace*{1em}{type}="{sub}"\hspace*{1em}{level}="{m}">}a tragi-comedie presented at the private\mbox{}\newline 
\hspace*{1em}\hspace*{1em}\hspace*{1em}\hspace*{1em}\hspace*{1em}\hspace*{1em}\hspace*{1em}\hspace*{1em} house in Salisbury Court by Her Majesties servants{</\textbf{title}>}\mbox{}\newline 
\hspace*{1em}\hspace*{1em}\hspace*{1em}\hspace*{1em}{<\textbf{imprint}>}\mbox{}\newline 
\hspace*{1em}\hspace*{1em}\hspace*{1em}\hspace*{1em}\hspace*{1em}{<\textbf{pubPlace}>}London{</\textbf{pubPlace}>}\mbox{}\newline 
\hspace*{1em}\hspace*{1em}\hspace*{1em}\hspace*{1em}\hspace*{1em}{<\textbf{publisher}>}H. Moseley{</\textbf{publisher}>}\mbox{}\newline 
\hspace*{1em}\hspace*{1em}\hspace*{1em}\hspace*{1em}\hspace*{1em}{<\textbf{date}\hspace*{1em}{when}="{1655}">}1655{</\textbf{date}>}\mbox{}\newline 
\hspace*{1em}\hspace*{1em}\hspace*{1em}\hspace*{1em}{</\textbf{imprint}>}\mbox{}\newline 
\hspace*{1em}\hspace*{1em}\hspace*{1em}\hspace*{1em}{<\textbf{extent}>}78 p.{</\textbf{extent}>}\mbox{}\newline 
\hspace*{1em}\hspace*{1em}\hspace*{1em}{</\textbf{monogr}>}\mbox{}\newline 
\hspace*{1em}\hspace*{1em}{</\textbf{biblStruct}>}\mbox{}\newline 
\hspace*{1em}{</\textbf{relatedItem}>}\mbox{}\newline 
{</\textbf{biblStruct}>}\end{shaded}\egroup\par \par
The {\itshape type} attribute should be used to indicate the relationship between the bibliographic item and any \hyperref[TEI.relatedItem]{<relatedItem>} it contains or points to. The relationships may be transitive (for example translatedAs or reprintedFrom) or non-transitive (for example otherEdition). The {\itshape subtype} attribute may be used to provide a more detailed classification, where this is appropriate. Some further examples follow: \par\bgroup\index{biblStruct=<biblStruct>|exampleindex}\index{monogr=<monogr>|exampleindex}\index{author=<author>|exampleindex}\index{title=<title>|exampleindex}\index{level=@level!<title>|exampleindex}\index{title=<title>|exampleindex}\index{type=@type!<title>|exampleindex}\index{editor=<editor>|exampleindex}\index{role=@role!<editor>|exampleindex}\index{imprint=<imprint>|exampleindex}\index{pubPlace=<pubPlace>|exampleindex}\index{publisher=<publisher>|exampleindex}\index{date=<date>|exampleindex}\index{relatedItem=<relatedItem>|exampleindex}\index{type=@type!<relatedItem>|exampleindex}\index{bibl=<bibl>|exampleindex}\index{author=<author>|exampleindex}\index{title=<title>|exampleindex}\index{level=@level!<title>|exampleindex}\index{publisher=<publisher>|exampleindex}\index{date=<date>|exampleindex}\exampleFont \begin{shaded}\noindent\mbox{}{<\textbf{biblStruct}>}\mbox{}\newline 
\hspace*{1em}{<\textbf{monogr}>}\mbox{}\newline 
\hspace*{1em}\hspace*{1em}{<\textbf{author}>}Tolkien, J.R.R.{</\textbf{author}>}\mbox{}\newline 
\hspace*{1em}\hspace*{1em}{<\textbf{title}\hspace*{1em}{level}="{m}">}Den hobbit{</\textbf{title}>}\mbox{}\newline 
\hspace*{1em}\hspace*{1em}{<\textbf{title}\hspace*{1em}{type}="{sub}">}aus dem Engleschen iwwersat{</\textbf{title}>}\mbox{}\newline 
\hspace*{1em}\hspace*{1em}{<\textbf{editor}\hspace*{1em}{role}="{translator}">}Henry Wickens{</\textbf{editor}>}\mbox{}\newline 
\hspace*{1em}\hspace*{1em}{<\textbf{imprint}>}\mbox{}\newline 
\hspace*{1em}\hspace*{1em}\hspace*{1em}{<\textbf{pubPlace}>}Esch-sur-Sûre{</\textbf{pubPlace}>}\mbox{}\newline 
\hspace*{1em}\hspace*{1em}\hspace*{1em}{<\textbf{publisher}>}Op der Lay S. àr. L{</\textbf{publisher}>}\mbox{}\newline 
\hspace*{1em}\hspace*{1em}\hspace*{1em}{<\textbf{date}>}2002{</\textbf{date}>}\mbox{}\newline 
\hspace*{1em}\hspace*{1em}{</\textbf{imprint}>}\mbox{}\newline 
\hspace*{1em}{</\textbf{monogr}>}\mbox{}\newline 
\hspace*{1em}{<\textbf{relatedItem}\hspace*{1em}{type}="{translatedFrom}">}\mbox{}\newline 
\hspace*{1em}\hspace*{1em}{<\textbf{bibl}>}\mbox{}\newline 
\hspace*{1em}\hspace*{1em}\hspace*{1em}{<\textbf{author}>}Tolkien, J.R.R.{</\textbf{author}>}\mbox{}\newline 
\hspace*{1em}\hspace*{1em}\hspace*{1em}{<\textbf{title}\hspace*{1em}{level}="{m}">}The Hobbit{</\textbf{title}>}.\mbox{}\newline 
\hspace*{1em}\hspace*{1em}{<\textbf{publisher}>}Collins{</\textbf{publisher}>}\mbox{}\newline 
\hspace*{1em}\hspace*{1em}\hspace*{1em}{<\textbf{date}>}1997{</\textbf{date}>}\mbox{}\newline 
\hspace*{1em}\hspace*{1em}{</\textbf{bibl}>}\mbox{}\newline 
\hspace*{1em}{</\textbf{relatedItem}>}\mbox{}\newline 
{</\textbf{biblStruct}>}\end{shaded}\egroup\par \noindent  In this example, a full bibliographic description of the edition used as source for the translation is provided within the content of the \hyperref[TEI.relatedItem]{<relatedItem>}. Alternatively this might be provided by means of a link, in which case the \hyperref[TEI.relatedItem]{<relatedItem>} would be empty: \par\bgroup\index{relatedItem=<relatedItem>|exampleindex}\index{type=@type!<relatedItem>|exampleindex}\index{target=@target!<relatedItem>|exampleindex}\exampleFont \begin{shaded}\noindent\mbox{}{<\textbf{relatedItem}\hspace*{1em}{type}="{translatedFrom}"\mbox{}\newline 
\hspace*{1em}{target}="{http://www.example.com/bibliography.xml\#TOLK97}"/>}\end{shaded}\egroup\par 
\paragraph[{Notes and Statement of Language}]{Notes and Statement of Language}\label{COBICON}\par
Explanatory notes about the publication of unusual items, the form of an item (e.g. \textit{[Score]} or \textit{[Microform]}), or its provenance (e.g. \textit{translation of ...}) may be tagged using the \hyperref[TEI.note]{<note>} element. The same element may be used for any descriptive annotation of a bibliographic entry in a database. 
\begin{sansreflist}
  
\item [\textbf{<note>}] (note) contains a note or annotation.
\end{sansreflist}
\par
For example: \par\bgroup\index{bibl=<bibl>|exampleindex}\index{author=<author>|exampleindex}\index{title=<title>|exampleindex}\index{level=@level!<title>|exampleindex}\index{title=<title>|exampleindex}\index{level=@level!<title>|exampleindex}\index{biblScope=<biblScope>|exampleindex}\index{note=<note>|exampleindex}\exampleFont \begin{shaded}\noindent\mbox{}{<\textbf{bibl}>}\mbox{}\newline 
\hspace*{1em}{<\textbf{author}>}Coombs, James H., Allen H. Renear,\mbox{}\newline 
\hspace*{1em}\hspace*{1em} and Steven J. DeRose.{</\textbf{author}>}\mbox{}\newline 
\hspace*{1em}{<\textbf{title}\hspace*{1em}{level}="{a}">}Markup Systems and the Future of Scholarly\mbox{}\newline 
\hspace*{1em}\hspace*{1em} Text Processing.{</\textbf{title}>}\mbox{}\newline 
\hspace*{1em}{<\textbf{title}\hspace*{1em}{level}="{j}">}Communications of the ACM{</\textbf{title}>}\mbox{}\newline 
\hspace*{1em}{<\textbf{biblScope}>}30.11 (November 1987): 933–947.{</\textbf{biblScope}>}\mbox{}\newline 
\hspace*{1em}{<\textbf{note}>}Classic polemic supporting descriptive over procedural\mbox{}\newline 
\hspace*{1em}\hspace*{1em} markup in scholarly work.{</\textbf{note}>}\mbox{}\newline 
{</\textbf{bibl}>}\end{shaded}\egroup\par \par
The \hyperref[TEI.textLang]{<textLang>} element may be used to record information about the languages used within a bibliographic item. 
\begin{sansreflist}
  
\item [\textbf{<textLang>}] (text language) describes the languages and writing systems identified within the bibliographic work being described, rather than its description.\hfil\\[-10pt]\begin{sansreflist}
    \item[@{\itshape mainLang}]
  (main language) supplies a code which identifies the chief language used in the bibliographic work.
    \item[@{\itshape otherLangs}]
  (other languages) one or more codes identifying any other languages used in the bibliographic work.
\end{sansreflist}  
\end{sansreflist}
 This element can take the form of a simple note such as: \par\bgroup\index{textLang=<textLang>|exampleindex}\exampleFont \begin{shaded}\noindent\mbox{}{<\textbf{textLang}>}Latin, with some glosses in Anglo-Saxon and French{</\textbf{textLang}>}\end{shaded}\egroup\par \noindent  However, it is generally recommended where feasible to use the {\itshape mainLang} attribute to record the chief language of the bibliographic item, and optionally the {\itshape otherLangs} to identify other languages used in the work. For example: \par\bgroup\index{textLang=<textLang>|exampleindex}\index{mainLang=@mainLang!<textLang>|exampleindex}\index{otherLangs=@otherLangs!<textLang>|exampleindex}\exampleFont \begin{shaded}\noindent\mbox{}{<\textbf{textLang}\hspace*{1em}{mainLang}="{la}"\hspace*{1em}{otherLangs}="{ang fr}">}Latin, with some glosses in Anglo-Saxon and French{</\textbf{textLang}>}\end{shaded}\egroup\par \par
The {\itshape mainLang} and {\itshape otherLangs} attributes should both provide language identifiers in the same form as used for {\itshape xml:lang} as described at \textit{\hyperref[CHSH]{vi.1\ Language Identification}}. Where additional detail is needed correctly to describe a language, or to discuss its deployment in a given text, this should be done using the \hyperref[TEI.langUsage]{<langUsage>} element in the TEI header, within which individual \hyperref[TEI.language]{<language>} elements document the languages used: see \textit{\hyperref[HD41]{2.4.2.\ Language Usage}}.\par
A description, in French, of a work predominantly in German, but also with some Latin might have a \hyperref[TEI.textLang]{<textLang>} like the following: \par\bgroup\index{textLang=<textLang>|exampleindex}\index{mainLang=@mainLang!<textLang>|exampleindex}\index{otherLangs=@otherLangs!<textLang>|exampleindex}\exampleFont \begin{shaded}\noindent\mbox{}{<\textbf{textLang}\hspace*{1em}{xml:lang}="{fr}"\hspace*{1em}{mainLang}="{de}"\mbox{}\newline 
\hspace*{1em}{otherLangs}="{la}">}allemand et latin{</\textbf{textLang}>}\end{shaded}\egroup\par \noindent  For more information about the use of \hyperref[TEI.textLang]{<textLang>} in manuscript descriptions see: \textit{\hyperref[mslangs]{10.6.6.\ Languages and Writing Systems}}.
\paragraph[{Order of Components within References}]{Order of Components within References}\label{COBICOO}\par
The order of elements in \hyperref[TEI.bibl]{<bibl>} elements is not constrained.\par
In \hyperref[TEI.biblStruct]{<biblStruct>} elements, the \hyperref[TEI.analytic]{<analytic>} element, if it occurs, must come first, followed by one or more \hyperref[TEI.monogr]{<monogr>} and \hyperref[TEI.series]{<series>} elements, which may appear intermingled (as long as a \hyperref[TEI.monogr]{<monogr>} element comes first), and then zero or more of the following in any order: \hyperref[TEI.note]{<note>}, \hyperref[TEI.witDetail]{<witDetail>}, \hyperref[TEI.idno]{<idno>}, \hyperref[TEI.ptr]{<ptr>}, \hyperref[TEI.ref]{<ref>}, \hyperref[TEI.relatedItem]{<relatedItem>}, and \hyperref[TEI.citedRange]{<citedRange>}. Within \hyperref[TEI.analytic]{<analytic>}, the title(s), author(s), editor(s), and other statements of responsibility may appear in any order; it is recommended that all forms of the title be given together. Within \hyperref[TEI.monogr]{<monogr>}, the author, editor, and statements of responsibility may either come first or else follow the monographic title(s). Following these, the elements listed below, if present, must appear in the following order: \begin{itemize}
\item \hyperref[TEI.note]{<note>}s on the publication (and \hyperref[TEI.meeting]{<meeting>} elements describing the conference, in the case of a proceedings volume)
\item \hyperref[TEI.edition]{<edition>} elements, each followed by any related \hyperref[TEI.editor]{<editor>} or \hyperref[TEI.respStmt]{<respStmt>} elements
\item \hyperref[TEI.imprint]{<imprint>}
\item \hyperref[TEI.biblScope]{<biblScope>}
\end{itemize}  Within \hyperref[TEI.imprint]{<imprint>}, the elements allowed may appear in any order.\par
Finally, within the \hyperref[TEI.series]{<series>} information in a \hyperref[TEI.biblStruct]{<biblStruct>}, the sequence of elements is not constrained.\par
If more detailed structuring of a bibliographic description is required, the \hyperref[TEI.biblFull]{<biblFull>} element should be used. This is not further described here, as its contents are essentially equivalent to those of the \hyperref[TEI.fileDesc]{<fileDesc>} element in the \hyperref[TEI.teiHeader]{<teiHeader>}, which is fully described in section \textit{\hyperref[HD2]{2.2.\ The File Description}}.
\subsubsection[{Bibliographic Pointers }]{Bibliographic Pointers }\label{COBIXR}\par
References which are pointers to bibliographic items, of whatever kind, should be treated in the same way as other cross-references (see section \textit{\hyperref[COXR]{3.7.\ Simple Links and Cross-References}}). As discussed in that section, cross-referencing within TEI texts is in general represented by means of \hyperref[TEI.ptr]{<ptr>} or \hyperref[TEI.ref]{<ref>} elements. A {\itshape target} attribute on these elements is used to supply an identifying value for the target of the cross-reference, which should be, in the case of bibliographic elements, a bibliographic reference of some kind. Where the form of the reference itself is unimportant, or may be reconstructed mechanically, or is not to be encoded, the \hyperref[TEI.ptr]{<ptr>} element is used, as in the following example: \par\bgroup\index{ptr=<ptr>|exampleindex}\index{target=@target!<ptr>|exampleindex}\exampleFont \begin{shaded}\noindent\mbox{}As shown above ({<\textbf{ptr}\hspace*{1em}{target}="{\#NEL80}"/>}) ...\end{shaded}\egroup\par \par
Where the form of the reference is important, or contains additional qualifying information which is to be kept but distinguished from the surrounding text, the \hyperref[TEI.ref]{<ref>} element should be used, as in the following example: \par\bgroup\index{ref=<ref>|exampleindex}\index{target=@target!<ref>|exampleindex}\exampleFont \begin{shaded}\noindent\mbox{}Nelson claims {<\textbf{ref}\hspace*{1em}{target}="{\#NEL80}">}(ibid, passim){</\textbf{ref}>} ...\end{shaded}\egroup\par \noindent  It may be important to distinguish between the short form of a bibliographic reference and some qualifying or additional information. The latter should not appear within the scope of the \hyperref[TEI.ref]{<ref>} element when this is the case, as for example in an application concerned to normalize bibliographic references: \par\bgroup\index{ref=<ref>|exampleindex}\index{target=@target!<ref>|exampleindex}\exampleFont \begin{shaded}\noindent\mbox{}Nelson claims ({<\textbf{ref}\hspace*{1em}{target}="{\#NEL80}">}Nelson [1980]{</\textbf{ref}>} pages 13–37)\mbox{}\newline 
 ...\end{shaded}\egroup\par \noindent  If it is desired to capture additional information like this in a short-form reference, then \hyperref[TEI.bibl]{<bibl>} may be used with the {\itshape corresp} attribute pointing to the full bibliographic reference: \par\bgroup\index{bibl=<bibl>|exampleindex}\index{corresp=@corresp!<bibl>|exampleindex}\index{biblScope=<biblScope>|exampleindex}\index{unit=@unit!<biblScope>|exampleindex}\index{from=@from!<biblScope>|exampleindex}\index{to=@to!<biblScope>|exampleindex}\exampleFont \begin{shaded}\noindent\mbox{}Nelson claims ({<\textbf{bibl}\hspace*{1em}{corresp}="{\#NEL80}">}Nelson [1980] pages {<\textbf{biblScope}\hspace*{1em}{unit}="{page}"\hspace*{1em}{from}="{13}"\hspace*{1em}{to}="{37}">}13–37{</\textbf{biblScope}>}\mbox{}\newline 
{</\textbf{bibl}>}) ...\end{shaded}\egroup\par \par
The \hyperref[TEI.ref]{<ref>} element may also be used to provide a reference to a copy of the bibliographic item itself, particularly if this is available online, as in the following example: \par\bgroup\index{biblStruct=<biblStruct>|exampleindex}\index{analytic=<analytic>|exampleindex}\index{author=<author>|exampleindex}\index{forename=<forename>|exampleindex}\index{surname=<surname>|exampleindex}\index{title=<title>|exampleindex}\index{level=@level!<title>|exampleindex}\index{ref=<ref>|exampleindex}\index{type=@type!<ref>|exampleindex}\index{monogr=<monogr>|exampleindex}\index{title=<title>|exampleindex}\index{level=@level!<title>|exampleindex}\index{imprint=<imprint>|exampleindex}\index{biblScope=<biblScope>|exampleindex}\index{unit=@unit!<biblScope>|exampleindex}\index{biblScope=<biblScope>|exampleindex}\index{unit=@unit!<biblScope>|exampleindex}\index{date=<date>|exampleindex}\exampleFont \begin{shaded}\noindent\mbox{}{<\textbf{biblStruct}>}\mbox{}\newline 
\hspace*{1em}{<\textbf{analytic}>}\mbox{}\newline 
\hspace*{1em}\hspace*{1em}{<\textbf{author}>}\mbox{}\newline 
\hspace*{1em}\hspace*{1em}\hspace*{1em}{<\textbf{forename}>}Suzana{</\textbf{forename}>}\mbox{}\newline 
\hspace*{1em}\hspace*{1em}\hspace*{1em}{<\textbf{surname}>}Sukovic{</\textbf{surname}>}\mbox{}\newline 
\hspace*{1em}\hspace*{1em}{</\textbf{author}>}\mbox{}\newline 
\hspace*{1em}\hspace*{1em}{<\textbf{title}\hspace*{1em}{level}="{a}">}Beyond the Scriptorium: The Role of the Library in Text\mbox{}\newline 
\hspace*{1em}\hspace*{1em}\hspace*{1em}\hspace*{1em} Encoding{</\textbf{title}>}\mbox{}\newline 
\hspace*{1em}\hspace*{1em}{<\textbf{ref}\hspace*{1em}{type}="{url}">}http://www.dlib.org/dlib/january02/sukovic/01sukovic.html{</\textbf{ref}>}\mbox{}\newline 
\hspace*{1em}{</\textbf{analytic}>}\mbox{}\newline 
\hspace*{1em}{<\textbf{monogr}>}\mbox{}\newline 
\hspace*{1em}\hspace*{1em}{<\textbf{title}\hspace*{1em}{level}="{j}">}D-Lib{</\textbf{title}>}\mbox{}\newline 
\hspace*{1em}\hspace*{1em}{<\textbf{imprint}>}\mbox{}\newline 
\hspace*{1em}\hspace*{1em}\hspace*{1em}{<\textbf{biblScope}\hspace*{1em}{unit}="{volume}">}8{</\textbf{biblScope}>}\mbox{}\newline 
\hspace*{1em}\hspace*{1em}\hspace*{1em}{<\textbf{biblScope}\hspace*{1em}{unit}="{issue}">}1{</\textbf{biblScope}>}\mbox{}\newline 
\hspace*{1em}\hspace*{1em}\hspace*{1em}{<\textbf{date}>}2002{</\textbf{date}>}\mbox{}\newline 
\hspace*{1em}\hspace*{1em}{</\textbf{imprint}>}\mbox{}\newline 
\hspace*{1em}{</\textbf{monogr}>}\mbox{}\newline 
{</\textbf{biblStruct}>}\end{shaded}\egroup\par \par
The \hyperref[TEI.ptr]{<ptr>} element may be used as a child element of \hyperref[TEI.biblStruct]{<biblStruct>} to refer to the online catalog record of this bibliographic item: \par\bgroup\index{biblStruct=<biblStruct>|exampleindex}\index{type=@type!<biblStruct>|exampleindex}\index{monogr=<monogr>|exampleindex}\index{author=<author>|exampleindex}\index{persName=<persName>|exampleindex}\index{forename=<forename>|exampleindex}\index{surname=<surname>|exampleindex}\index{title=<title>|exampleindex}\index{level=@level!<title>|exampleindex}\index{imprint=<imprint>|exampleindex}\index{date=<date>|exampleindex}\index{when=@when!<date>|exampleindex}\index{pubPlace=<pubPlace>|exampleindex}\index{publisher=<publisher>|exampleindex}\index{extent=<extent>|exampleindex}\index{ptr=<ptr>|exampleindex}\index{type=@type!<ptr>|exampleindex}\index{target=@target!<ptr>|exampleindex}\exampleFont \begin{shaded}\noindent\mbox{}{<\textbf{biblStruct}\hspace*{1em}{type}="{manifestation}">}\mbox{}\newline 
\hspace*{1em}{<\textbf{monogr}>}\mbox{}\newline 
\hspace*{1em}\hspace*{1em}{<\textbf{author}>}\mbox{}\newline 
\hspace*{1em}\hspace*{1em}\hspace*{1em}{<\textbf{persName}>}\mbox{}\newline 
\hspace*{1em}\hspace*{1em}\hspace*{1em}\hspace*{1em}{<\textbf{forename}>}Germain{</\textbf{forename}>}\mbox{}\newline 
\hspace*{1em}\hspace*{1em}\hspace*{1em}\hspace*{1em}{<\textbf{surname}>}Brice{</\textbf{surname}>}\mbox{}\newline 
\hspace*{1em}\hspace*{1em}\hspace*{1em}{</\textbf{persName}>}\mbox{}\newline 
\hspace*{1em}\hspace*{1em}{</\textbf{author}>}\mbox{}\newline 
\hspace*{1em}\hspace*{1em}{<\textbf{title}\hspace*{1em}{level}="{m}">}Description de la ville de Paris et de tout ce qu’elle contient de plus remarquable, par Germain Brice ; enrichie d’un nouveau plan et de figures dessinées et gravées correctement. 7e édition, revue et augmentée par l’auteur{</\textbf{title}>}\mbox{}\newline 
\hspace*{1em}\hspace*{1em}{<\textbf{imprint}>}\mbox{}\newline 
\hspace*{1em}\hspace*{1em}\hspace*{1em}{<\textbf{date}\hspace*{1em}{when}="{1717}">}1717{</\textbf{date}>}\mbox{}\newline 
\hspace*{1em}\hspace*{1em}\hspace*{1em}{<\textbf{pubPlace}>}Paris{</\textbf{pubPlace}>}\mbox{}\newline 
\hspace*{1em}\hspace*{1em}\hspace*{1em}{<\textbf{publisher}>}F. Fournier{</\textbf{publisher}>}\mbox{}\newline 
\hspace*{1em}\hspace*{1em}{</\textbf{imprint}>}\mbox{}\newline 
\hspace*{1em}\hspace*{1em}{<\textbf{extent}>}In-12{</\textbf{extent}>}\mbox{}\newline 
\hspace*{1em}{</\textbf{monogr}>}\mbox{}\newline 
\hspace*{1em}{<\textbf{ptr}\hspace*{1em}{type}="{catBnf}"\mbox{}\newline 
\hspace*{1em}\hspace*{1em}{target}="{http://catalogue.bnf.fr/ark:/12148/cb30160624f/}"/>}\mbox{}\newline 
{</\textbf{biblStruct}>}\end{shaded}\egroup\par 
\subsubsection[{Relationship to Other Bibliographic Schemes}]{Relationship to Other Bibliographic Schemes}\label{COBIOT}\par
The bibliographic tagging defined here can capture the distinctions required by most bibliographic encoding systems; for the benefit of users of some commonly used systems, the following lists of equivalences are offered, showing the relationship of the markup defined here to the fields defined for bibliographic records in the Scribe, BibTeX, and ProCite systems. \par
Listed below are the equivalences between the various bibliographic fields defined for use in the Scribe and BibTeX systems of bibliographic databases and the elements defined in this module.\footnote{The BibTeX scheme is intentionally compatible with that of Scribe, although it omits some fields used by Scribe. Hence only one list of fields is given here.} Elements and structures available in the module defined here which have no analogues in Scribe and BibTeX are not noted. \begin{description}

\item[{address}]tag as \hyperref[TEI.placeName]{<placeName>} or \hyperref[TEI.address]{<address>}
\item[{annote}]tag as \hyperref[TEI.note]{<note>}
\item[{author}]tag as \hyperref[TEI.author]{<author>}
\item[{booktitle}]tag as <title level="m"> or \hyperref[TEI.title]{<title>} within \hyperref[TEI.monogr]{<monogr>}
\item[{chapter}]tag as <biblScope unit="chap">
\item[{date}]used only to record date entry was made in the bibliographic database; not supported
\item[{edition}]tag as \hyperref[TEI.edition]{<edition>}
\item[{editor}]tag as \hyperref[TEI.editor]{<editor>} or \hyperref[TEI.respStmt]{<respStmt>}
\item[{editors}]tag as multiple \hyperref[TEI.editor]{<editor>} or \hyperref[TEI.respStmt]{<respStmt>} elements
\item[{fullauthor}]use the \hyperref[TEI.reg]{<reg>} element, possibly inside a \hyperref[TEI.choice]{<choice>} element, inside either an \hyperref[TEI.author]{<author>} or \hyperref[TEI.name]{<name>}
\item[{fullorganization}]use the \hyperref[TEI.reg]{<reg>} element, possibly inside a \hyperref[TEI.choice]{<choice>} element, inside a <name type="org">
\item[{howpublished}]tag as \hyperref[TEI.note]{<note>}, possibly using the form <note place="inline">
\item[{institution}]used only for issuer of technical reports; tag as \hyperref[TEI.publisher]{<publisher>}
\item[{journal}]tag as <title level="j"> or \hyperref[TEI.title]{<title>} within \hyperref[TEI.monogr]{<monogr>}
\item[{key}]used to specify an alternate sort key for the bibliographic item, for use instead of author's or editor's name; not supported
\item[{meeting}]tag as \hyperref[TEI.meeting]{<meeting>} or as \hyperref[TEI.note]{<note>}
\item[{month}]use \hyperref[TEI.date]{<date>}; if the date is not in a trivially parseable form, use the {\itshape when} attribute to provide a normalized equivalent in one of the format from \textit{XML Schema Part 2: Datatypes Second Edition}
\item[{note}]tag as \hyperref[TEI.note]{<note>}
\item[{number}]tag as <biblScope unit="issue"> or <biblScope unit="number">; for technical report numbers, use <idno type="docno">
\item[{organization}]used only for sponsor of conference; use <name type="org"> within \hyperref[TEI.respStmt]{<respStmt>} within \hyperref[TEI.meeting]{<meeting>} element
\item[{pages}]tag as <biblScope unit="pp">
\item[{publisher}]tag as \hyperref[TEI.publisher]{<publisher>}
\item[{school}]used only for institutions at which thesis work is done; tag as \hyperref[TEI.publisher]{<publisher>}
\item[{series}]tag as <title level="s"> or \hyperref[TEI.title]{<title>} within \hyperref[TEI.series]{<series>}
\item[{title}]tag as \hyperref[TEI.title]{<title>} in appropriate context or with appropriate {\itshape level} value
\item[{volume}]tag as <biblScope unit="volume">
\item[{year}]tag as \hyperref[TEI.date]{<date>}; if the date is not in a trivially parseable form, use the {\itshape when} attribute to provide an ISO-format equivalent
\end{description} 
\subsection[{Passages of Verse or Drama}]{Passages of Verse or Drama}\label{CODV}\par
The following elements are included in the core module for the convenience of those encoding texts which include mixtures of prose, verse and drama. 
\begin{sansreflist}
  
\item [\textbf{<l>}] (verse line) contains a single, possibly incomplete, line of verse.
\item [\textbf{<lg>}] (line group) contains one or more verse lines functioning as a formal unit, e.g. a stanza, refrain, verse paragraph, etc.
\item [\textbf{<sp>}] (speech) contains an individual speech in a performance text, or a passage presented as such in a prose or verse text.
\item [\textbf{<speaker>}] contains a specialized form of heading or label, giving the name of one or more speakers in a dramatic text or fragment.
\item [\textbf{<stage>}] (stage direction) contains any kind of stage direction within a dramatic text or fragment.
\end{sansreflist}
\par
Full details of other, more specialized, elements for the encoding of texts which are predominantly verse or drama are described in the appropriate chapter of part three (for verse, see the verse base described in chapter \textit{\hyperref[VE]{6.\ Verse}}; for performance texts, see the drama base described in chapter \textit{\hyperref[DR]{7.\ Performance Texts}}). In this section, we describe only the elements listed above, all of which can appear in any text, whichever of the three modes prose, verse, or drama may predominate in it.
\subsubsection[{Core Tags for Verse}]{Core Tags for Verse}\label{COVE}\par
Like other written texts, verse texts or poems may be hierarchically subdivided, for example into books or cantos. These structural subdivisions should be encoded using the general purpose \hyperref[TEI.div]{<div>} or \hyperref[TEI.div1]{<div1>} (etc.) elements described below in chapters \textit{\hyperref[DS]{4.\ Default Text Structure}} and \textit{\hyperref[VE]{6.\ Verse}}. The fundamental unit of a verse text is the verse line rather than the paragraph, however.\par
The \hyperref[TEI.l]{<l>} element is used to mark up verse lines, that is metrical rather than typographic lines. In some modern or free verse, it may be hard to decide whether the typographic line is to be regarded as a verse line or not, but the distinction is quite clear for verse following regular metrical patterns. Where a metrical line is interrupted by a typographic line break, the encoder may choose to ignore the fact entirely or to use the empty \hyperref[TEI.lb]{<lb>} (line break) element discussed in \textit{\hyperref[CORS]{3.11.\ Reference Systems}}. By convention, the start of a metrical line implies the start of a typographic line; hence there is no need to introduce an \hyperref[TEI.lb]{<lb>} tag at the start of every \hyperref[TEI.l]{<l>} element, but only at places where a new typographic line starts within a metrical line, as in the following example: \par\bgroup\index{l=<l>|exampleindex}\index{lb=<lb>|exampleindex}\index{l=<l>|exampleindex}\index{lb=<lb>|exampleindex}\index{l=<l>|exampleindex}\index{lb=<lb>|exampleindex}\index{l=<l>|exampleindex}\index{l=<l>|exampleindex}\exampleFont \begin{shaded}\noindent\mbox{}{<\textbf{l}>}Of Mans First Disobedience, and{<\textbf{lb}/>} the Fruit{</\textbf{l}>}\mbox{}\newline 
{<\textbf{l}>}Of that Forbidden Tree, whose{<\textbf{lb}/>} mortal tast{</\textbf{l}>}\mbox{}\newline 
{<\textbf{l}>}Brought Death into the World,{<\textbf{lb}/>} and all our woe,{</\textbf{l}>}\mbox{}\newline 
{<\textbf{l}>}With loss of Eden, till one greater Man{</\textbf{l}>}\mbox{}\newline 
{<\textbf{l}>}Restore us, and regain the blissful Seat...{</\textbf{l}>}\end{shaded}\egroup\par \noindent  In the original copy text, the presence of an ornamental capital at the start of the poem means that the measure is not wide enough to print the first four lines on four lines; instead each metrical line occupies two typographic lines, with a break at the point indicated. Note that this encoding makes no attempt to preserve information about the whitespace or indentation associated with either kind of line; if regarded as essential, this information would be recorded using the {\itshape rend} or {\itshape rendition} attributes discussed in \textit{\hyperref[STGA]{1.3.1.1.\ Global Attributes}}.\par
The \hyperref[TEI.l]{<l>} element should not be used to represent typographic lines in non-verse materials: if the line-breaking points in a prose text are considered important for analysis, they should be marked with the \hyperref[TEI.lb]{<lb>} element. Alternatively, a neutral segmentation element such as \hyperref[TEI.seg]{<seg>} or \hyperref[TEI.ab]{<ab>} may be used; see further discussion of these elements in chapter \textit{\hyperref[SA]{16.\ Linking, Segmentation, and Alignment}}. The \hyperref[TEI.l]{<l>} element is a member of the \textsf{model.lLike} class, which is a subclass of the \textsf{model.divPart} class, along with elements from the \textsf{model.pLike} (paragraph-like) class.\par
In some verse forms, regular groupings of lines are regarded as units of some kind, often identified by a regular verse scheme. In stichic verse and couplets, groups of lines analogous to paragraphs are often indicated by indentation. In other verse forms, lines are grouped into irregular sequences indicated simply by whitespace. The \hyperref[TEI.lg]{<lg>} or line group element may be used to mark any such grouping of elements from the \textsf{model.lLike} class. As a member of the \textsf{att.typed} class, the \hyperref[TEI.lg]{<lg>} element bears the following attributes: 
\begin{sansreflist}
  
\item [\textbf{att.typed}] provides attributes which can be used to classify or subclassify elements in any way.\hfil\\[-10pt]\begin{sansreflist}
    \item[@{\itshape type}]
  characterizes the element in some sense, using any convenient classification scheme or typology.
    \item[@{\itshape subtype}]
  (subtype) provides a sub-categorization of the element, if needed
\end{sansreflist}  
\end{sansreflist}
 which may be used to further categorize the line group where this is felt desirable, as in the following example. This example also demonstrates the {\itshape rend} attribute to indicate whether or not a line is indented. \par\bgroup\index{lg=<lg>|exampleindex}\index{l=<l>|exampleindex}\index{l=<l>|exampleindex}\index{rend=@rend!<l>|exampleindex}\index{l=<l>|exampleindex}\index{l=<l>|exampleindex}\index{rend=@rend!<l>|exampleindex}\index{l=<l>|exampleindex}\index{rend=@rend!<l>|exampleindex}\index{l=<l>|exampleindex}\index{lg=<lg>|exampleindex}\index{n=@n!<lg>|exampleindex}\index{type=@type!<lg>|exampleindex}\index{l=<l>|exampleindex}\index{l=<l>|exampleindex}\index{l=<l>|exampleindex}\index{l=<l>|exampleindex}\exampleFont \begin{shaded}\noindent\mbox{}{<\textbf{lg}>}\mbox{}\newline 
\hspace*{1em}{<\textbf{l}>}Come fill up the Glass,{</\textbf{l}>}\mbox{}\newline 
\hspace*{1em}{<\textbf{l}\hspace*{1em}{rend}="{indent}">}Round, round let it pass,{</\textbf{l}>}\mbox{}\newline 
\hspace*{1em}{<\textbf{l}>}'Till our Reason be lost in our Wine:{</\textbf{l}>}\mbox{}\newline 
\hspace*{1em}{<\textbf{l}\hspace*{1em}{rend}="{indent}">}Leave Conscience's Rules{</\textbf{l}>}\mbox{}\newline 
\hspace*{1em}{<\textbf{l}\hspace*{1em}{rend}="{indent}">}To Women and Fools,{</\textbf{l}>}\mbox{}\newline 
\hspace*{1em}{<\textbf{l}>}This only can make us divine.{</\textbf{l}>}\mbox{}\newline 
{</\textbf{lg}>}\mbox{}\newline 
{<\textbf{lg}\hspace*{1em}{n}="{Chorus}"\hspace*{1em}{type}="{refrain}">}\mbox{}\newline 
\hspace*{1em}{<\textbf{l}>}Then a Mohock, a Mohock I'll be,{</\textbf{l}>}\mbox{}\newline 
\hspace*{1em}{<\textbf{l}>}No Laws shall restrain{</\textbf{l}>}\mbox{}\newline 
\hspace*{1em}{<\textbf{l}>}Our Libertine Reign,{</\textbf{l}>}\mbox{}\newline 
\hspace*{1em}{<\textbf{l}>}We'll riot, drink on, and be free.{</\textbf{l}>}\mbox{}\newline 
{</\textbf{lg}>}\end{shaded}\egroup\par \par
For some kinds of analysis, it may be useful to identify different kinds of line group within the same piece of verse. Such line groups may self-nest, in much the same way as the un-numbered \hyperref[TEI.div]{<div>} element described in chapter \textit{\hyperref[DS]{4.\ Default Text Structure}}. For example: \par\bgroup\index{lg=<lg>|exampleindex}\index{type=@type!<lg>|exampleindex}\index{lg=<lg>|exampleindex}\index{type=@type!<lg>|exampleindex}\index{l=<l>|exampleindex}\index{l=<l>|exampleindex}\index{name=<name>|exampleindex}\index{l=<l>|exampleindex}\index{l=<l>|exampleindex}\index{l=<l>|exampleindex}\index{l=<l>|exampleindex}\index{l=<l>|exampleindex}\index{l=<l>|exampleindex}\index{lg=<lg>|exampleindex}\index{type=@type!<lg>|exampleindex}\index{l=<l>|exampleindex}\index{l=<l>|exampleindex}\index{l=<l>|exampleindex}\index{l=<l>|exampleindex}\index{name=<name>|exampleindex}\index{l=<l>|exampleindex}\index{name=<name>|exampleindex}\index{l=<l>|exampleindex}\index{rend=@rend!<l>|exampleindex}\exampleFont \begin{shaded}\noindent\mbox{}{<\textbf{lg}\hspace*{1em}{type}="{sonnet}">}\mbox{}\newline 
\hspace*{1em}{<\textbf{lg}\hspace*{1em}{type}="{octet}">}\mbox{}\newline 
\hspace*{1em}\hspace*{1em}{<\textbf{l}>}Thus speaks the Muse, and bends her brow severe:—{</\textbf{l}>}\mbox{}\newline 
\hspace*{1em}\hspace*{1em}{<\textbf{l}>}“Did I, {<\textbf{name}>}Lætitia{</\textbf{name}>}, lend my choicest lays,{</\textbf{l}>}\mbox{}\newline 
\hspace*{1em}\hspace*{1em}{<\textbf{l}>}And crown thy youthful head with freshest bays,{</\textbf{l}>}\mbox{}\newline 
\hspace*{1em}\hspace*{1em}{<\textbf{l}>}That all the' expectance of thy full-grown year{</\textbf{l}>}\mbox{}\newline 
\hspace*{1em}\hspace*{1em}{<\textbf{l}>}Should lie inert and fruitless? O revere{</\textbf{l}>}\mbox{}\newline 
\hspace*{1em}\hspace*{1em}{<\textbf{l}>}Those sacred gifts whose meed is deathless praise,{</\textbf{l}>}\mbox{}\newline 
\hspace*{1em}\hspace*{1em}{<\textbf{l}>}Whose potent charms the' enraptured soul can raise{</\textbf{l}>}\mbox{}\newline 
\hspace*{1em}\hspace*{1em}{<\textbf{l}>}Far from the vapours of this earthly sphere!{</\textbf{l}>}\mbox{}\newline 
\hspace*{1em}{</\textbf{lg}>}\mbox{}\newline 
\hspace*{1em}{<\textbf{lg}\hspace*{1em}{type}="{sestet}">}\mbox{}\newline 
\hspace*{1em}\hspace*{1em}{<\textbf{l}>}Seize, seize the lyre! resume the lofty strain!{</\textbf{l}>}\mbox{}\newline 
\hspace*{1em}\hspace*{1em}{<\textbf{l}>}'T is time, 't is time! hark how the nations round{</\textbf{l}>}\mbox{}\newline 
\hspace*{1em}\hspace*{1em}{<\textbf{l}>}With jocund notes of liberty resound,—{</\textbf{l}>}\mbox{}\newline 
\hspace*{1em}\hspace*{1em}{<\textbf{l}>}And thy own {<\textbf{name}>}Corsica{</\textbf{name}>} has burst her chain!{</\textbf{l}>}\mbox{}\newline 
\hspace*{1em}\hspace*{1em}{<\textbf{l}>}O let the song to {<\textbf{name}>}Britain's{</\textbf{name}>} shores rebound,{</\textbf{l}>}\mbox{}\newline 
\hspace*{1em}\hspace*{1em}{<\textbf{l}\hspace*{1em}{rend}="{indent(-1)}">}Where Freedom's once-loved voice is heard,\mbox{}\newline 
\hspace*{1em}\hspace*{1em}\hspace*{1em}\hspace*{1em} alas! in vain.”{</\textbf{l}>}\mbox{}\newline 
\hspace*{1em}{</\textbf{lg}>}\mbox{}\newline 
{</\textbf{lg}>}\end{shaded}\egroup\par \par
It is often the case that verse line boundaries conflict with the boundaries of other structural elements. In the following example, the single verse line ‘A Workeman in't... welcome’ is interrupted by a stage direction: \par\bgroup\index{l=<l>|exampleindex}\index{name=<name>|exampleindex}\index{l=<l>|exampleindex}\index{l=<l>|exampleindex}\index{l=<l>|exampleindex}\index{l=<l>|exampleindex}\index{part=@part!<l>|exampleindex}\index{stage=<stage>|exampleindex}\index{l=<l>|exampleindex}\index{part=@part!<l>|exampleindex}\exampleFont \begin{shaded}\noindent\mbox{}{<\textbf{l}>}Thou fumblest {<\textbf{name}>}Eros{</\textbf{name}>}, and my Queenes a Squire{</\textbf{l}>}\mbox{}\newline 
{<\textbf{l}>}More tight at this, then thou: Dispatch. O Loue,{</\textbf{l}>}\mbox{}\newline 
{<\textbf{l}>}That thou couldst see my Warres to day, and knew'st{</\textbf{l}>}\mbox{}\newline 
{<\textbf{l}>}The Royall Occupation, thou should'st see{</\textbf{l}>}\mbox{}\newline 
{<\textbf{l}\hspace*{1em}{part}="{I}">}A Workeman in't.{</\textbf{l}>}\mbox{}\newline 
{<\textbf{stage}>}Enter an Armed Soldier.{</\textbf{stage}>}\mbox{}\newline 
{<\textbf{l}\hspace*{1em}{part}="{F}">}Good morrow to thee, welcome.{</\textbf{l}>}\end{shaded}\egroup\par \noindent  In this encoding, the {\itshape part} attribute is used, as with \hyperref[TEI.div]{<div>}, to indicate that the last two \hyperref[TEI.l]{<l>} elements should be regarded as the initial and final parts of a single line, rather than as two lines.\par
The same technique may be used where verse lines are collected together into units such as verse paragraphs: \par\bgroup\index{lg=<lg>|exampleindex}\index{n=@n!<lg>|exampleindex}\index{type=@type!<lg>|exampleindex}\index{l=<l>|exampleindex}\index{l=<l>|exampleindex}\index{l=<l>|exampleindex}\index{part=@part!<l>|exampleindex}\index{lg=<lg>|exampleindex}\index{type=@type!<lg>|exampleindex}\index{n=@n!<lg>|exampleindex}\index{l=<l>|exampleindex}\index{part=@part!<l>|exampleindex}\index{l=<l>|exampleindex}\index{l=<l>|exampleindex}\exampleFont \begin{shaded}\noindent\mbox{}{<\textbf{lg}\hspace*{1em}{n}="{6}"\hspace*{1em}{type}="{para}">}\mbox{}\newline 
\textit{<!-- ... -->}\mbox{}\newline 
\hspace*{1em}{<\textbf{l}>}Unprofitably travelling toward the grave,{</\textbf{l}>}\mbox{}\newline 
\hspace*{1em}{<\textbf{l}>}Like a false steward who hath much received{</\textbf{l}>}\mbox{}\newline 
\hspace*{1em}{<\textbf{l}\hspace*{1em}{part}="{I}">}And renders nothing back.{</\textbf{l}>}\mbox{}\newline 
{</\textbf{lg}>}\mbox{}\newline 
{<\textbf{lg}\hspace*{1em}{type}="{para}"\hspace*{1em}{n}="{7}">}\mbox{}\newline 
\hspace*{1em}{<\textbf{l}\hspace*{1em}{part}="{F}">}Was it for this{</\textbf{l}>}\mbox{}\newline 
\hspace*{1em}{<\textbf{l}>}That one, the fairest of all rivers, loved{</\textbf{l}>}\mbox{}\newline 
\hspace*{1em}{<\textbf{l}>}To blend his murmurs with my nurse's song,{</\textbf{l}>}\mbox{}\newline 
\textit{<!-- ... -->}\mbox{}\newline 
{</\textbf{lg}>}\end{shaded}\egroup\par \par
The {\itshape part} attribute may also be attached to an \hyperref[TEI.lg]{<lg>} element to indicate that it is incomplete, for example because it forms part of a group that is divided between two speakers, as in the following example: \par\bgroup\index{sp=<sp>|exampleindex}\index{speaker=<speaker>|exampleindex}\index{lg=<lg>|exampleindex}\index{type=@type!<lg>|exampleindex}\index{part=@part!<lg>|exampleindex}\index{l=<l>|exampleindex}\index{l=<l>|exampleindex}\index{sp=<sp>|exampleindex}\index{speaker=<speaker>|exampleindex}\index{lg=<lg>|exampleindex}\index{type=@type!<lg>|exampleindex}\index{part=@part!<lg>|exampleindex}\index{l=<l>|exampleindex}\index{l=<l>|exampleindex}\exampleFont \begin{shaded}\noindent\mbox{}{<\textbf{sp}>}\mbox{}\newline 
\hspace*{1em}{<\textbf{speaker}>}First Voice{</\textbf{speaker}>}\mbox{}\newline 
\hspace*{1em}{<\textbf{lg}\hspace*{1em}{type}="{stanza}"\hspace*{1em}{part}="{I}">}\mbox{}\newline 
\hspace*{1em}\hspace*{1em}{<\textbf{l}>}But why drives on that ship so fast{</\textbf{l}>}\mbox{}\newline 
\hspace*{1em}\hspace*{1em}{<\textbf{l}>}Withouten wave or wind?{</\textbf{l}>}\mbox{}\newline 
\hspace*{1em}{</\textbf{lg}>}\mbox{}\newline 
{</\textbf{sp}>}\mbox{}\newline 
{<\textbf{sp}>}\mbox{}\newline 
\hspace*{1em}{<\textbf{speaker}>}Second Voice{</\textbf{speaker}>}\mbox{}\newline 
\hspace*{1em}{<\textbf{lg}\hspace*{1em}{type}="{stanza}"\hspace*{1em}{part}="{F}">}\mbox{}\newline 
\hspace*{1em}\hspace*{1em}{<\textbf{l}>}The air is cut away before,{</\textbf{l}>}\mbox{}\newline 
\hspace*{1em}\hspace*{1em}{<\textbf{l}>}And closes from behind.{</\textbf{l}>}\mbox{}\newline 
\hspace*{1em}{</\textbf{lg}>}\mbox{}\newline 
{</\textbf{sp}>}\end{shaded}\egroup\par \par
For alternative methods of aligning groups of lines which do not form simple hierarchic groups, or which are discontinuous, see the more detailed discussion in chapter \textit{\hyperref[SA]{16.\ Linking, Segmentation, and Alignment}}. For discussion of other elements and attributes specific to the encoding of verse, see chapter \textit{\hyperref[VE]{6.\ Verse}}.
\subsubsection[{Core Tags for Drama}]{Core Tags for Drama}\label{CODR}\par
Like other written texts, dramatic and other \textit{performance texts} such as cinema or TV scripts are often hierarchically organized, for example into acts and scenes. These structural subdivisions should be encoded using the general purpose \hyperref[TEI.div]{<div>} or \hyperref[TEI.div1]{<div1>} (etc.) elements described below in chapters \textit{\hyperref[DS]{4.\ Default Text Structure}} and \textit{\hyperref[DR]{7.\ Performance Texts}}. Within these divisions, the body of a performance text typically consists of \textit{speeches}, often prefixed by a phrase indicating who is speaking, and occasionally interspersed with stage directions of various kinds.\par
In the following simple example, each speech consists of a single paragraph: \par\bgroup\index{div2=<div2>|exampleindex}\index{n=@n!<div2>|exampleindex}\index{type=@type!<div2>|exampleindex}\index{head=<head>|exampleindex}\index{stage=<stage>|exampleindex}\index{type=@type!<stage>|exampleindex}\index{sp=<sp>|exampleindex}\index{speaker=<speaker>|exampleindex}\index{p=<p>|exampleindex}\index{sp=<sp>|exampleindex}\index{speaker=<speaker>|exampleindex}\index{p=<p>|exampleindex}\index{sp=<sp>|exampleindex}\index{speaker=<speaker>|exampleindex}\index{p=<p>|exampleindex}\exampleFont \begin{shaded}\noindent\mbox{}{<\textbf{div2}\hspace*{1em}{n}="{I.2}"\hspace*{1em}{type}="{scene}">}\mbox{}\newline 
\hspace*{1em}{<\textbf{head}>}Scene 2.{</\textbf{head}>}\mbox{}\newline 
\hspace*{1em}{<\textbf{stage}\hspace*{1em}{type}="{setting}">}Peachum, Filch.{</\textbf{stage}>}\mbox{}\newline 
\hspace*{1em}{<\textbf{sp}>}\mbox{}\newline 
\hspace*{1em}\hspace*{1em}{<\textbf{speaker}>}FILCH.{</\textbf{speaker}>}\mbox{}\newline 
\hspace*{1em}\hspace*{1em}{<\textbf{p}>}Sir, Black Moll hath sent word her Trial comes on in\mbox{}\newline 
\hspace*{1em}\hspace*{1em}\hspace*{1em}\hspace*{1em} the Afternoon, and she hopes you will order Matters\mbox{}\newline 
\hspace*{1em}\hspace*{1em}\hspace*{1em}\hspace*{1em} so as to bring her off.{</\textbf{p}>}\mbox{}\newline 
\hspace*{1em}{</\textbf{sp}>}\mbox{}\newline 
\hspace*{1em}{<\textbf{sp}>}\mbox{}\newline 
\hspace*{1em}\hspace*{1em}{<\textbf{speaker}>}PEACHUM.{</\textbf{speaker}>}\mbox{}\newline 
\hspace*{1em}\hspace*{1em}{<\textbf{p}>}Why, she may plead her Belly at worst; to my\mbox{}\newline 
\hspace*{1em}\hspace*{1em}\hspace*{1em}\hspace*{1em} Knowledge she hath taken care of that Security.\mbox{}\newline 
\hspace*{1em}\hspace*{1em}\hspace*{1em}\hspace*{1em} But, as the Wench is very active and industrious,\mbox{}\newline 
\hspace*{1em}\hspace*{1em}\hspace*{1em}\hspace*{1em} you may satisfy her that I'll soften the Evidence.{</\textbf{p}>}\mbox{}\newline 
\hspace*{1em}{</\textbf{sp}>}\mbox{}\newline 
\hspace*{1em}{<\textbf{sp}>}\mbox{}\newline 
\hspace*{1em}\hspace*{1em}{<\textbf{speaker}>}FILCH.{</\textbf{speaker}>}\mbox{}\newline 
\hspace*{1em}\hspace*{1em}{<\textbf{p}>}Tom Gagg, sir, is found guilty.{</\textbf{p}>}\mbox{}\newline 
\hspace*{1em}{</\textbf{sp}>}\mbox{}\newline 
{</\textbf{div2}>}\end{shaded}\egroup\par \par
In the following example, each speech consists of a sequence of verse lines, some of them being marked as metrically incomplete: \par\bgroup\index{div1=<div1>|exampleindex}\index{n=@n!<div1>|exampleindex}\index{type=@type!<div1>|exampleindex}\index{head=<head>|exampleindex}\index{div2=<div2>|exampleindex}\index{n=@n!<div2>|exampleindex}\index{type=@type!<div2>|exampleindex}\index{head=<head>|exampleindex}\index{stage=<stage>|exampleindex}\index{rend=@rend!<stage>|exampleindex}\index{sp=<sp>|exampleindex}\index{speaker=<speaker>|exampleindex}\index{l=<l>|exampleindex}\index{part=@part!<l>|exampleindex}\index{sp=<sp>|exampleindex}\index{speaker=<speaker>|exampleindex}\index{l=<l>|exampleindex}\index{sp=<sp>|exampleindex}\index{speaker=<speaker>|exampleindex}\index{l=<l>|exampleindex}\index{part=@part!<l>|exampleindex}\index{sp=<sp>|exampleindex}\index{speaker=<speaker>|exampleindex}\index{l=<l>|exampleindex}\index{part=@part!<l>|exampleindex}\index{sp=<sp>|exampleindex}\index{speaker=<speaker>|exampleindex}\index{l=<l>|exampleindex}\index{part=@part!<l>|exampleindex}\index{sp=<sp>|exampleindex}\index{speaker=<speaker>|exampleindex}\index{l=<l>|exampleindex}\index{sp=<sp>|exampleindex}\index{speaker=<speaker>|exampleindex}\index{l=<l>|exampleindex}\index{sp=<sp>|exampleindex}\index{speaker=<speaker>|exampleindex}\index{l=<l>|exampleindex}\index{l=<l>|exampleindex}\index{part=@part!<l>|exampleindex}\exampleFont \begin{shaded}\noindent\mbox{}{<\textbf{div1}\hspace*{1em}{n}="{I}"\hspace*{1em}{type}="{Act}">}\mbox{}\newline 
\hspace*{1em}{<\textbf{head}>}ACT I{</\textbf{head}>}\mbox{}\newline 
\hspace*{1em}{<\textbf{div2}\hspace*{1em}{n}="{1}"\hspace*{1em}{type}="{Scene}">}\mbox{}\newline 
\hspace*{1em}\hspace*{1em}{<\textbf{head}>}SCENE I{</\textbf{head}>}\mbox{}\newline 
\hspace*{1em}\hspace*{1em}{<\textbf{stage}\hspace*{1em}{rend}="{italic}">}Enter Barnardo and Francisco,\mbox{}\newline 
\hspace*{1em}\hspace*{1em}\hspace*{1em}\hspace*{1em} two Sentinels, at several doors{</\textbf{stage}>}\mbox{}\newline 
\hspace*{1em}\hspace*{1em}{<\textbf{sp}>}\mbox{}\newline 
\hspace*{1em}\hspace*{1em}\hspace*{1em}{<\textbf{speaker}>}Barn{</\textbf{speaker}>}\mbox{}\newline 
\hspace*{1em}\hspace*{1em}\hspace*{1em}{<\textbf{l}\hspace*{1em}{part}="{Y}">}Who's there?{</\textbf{l}>}\mbox{}\newline 
\hspace*{1em}\hspace*{1em}{</\textbf{sp}>}\mbox{}\newline 
\hspace*{1em}\hspace*{1em}{<\textbf{sp}>}\mbox{}\newline 
\hspace*{1em}\hspace*{1em}\hspace*{1em}{<\textbf{speaker}>}Fran{</\textbf{speaker}>}\mbox{}\newline 
\hspace*{1em}\hspace*{1em}\hspace*{1em}{<\textbf{l}>}Nay, answer me. Stand and unfold yourself.{</\textbf{l}>}\mbox{}\newline 
\hspace*{1em}\hspace*{1em}{</\textbf{sp}>}\mbox{}\newline 
\hspace*{1em}\hspace*{1em}{<\textbf{sp}>}\mbox{}\newline 
\hspace*{1em}\hspace*{1em}\hspace*{1em}{<\textbf{speaker}>}Barn{</\textbf{speaker}>}\mbox{}\newline 
\hspace*{1em}\hspace*{1em}\hspace*{1em}{<\textbf{l}\hspace*{1em}{part}="{I}">}Long live the King!{</\textbf{l}>}\mbox{}\newline 
\hspace*{1em}\hspace*{1em}{</\textbf{sp}>}\mbox{}\newline 
\hspace*{1em}\hspace*{1em}{<\textbf{sp}>}\mbox{}\newline 
\hspace*{1em}\hspace*{1em}\hspace*{1em}{<\textbf{speaker}>}Fran{</\textbf{speaker}>}\mbox{}\newline 
\hspace*{1em}\hspace*{1em}\hspace*{1em}{<\textbf{l}\hspace*{1em}{part}="{M}">}Barnardo?{</\textbf{l}>}\mbox{}\newline 
\hspace*{1em}\hspace*{1em}{</\textbf{sp}>}\mbox{}\newline 
\hspace*{1em}\hspace*{1em}{<\textbf{sp}>}\mbox{}\newline 
\hspace*{1em}\hspace*{1em}\hspace*{1em}{<\textbf{speaker}>}Barn{</\textbf{speaker}>}\mbox{}\newline 
\hspace*{1em}\hspace*{1em}\hspace*{1em}{<\textbf{l}\hspace*{1em}{part}="{F}">}He.{</\textbf{l}>}\mbox{}\newline 
\hspace*{1em}\hspace*{1em}{</\textbf{sp}>}\mbox{}\newline 
\hspace*{1em}\hspace*{1em}{<\textbf{sp}>}\mbox{}\newline 
\hspace*{1em}\hspace*{1em}\hspace*{1em}{<\textbf{speaker}>}Fran{</\textbf{speaker}>}\mbox{}\newline 
\hspace*{1em}\hspace*{1em}\hspace*{1em}{<\textbf{l}>}You come most carefully upon your hour.{</\textbf{l}>}\mbox{}\newline 
\hspace*{1em}\hspace*{1em}{</\textbf{sp}>}\mbox{}\newline 
\hspace*{1em}\hspace*{1em}{<\textbf{sp}>}\mbox{}\newline 
\hspace*{1em}\hspace*{1em}\hspace*{1em}{<\textbf{speaker}>}Barn{</\textbf{speaker}>}\mbox{}\newline 
\hspace*{1em}\hspace*{1em}\hspace*{1em}{<\textbf{l}>}'Tis now struck twelve. Get thee to bed, Francisco.{</\textbf{l}>}\mbox{}\newline 
\hspace*{1em}\hspace*{1em}{</\textbf{sp}>}\mbox{}\newline 
\hspace*{1em}\hspace*{1em}{<\textbf{sp}>}\mbox{}\newline 
\hspace*{1em}\hspace*{1em}\hspace*{1em}{<\textbf{speaker}>}Fran{</\textbf{speaker}>}\mbox{}\newline 
\hspace*{1em}\hspace*{1em}\hspace*{1em}{<\textbf{l}>}For this relief much thanks. 'Tis bitter cold,{</\textbf{l}>}\mbox{}\newline 
\hspace*{1em}\hspace*{1em}\hspace*{1em}{<\textbf{l}\hspace*{1em}{part}="{I}">}And I am sick at heart.{</\textbf{l}>}\mbox{}\newline 
\hspace*{1em}\hspace*{1em}{</\textbf{sp}>}\mbox{}\newline 
\hspace*{1em}{</\textbf{div2}>}\mbox{}\newline 
{</\textbf{div1}>}\end{shaded}\egroup\par \par
In some cases, as here in the First Quarto of \textit{Hamlet}, the printed speaker attributions need to be supplemented by use of the {\itshape who} attribute; again, the lines are marked as complete or incomplete: \par\bgroup\index{stage=<stage>|exampleindex}\index{add=<add>|exampleindex}\index{place=@place!<add>|exampleindex}\index{name=<name>|exampleindex}\index{name=<name>|exampleindex}\index{sp=<sp>|exampleindex}\index{who=@who!<sp>|exampleindex}\index{speaker=<speaker>|exampleindex}\index{l=<l>|exampleindex}\index{part=@part!<l>|exampleindex}\index{sp=<sp>|exampleindex}\index{who=@who!<sp>|exampleindex}\index{speaker=<speaker>|exampleindex}\index{l=<l>|exampleindex}\index{part=@part!<l>|exampleindex}\index{sp=<sp>|exampleindex}\index{who=@who!<sp>|exampleindex}\index{speaker=<speaker>|exampleindex}\index{l=<l>|exampleindex}\index{sp=<sp>|exampleindex}\index{who=@who!<sp>|exampleindex}\index{speaker=<speaker>|exampleindex}\index{l=<l>|exampleindex}\index{l=<l>|exampleindex}\index{sp=<sp>|exampleindex}\index{who=@who!<sp>|exampleindex}\index{speaker=<speaker>|exampleindex}\index{l=<l>|exampleindex}\index{part=@part!<l>|exampleindex}\index{stage=<stage>|exampleindex}\exampleFont \begin{shaded}\noindent\mbox{}{<\textbf{stage}>}Enter two Centinels.\mbox{}\newline 
{<\textbf{add}\hspace*{1em}{place}="{margin}">}Now call'd {<\textbf{name}\hspace*{1em}{xml:id}="{barnardo}">}Bernardo{</\textbf{name}>} \&\mbox{}\newline 
\hspace*{1em}{<\textbf{name}\hspace*{1em}{xml:id}="{francisco}">}Francesco{</\textbf{name}>}.{</\textbf{add}>}\mbox{}\newline 
{</\textbf{stage}>}\mbox{}\newline 
{<\textbf{sp}\hspace*{1em}{who}="{\#francisco}">}\mbox{}\newline 
\hspace*{1em}{<\textbf{speaker}>}1.{</\textbf{speaker}>}\mbox{}\newline 
\hspace*{1em}{<\textbf{l}\hspace*{1em}{part}="{Y}">}Stand: who is that?{</\textbf{l}>}\mbox{}\newline 
{</\textbf{sp}>}\mbox{}\newline 
{<\textbf{sp}\hspace*{1em}{who}="{\#barnardo}">}\mbox{}\newline 
\hspace*{1em}{<\textbf{speaker}>}2.{</\textbf{speaker}>}\mbox{}\newline 
\hspace*{1em}{<\textbf{l}\hspace*{1em}{part}="{Y}">}Tis I.{</\textbf{l}>}\mbox{}\newline 
{</\textbf{sp}>}\mbox{}\newline 
{<\textbf{sp}\hspace*{1em}{who}="{\#francisco}">}\mbox{}\newline 
\hspace*{1em}{<\textbf{speaker}>}1.{</\textbf{speaker}>}\mbox{}\newline 
\hspace*{1em}{<\textbf{l}>}O you come most carefully vpon your watch,{</\textbf{l}>}\mbox{}\newline 
{</\textbf{sp}>}\mbox{}\newline 
{<\textbf{sp}\hspace*{1em}{who}="{\#barnardo}">}\mbox{}\newline 
\hspace*{1em}{<\textbf{speaker}>}2.{</\textbf{speaker}>}\mbox{}\newline 
\hspace*{1em}{<\textbf{l}>}And if you meete Marcellus and Horatio,{</\textbf{l}>}\mbox{}\newline 
\hspace*{1em}{<\textbf{l}>}The partners of my watch, bid them make haste.{</\textbf{l}>}\mbox{}\newline 
{</\textbf{sp}>}\mbox{}\newline 
{<\textbf{sp}\hspace*{1em}{who}="{\#francisco}">}\mbox{}\newline 
\hspace*{1em}{<\textbf{speaker}>}1.{</\textbf{speaker}>}\mbox{}\newline 
\hspace*{1em}{<\textbf{l}\hspace*{1em}{part}="{Y}">}I will: See who goes there.{</\textbf{l}>}\mbox{}\newline 
{</\textbf{sp}>}\mbox{}\newline 
{<\textbf{stage}>}Enter Horatio and Marcellus.{</\textbf{stage}>}\end{shaded}\egroup\par \par
By contrast with the preceding examples, the following encodes an early printed edition without making any assumption about which parts are prose or verse: \par\bgroup\index{div1=<div1>|exampleindex}\index{n=@n!<div1>|exampleindex}\index{type=@type!<div1>|exampleindex}\index{div2=<div2>|exampleindex}\index{n=@n!<div2>|exampleindex}\index{type=@type!<div2>|exampleindex}\index{head=<head>|exampleindex}\index{rend=@rend!<head>|exampleindex}\index{stage=<stage>|exampleindex}\index{rend=@rend!<stage>|exampleindex}\index{type=@type!<stage>|exampleindex}\index{sp=<sp>|exampleindex}\index{speaker=<speaker>|exampleindex}\index{p=<p>|exampleindex}\index{sp=<sp>|exampleindex}\index{speaker=<speaker>|exampleindex}\index{p=<p>|exampleindex}\index{sp=<sp>|exampleindex}\index{speaker=<speaker>|exampleindex}\index{p=<p>|exampleindex}\index{stage=<stage>|exampleindex}\index{type=@type!<stage>|exampleindex}\index{stage=<stage>|exampleindex}\index{type=@type!<stage>|exampleindex}\index{sp=<sp>|exampleindex}\index{speaker=<speaker>|exampleindex}\index{p=<p>|exampleindex}\exampleFont \begin{shaded}\noindent\mbox{}{<\textbf{div1}\hspace*{1em}{n}="{I}"\hspace*{1em}{type}="{act}">}\mbox{}\newline 
\hspace*{1em}{<\textbf{div2}\hspace*{1em}{n}="{1}"\hspace*{1em}{type}="{scene}">}\mbox{}\newline 
\hspace*{1em}\hspace*{1em}{<\textbf{head}\hspace*{1em}{rend}="{italic}">}Actus primus, Scena prima.{</\textbf{head}>}\mbox{}\newline 
\hspace*{1em}\hspace*{1em}{<\textbf{stage}\hspace*{1em}{rend}="{italic}"\hspace*{1em}{type}="{setting}">}A tempestuous\mbox{}\newline 
\hspace*{1em}\hspace*{1em}\hspace*{1em}\hspace*{1em} noise of Thunder and Lightning heard: Enter\mbox{}\newline 
\hspace*{1em}\hspace*{1em}\hspace*{1em}\hspace*{1em} a Ship-master, and a Boteswaine.{</\textbf{stage}>}\mbox{}\newline 
\hspace*{1em}\hspace*{1em}{<\textbf{sp}>}\mbox{}\newline 
\hspace*{1em}\hspace*{1em}\hspace*{1em}{<\textbf{speaker}>}Master.{</\textbf{speaker}>}\mbox{}\newline 
\hspace*{1em}\hspace*{1em}\hspace*{1em}{<\textbf{p}>}Bote-swaine.{</\textbf{p}>}\mbox{}\newline 
\hspace*{1em}\hspace*{1em}{</\textbf{sp}>}\mbox{}\newline 
\hspace*{1em}\hspace*{1em}{<\textbf{sp}>}\mbox{}\newline 
\hspace*{1em}\hspace*{1em}\hspace*{1em}{<\textbf{speaker}>}Botes.{</\textbf{speaker}>}\mbox{}\newline 
\hspace*{1em}\hspace*{1em}\hspace*{1em}{<\textbf{p}>}Heere Master: What cheere?{</\textbf{p}>}\mbox{}\newline 
\hspace*{1em}\hspace*{1em}{</\textbf{sp}>}\mbox{}\newline 
\hspace*{1em}\hspace*{1em}{<\textbf{sp}>}\mbox{}\newline 
\hspace*{1em}\hspace*{1em}\hspace*{1em}{<\textbf{speaker}>}Mast.{</\textbf{speaker}>}\mbox{}\newline 
\hspace*{1em}\hspace*{1em}\hspace*{1em}{<\textbf{p}>}Good: Speake to th' Mariners: fall\mbox{}\newline 
\hspace*{1em}\hspace*{1em}\hspace*{1em}\hspace*{1em}\hspace*{1em}\hspace*{1em} too't, yarely, or we run our selues a ground,\mbox{}\newline 
\hspace*{1em}\hspace*{1em}\hspace*{1em}\hspace*{1em}\hspace*{1em}\hspace*{1em} bestirre, bestirre. {<\textbf{stage}\hspace*{1em}{type}="{move}">}Exit.{</\textbf{stage}>}\mbox{}\newline 
\hspace*{1em}\hspace*{1em}\hspace*{1em}{</\textbf{p}>}\mbox{}\newline 
\hspace*{1em}\hspace*{1em}{</\textbf{sp}>}\mbox{}\newline 
\hspace*{1em}\hspace*{1em}{<\textbf{stage}\hspace*{1em}{type}="{move}">}Enter Mariners.{</\textbf{stage}>}\mbox{}\newline 
\hspace*{1em}\hspace*{1em}{<\textbf{sp}>}\mbox{}\newline 
\hspace*{1em}\hspace*{1em}\hspace*{1em}{<\textbf{speaker}>}Botes.{</\textbf{speaker}>}\mbox{}\newline 
\hspace*{1em}\hspace*{1em}\hspace*{1em}{<\textbf{p}>}Heigh my hearts, cheerely, cheerely my harts: yare,\mbox{}\newline 
\hspace*{1em}\hspace*{1em}\hspace*{1em}\hspace*{1em}\hspace*{1em}\hspace*{1em} yare: Take in the toppe-sale: Tend to th' Masters whistle:\mbox{}\newline 
\hspace*{1em}\hspace*{1em}\hspace*{1em}\hspace*{1em}\hspace*{1em}\hspace*{1em} Blow till thou burst thy winde, if roome e-nough.{</\textbf{p}>}\mbox{}\newline 
\hspace*{1em}\hspace*{1em}{</\textbf{sp}>}\mbox{}\newline 
\hspace*{1em}{</\textbf{div2}>}\mbox{}\newline 
{</\textbf{div1}>}\end{shaded}\egroup\par \par
The \hyperref[TEI.sp]{<sp>} and \hyperref[TEI.stage]{<stage>} elements should also be used to mark parts of a text otherwise in prose which are presented as if they were dialogue in a play. The following example is taken from a 19th century novel in which passages of narrative and passages of dialogue are mixed within the same chapter: \par\bgroup\index{sp=<sp>|exampleindex}\index{speaker=<speaker>|exampleindex}\index{p=<p>|exampleindex}\index{sp=<sp>|exampleindex}\index{speaker=<speaker>|exampleindex}\index{p=<p>|exampleindex}\index{sp=<sp>|exampleindex}\index{speaker=<speaker>|exampleindex}\index{p=<p>|exampleindex}\index{p=<p>|exampleindex}\exampleFont \begin{shaded}\noindent\mbox{}{<\textbf{sp}>}\mbox{}\newline 
\hspace*{1em}{<\textbf{speaker}>}The reverend Doctor Opimian{</\textbf{speaker}>}\mbox{}\newline 
\hspace*{1em}{<\textbf{p}>}I do not think I have named a single unpresentable fish.{</\textbf{p}>}\mbox{}\newline 
{</\textbf{sp}>}\mbox{}\newline 
{<\textbf{sp}>}\mbox{}\newline 
\hspace*{1em}{<\textbf{speaker}>}Mr Gryll{</\textbf{speaker}>}\mbox{}\newline 
\hspace*{1em}{<\textbf{p}>}Bream, Doctor: there is not much to be said for bream.{</\textbf{p}>}\mbox{}\newline 
{</\textbf{sp}>}\mbox{}\newline 
{<\textbf{sp}>}\mbox{}\newline 
\hspace*{1em}{<\textbf{speaker}>}The Reverend Doctor Opimian{</\textbf{speaker}>}\mbox{}\newline 
\hspace*{1em}{<\textbf{p}>}On the contrary, sir, I think there is much to be said for him.\mbox{}\newline 
\hspace*{1em}\hspace*{1em} In the first place ...{</\textbf{p}>}\mbox{}\newline 
\hspace*{1em}{<\textbf{p}>}Fish, Miss Gryll — I could discourse to you on fish by the\mbox{}\newline 
\hspace*{1em}\hspace*{1em} hour: but for the present I will forbear ...{</\textbf{p}>}\mbox{}\newline 
{</\textbf{sp}>}\end{shaded}\egroup\par \noindent  \par\bgroup\index{sp=<sp>|exampleindex}\index{speaker=<speaker>|exampleindex}\index{stage=<stage>|exampleindex}\index{p=<p>|exampleindex}\index{q=<q>|exampleindex}\index{title=<title>|exampleindex}\index{q=<q>|exampleindex}\index{p=<p>|exampleindex}\exampleFont \begin{shaded}\noindent\mbox{}{<\textbf{sp}>}\mbox{}\newline 
\hspace*{1em}{<\textbf{speaker}>}Lord Curryfin{</\textbf{speaker}>}\mbox{}\newline 
\hspace*{1em}{<\textbf{stage}>}(after a pause).{</\textbf{stage}>}\mbox{}\newline 
\hspace*{1em}{<\textbf{p}>}\mbox{}\newline 
\hspace*{1em}\hspace*{1em}{<\textbf{q}>}Mass{</\textbf{q}>} as the second grave-digger says\mbox{}\newline 
\hspace*{1em}\hspace*{1em} in {<\textbf{title}>}Hamlet{</\textbf{title}>}, {<\textbf{q}>}I cannot tell.{</\textbf{q}>}\mbox{}\newline 
\hspace*{1em}{</\textbf{p}>}\mbox{}\newline 
{</\textbf{sp}>}\mbox{}\newline 
{<\textbf{p}>}A chorus of laughter dissolved the sitting.{</\textbf{p}>}\end{shaded}\egroup\par 
\subsection[{Overview of the Core Module }]{Overview of the Core Module }\label{COOV}\par
All the elements described in this chapter are provided by the \textsf{core} module. \begin{description}

\item[{Module core: Elements common to all TEI documents}]\hspace{1em}\hfill\linebreak
\mbox{}\\[-10pt] \begin{itemize}
\item {\itshape Elements defined}: \hyperref[TEI.abbr]{abbr} \hyperref[TEI.add]{add} \hyperref[TEI.addrLine]{addrLine} \hyperref[TEI.address]{address} \hyperref[TEI.analytic]{analytic} \hyperref[TEI.author]{author} \hyperref[TEI.bibl]{bibl} \hyperref[TEI.biblScope]{biblScope} \hyperref[TEI.biblStruct]{biblStruct} \hyperref[TEI.binaryObject]{binaryObject} \hyperref[TEI.cb]{cb} \hyperref[TEI.choice]{choice} \hyperref[TEI.cit]{cit} \hyperref[TEI.citedRange]{citedRange} \hyperref[TEI.corr]{corr} \hyperref[TEI.date]{date} \hyperref[TEI.del]{del} \hyperref[TEI.desc]{desc} \hyperref[TEI.distinct]{distinct} \hyperref[TEI.divGen]{divGen} \hyperref[TEI.editor]{editor} \hyperref[TEI.email]{email} \hyperref[TEI.emph]{emph} \hyperref[TEI.expan]{expan} \hyperref[TEI.foreign]{foreign} \hyperref[TEI.gap]{gap} \hyperref[TEI.gb]{gb} \hyperref[TEI.gloss]{gloss} \hyperref[TEI.graphic]{graphic} \hyperref[TEI.head]{head} \hyperref[TEI.headItem]{headItem} \hyperref[TEI.headLabel]{headLabel} \hyperref[TEI.hi]{hi} \hyperref[TEI.imprint]{imprint} \hyperref[TEI.index]{index} \hyperref[TEI.item]{item} \hyperref[TEI.l]{l} \hyperref[TEI.label]{label} \hyperref[TEI.lb]{lb} \hyperref[TEI.lg]{lg} \hyperref[TEI.list]{list} \hyperref[TEI.listBibl]{listBibl} \hyperref[TEI.measure]{measure} \hyperref[TEI.measureGrp]{measureGrp} \hyperref[TEI.media]{media} \hyperref[TEI.meeting]{meeting} \hyperref[TEI.mentioned]{mentioned} \hyperref[TEI.milestone]{milestone} \hyperref[TEI.monogr]{monogr} \hyperref[TEI.name]{name} \hyperref[TEI.note]{note} \hyperref[TEI.noteGrp]{noteGrp} \hyperref[TEI.num]{num} \hyperref[TEI.orig]{orig} \hyperref[TEI.p]{p} \hyperref[TEI.pb]{pb} \hyperref[TEI.postBox]{postBox} \hyperref[TEI.postCode]{postCode} \hyperref[TEI.ptr]{ptr} \hyperref[TEI.pubPlace]{pubPlace} \hyperref[TEI.publisher]{publisher} \hyperref[TEI.q]{q} \hyperref[TEI.quote]{quote} \hyperref[TEI.rb]{rb} \hyperref[TEI.ref]{ref} \hyperref[TEI.reg]{reg} \hyperref[TEI.relatedItem]{relatedItem} \hyperref[TEI.resp]{resp} \hyperref[TEI.respStmt]{respStmt} \hyperref[TEI.rs]{rs} \hyperref[TEI.rt]{rt} \hyperref[TEI.ruby]{ruby} \hyperref[TEI.said]{said} \hyperref[TEI.series]{series} \hyperref[TEI.sic]{sic} \hyperref[TEI.soCalled]{soCalled} \hyperref[TEI.sp]{sp} \hyperref[TEI.speaker]{speaker} \hyperref[TEI.stage]{stage} \hyperref[TEI.street]{street} \hyperref[TEI.teiCorpus]{teiCorpus} \hyperref[TEI.term]{term} \hyperref[TEI.textLang]{textLang} \hyperref[TEI.time]{time} \hyperref[TEI.title]{title} \hyperref[TEI.unclear]{unclear} \hyperref[TEI.unit]{unit}
\item {\itshape Classes defined}: \hyperref[TEI.att.milestoneUnit]{att.milestoneUnit}
\end{itemize} 
\end{description}  The selection and combination of modules to form a TEI schema is described in \textit{\hyperref[STIN]{1.2.\ Defining a TEI Schema}}.